% !TeX program = xelatex
% Standalone English volume. Use XeLaTeX until the references stabilize.
\documentclass[a4paper,11pt]{article}
\usepackage[T1]{fontenc}
\usepackage{lmodern}
\usepackage[english]{babel}
\usepackage[left=26mm,right=26mm,top=25mm,bottom=25mm,headheight=15pt,headsep=7mm,footskip=11mm]{geometry}
\usepackage{amsmath,amssymb,mathtools,amsthm}
\newtheorem{proposition}{Proposition}[section]
\newtheorem{lemma}[proposition]{Lemma}
\newtheorem{theorem}[proposition]{Theorem}
\newtheorem{corollary}[proposition]{Corollary}
\newtheorem{maintheorem}[proposition]{Main Theorem}
\theoremstyle{definition}
\newtheorem{definition}[proposition]{Definition}
\theoremstyle{remark}
\newtheorem{remark}[proposition]{Remark}
\makeatletter
\renewcommand*\l@section{\@dottedtocline{1}{0em}{2.4em}}
\renewcommand*\l@subsection{\@dottedtocline{2}{1.45em}{3.45em}}
\renewcommand*\l@subsubsection{\@dottedtocline{3}{3.8em}{4.55em}}
\makeatother
\usepackage{xcolor,fancyhdr,titlesec,needspace,xurl}
\usepackage{longtable,multicol}
\usepackage{hyperref}
\definecolor{ink}{HTML}{1A3542}
\definecolor{muted}{HTML}{4F606A}
\hypersetup{unicode=true,colorlinks=true,linkcolor=ink,filecolor=ink,urlcolor=ink,citecolor=ink,bookmarksnumbered=true,
 pdfdisplaydoctitle=true,pdflang={en},
 pdftitle={The generalized Riemann hypothesis for Dirichlet L-functions. Edition 930, Volume XXV},
 pdfsubject={A natural-language reconstruction of the proposed proof},
 pdfauthor={A mathematical reconstruction}}
\titleformat{\section}{\Large\bfseries\color{ink}}{\thesection}{.65em}{}
\titlespacing*{\section}{0pt}{.8em}{.65em}
\setlength{\parindent}{1.2em}
\linespread{1.045}
\setlength{\parskip}{.22em plus .05em minus .03em}
\setlength{\emergencystretch}{2em}
\setcounter{topnumber}{2}
\setcounter{tocdepth}{3}\setcounter{secnumdepth}{3}
\widowpenalty=4500
\clubpenalty=4500
\pagestyle{fancy}\fancyhf{}
\fancyhead[L]{\small\color{muted}Dirichlet GRH \textperiodcentered\ Volume I}
\fancyhead[R]{\small\color{muted}Dirichlet zeros and power intervals}
\fancyfoot[C]{\small\thepage}
\renewcommand{\headrulewidth}{.3pt}
\newcommand{\lead}[1]{\par\smallskip\noindent{\bfseries\color{ink}#1.}\enspace}
\newcommand{\sour}[2]{\unskip\penalty150\hspace{.35em}{\small\hyperlink{ref#1}{[#2]}}}
\titleformat{\subsection}{\large\bfseries}{\thesubsection}{.6em}{}
\titlespacing*{\subsection}{0pt}{1.15em plus .15em}{.4em}
\setlength{\abovedisplayskip}{9pt plus 2pt minus 3pt}
\setlength{\belowdisplayskip}{9pt plus 2pt minus 3pt}
\newcommand{\RR}{\mathbb R}\newcommand{\CC}{\mathbb C}\newcommand{\QQ}{\mathbb Q}
\newcommand{\ZZ}{\mathbb Z}\newcommand{\ord}{\operatorname{ord}}
\newcommand{\Frob}{\operatorname{Frob}}\newcommand{\Gal}{\operatorname{Gal}}
\newcommand{\len}{\operatorname{length}}\newcommand{\Res}{\operatorname*{Res}}
\newcommand{\wmean}{\operatorname{av}}\newcommand{\Ker}{\operatorname{ker}}
\newcommand{\Lh}{\mathcal L}\newcommand{\Gg}{\mathcal G}
\newcommand{\real}{\operatorname{Re}}\newcommand{\imag}{\operatorname{Im}}
\newcommand{\abs}[1]{\left|#1\right|}\newcommand{\eps}{\varepsilon}
\newcommand{\anchoritem}[2]{\par\Needspace{2.8\baselineskip}\hypertarget{ref#1}{}\noindent{\bfseries\color{ink}#2.}\enspace}
\newcommand{\srcloc}[2]{{\small\nolinkurl{#1.lean}}, line~#2}
\numberwithin{equation}{section}
\raggedbottom

\newcommand{\Sel}{\mathcal S}
\newcommand{\Cl}{\operatorname{Cl}}
\newcommand{\Aut}{\operatorname{Aut}}
\newcommand{\GL}{\operatorname{GL}}
\newcommand{\OO}{\mathcal O}
\hypersetup{bookmarksdepth=3}


\addto\captionsenglish{\renewcommand{\partname}{Volume}}


\makeatletter
\expandafter\def\csname r@cor:oddwindow677\endcsname{{11.30}{27}{}{proposition.11.30}{Dirichlet-GRH-v930-Volume-II.pdf}}
\expandafter\def\csname r@lem:covol127\endcsname{{92.23}{19}{}{proposition.92.23}{Dirichlet-GRH-v930-Volume-XVIII.pdf}}
\expandafter\def\csname r@lem:mass209\endcsname{{93.7}{26}{}{proposition.93.7}{Dirichlet-GRH-v930-Volume-XVIII.pdf}}
\expandafter\def\csname r@lem:terminalintegral867\endcsname{{84.13}{90}{}{proposition.84.13}{Dirichlet-GRH-v930-Volume-XVI.pdf}}
\expandafter\def\csname r@seriesguide\endcsname{{1.1}{4}{}{subsection.1.1}{Dirichlet-GRH-v930-Volume-I.pdf}}
\expandafter\def\csname r@thm:commonsource806\endcsname{{73.8}{81}{}{proposition.73.8}{Dirichlet-GRH-v930-Volume-XIV.pdf}}
\expandafter\def\csname r@thm:longframe867\endcsname{{84.14}{91}{}{proposition.84.14}{Dirichlet-GRH-v930-Volume-XVI.pdf}}
\expandafter\def\csname r@thm:norm91\endcsname{{94.9}{44}{}{proposition.94.9}{Dirichlet-GRH-v930-Volume-XVIII.pdf}}
\expandafter\def\csname r@thm:oddinteger677\endcsname{{11.29}{26}{}{proposition.11.29}{Dirichlet-GRH-v930-Volume-II.pdf}}
\makeatother
\begin{document}
\setcounter{part}{24}
\fancyhead[L]{\small\color{muted}Dirichlet GRH \textperiodcentered\ Volume XXV}
\fancyhead[R]{\small\color{muted}Rational flags and the second moment}
\thispagestyle{plain}\begin{center}
{\LARGE\bfseries The generalized Riemann hypothesis\\[2pt]for Dirichlet \(L\)-functions}\\[6pt]
{\small 27 November 2026 \quad\textperiodcentered\quad Edition 930}\end{center}
\part{Rational flags and the second moment of Bessel periods}\label{part:XXV}
\begin{abstract}\noindent We count pairs of rational flags, separate the open and closed orbital integrals and construct the real test vector near a stationary point. The resulting Poincar\'e series give a second moment on the marked classes. Its terminal norm constraint leads to the local correlations in the next volume.\end{abstract}
\setcounter{section}{129}
\noindent The \hyperref[seriesguide]{guide in Volume~I} lists all twenty-seven volumes.
\begingroup\small\setlength{\parskip}{0pt}\tableofcontents\endgroup
\clearpage

\section{Rational flags at the two marked primes}\label{nav:two-mark-flags}\label{nav:volume-XXV}

\subsection{The two level projections on definite automorphic forms}
\label{sec:levelmass678}

Consider the family \(E=\QQ(\sqrt{-d},\mu_{q_1})\) of
Corollary~\ref{cor:oddwindow677}, and put \(F=E^+\),
\(f=[F:\QQ]\). Write \(r\) for the rank of the unitary group
in this family. Its two distinct level primes \(p_0,p_1\) lie
in \([b^{10},b^{11}]\) and are one modulo \(Q^2\), where
\(Q=(q_0d)^{s+6}\) is a multiple of \(q_1d\). Thus each splits
completely in \(E\), by its action on cyclotomic roots. At either
prime there are \(f\) places of \(F\), with residue field
\(\mathbb F_{p_i}\); choose one of the two places of \(E\)
over each. The two residue groups are
\[
       H_i=\mathrm{GL}_r(\mathbb F_{p_i})^f,\qquad H=H_0\times H_1 .
\]

For a field \(\mathbb F_q\) and \(r\ge2\), call a monic
degree-\(r\) polynomial \emph{allowed} if it is squarefree,
has nonzero constant term, and has no root in
\(\{a\in\mathbb F_q^\times:a^Q=1\}\). Call it \emph{primitive}
if it is the minimal polynomial of a generator of
\(\mathbb F_{q^r}^\times\). Such a polynomial has degree \(r\)
and is allowed. Green's parametrization attaches distinct
irreducible characters to the squarefree degree-\(r\) polynomials
of nonzero constant term. In this parametrization one fixes
compatible generators for the character groups of the finite
tori. For a polynomial with irreducible factor degrees
\(d_1,\ldots,d_k\), take the one-box partition on each
corresponding Frobenius orbit. The resulting character \(\chi_c\)
has degree
\[
             \chi_c(1)=
        \frac{\prod_{j=1}^{r}(q^j-1)}{\prod_{j=1}^{k}(q^{d_j}-1)} .
\]
This is the squarefree case of Green's degree and classification
theorems.~\cite[\S8, Thms.~13--14]{XXV-XX-XX-green678}

In \(\mathbb C[H_i]\) let \(e_i\) be the sum of central
character idempotents for the \(f\)-tuples of allowed labels
having at least one primitive member. Write \(e=e_0e_1\), viewing
each factor in \(\mathbb C[H]\). We compute the fraction
\(\delta_i\) that \(e_i\) discards in the regular representation
and then evaluate its rank on the definite automorphic space.

\begin{proposition}[Plancherel mass at one level prime]
\leavevmode\label{prop:levelmass678}
Let \(q=p_i\), \(r\ge2\), \(q\ge3\), and let \(A\) be the
sum of \(\chi_c(1)^2/|\mathrm{GL}_r(\mathbb F_q)|\) over allowed
\(c\). Put \(n=q^r-1\). Then
\[
 P=\frac{\varphi(n)}{rn}\prod_{j=1}^{r-1}(1-q^{-j}),
 \qquad
 \delta_i=1-A^f+(A-P)^f .
\]
In particular
\[
 \delta_i\le f\left(\frac{2+\gcd(Q,q-1)}q+\frac1{q-1}\right)
                +\exp\!\left(-\frac{f}{4r(1+r\log q)}\right).
\]
For every sufficiently large node of the stated family,
\(\delta_0+\delta_1\le b^{-8}\), uniformly for its prescribed
level primes.
\end{proposition}
\begin{proof}
Counting ordered bases gives
\(|\mathrm{GL}_r(\mathbb F_q)|=q^{r(r-1)/2}\prod_{j=1}^r(q^j-1)\).
Consequently each squarefree label has mass
\[
 \frac{\chi_c(1)^2}{|\mathrm{GL}_r(\mathbb F_q)|}
   =q^{-r}\frac{\prod_{j=1}^r(1-q^{-j})}
                  {\prod_{j=1}^k(1-q^{-d_j})^2}
   \ge q^{-r}\Big(1-\frac1{q-1}\Big).
\]
There are \(\varphi(n)/r\) primitive labels. Indeed a generator
of the multiplicative group cannot lie in a smaller subfield,
and its \(r\) Frobenius conjugates have the same minimal
polynomial. Substitution of \(k=1,d_1=r\) in the exact mass
formula gives their total mass \(P\).

Among the \(q^r\) monic polynomials, exactly \(q^{r-1}\) have
a repeated factor. To see this, unique factorization gives
\(\prod_L(1+z^{\deg L})=(1-qz^2)/(1-qz)\) for monic irreducibles
\(L\), by taking the quotient of the two Euler products for all
monic polynomials at \(z\) and \(z^2\). Each prescribed root
excludes \(q^{r-1}\) polynomials. There are
\(1+\gcd(Q,q-1)\) excluded roots including zero. Taking the
union bound and the lower mass per allowed label gives
\[
        1-A\le (2+\gcd(Q,q-1))/q+1/(q-1).
\]
This follows directly when the root-count lower bound is
positive; otherwise the displayed upper bound is at least one.
The ranks of distinct character idempotents add, and
Plancherel masses multiply in a direct product. All labels
being allowed has probability \(A^f\); all being allowed with
none primitive has probability \((A-P)^f\). These events give
the asserted equality for \(\delta_i\).

For any positive integer \(n\),
\[
 \frac n{\varphi(n)}
   =\prod_{p\mid n}\frac{1+p^{-1}}{1-p^{-2}}
   \le \zeta(2)\sum_{d\mid n}d^{-1}\le 2(1+\log n).
\]
Also \(\prod_{j=1}^{r-1}(1-q^{-j})\ge 1-1/(q-1)\ge1/2\).
Thus \(P\ge1/[4r(1+r\log q)]\).
Using \(1-A^f\le f(1-A)\) and \((A-P)^f\le e^{-fP}\)
proves the inequality.

We verify both terms on the specified fields. Since \(E\)
contains \(\QQ(\mu_{2^s})\) and lies in \(\QQ(\mu_{q_1d})\),
Corollary~\ref{cor:oddwindow677} gives
\[
         x^{11}/4\le 2^{s-2}\le f
                   \ll_{q_0,d}x^{11}\log(2x).
\]
Here \(2^s=m^{22}\), \(m\ge\sqrt{x}\). Also \(s=O(\log x)\),
\(r\le3x^5\), \((\log x)^6\le\log b\le x^{1/100}\), and
\(10\log b\le\log q\le11\log b\).
Thus \(Q\) grows at most polynomially in \(x\) for fixed
\(q_0,d\), whereas \(b\) exceeds every such polynomial
eventually. The first term is at most \(b^{-9}\) for large
\(x\). For the second, \(r(1+r\log q)\log b\ll x^{10+2/100}\);
hence \(f\ge 36r(1+r\log q)\log b\) for all sufficiently
large \(x\). Each \(\delta_i\le2b^{-9}\), proving the claim.
\end{proof}

\begin{theorem}[Dimension and integral index on the definite endpoint]
\leavevmode\label{thm:definitelevels678}
Let the unitary group of rank \(r\) over \(F\) be definite at
all real places. Take the standard compact level at the two
split primes and any fixed compact open level away from them.
Let \(K\) impose principal reduction at both \(p_0,p_1\).
For algebraic automorphic forms \(V\) of level \(K\) and any
fixed finite-dimensional archimedean weight, one has
\[
 \dim (1-e_i)V=\delta_i\dim V,\qquad
 \dim(1-e)V\le(\delta_0+\delta_1)\dim V\le b^{-8}\dim V
\]
at sufficiently large nodes. The idempotent \(e\) lies in
\(|H|^{-1}\ZZ[H]\). If \(\Lambda_\ell\) is any \(H\)-stable
full lattice in a characteristic-zero model of \(V\) over a
finite extension of \(\QQ_\ell\), put
\[
 I_\ell=e\Lambda_\ell/(\Lambda_\ell\cap eV_\ell),\qquad
 d_\ell=\text{ramification index of the coefficient field}.
\]
For arbitrary such coefficient fields and lattices at each prime,
\[
  \sum_\ell\frac{\len_{\mathcal O_\ell}I_\ell}{d_\ell}\log\ell
       \le b^{-8}\dim V\,\log |H|,\qquad
  \log|H|\ll_{q_0,d}x^{21}\log(2x)\log b .
\]
\end{theorem}
\begin{proof}
Let \(K^{(i)}\) restore the full compact group at \(p_i\) while
retaining principal reduction at the other prime. Reduction of
invertible integral matrices is onto, by lifting their entries
and using the determinant. Hence \(K^{(i)}/K=H_i\).
Its right action on the finite level-\(K\) class set is free.
Indeed, for a finite adelic representative \(g\), the group
of rational points lying in \(gK^{(i)}g^{-1}\) is finite.
The diagonal rational points are closed and discrete in the
full adelic group, and both the indicated finite-adelic set and
the entire archimedean group are compact. At a place over
\(p_{1-i}\), an element of this finite group is conjugate into
\(1+p_{1-i}M_r(\ZZ_{p_{1-i}})\). This group is torsion free for
odd \(p=p_{1-i}\). In fact, if \(C=1+p^aB\) with \(a\ge1\)
and \(B\) nonzero modulo \(p\), the first nonzero term of
\(C^u-1\) has valuation \(a\) when \(p\nmid u\), and that of
\(C^p-1\) has valuation \(a+1\) by the binomial formula.
Iteration rules out every finite order. Thus the rational
element is the identity at one completion, and hence is the
identity. If a lift in \(K^{(i)}\) fixes the double coset of
\(g\), this trivial intersection puts the lift in \(K\),
proving freeness.

The level-\(K\) forms have the full archimedean weight space as
fiber over each double coset, since its stabilizer is also
trivial. The finite class set here is the usual class set of
positive-definite Hermitian lattices at a compact open level.
On a free \(H_i\)-orbit identify each fiber with the fiber over
one base point by translation. The sections over this orbit
then form copies of the regular representation of \(H_i\).
The coefficient at one of a group algebra projection is its
normalized trace on that representation. This proves the
single-prime rank formula. Since \(e_0,e_1\) are commuting
orthogonal projections (choose an invariant inner product),
the kernel of their product is the sum of their kernels.
Proposition~\ref{prop:levelmass678} gives the joint bound.

For an irreducible character \(\chi\) the coefficient of the
central idempotent at \(h\) is
\(\chi(1)\overline{\chi(h)}/|H|\), with \(H\) here the whole
product and \(\chi\) a tensor product character. Its numerator
is an algebraic integer. The indicated set of characters is
stable under every automorphism of the character-value
cyclotomic field. Indeed Green's formula takes a Galois power
on the roots of unity to the same power on the torus character
labels. This permutes Frobenius orbits and preserves their
length, primitive order, and the property of being a
degree-one root of order dividing \(Q\). Thus summing the
numerators gives rational algebraic integers, as asserted.
In particular both \(|H|e\) and \(|H|(1-e)\) preserve every
\(H\)-stable lattice.

The map \(\lambda\mapsto e\lambda\) identifies \(I_\ell\) with
the quotient of \(\Lambda_\ell\) by the sum of its intersections
with the two eigenspaces. Indeed it is onto, and \(e\lambda\)
lies in the lattice precisely when both components of \(\lambda\)
do. Using \(1-e\) instead gives the same quotient. It follows
that \(I_\ell\) is killed by \(|H|\) and can be generated by at
most \(\dim(1-e)V\) elements, since the image of the lattice
under \(1-e\) is free of that rank over the coefficient
valuation ring. The elementary divisor theorem bounds its
normalized length by \(\dim(1-e)V\,v_\ell(|H|)\). Sum and use
integer factorization. Finally,
\(\log |H|\le fr^2\log(p_0p_1)\) by counting all matrices, and
the estimates for \(f,r,p_i\) above give the asserted order.
\end{proof}

The rank computation uses the other prime's level each time
it recovers a regular representation. It therefore applies to
the definite space with both markings on the same field as
Corollary~\ref{cor:oddwindow677}. Estimating the norm of a
given projected Poincar\'e series further requires its mass in
these isotypic components. Projection to the global subspace
of theta lifts also introduces the congruences among
automorphic representations.


\subsection{The finite Bessel character at the two level primes}
\label{sec:finitebessel679}

We compute the character average that occurs at the split
places in the Bessel period of the shorter construction. Its
Fourier transform counts vectors for which the first powers of
a matrix are linearly independent. This count controls the
effect of the finite level projector on the Bessel character
itself.

Let \(V\) be an \(r\)-dimensional space over \(k=\mathbb F_q\).
Fix \(a\ge1\) and \(c=r-2a\ge1\), and choose ordered blocks of
sizes \(1,\ldots,1,c,1,\ldots,1\), with \(a\) lines on each
side. Let \(\mathfrak p\) be the block upper triangular algebra,
\(\mathfrak u\) its strictly block upper triangular ideal,
\(U=1+\mathfrak u\), and \(G=\mathrm{GL}(V)\).
Fix a nontrivial additive character \(\vartheta\) of \(k\).
Choose \(D\) with blocks immediately below the diagonal as
follows. On consecutive lines in each arm the map is the
identity in the specified coordinates. The map from the last
left line to the middle space is a vector \(w\), and the map
from the middle space to the first right line is a functional
\(u\). Suppose \(\gamma=u(w)\ne0\). Set its other blocks to
zero. Nonzero coefficients along either arm can be put into
this form by rescaling the line bases, which preserves
nonvanishing of the middle pairing. Define
\[
  \psi(1+X)=\vartheta(\operatorname{tr}(DX))\quad(X\in\mathfrak u).
\]
It is a group character since \(\operatorname{tr}(DXY)=0\)
for \(X,Y\in\mathfrak u\). Extend it by zero to \(G\), and let
\(\psi^\#(h)=|G|^{-1}\sum_{g\in G}\psi(g^{-1}hg)\).

\begin{theorem}[The Fourier transform of the finite Bessel character]
\leavevmode\label{thm:besselfourier679}
Set \(F(X)=\psi^\#(1+X)\) on matrices where \(1+X\) is
invertible and set it to zero elsewhere. Use the transform
\[
       \widehat F(A)=\sum_{X\in\operatorname{End}(V)}
           F(X)\overline{\vartheta(\operatorname{tr}(AX))}.
\]
Write \(k_A\) for the fraction of vectors \(v\in V\) for which
\(v,Av,\ldots,A^{2a}v\) are independent. Then
\[
     \widehat F(A)=k_A/B_{q,r,a},\qquad
     B_{q,r,a}=\prod_{j=0}^{2a}(1-q^{j-r}),\qquad
     \sum_{\substack{h\in G\\h\ne1}}|\psi^\#(h)|^2
                      \le B_{q,r,a}^{-1}-1 .
\]
\end{theorem}
\begin{proof}
Before averaging over conjugates, the function on matrices is
\(\vartheta(\operatorname{tr}(DX))1_{\mathfrak u}(X)\).
Each \(1+X\) in its support is invertible. The trace pairing
identifies the annihilator of \(\mathfrak u\) with \(\mathfrak p\):
testing one entry of a strict upper block tests the transposed
lower block. Orthogonality of additive characters therefore
gives transform \(|\mathfrak u|1_{D+\mathfrak p}\).
Fourier transformation commutes with conjugation. Hence
\(\widehat F(A)\) is \(|\mathfrak u|\) times the probability
that the matrix of \(A\) in a uniformly chosen basis belongs
to \(D+\mathfrak p\).

We count that probability. Record from a basis its first \(a\)
columns \(v_1,\ldots,v_a\) and its last \(a\) coordinate
functionals \(\lambda_1,\ldots,\lambda_a\), starting with the
last functional and proceeding backwards. Write \(L_j\) for
the span of the first \(j\) columns, \(R_j\) for the span of
the first \(j\) of these functionals, and \(W=\bigcap_i\ker
\lambda_i\). Then \(L_a\subset W\) and \(C=W/L_a\) has
dimension \(c\). The basis identifies the original middle
space with \(C\), and thereby gives \(w'\in C\) and \(u'\in C^*\)
with \(u'(w')=\gamma\). Every tuple of independent frames with
these annihilations and every such pair in \(C\) has the same
number of basis lifts. Indeed one may choose a basis of the
kernel of \(u'\), complement it by \(w'\), lift a resulting
middle basis to \(W\), and complete it by lifts of the basis
dual to the \(\lambda_i\). Equivalently \(G\) acts transitively
on these data and the map from bases is equivariant.

The condition on the blocks of \(A\) in this basis says exactly
\[
 \begin{split}
 Av_j-v_{j+1}&\in L_j,\qquad
       \lambda_j A-\lambda_{j+1}\in R_j \quad(1\le j<a),\\
 Av_a&\in W,\qquad [Av_a]=w'\ \text{in } C,\\
 (\lambda_a A)|_W&=u'\circ(W\longrightarrow C).
 \end{split}
\]
For example the first line fixes one subdiagonal entry and
kills all entries below it in the corresponding column or
to its left in the corresponding row. The two last conditions
fix the two blocks adjacent to \(C\) and kill the remaining
blocks below them. Other entries belong to \(\mathfrak p\).

Put \(v=v_1,\lambda=\lambda_1\). The first line makes \(v_j\)
a monic polynomial of degree \(j-1\) in \(A\) applied to \(v\).
It gives the analogous monic polynomial on the right of
\(\lambda\) for \(\lambda_j\). The equations
\(\lambda_i(v_j)=0\) successively give
\(\lambda A^h v=0\) through \(h=2a-2\), since every such
\(h\) equals \(i+j-2\) for \(1\le i,j\le a\).
The condition \(Av_a\in W\) gives the next vanishing. Finally
\(\lambda_a A^2v_a=u'(w')=\gamma\) gives
\[
       \lambda A^h v=0\quad(0\le h<2a),\qquad
       \lambda A^{2a}v=\gamma .
\]
Conversely these moment equations force independence of
\(v,Av,\ldots,A^{2a}v\), because a relation ending at an
earlier power would express \(A^{2a}v\) in terms of earlier
powers. The same argument works for
\(\lambda,\lambda A,\ldots,\lambda A^{2a}\).
Choose \(v_{j+1}\) freely in \(Av_j+L_j\) and choose
\(\lambda_{j+1}\) freely in \(\lambda_j A+R_j\). The moment
equations give all required annihilations, put \(Av_a\) in
\(W\), and make \(\lambda_a A\) vanish on \(L_a\).
Its induced functional on \(C\) pairs with \([Av_a]\) to
\(\gamma\), so \(u',w'\) are determined. Each solution
\((v,\lambda)\) gives exactly \(q^{a(a-1)}\) tuples, by
multiplying the two sets of choices of \(q^j\) at each step.

For each of the \(q^r k_A\) independent starting vectors
there are exactly \(q^{r-2a-1}\) functionals satisfying the
\(2a+1\) moment equations. For the other vectors there are
none. Thus \(\widehat F(A)\) is a constant times \(k_A\),
with the constant independent of \(A\). Its average over all
matrices is \(F(0)=1\), by character orthogonality.
For uniform independent \(A,v\), the probability that \(v\ne0\)
is \(1-q^{-r}\). Given \(j\) independent successive vectors
and the previous values of \(A\), the next value of \(A\) on
the last of them is uniform on \(V\): a linear map prescribed
on the first \(j-1\) independent vectors extends in equally
many ways for each image of the \(j\)-th. Its chance of missing
their span is \(1-q^{j-r}\). This computes the average of
\(k_A\) as \(B_{q,r,a}\), proving the transform formula.
Since \(0\le k_A\le1\), the average squared transform is at
most \(B_{q,r,a}^{-1}\). Finite Fourier Plancherel gives
the sum of \(|F(X)|^2\) as that average. Removing \(X=0\)
subtracts exactly one and proves the estimate.
\end{proof}

\begin{corollary}[Mass of the projected finite Bessel character]
\leavevmode\label{cor:markedbessel679}
On the fields and levels of Proposition~\ref{prop:levelmass678},
put \(a=m/2-1\), \(c=r-2a\), and
\(K_x=8f(p_0^{-c}+p_1^{-c})\). Let \(\Psi\) be the product of
the reduced Bessel characters extended by zero on \(H\).
For all sufficiently large \(x\),
\[
 \sum_{s\in H,\ s\ne1}|\Psi^\#(s)|^2\le K_x.
\]
In the representation induced to \(H\) from the character of
its Bessel unipotent subgroup, the fraction of the dimension
discarded by \(e\) is
\[
        \delta+O\big(\sqrt{\delta(1-\delta)K_x}\big),
        \qquad \delta=1-(1-\delta_0)(1-\delta_1).
\]
The implied constant is one. The same
error bound with \(\delta_i\) holds for the rank of
\(1-e_i\). In particular all these fractions are \(O(b^{-9})\).
\end{corollary}
\begin{proof}
We check the middle pairing, since its nonvanishing is used
in the moment calculation. At a marked prime the split algebra
is \(\QQ_p\times\QQ_p\) with involution interchanging factors.
Write the hyperbolic lines in the order
\(e_1,\ldots,e_a\), the middle space, \(e_a^*,\ldots,e_1^*\).
If \(g\) is the first matrix, the second equals
\(J^{-1}(g^{-1})^tJ\), where \(J\) pairs the opposite lines
and has a symmetric, unimodular middle block \(T\).
In the lattice basis the distinguished middle vector \(z\)
satisfies \(z^tTz=-N\), where \(N>p_0,p_1\) is the separate
prime level. Indeed the middle basis replaces one unit vector
\(x_j\) of the form \(\operatorname{diag}(1,\ldots,1,-N)\)
by \(x_j/\overline\pi+z/N\), with \(\pi\overline\pi=N\).
Its pair with \(z\) is \(-1\), its norm is zero, and the
two-by-two middle block has determinant \(-1\).
The reduction in the first matrix is the entire strictly
block upper unipotent group, since entries can be lifted freely
over \(\ZZ_p\) and the second matrix is then determined.

The additive character reads the consecutive entries along
\(e_1,\ldots,e_a,z\) in both components, with a unit coefficient
in each. In fact the two completions are \(\QQ_p\), the
different is a unit there, and division by \(p_0p_1\) has
valuation \(-1\) in each completion. In the first matrix this
gives nonzero coefficients along the left arm and the last
row evaluated at \(z\). In the second matrix the inverse
negates each first off-block entry, since every higher power
of \(g-1\) has zero in that block. Transposition and \(J\)
therefore give the reversed arm in the first matrix and, at its
central end, the functional \(-z^tT\). The two central blocks
have pairing a nonzero scalar times \(z^tTz=-N\bmod p\).
This verifies precisely the hypothesis of
Theorem~\ref{thm:besselfourier679} for the actual reduced
character. Reductions in all factors are independent: take
the specified integral unipotent lift at each split place
and the identity away from the marked primes. Thus the
extended adelic character on its residue image is the stated
product.

Conjugation averages and sums of squared absolute values
factor over the product. The full sum is at most
\(\prod_{i=0}^1 B_{p_i,r,a}^{-f}\), and the identity contributes
one. Since \(p_i\ge3\), summing
\(-\log(1-y)\le2y\) over \(y=p_i^{j-r}\), \(0\le j\le2a\),
bounds the log of the product by
\(4f(p_0^{-c}+p_1^{-c})\). This quantity tends to zero by
the scale bounds of Proposition~\ref{prop:levelmass678},
since \(c=r-m+2\ge x^5-2\sqrt x\).
Using \(e^v-1\le2v\) for small \(v\ge0\) gives \(K_x\).

The induced character divided by its degree is exactly
\(\Psi^\#\), by the character formula for induction from a
subgroup. Let \(d=1-e=\sum_s d_s s\). Its identity
coefficient is \(\delta\) by the product formula for the
regular trace. Since \(d=d^*=d^2\),
\(\sum_{s\ne1}|d_s|^2=\delta-\delta^2\), by taking the
coefficient at one in \(d^*d\). The normalized trace of \(d\)
on the induced space is \(\sum_s d_s\Psi^\#(s)\). Its identity
term is \(\delta\) and Cauchy--Schwarz bounds the rest by
\(\sqrt{\delta(1-\delta)K_x}\). Taking \(d=1-e_i\) gives the
single-prime assertion. Finally \(\delta\le4b^{-9}\) by the
proof of Proposition~\ref{prop:levelmass678}, and \(K_x\)
is smaller than any fixed negative power of \(b\) on this
family.
\end{proof}

These bounds and the odd-cocycle estimate can be used on the
same entire interval in this construction. More precisely,
fix \(q_0,d\), take \(V_0\le x\le X=V_0^{3/2}\), and put
\[
 s_*=\lceil 22\log_2(2\sqrt X)\rceil,\quad
 N_*=3q_0d\,l_0(s_*)2^{s_*},\qquad t_*=\lceil 15X^5\rceil ,
\]
where \(l_0(s)\) is the least power of three at least \(27s\).
For every possible tower exponent \(1\le s\le s_*\), put
\(E_s=\QQ(\sqrt{-d},\mu_{3q_0l_0(s)2^s})\). Then one integer
\[
             J_I=\mathcal J_{N_*,t_*}\prod_{1\le s\le s_*}h_{E_s}
\]
controls the odd-cocycle lengths of Theorem~\ref{thm:oddinteger677},
including weight zero, on every field used in the finite
Bessel comparison above. Indeed \(2^s=m^{22}\le(2\sqrt X)^{22}\)
at every node, and \(l_0(s)\mid l_0(s_*)\) because both are
powers of three. Hence all \(E_s\subseteq\QQ(\mu_{N_*})\) and
all needed weights lie below \(t_*\). Each
\(\log h_{E_s}\ll N_*\log(2N_*)\) by the class-number estimate
in Corollary~\ref{cor:oddwindow677}. Since \(s_*=O(\log(2X))\)
and \(N_*\ll_{q_0,d}X^{11}\log(2X)\), the theorem gives
\(\log J_I\ll_{q_0,d}X^{54}\log^5(2X)\). The finite projector
and its character estimate at a node use that node's own two
places and field \(E_s\). All the preceding bounds for them
hold throughout the interval when \(V_0\) is sufficiently large.

This evaluates the finite character in the Bessel term of the
period. Extending it to the whole Poincar\'e Gram matrix
requires bounds for rational elements outside the Bessel
subgroup in the unfolded integral.


\subsection{Rational orbital cancellation and the two frame filtrations}
\label{sec:rationalflags680}

Retain the shorter construction and its levels from
Corollary~\ref{cor:markedbessel679}. Work at a single node \(x\) of
the interval there. Write \(F=E^+\), \(a=m/2-1\), and \(G=\mathrm U(h)\)
over \(F\). Use a Hermitian form conjugate linear in its first
variable. The space over \(E\) is
\[
 V=L\oplus (X_0\perp Ez)\oplus L^\vee,\qquad
 L=\langle e_1,\ldots,e_a\rangle,\qquad h(z,z)=-N .
\]
The first and last summands are opposite totally isotropic spaces,
\(X_0\) is totally positive definite, and \(N>0\). Put \(n=a+1\),
\(f_i=e_i\) for \(i\le a\), \(f_n=z\), and
\(T_i=\langle f_1,\ldots,f_i\rangle\), with \(T_0=0\) and \(T=T_n\).
Let \(U\) be the unipotent radical of the stabilizer of the flag
\(T_1,\ldots,T_a\), and set \(S=U\rtimes\mathrm U(X_0)\), with
\(\mathrm U(X_0)\) acting trivially on the other summands. Write
\(\sigma(u)\) for the sum of the entries in positions \((i,i+1)\)
in the action on the frame, \(1\le i\le a\). In the final entry
one applies the last row on \(L\) to the vector \(z\). This is
additive on \(U\), since a product of two strictly upper block
matrices has zero in every first off-diagonal block. Put
\(\alpha=-1/(M\delta_E)\), with \(\delta_E\) the chosen different
generator and \(M=p_0p_1\). The finite character is the standard
finite adelic additive character of \(E\) applied to
\(\alpha\sigma(u)\).

\begin{lemma}[A simultaneous endomorphism for two rational characters]
\leavevmode\label{lem:rationalshift680}
For \(\gamma\in G(F)\), set \(H=U\cap\gamma U\gamma^{-1}\).
Suppose the finite character takes equal values at \(u\) and
\(\gamma^{-1}u\gamma\) for every \(u\in H(F)\). There is an
\(E\)-linear endomorphism \(A\) of \(V\) such that, for \(1\le i<n\),
\[
 A f_i-f_{i+1}\in T_i,\qquad
 A\gamma f_i-\gamma f_{i+1}\in\gamma T_i .
\]
\end{lemma}
\begin{proof}
Let \(\mathfrak n=\operatorname{Lie}U\) over \(F\). Matrix logarithm
and exponential are inverse finite polynomials on strict block
unipotents and nilpotents in characteristic zero. They commute
with conjugation. For the Hermitian adjoint \(u^*=u^{-1}\) amounts
to \((\log u)^*=-\log u\). Thus the logarithms of \(H(F)\) are
exactly the intersection of \(\mathfrak n\) and
\(\gamma\mathfrak n\gamma^{-1}\) on \(F\)-points. Also
\(\sigma(\exp X)=\sigma(X)\), using the same entries on the Lie
algebra. On rational arguments the finite character is the inverse
of its archimedean factor, the exponential of the absolute trace
times \(2\pi i\), up to the common choice of sign. Consequently
\[
 \lambda(X)=\operatorname{Tr}_{E/F}(\alpha\sigma(X))
\]
and its conjugate by \(\gamma\), given by
\(X\mapsto\lambda(\gamma^{-1}X\gamma)\), agree on the intersection.
Indeed if their difference were \(d\ne0\) at \(X\), evaluating at
\(X/(2[F:\mathbb Q]d)\) would give absolute trace \(1/2\).

The two functionals now define one \(F\)-linear functional on the
sum of their domains. Extend it to the unitary Lie algebra by
extending a basis, and extend scalars to \(E\). This last Lie
algebra becomes \(\operatorname{End}_E(V)\). For instance each
endomorphism is the sum of a skew-adjoint endomorphism and an
imaginary scalar times a skew-adjoint endomorphism, by taking its
skew-adjoint and self-adjoint parts. The scalar extension of
\(\mathfrak n\) is the entire strict upper block space: the
adjoint preserves that space in the ordering with reversed dual
lines, so the same decomposition applies. Represent the resulting
functional by \(\operatorname{tr}(B\,\cdot)\) using the
nondegenerate ordinary matrix trace pairing.

On a skew-adjoint matrix, \(\lambda\) is the sum of
\(\alpha\sigma(X)\) and its conjugate. After scalar extension the
conjugate term uses only the opposite arm: conjugation on the
entries of a skew-adjoint matrix is obtained by negative
transposition and the Hermitian Gram matrix, which interchanges
the two outer flags and preserves the middle block. Thus it has
zero coefficient on each entry from a later block into a left
line. Testing an arbitrary entry of this kind in row \(i\)
specifies precisely the part of \(Bf_i\) below that line.
The only prescribed nonzero part is \(\alpha f_{i+1}\), including
\(\alpha z\) for the last left line. This gives
\(Bf_i-\alpha f_{i+1}\in T_i\). Applying the identical argument
in the \(\gamma\)-translated basis gives the second relation for
\(B\). Take \(A=\alpha^{-1}B\).
\end{proof}

\begin{theorem}[Ranks and minors of agreeing rational frames]
\leavevmode\label{thm:rationalflags680}
Write \(C=(h(f_i,\gamma f_l))_{i,l\le n}\), let \(D_i\) be its
leading minor of order \(i\), set \(D_0=1\), and put
\(\Delta=D_n\overline{D_n}-N^2D_a\overline{D_a}\).
If \(T\cap\gamma T=0\), then \(\Delta\) is totally positive and
\(D_n\ne0\). Otherwise \(\Delta=0\). In the latter case, under
the character agreement of Lemma~\ref{lem:rationalshift680},
there is a unique \(0\le j\le a\) with
\[
 \dim_E(T\cap\gamma T)=n-j,\qquad
 \dim_E(L\cap\gamma L)=a-j,\qquad D_j\ne0.
\]
There is one \(\zeta\in E^\times\), \(\zeta\bar\zeta=1\), such
that for each \(j<i\le n\) a common vector in
\(T_i\cap\gamma T_i\) has top coefficient \(\zeta\) in the first
frame and \(1\) in the second. None of the lower levels \(i\le j\)
has a nonzero common vector.
\end{theorem}
\begin{proof}
The joint Gram matrix of the \(2n\) columns has off-diagonal
blocks \(C,C^*\). Both diagonal blocks are zero apart from their
last diagonal entry \(-N\). Its determinant is
\((-1)^n\Delta\). Indeed the terms using neither diagonal entry
give \((-1)^nD_n\overline{D_n}\), those using both give
\((-1)^{n-1}N^2 D_a\overline{D_a}\), and using exactly one leaves
unequal numbers of rows and columns to match across the two
blocks. A common vector gives linearly dependent columns. When
the columns are independent, the span is nondegenerate at each
real place. To see this, \(T^\perp=L\oplus X_0\) is positive
semidefinite with nullspace \(L\). A vector in the radical of
the joint span is orthogonal to both frames and has norm zero,
so lies in \(L\cap\gamma L=0\); scalar extension preserves this
intersection. The span contains \(z\) and a totally isotropic
\(a\)-space orthogonal to \(z\). Its negative index is at least
\(a+1=n\), and the ambient negative index is exactly \(n\).
The determinant therefore has sign \((-1)^n\) at every real
place. This proves the assertions about \(\Delta\) and \(D_n\).

For the remaining assertions write \(W=T\cap\gamma T\).
For \(0\ne w\in W\) take the larger \(i\) of its two top
indices. Write \(u,v\) for its two coefficients at \(i\),
temporarily allowing one to be zero. Applying \(A^{n-i}\)
takes the two expressions into the top frame level, with
coefficients \(u,v\) at \(z,\gamma z\), by
Lemma~\ref{lem:rationalshift680}. The norm of this one vector
is both \(-N u\bar u\) and \(-N v\bar v\).
It follows that both coefficients are nonzero. In particular
the two filtrations induced on \(W\) coincide. In their common
level \(W_i\), the map to the top coefficient in either frame
has kernel \(W_{i-1}\). Each nonzero successive quotient is
therefore one-dimensional; its pair of coefficients spans a
line generated by \((\zeta,1)\) with \(\zeta\bar\zeta=1\).
The map \(A\) takes such a pair unchanged to the next level
for \(i<n\). All nonzero quotients thus occur consecutively at
the top and give the same \(\zeta\). Counting the quotients in
\(W_n\) and \(W_a\) proves the two dimension formulas.

We have \(T=T_j\oplus W\) and \(\gamma T=\gamma T_j\oplus W\).
If \(y\in\gamma T_j\) is orthogonal to \(T_j\), it is also
orthogonal to \(W\subset\gamma T\), since
\(\gamma T_j\subset\gamma L\). Hence \(y\in T^\perp\). Its norm
is zero, so \(y\in L\) by positivity. Now
\(y\in W\cap\gamma T_j=0\). The pairing of the two \(j\)-spaces
is nonsingular, which proves \(D_j\ne0\).
\end{proof}

\begin{proposition}[Cancellation over the stabilizer of a rational term]
\leavevmode\label{prop:stabilizercancel680}
For class representatives \(h_c,h_{c'}\in\mathrm U(X_0)(\mathbb A_f)\)
write \(\Gamma_c=U(F)\cap h_c K(M) h_c^{-1}\), where \(K(M)\)
is the principal level in the integral stabilizer. Put
\(\Gamma=\Gamma_c\cap\gamma\Gamma_{c'}\gamma^{-1}\).
If the two rational characters disagree, the integral of the
Bessel orbital kernel over every fiber
\(\Gamma\backslash H(F\otimes_{\mathbb Q}\mathbb R)\) is zero.
In particular every absolutely convergent full or cut orbital
integral for this term vanishes.
\end{proposition}
\begin{proof}
Write \(H_\infty=H(F\otimes_{\mathbb Q}\mathbb R)\) and use a
similar notation for \(G_\infty\). In the unfolding, the two
principal congruence groups act on \(\gamma\) by left and right
multiplication; its stabilizer consists of the pairs
\((u,\gamma^{-1}u\gamma)\) with \(u\in\Gamma\).
The finite character is trivial on \(\Gamma_c\).
Indeed the frame is fixed pointwise by \(h_c\), and for \(u\in U\)
its first off-diagonal frame entries are unchanged by conjugation
by \(h_c\). These entries are integral, and at the levels are
divisible by \(M\); division by \(\delta_E\) therefore puts their
trace in the conductor of the finite character. The same argument
applies to \(c'\).

The quotient \(\Gamma\backslash H_\infty\) is compact. Here are
details for this particular type of arithmetic subgroup. Restrict
scalars to \(\mathbb Q\) and choose a rational flag of Lie ideals
with successive central, one-dimensional quotients, successively
taking rational central lines in nilpotent quotients. Exponential
coordinates in a basis adapted to this flag give real coordinates
on \(H_\infty\); in particular each successive central direction
has group law addition on its quotient. For each rational basis
element \(X\), \(\exp(dX)\in\Gamma\) for a sufficiently divisible
positive integer \(d\). This follows from the finite exponential
polynomial: outside finitely many primes its coefficients and
both class changes are integral, and at the exceptional primes
it tends to the identity as the local divisibility of \(d\)
increases. Reducing one coordinate at a time, starting at the
outermost quotient, gives representatives in a product of compact
intervals in exponential product coordinates. The subgroup is
discrete, since its rational matrix entries lie in fixed fractional
integer lattices at all finite places. This proves compactness and
gives a finite invariant quotient measure.

For either Bessel function \(W\), its defining equivariance on
the left by the real unipotent is multiplication by the
archimedean character \(\chi_\infty\). For the full vector this
is the character in the Bessel integral (translate its unipotent
integration variable); the cut vector is multiplied by a function
constant on left \(U_\infty\)-cosets. On a rational unipotent
the character is the inverse of its finite factor, by triviality
of the adelic additive character on \(E\). It follows that
\(K_\gamma(y)=W(y)\otimes W(\gamma^{-1}y)^*\) transforms under
\(h\in H_\infty\) by the character
\(\eta(h)=\chi_\infty(h)/\chi_\infty(\gamma^{-1}h\gamma)\).
It is trivial on \(\Gamma\) and nontrivial by rational
disagreement. For fixed \(y\), right invariance of measure on
\(\Gamma\backslash H_\infty\) shows that the integral of
\(K_\gamma(hy)\) there is unchanged by multiplication by any
\(\eta(h_0)\). Choosing \(\eta(h_0)\ne1\) proves the vanishing.
Finally integrate the fiber integral on \(H_\infty\backslash
G_\infty\) using quotient Haar measure. Both groups are
unimodular. This gives zero with compact cutoffs on the outer
quotient, and for every absolutely convergent orbital integral.
\end{proof}

\begin{corollary}[Elimination of the isotropic parabolic on a level slice]
\leavevmode\label{cor:parabolicslice680}
Suppose \(h_c^{-1}\gamma h_{c'}\) is integral at all finite
places and that, in the adapted basis at one of the split marks,
its first diagonal entry reduces to \(1\). Every rational term
preserving \(L\) either belongs to \(S(F)\) or has the
stabilizer cancellation of Proposition~\ref{prop:stabilizercancel680}.
In particular this holds on either slice obtained by requiring
the residue to be in the Bessel image above \(p_i\).
\end{corollary}
\begin{proof}
For an agreeing term with \(\gamma L=L\) the preceding theorem
has \(j=0\), since \(a\ge1\). Starting at \(i=1\), its common
vectors with coefficient pair \((\zeta,1)\) prove inductively
that \(\gamma T_i=T_i\) and \(\gamma f_i-\zeta f_i\in T_{i-1}\),
also at \(i=n\). The unitary element \(\zeta^{-1}\gamma\) lies
in the flag parabolic. In its Levi quotient the diagonal on the
left lines is one, and the middle unitary transformation fixes
\(z\); hence it preserves \(z^\perp=X_0\) inside the middle
space. Removing that middle action leaves the unipotent radical.
Thus \(\gamma\in\zeta S(F)\).

Both class representatives fix \(e_1\). Its line is a primitive
lattice summand, so integrality at every finite place implies
that the eigenvalue \(\zeta\) is integral there. Since
\(\zeta\bar\zeta=1\), all complex conjugates of \(\zeta\) have
absolute value one. Such an algebraic integer is a root of
unity: the monic polynomials of its positive powers have bounded
degree and bounded integer coefficients, so two of the powers
coincide. At the selected component \(E_w=\mathbb Q_{p_i}\),
its residue equals the indicated first diagonal entry. The group
\(1+p_i\mathbb Z_{p_i}\) is torsion-free for odd \(p_i\), by
the convergent injective \(p_i\)-adic logarithm there. Therefore
\(\zeta=1\) and \(\gamma\in S(F)\). All disagreeing terms have
zero fiber integral by the proposition. Every element of the
Bessel image has first diagonal entry one, as required.
\end{proof}

On these slices the agreeing terms outside \(S(F)\) therefore
have either \(1\le j\le a\) or \(\Delta\ne0\). Quantitative control
requires their multiplicities and the sizes of their integrals.


\subsection{Divisors of the frame discriminant at the marked primes}
\label{sec:discriminant681}

Continue with the two frames of Theorem~\ref{thm:rationalflags680}.
We work with pairs of classes in the actual level. Write \(f=[F:\mathbb Q]\)
in exponents, and write \(\mathcal F=(f_1,\ldots,f_n)\) for the
frame matrix. A rational term \(\gamma\) occurs between \(c,c'\)
only if \(g=h_c^{-1}\gamma h_{c'}\) carries the standard finite
lattice to itself. In these coordinates the two integral frames
are \(\mathcal F,g\mathcal F\). Indeed both class representatives
fix the original frame, and the change to the lattice basis
preserves its columns. Their Hermitian Gram matrices are
therefore the rational Gram matrices of the preceding subsection.
The lattice form is integral, including at \(N\): the altered
middle basis there has an integral two-by-two block with diagonal
\(0,-N\) and off-diagonal entries of norm one. In particular
\(C\) has entries in \(\mathcal O_E\) and \(\Delta\in\mathcal O_F\).

\begin{theorem}[Divisibility by the indices of the two frames]
\leavevmode\label{thm:discriminant681}
Suppose \(\Delta\ne0\). At a split place \(v\mid p_0p_1\) of
\(F\), let \(w,\bar w\) be the places of \(E\) over \(v\). Form the
matrix \(B_w=(\mathcal F,g\mathcal F)_w\) in the integral lattice
at \(w\). Let \(\ell_w\) be the length of the saturation of its
column lattice modulo that lattice, and let \(d_w\) be the
dimension of the intersection of the two reduced frame spans.
With normalized integral valuations one has
\[
 \operatorname{ord}_v\Delta\ \ge\ \ell_w+\ell_{\bar w}
      \ \ge\ d_w+d_{\bar w}.
\]
In particular on the slice where the reduction belongs to the
Bessel image at all places above \(p_i\),
\[
 p_i^{\,2nf}\mid {\rm N}_{F/\mathbb Q}(\Delta),\qquad
 \big|{\rm N}_{F/\mathbb Q}(\Delta)\big|^{-s}
              \le p_i^{-2nfs}\quad(s\ge0).
\]
If both slices apply, the two divisors multiply. On the same
slice, for any rational cell, \(p_i^j\mathcal O_E\) contains
\(D_j\) for \(0\le j\le a\).
\end{theorem}
\begin{proof}
The two frames each have independent reduction: the first consists
of primitive basis columns and the second is carried there by a
lattice automorphism. Their joint reduction thus has rank
\(2n-d_w\). The columns are independent in characteristic zero
by Theorem~\ref{thm:rationalflags680}. Smith normal form over
the local discrete valuation ring shows that the least valuation
of any maximal minor of \(B_w\) is \(\ell_w\), the sum of its
elementary divisor exponents. Exactly \(d_w\) of those exponents
are positive. Write \(J_w\) for the integral matrix pairing the
two components of the ambient lattice. The joint Gram over \(w\)
is \(B_{\bar w}^t J_w B_w\). Expanding its determinant twice by
Cauchy--Binet shows that each summand contains a maximal minor
of each \(B\) and an integral minor of \(J_w\). Their determinant
is \((-1)^n\Delta\), which proves the inequalities. This argument
retains nontrivial elementary divisors at arbitrary powers of
the mark.

If the reduction is in the Bessel image then it preserves the
whole frame span in both components, so that \(d_w=d_{\bar w}=n\).
There are \(f\) places of \(F\) over the completely split prime
\(p_i\), each of residue degree one. Taking the field norm gives
the divisibility assertion; total positivity makes the inverse
power estimate immediate. For the last assertion each entry of
the leading \(j\)-square Gram block vanishes modulo every \(w\)
over \(p_i\): a vector in the first isotropic arm pairs to zero
with the whole frame. Expand the determinant at each such place.
\end{proof}

\begin{proposition}[The positive vector measured by the discriminant]
\leavevmode\label{prop:normalvector681}
Put \(u=\operatorname{adj}(C)e_n^{\rm std}\), where
\(e_n^{\rm std}\) is the last coordinate vector in \(E^n\).
Let \(t=\gamma\mathcal F u\), and let \(q\) be the \(X_0\)-component
of \(t\) in the rational orthogonal middle decomposition. Then
\[
                 N h(q,q)=\Delta .
\]
At every marked component where the reduction preserves the
frame span, the vector \(q\) belongs to the \(p_i^{a+1}\)-multiple
of the local positive lattice of class \(c\), provided the
reduction preserves the span at its paired component as well.
\end{proposition}
\begin{proof}
The identity \(Cu=D_n e_n^{\rm std}\) gives \(t\in L^\perp\)
and \(h(z,t)=D_n\). Thus \(t=l+q-(D_n/N)z\), where \(l\in L\).
On the other hand the last coordinate of \(u\) is \(D_a\),
so \(h(t,t)=-N D_a\overline{D_a}\) by unitarity. Equating the
two evaluations gives the formula at every embedding.
In the marked lattice basis \(q\) is obtained by projecting
\(g\mathcal F u\) to its positive part, then carrying it back
by the class representative. The projection is integral there:
the only middle change of basis has denominators dividing
\(N\), a unit at the marks. The components of \(u\) are signed
\(a\)-minors of the first \(a\) rows of \(C\), whose entries
vanish modulo \(p_i\) when the frame span is preserved in the
specified components. Hence \(p_i^a\) divides \(u\).
Each positive projection of a column of \(g\mathcal F\) vanishes
modulo \(p_i\) also. Multiplication yields the claimed extra
factor of \(p_i\).
\end{proof}

We next allow the two reduced frame spans to share only one
vector. In this situation there is a further divisor that can
be retained after removing the elementary divisors of both
isotropic arms. For a local lattice \(D\), the saturation of a
submodule is its fraction-field span intersected with \(D\).

\begin{theorem}[The two terminal columns at a meeting place]
\leavevmode\label{thm:terminalmeet681}
Assume \(\Delta\ne0\). At each marked component \(w\) form
\(B_{0,w}\) from the two isotropic arms, omitting the two last
columns of the frames. Write \(S_{0,w}\) for the saturation of
its columns and let \(\ell_{0,w}\) be their index length in
\(S_{0,w}\). In the free quotient by \(S_{0,w}\), let \(b,b'\)
be the images of the two terminal columns. The ideal generated
by all determinants
\[
 \det\begin{pmatrix}t_1(b)&t_1(b')\\t_2(b)&t_2(b')\end{pmatrix}
\]
as \(t_1,t_2\) range over integral linear forms on the quotient
has valuation \(\ell_w-\ell_{0,w}\).
Suppose the two finite characters agree on the intersection
of their conjugate reduced unipotent groups. Then
\[
             \ell_w-\ell_{0,w}\ge 1\qquad(d_w>0).
\]
Consequently the product of the norms of all meeting components
divides the absolute finite index given by any such determinant
chosen locally at every finite place, if the forms and columns
are integral everywhere and the determinant is an idele.
\end{theorem}
\begin{proof}
Choose bases adapted to \(S_{0,w}\) and its direct complement.
The exterior product of the \(2a\) isotropic columns is a
uniformizer to the power \(\ell_{0,w}\) times a primitive generator
of the top exterior power of \(S_{0,w}\), up to a unit. Wedging
the two terminal columns into it uses only their images in the
free quotient. Thus the minimum valuation of the maximal minors
of \(B_w\) is \(\ell_{0,w}\) plus the minimum valuation of a
two-row minor of \((b,b')\). Arbitrary pairs of forms generate
the same ideal as these minors, by determinant expansion. This
proves the equality of valuations.

For the remaining assertion work over \(\mathbb F_{p_i}\) at
\(w\), writing \(s\) for the local matrix. In this component
the unipotent is \(1+\mathfrak n\), with \(\mathfrak n\) the full
strict block upper algebra. Its intersection with its conjugate
is \(1+(\mathfrak n\cap s\mathfrak n s^{-1})\). As in the proof of
Corollary~\ref{cor:markedbessel679}, the two-component character
in this single matrix realization is the additive exponential
of a linear form on \(X=u-1\); its consecutive coefficients
on the left frame are the same nonzero scalar. This description
holds with \(w,\bar w\) interchanged, using the other component
of the additive character on the left. A character equality
on the global reduced intersection gives one on each factor
by taking identity components at all other places.
Over the prime field the additive exponential is injective as
a map to the \(p_i\)-th roots of unity. The two linear forms
therefore agree on the intersection. Extend them to one
functional on the full matrix algebra and represent it by
ordinary matrix trace. Dividing its representing matrix by
the left coefficient gives a matrix \(A\) with
\[
 A f_i\equiv f_{i+1}\pmod {T_i},\qquad
 A s f_i\equiv s f_{i+1}\pmod {s T_i}\quad(i<n)
\]
over the residue field. The prescribed lower blocks give
these formulas just as in Lemma~\ref{lem:rationalshift680}.
This finite-field proof uses \(u-1\) and holds independently
of denominators in a matrix logarithm.

If the spans meet, write a nonzero common vector in the two
frames. Let \(i\) be the largest index at which at least one
coefficient is nonzero. Apply \(A^{n-i}\). The two coefficients
at the terminal positions are now the former coefficients at
\(i\), so they are not both zero. The resulting equality, modulo
the span of both isotropic arms, is a nontrivial relation between
the terminal columns. It remains a relation modulo the reduction
of the larger saturated lattice \(S_{0,w}\). Hence every
two-row minor there vanishes. The index of an integral finite
idele is the product of residue norms to its local valuations.
Taking one factor from every meeting component proves the last
assertion.
\end{proof}

The two estimates apply at each \(x\) on the interval in
Corollary~\ref{cor:markedbessel679}, with its own \(E,M\) and
class lattices. They constrain an individual term and its
terminal determinant. The second moment involves products of
two amplitudes summed over all pairs in a residue class.

\section{Open and closed orbital integrals}\label{nav:marked-open-closed}

\subsection{The exact real measure over a pair of frames}
\label{sec:realgram682}

At a real place normalize \(z\) to have norm \(-1\).
Write \(n=a+1\), \(k=\dim X_0\ge1\) and \(r=2a+k+1\).
In coordinates \((u,x,t,v)\) of sizes \(a,k,1,a\) the Hermitian
form has matrix \(J\) with quadratic form
\(u^*v+v^*u+x^*x-|t|^2\). Use the matching order of the dual
lines here. Let \(\mathcal F=(e_1,\ldots,e_a,z)\) and \(G=\mathrm U(J)\).
For \(g\in G\) put \(Z(g)=\mathcal F^*Jg\mathcal F\). For any
\(n\)-square complex matrix \(Z\) set
\[
 \delta(Z)=|\det Z|^2-|\det Z_{[1,a],[1,a]}|^2,\qquad
 \Omega=\{Z:\delta(Z)>0\}.
\]
Write \(S_j=2\pi^j/(j-1)!\) for the area of the unit sphere in
\(\mathbb C^j\) and \(V_j=\prod_{l=1}^j S_l\).
Lebesgue measure on complex entries is \(d\Re\,d\Im\). On
Hermitian matrices use the diagonal real and the strict upper
complex entries, and on skew-Hermitian matrices use the imaginary
parts of the diagonal and the same strict upper entries.

We specify Haar measures. On \(G\) use the invariant measure
that at the identity is \(1/V_r\) times Lebesgue measure in the
skew-Hermitian tangent coordinates \(J\,dg\). This is the
compact-dual normalization: for \(J=1\) it gives probability
measure on \(\mathrm U(r)\), as verified below. Write \(P\)
for the subgroup fixing every column of \(\mathcal F\). Its
elements are the transformations
\[
 (u,x,t,v)\longmapsto
 \big(u-B^*h x+(K-B^*B/2)v,\ h x+Bv,\ t,\ v\big),
\]
where \(h\in\mathrm U(k)\), \(B\in\mathbb C^{k\times a}\) and
\(K^*=-K\). The formula follows by pairing the proposed images
of the positive and dual columns with the fixed frame and with
one another. Direct substitution also verifies the isometry.
Give \(P\) the measure \(ds=dB\,dK\,dh\) with \(dh\) probability
Haar. This is left and right invariant: multiplication translates
the central \(K\) variable by an expression in \(B\), and the
linear actions on \(B\) are unitary or the identity.

\begin{theorem}[Disintegration on the open set of frame products]
\leavevmode\label{thm:realgram682}
There are measurable representatives \(g_Z\in G\) with
\(Z(g_Z)=Z\) for \(Z\in\Omega\). For any nonnegative measurable
function \(\varphi\) on \(G\),
\[
 \begin{split}
 \int_{Z(g)\in\Omega}\varphi(g)\,dg
  =\frac{V_k S_k}{V_r}\int_\Omega \delta(Z)^{k-1}
    \int_P\int_P \varphi(s' g_Z s)\,ds\,ds'\,dZ .
 \end{split}
\]
Both sides may be infinite. The inner integral is independent
of the representative of the indicated frame products.
\end{theorem}
\begin{proof}
We first compute directly on the set of frames. For a complex
\(r\)-by-\(n\) matrix \(Y\) of independent columns, the map
\(Y\mapsto Y^*JY\) is a submersion, even at a singular Gram
value. In fact \(Y^*J\) is surjective and one can solve
\(Y^*J\,dY=H/2\) for any Hermitian variation \(H\). Let \(d\mu\)
be the fiber density at \(Y^*JY=-e_n e_n^*\), given in
submersion coordinates by dividing Lebesgue measure \(dY\) by
the specified measure on Hermitian matrices. Such coordinates
define the same density on overlaps by the change of variables
formula. It is invariant under \(G\), since left multiplication
preserves both measures and the Gram value.

Write \(Z=\mathcal F^*JY\); it consists of the \(a\) dual rows
of \(Y\) and the negative of its \(z\)-row. If \(Z\) is invertible
pass to \(Y'=YZ^{-1}\), with dual rows \((I_a,0)\) and \(z\)-row
\((0,-1)\). Let \((\eta,\xi)\) be the last row of \(Z^{-1}\),
let \((A,p)\) be the top rows of \(Y'\), and \((B,x)\) its
\(X_0\)-rows. Here \(\eta\) is a row and \(x\) a single column.
The Gram equations for \(Y'\) are precisely
\[
 \begin{split}
 A+A^*+B^*B&=-\eta^*\eta,\\
 p+B^*x&=-\eta^*\xi,\\
 x^*x&=1-|\xi|^2 .
 \end{split}
\]
Indeed the required Gram is the negative outer product of the
last row of \(Z^{-1}\). The right side of the last equation is
positive on \(\Omega\), since \(\xi=\det Z_{[1,a],[1,a]}/\det Z\).

There are \(a+k\) free complex rows before imposing the Gram
equations. Their change from \(Y'\) to \(Y\) has real Jacobian
\(|\det Z|^{2(a+k)}\) with \(Z\) held fixed. On Hermitian
matrices the change of prescribed Gram has Jacobian
\(|\det Z|^{2n}\). For a diagonal matrix with positive singular
values \(s_i\), the diagonal factors are \(s_i^2\) and the
complex off-diagonal factors contribute \((s_i s_j)^2\);
unitary factors preserve the Hermitian Lebesgue measure.
Dividing the two Jacobians yields \(|\det Z|^{2(k-1)}\).
Solving for \(A+A^*\) leaves its skew part \(K=(A-A^*)/2\)
with factor \(2^{-a}\) from the real diagonal variables. Each
off-diagonal pair has Jacobian one, as one checks on the sum
and half-difference of \(A_{ij}\) and \(\overline{A_{ji}}\).
Solving for \(p\) has Jacobian one. Polar coordinates for \(x\)
with squared radius \(1-|\xi|^2\) give
\(S_k(1-|\xi|^2)^{k-1}/2\) times uniform probability on the unit
sphere. These calculations give the frame density on \(\Omega\)
in the coordinates \(Z,B,K,\omega=x/\|x\|\):
\[
             d\mu=2^{-n}S_k\,\delta(Z)^{k-1}
                         \,dZ\,dB\,dK\,d\sigma(\omega).
\]

For a fixed \(Z\), take \(B=K=0\) and a fixed unit \(\omega_0\)
to construct \(Y_Z=Y'Z\) by the displayed equations. This is
a frame of the required Gram, with independent columns since
\(\mathcal F^*JY_Z=Z\). Such a frame is the image of \(\mathcal F\)
by an isometry: inside the orthogonal complement of its negative
line choose duals to the \(a\) isotropic columns, then subtract
one half of their mutual Gram matrix times the isotropic columns
to make the duals isotropic. Complete the hyperbolic pairs with
an orthonormal basis of their positive complement inside that
orthogonal complement. These choices can be made measurably
by using the first successful basis minors and Gram--Schmidt.
They yield \(g_Z\mathcal F=Y_Z\). Applying the formula for \(P\)
to \(Y_Z\) gives exactly the free \(B,K\) and
\(\omega=h\omega_0\) above. The image of \(dh\) is sphere
probability by rotation invariance. Consequently, in the frame
density the integral over these three coordinates is the
integral over the left copy of \(P\).

It remains to compute the factor for lifting frames to \(G\).
In the space of complex square matrices let \(d\nu_r\) be the
fiber density of \(g^*Jg=J\) by the same Gram construction.
It is invariant, and at the identity its density is \(2^{-r}\)
in the specified tangent coordinates: the derivative is the
Hermitian part \(J\,dg+(J\,dg)^*\), with unit Jacobian off the
diagonal and a factor two on each real diagonal. Multiplication
by \(J\) on the ambient entries preserves Lebesgue measure.
Thus \(dg=2^r d\nu_r/V_r\). For the positive identity form,
one can integrate its square Gram fiber by adding orthonormal
columns successively. Orthogonality fixes complex coordinates
with unit Jacobian; the remaining unit length equation in
complex dimension \(j\) contributes \(S_j/2\) by polar
coordinates. The mass of that fiber is \(V_r/2^r\), which
verifies the stated normalization.

Disintegrate the square Gram fiber first in the columns of the
frame and then in the remaining columns. Its conditional fiber
over \(\mathcal F\) can be computed there, since translation
by any isometry has complex determinant of absolute value one.
The equations pairing the remaining columns with \(\mathcal F\)
fix their dual rows and \(z\)-rows with unit Jacobian. For the
\(a\) dual columns the remaining rows are \(K-B^*B/2,B\),
with factor \(2^{-a}\) from their zero Gram as before.
For the \(k\) positive columns, pairing with these duals fixes
their top rows as \(-B^*h\) with unit Jacobian; their remaining
Gram equation is \(h^*h=I_k\). Its density is \(V_k\,dh/2^k\)
by the compact calculation. The ordering of these equations
only permutes real or complex coordinates, so the full
factor along the right copy of \(P\) in \(d\nu_r\) is
\(2^{-a-k}V_k\,ds\). These successive changes of variables
may first be made in local submersion charts for continuous
compactly supported functions and then extended to nonnegative
measurable functions by uniqueness of the resulting measures.
Combining them multiplies the frame density by
\(2^{r-a-k}V_k/V_r\). Since \(r-a-k=n\), this proves the formula.
The same coordinates show that two representatives differ by
left and right factors in \(P\); Haar invariance proves the
last assertion.
\end{proof}

\begin{corollary}[The density at a rational term on the marked slices]
\leavevmode\label{cor:realmarked682}
For a rational term with \(\Delta\ne0\), write \(Z_v\) for the
frame products at each normalized real place of the shorter
construction. Set \(\kappa_v=|M\Phi_v(\delta_E)|/(2\pi)\) and
\(b_{j,v}=\kappa_v^{2(n-j)}/N\) for \(1\le j\le n\).
The coefficient of the two Haar integrals in the product of the
open-cell densities at \((Z_v)_v\) is
\[
 \left(\frac{V_kS_k}{V_r}\right)^f
 \left(\prod_{v,j} b_{j,v}\right)^{2(k-1)}
             {\rm N}_{F/\mathbb Q}(\Delta)^{k-1}.
\]
For terms on the \(p_i\)-slice, its reciprocal is bounded above
by
\[
 \left(\frac{V_r}{V_kS_k}\right)^f
 \left(\prod_{v,j} b_{j,v}\right)^{-2(k-1)}
                 p_i^{-2nf(k-1)} .
\]
\end{corollary}
\begin{proof}
The real frame is obtained from the rational frame by scaling
the last vector by \(1/\sqrt N\) and the lengths along consecutive
left lines by successive additional factors \(\kappa_v\).
In the group coordinates this is the diagonal substitution
making the middle norm \(-1\) and the simple-root character
have unit coefficient in the real wave. Indeed the archimedean
coefficient in the wave \(\exp(s-\bar s)\) has absolute value
\(2\pi/|M\Phi_v(\delta_E)|=\kappa_v^{-1}\); the ratios of adjacent
frame scales are \(\kappa_v\). Phases of basis vectors cancel in
absolute determinants. Thus each leading Gram determinant of
order \(l\) has its absolute value multiplied by
\(\prod_{j\le l} b_{j,v}\). Since \(b_{n,v}=1/N\), subtraction
of their squared absolute values at orders \(n\) and \(a\)
gives \(\delta(Z_v)=(\prod_j b_{j,v})^2\,\Phi_v(\Delta)>0\).
Multiply the densities of Theorem~\ref{thm:realgram682} and
use Theorem~\ref{thm:discriminant681} for the reciprocal bound.
\end{proof}

The density describes integration in the continuous Gram
variable. Estimating a period at a specified rational value
requires control of the function being integrated there.


\subsection{Pointwise restriction at a prescribed real Gram matrix}
\label{sec:gramtrace683}

We pass from integration in \(Z\) to evaluation at a single \(Z\).
Keep \(G,P,\Omega,n,k,r\) and the Haar measures of
Theorem~\ref{thm:realgram682}. Write \(S\) for the real Bessel
subgroup from Section~\ref{sec:rationalflags680}. Its action
on the ordered frame has image the group \(N_n\) of complex
upper unitriangular matrices and kernel \(P\). Indeed the
upper unitriangular changes among the first \(a\) lines extend
by their inverse adjoints on the dual lines. The possible
changes \(z\mapsto z+l,\ l\in L\), are all realized by the
unipotent radical: for a middle-space column map \(D\) from
the dual lines into \(\mathbb Cz\), the middle coordinates
change by \(Dv\), and the top line coordinates by
\(-D^*J_{\rm mid}(x,t)-D^*J_{\rm mid}Dv/2\). This preserves the
form and the linear map \(D\mapsto-D^*J_{\rm mid}z\) ranges
over \(L\). We use Lebesgue measure on the strict upper entries
of \(N_n\), which is Haar in these triangular coordinates.
Together with \(ds\) on \(P\) this specifies Haar on \(S\);
equivalently it is Haar on its unipotent radical times
probability Haar on \(\mathrm U(k)\). All quotient measures
below use these choices.

Put \(c_G=V_k S_k/V_r\) and let \(d=2n^2\), the real dimension of
the Gram variable. Use operator norm for complex matrices. For
\(Z_0\in\Omega\) write \(\xi_0=(Z_0^{-1})_{nn}\). A real number
\(P_0\) will satisfy
\[
 P_0\ge\max\{8r,\,\|Z_0\|,\,\|Z_0^{-1}\|,
                   \,(1-|\xi_0|^2)^{-1/2}\}.
\]
Use the complex entry coordinates on matrices, regarded as
\(d\) real coordinates. Let \(Q\) be the cube centered at \(Z_0\)
with coordinate half-width \(\epsilon=P_0^{-16}/4\).

\begin{lemma}[A section with quantitative control in a Gram cube]
\leavevmode\label{lem:gramchart683}
On \(Q\) one can choose the representatives \(g_Z\) of
Theorem~\ref{thm:realgram682} real analytically, fixing a positive
space of dimension \(k-1\) pointwise. On \(Q\),
\(\delta(Z)^{k-1}\ge\delta(Z_0)^{k-1}/2\).
Both \(g_Z\) and \(g_Z^{-1}\) have analytic continuations of
operator norm at most \(P_0^{32}\) to the polydisc of radius
\(P_0^{-16}\) in the complexified real coordinates around \(Z_0\).
\end{lemma}
\begin{proof}
Use \(Y=Y_Z\) constructed in Theorem~\ref{thm:realgram682},
with the sphere direction a fixed positive coordinate vector.
This takes values in the \(2n\)-space consisting of the two
isotropic arms, \(z\), and this one positive direction. Write
\(Y_a=(Y_1,\ldots,Y_a)\) and put
\[
 U=\mathcal F(Z^*)^{-1},\qquad w=Z^{-1}e_n,\qquad
 D=U_{[1,a]}+\tfrac12Y_a w_{[1,a]}w_{[1,a]}^* ,
\]
where the subscript on \(U\) selects its first \(a\) columns.
Since \(Y^*JU=I\) and \(U^*JU=-ww^*\), the columns of \(D\)
are isotropic duals to \(Y_a\) and are orthogonal to \(Y_n\).
Set \(t=U_n+Y_n\). This is orthogonal to every column of \(Y\)
and has norm for \(J\) equal to \(1-|\xi|^2\), where
\(\xi=(Z^{-1})_{nn}\). Subtracting \(Y_a D^*Jt\) makes \(t\)
orthogonal to \(D\) and preserves its norm. Send the source frame
to \(Y\), the duals to \(D\), and the distinguished positive
unit vector to \((t-Y_aD^*Jt)/(1-|\xi|^2)^{1/2}\). Extend by the
identity on the other \(k-1\) positive directions. The stated
Gram relations give an isometry and specify \(g_Z\).

For the analytic estimates write \(p=P_0\). Extend both \(Z\)
and its entrywise conjugate using the same complexified real
coordinates, with opposite signs on the imaginary part.
Perturbations of either matrix have operator norm at most
\(2np^{-16}\). The inverse perturbation series bounds the new
inverse norms by \(2p\), and their changes by \(4np^2p^{-16}\).
In particular the change in \(1-\xi\bar\xi\), relative to its
central value, is at most \(12np^4p^{-16}\): at the center
\(|\xi_0|<1\) and \(1-|\xi_0|^2\ge p^{-2}\). Its analytic square
root is thus single valued and its reciprocal has norm at most
\(2p\). On the real cube the variation of
\(\log|\det Z/\det Z_0|\) is at most \(4n^2p\,p^{-16}\), by
the convergent matrix logarithm of \(I+Z_0^{-1}(Z-Z_0)\).
In \(\log(\delta(Z)/\delta(Z_0))\) this term occurs twice, plus
the logarithm of the relative change in \(1-|\xi|^2\).
Multiplication by \(k-1\le p\) gives absolute value less than
\(\log 2\), since \(p\ge8r\). This proves the lower bound.

For clarity, on the complex polydisc the free blocks of \(YZ^{-1}\)
are given by the previous Gram equations with \(B=K=0\).
The row \(\eta\) has length at most \(2p\), as does its continued
conjugate. The formulas there give \(\|Y\|\le p^7\).
The displayed formulas here then give \(\|U\|\le p^2\),
\(\|D\|\le p^{12}\), \(\|t\|\le p^8\), and
\(\|Y_aD^*Jt\|\le p^{27}\). With \(p\ge8r\) these bounds give
\(\|g_Z\|\le p^{32}\), also for the continued conjugate matrix.
The analytic continuation of \(g_Z^{-1}\) is the continuation
of \(Jg_Z^*J\) and satisfies the same bound.
\end{proof}

\begin{lemma}[The quotient measure in the Gram cube]
\leavevmode\label{lem:gramquotient683}
For every nonnegative measurable left \(S\)-invariant function
\(v\) and every measurable \(P'\subset P\),
\[
 \int_Q\int_{P'} v(g_Zs)\,ds\,dZ
       \le \frac{2}{c_G\,\delta(Z_0)^{k-1}}
                    \int_{S\backslash G}v(g)\,d\dot g .
\]
\end{lemma}
\begin{proof}
Theorem~\ref{thm:realgram682}, disintegrated on the left by \(P\),
gives the measure on \(P\backslash G\) over the open set as
the image of \(c_G\delta(Z)^{k-1}dZ\,ds\) under
\((Z,s)\mapsto P g_Zs\). To verify this for a compactly
supported function on the quotient one can multiply its lift
by a nonnegative function of left-orbit integral one, using
local sections and a partition of unity. The left integral in
the theorem removes that factor. Monotone exhaustion gives the
general measure identity. Use Lemma~\ref{lem:gramchart683} and
enlarge \(P'\) to \(P\). It suffices to show that the part with
\(Z\in Q\) of each \(P\backslash S\)-fiber has measure at most one.

On a fiber meeting the cube choose a representative \(g_1\) with
\(Z_1=Z(g_1)\in Q\). For \(h\in S\),
\(Z(hg_1)=((h^{-1})|_T)^* Z_1\), since \(h\) preserves the
form. Write \(A=((h^{-1})|_T)^*\), a lower unitriangular matrix.
If \(AZ_1\in Q\), then
\(\|A-I\|\le\|AZ_1-Z_1\|\,\|Z_1^{-1}\|
                 \le 8n\epsilon P_0<1/4\).
Haar on \(P\backslash S\) is Lebesgue measure in the entries
of \(A\): inversion on a unipotent triangular group preserves
Haar, and adjoint exchanges upper and lower coordinates with
unit absolute Jacobian. Each available complex entry is in a
disc of area less than one. The volume of their product is
at most one. Integrating over \(S\backslash G\) proves the lemma.
\end{proof}

Let \(X_1,\ldots,X_{r^2}\) be a real Hilbert--Schmidt orthonormal
basis of \(\operatorname{Lie}G\), with the real inner product
\(\Re\operatorname{tr}(X^*Y)\). Let \(R_X W(g)\) be the
derivative of \(W(g\exp(tX))\) at zero, and \(R_{\bf i}\) the
successive derivatives in a word \({\bf i}\). Suppose \(W\) is a
smooth, finite-dimensional Hilbert-space-valued function, with
\(W(hg)=\tau(h)W(g)\) for a unitary representation \(\tau\) of \(S\).
For an integrable compactly supported scalar \(\psi\) set
\(T_\psi W(g)=\int_G\psi(y)W(gy)\,dy\). Right derivatives of
\(W\) and the corresponding convolutions have the same left
equivariance, so their squared norms descend to \(S\backslash G\).

\begin{theorem}[A restriction estimate for localized Bessel convolutions]
\leavevmode\label{thm:gramrestriction683}
Put \(p=P_0\), \(P'=\{s\in P:\|s\|\le p^{20}\}\), and take
\(\phi\in L^2(G)\) supported on \(\|y\|\le p^{20}\).
For every word of length at most \(d\), suppose
\[
       \|T_\psi(R_{\bf i}W)\|_{L^2(S\backslash G)}
              \le A p^{|{\bf i}|}\|\psi\|_{L^2(G)}
\]
for all \(\psi\) supported on the same ball. Then
\[
 \int_{P'}\|T_\phi W(g_{Z_0}s)\|^2\,ds
       \le \frac{p^{1000n^2}}{c_G\,\delta(Z_0)^{k-1}}\,
                      A^2\|\phi\|_{L^2(G)}^2 .
\]
There are no differentiability assumptions on the scalar test
function \(\phi\).
\end{theorem}
\begin{proof}
The ball in \(G\) is compact: a unitary matrix for \(J\)
satisfies \(g^{-1}=Jg^*J\), and \(G\) is closed in matrix space.
The same holds for \(P'\). In particular \(\phi\) has finite
\(L^1\)-norm. Differentiation in \(Z\) under its integral is
valid by local uniform boundedness of the derivatives of \(W\)
on compact sets.

We bound a mixed derivative in any \(l\le d\) distinct real
coordinates \(Z_j\). Write \(B_j(Z)=g_Z^{-1}\partial_j g_Z\).
Differentiating \(W(g_Zs y)\) uses the Lie element
\(\operatorname{Ad}(y^{-1})\operatorname{Ad}(s^{-1})B_j(Z)\)
on its right. We describe the bounds after iteration, including
the coefficients that depend on \(y\). Express the adjoint
maps and \(B_j\) in the fixed orthonormal basis. Each matrix
entry of \(\operatorname{Ad}(y^{-1})/p^{40}\) and of
\(\operatorname{Ad}(s^{-1})/p^{40}\) has modulus at most one.
This follows from \(\|g^{-1}Xg\|_{\rm HS}\le
\|g\|^2\|X\|_{\rm HS}\) on the indicated balls. In each
differentiated integral absorb the entries depending on \(y\),
after division by \(p^{40}\), into \(\phi\). The resulting
weights are \(\phi\) times products of such entries; each has
\(L^2\)-norm at most \(\|\phi\|_2\) and is independent of \(s,Z\).

More explicitly the mixed derivative is a sum of terms
\(c(Z,s)\,T_\psi(R_{\bf i}W)(g_Zs)\), with \(|{\bf i}|\le l\).
The sum of the suprema of \(|c|\) on \(Q\times P'\) is at
most \(p^{180l}\). Here are estimates establishing this bound.
On the polydisc of radius \(3p^{-16}/4\), Cauchy's inequality
and Lemma~\ref{lem:gramchart683} give
\(\|B_j\|\le4p^{80}\); each Lie coefficient, extended
complex linearly from the real algebra, is bounded by \(p^{84}\).
For the latter assertion one may choose the orthonormal basis
as \(J\) times the usual orthonormal skew-Hermitian basis.
Its coordinates complexify to trace pairings of Hilbert--Schmidt
norm one; any other real orthonormal basis is related by an
orthogonal change of coordinates. A derivative of an existing
coefficient costs at most \(4(d+1)p^{16}\le p^{20}\), by shrinking
the polydisc radius by \(p^{-16}/(4(d+1))\) at each step.
When the derivative hits \(W\), the expansion uses at most
\(r^6\) summands, from the three Lie indices, of size at most
\(p^{84+80}\) times the old coefficient. The factors depending
on \(s\) stay in the coefficient and depend on no Gram variables.
Thus the sum of suprema increases by at most
\(p^{20}+r^6 p^{164}\le p^{180}\) per derivative. The final
polydisc still contains \(Q\). This proves the indicated expansion
and estimate by induction, starting from coefficient one.

Take the \(L^2(Q\times P')\)-norm. For each term,
Lemma~\ref{lem:gramquotient683}, applied to the squared norm
of its convolution, bounds that norm after removal of \(c\) by
\[
           \left(\frac{2}{c_G\delta(Z_0)^{k-1}}\right)^{1/2}
                      A p^l\|\phi\|_2 .
\]
The triangle inequality gives the same expression times
\(p^{180l}\) for the whole mixed derivative.

Finally, on an interval of length \(2\epsilon\), evaluation
at its center is bounded in norm by
\((2\epsilon)^{-1}\int\|u(t)\|\,dt+\int\|u'(t)\|\,dt\),
by the fundamental theorem of calculus. Apply this successively
in all \(d\) coordinates with values in \(L^2(P')\); equivalently
one may iterate the averaging and signed integration operators
of the one-variable identity. Cauchy--Schwarz bounds the result
by \((2\epsilon)^{-d/2}\) times the sum, over coordinate subsets,
of \((2\epsilon)^l\) times the corresponding mixed derivative
norm. Since \(2\epsilon=p^{-16}/2\), \(l\le d=2n^2\) and
\(p\ge8r\), the scalar factor after squaring is at most
\[
  2(2\epsilon)^{-d}\big(2^d p^{181d}\big)^2
       \le p^{1000n^2}.
\]
This proves the claim.
\end{proof}

The estimate evaluates a real convolution at the prescribed
Gram value, whereas Theorem~\ref{thm:realgram682} integrates
over Gram values. Corollary~\ref{cor:realmarked682} supplies
the product of its inverse densities at the actual rational
terms. The remaining factor \(A\) concerns the specific
archimedean Bessel vector and its discrete-series norm.


\subsection{The formal-degree bound for the real Bessel average}
\label{sec:besselschu684}

We evaluate the quotient-norm constant in
Theorem~\ref{thm:gramrestriction683} for the real function
that enters the shorter construction. All measures are those
of Theorem~\ref{thm:realgram682} and Section~\ref{sec:gramtrace683}.
In this subsection \(K\) is a maximal compact subgroup of
\(G=\mathrm U(a+k,a+1)\). Let \(\pi\) be its discrete series of
Harish-Chandra parameter \(\lambda\). Write \(d_\pi\) for its
formal degree with respect to \(dg\). Inner products are
conjugate linear in the first variable. Let \(w\) be its unit
lowest-\(K\)-type coefficient vector, \(x\in G\), and
\(b=\pi(x)w\). Let \(\tau\) be the unitary representation
of \(S\) on the finite-dimensional value space \(V_\tau\),
given by the simple-root character on the unipotent radical
and the chosen type on \(\mathrm U(k)\). For \(v\in V_\tau\)
and a \(K\)-finite vector \(u\) of \(\pi\), set
\[
 c_u(g)=\langle b,\pi(g)u\rangle,\qquad
 W_u(g)=\int_S c_u(sg)\tau(s)^{-1}v\,ds .
\]
In particular \(c_w(g)=\langle w,\pi(x^{-1}g)w\rangle\);
the freedom to choose the real point \(x\) is part of the
Bessel average.

\begin{proposition}[Orthogonality after averaging a discrete-series coefficient]
\leavevmode\label{prop:averagedschur684}
Suppose \(|\lambda_i-\lambda_j|>r-1\) for every noncompact
root \(e_i-e_j\). Then \(W_u\) exists, is smooth and bounded,
and belongs in quotient norm to \(L^1\cap L^2(S\backslash G)\).
It transforms on the left by \(\tau\). For some real \(m\ge0\)
one has
\[
 \begin{split}
 m&=\int_S\langle b,\pi(s)b\rangle\,
                  \langle v,\tau(s)^{-1}v\rangle\,ds,\\
 \langle W_u,W_{u'}\rangle_{S\backslash G}
         &=\frac{m}{d_\pi}\langle u,u'\rangle.
 \end{split}
\]
In particular \(m=d_\pi\|W_w\|_2^2\). For every compactly
supported \(\psi\in L^2(G)\),
\[
    \|T_\psi W_u\|_{L^2(S\backslash G)}
          \le \frac{\|W_w\|_2}{\sqrt{d_\pi}}\,
                        \|u\|\,\|\psi\|_2 .
\]
\end{proposition}
\begin{proof}
We use the discrete-series integrability criterion in
\cite{XXV-XX-XX-milicic684}. For root \(\alpha\) the threshold is
\(\tfrac14\sum_{\beta}|\langle\alpha,\beta\rangle|\) in the
same metric as the pairing with \(\lambda\). In type \(A_{r-1}\)
take roots of squared length two. The two roots \(\pm\alpha\)
contribute four; \(4(r-2)\) further roots each contribute one.
This gives \(r-1\). The criterion for the semisimple derived
group applies here since multiplication by the compact center
is a finite covering and the absolute values of coefficients
are invariant in that central direction. Every \(K\)-finite
coefficient is consequently in \(L^1(G)\), including its
derivatives: applying any Lie algebra word to a \(K\)-finite
vector leaves it \(K\)-finite by covariance under the finite
dimensional adjoint representation. Fixed left or right
translations have the same property of integrability of all
derivatives. For example a left derivative after left translation
is the translation of the derivative conjugated by its group
element. This permits \(b=\pi(x)w\) at arbitrary \(x\).

Here is the point that allows integration on the noncompact
subgroup \(S\). If \(c\) is such a coefficient or translate,
then \(\sup_g\int_S|c(sg)|ds<\infty\). Take a small coordinate
box transversal to \(S\) in \(G/S\) at the identity and a smooth
section \(t\mapsto z(t)\) with \(z(0)=1\). Evaluation at zero of
\(c(z(t)sg)\) is bounded by a fixed sum of integrals of absolute
mixed derivatives on this box, by the iterated one-dimensional
formula used in Theorem~\ref{thm:gramrestriction683}. These
derivatives are linear combinations with bounded coefficients
of left-translation derivatives of \(c\) at \(z(t)sg\).
The coordinates \(z(t)s\) disintegrate Haar as \(j(t)\,dt\,ds\)
with smooth positive \(j\), bounded below on the closed box.
Right translation by \(g\) preserves Haar. Thus integrating
this point bound in \(s\) is bounded by a finite sum of global
\(L^1(G)\)-norms, independent of \(g\). It applies to the
derivatives also. Differentiation of \(W_u\) along right
translations now gives \(W_{d\pi(X)u}\); for example the
integrated Taylor remainder is bounded by a constant times
the square of the increment using the same bound on the second
derivative. The argument iterates. It proves smoothness and
boundedness; changing the integration variable gives the
equivariance. The quotient \(L^1\)-norm is at most
\(\|v\|\|c_u\|_1\) by disintegration. Boundedness gives \(L^2\).

Expand the inner product of the two averages, write the second
integration variable in \(S\) as \(hs\), and unfold \(sg=t\)
over \(G\). Absolute interchange follows from
\[
 \int_G |c_u(t)|\int_S|c_{u'}(ht)|\,dh\,dt<\infty
\]
by the uniform restriction estimate just proved. Schur
orthogonality, in the normalization defining \(d_\pi\), gives
\[
 \int_G\overline{c_u(t)}\,c_{u'}(ht)\,dt
     = d_\pi^{-1}\langle b,\pi(h)b\rangle\,\langle u,u'\rangle .
\]
The value-space inner product is
\(\langle v,\tau(h)^{-1}v\rangle\). This proves the formulas
and the nonnegativity assertion by taking \(u=u'=w\).
The same computation holds for finite sums of translates of
\(K\)-finite vectors, since the coefficient bounds used for
Fubini still apply. Hence \(u\mapsto W_u\) extends by density
to a scaled isometry of norm \(\|W_w\|_2\) (the zero map if
\(m=0\)), and intertwines \(\pi\) with right translation.
Consequently \(T_\psi W_u\) is the image of
\(\int\psi(y)\pi(y)u\,dy\). Testing this vector against a unit
vector and using Cauchy--Schwarz and Schur orthogonality gives
norm at most \(\|u\|\|\psi\|_2/\sqrt{d_\pi}\).
\end{proof}

\begin{lemma}[Derivatives of the lowest compact type]
\leavevmode\label{lem:casimirs684}
Suppose \(\max_i|\lambda_i|\le L\), and let the real Lie basis
be as in Theorem~\ref{thm:gramrestriction683}. For any word
of length \(s\le d=2n^2\),
\[
       \|d\pi(X_{i_s})\cdots d\pi(X_{i_1})w\|
           \le \big[\,4\sqrt r\,(L+3r+d+1)\,\big]^s .
\]
\end{lemma}
\begin{proof}
Conjugate the form to a diagonal sign matrix by a unitary
change of basis, preserving the Hilbert--Schmidt metric.
In the Cartan decomposition choose orthonormal bases of the
compact and noncompact parts. For \(\Re\operatorname{tr}(XY)\)
the corresponding Casimir is
\(\Omega_G=\sum_{\mathfrak p}X^2-\sum_{\mathfrak k}X^2\);
write \(C_K=-\sum_{\mathfrak k}X^2\). Differentiated unitarity
on smooth \(K\)-finite vectors gives
\[
 \sum_i\|d\pi(X_i)u\|^2
      =\langle u,(2C_K-\Omega_G)u\rangle .
\]
The sum is unchanged under real orthogonal change of basis.
The Harish-Chandra parameter gives the usual Casimir scalar
\(\|\lambda\|^2-\|\rho\|^2\), with roots of squared length two.
On a compact highest weight \(\mu\), \(C_K\) has scalar
\(\|\mu+\rho_c\|^2-\|\rho_c\|^2\). These formulas use the
trace form, including the compact center. Equivalently, in
the complex matrix basis the quadratic element is
\(\sum_{ij}E_{ij}E_{ji}\); commuting opposite root vectors
into triangular order adds \(2\rho\) to its linear term,
which gives the indicated shifted-square formula. The compact
calculation uses the same root pairs within the compact blocks.

The lowest compact type has highest weight
\(\lambda+\rho_n-\rho_c\), by the discrete-series lowest-type
theorem \cite{XXV-XX-XX-blattner684}; each coordinate differs from
the corresponding \(\lambda_i\) by at most \(r\).
All weights of this compact type have coordinates bounded in
absolute value by \(L+r\): within either unitary block the
weights lie in the convex hull of the permutations of its
highest weight, by the highest-weight theorem. Lie algebra
application is a \(K\)-map from the tensor product with the
adjoint representation. Each weight of that representation is
zero or a root, and changes a coordinate by at most one.
After \(t\) applications every possible compact highest weight
has coordinates bounded by \(L+r+t\). Since those of
\(\rho,\rho_c\) are at most \(r\), inserting these bounds in
the displayed Casimir identity bounds each next derivative by
\(4\sqrt r(L+3r+t+1)\|u\|\), also for a sum of compact types
by their orthogonality. Iteration proves the lemma.
\end{proof}

\begin{theorem}[Pointwise real estimate with formal degree]
\leavevmode\label{thm:archbessel684}
Take the real Bessel average \(W=W_w\) above, or a scalar multiple,
with a parameter satisfying the gap in
Proposition~\ref{prop:averagedschur684}. With the Gram and support
conditions of Theorem~\ref{thm:gramrestriction683}, choose also
\(P_0\ge4\sqrt r(L+3r+d+1)\). Then
\[
 \int_{P'}\|T_\phi W(g_{Z_0}s)\|^2\,ds
  \le \frac{P_0^{1000n^2}}{c_G\delta(Z_0)^{k-1}}\,
             \frac{\|W\|_2^2}{d_\pi}\,\|\phi\|_2^2 .
\]
For \(W\) scaled to squared quotient norm \(\dim V_\tau\), the
numerator of the formal-degree factor is \(\dim V_\tau\).
\end{theorem}
\begin{proof}
Apply Proposition~\ref{prop:averagedschur684} with \(u\) a word
in \(d\pi(X_i)\) applied to \(w\). Its Bessel average is the
corresponding right derivative of \(W\). The derivative lemma
gives the hypothesis of Theorem~\ref{thm:gramrestriction683}
with \(A=\|W\|_2/\sqrt{d_\pi}\), including the case \(W=0\).
The construction uses the doubled parameter \(2\lambda\); its
separation for the undoubled entries is at least \(500r^3\),
which exceeds \(r-1\). The test and the Gram matrix use the
normalizations of the preceding theorem. Rescaling multiplies
both squared norms by the same factor.
\end{proof}

This estimate holds for every translate \(x\). The identity for
\(m\) in Proposition~\ref{prop:averagedschur684} identifies the
positive local period that must be bounded below to choose and
control a nonzero normalized vector in a prescribed range of \(x\).


\subsection{One open-cell pairing at all infinite places}
\label{sec:realpair685}

We estimate the real integral associated with a single rational
term of nonzero discriminant. The estimate keeps the cut error
as an integral with its original weight. We begin at one real
place, retaining \(G,S,P,g_Z,c_G\) from the preceding subsections.
A point \(y\) lies in the frame box of radius \(\rho\) if
\(\|y^{-1} f_j\|\le\rho\) for all normalized frame vectors.
Set \(\ell(y)=\inf_{h\in S}(\|hy\|+1)\). For \(Z\in\Omega\)
with \(|Z_{ij}|\le\rho^2\) and \(\rho\ge8r\), put
\[
             H=\frac{\rho^{9n}}{\delta(Z)},\qquad T=H^{140n}.
\]

\begin{lemma}[Bounding a relative position by the two frames]
\leavevmode\label{lem:pairbox685}
If \(y\) is in the frame box and \(\ell(g_Z y)\le\rho\), then
\(\|y\|\le T\) for the section of Lemma~\ref{lem:gramchart683}.
\end{lemma}
\begin{proof}
We have \(\|Z\|\le n\rho^2\le\rho^3\),
\(\delta(Z)\le|\det Z|^2\le\rho^{6n}\), and \(H\ge\rho^{3n}\).
The adjugate bound yields
\(\|Z^{-1}\|\le\rho^{3(n-1)}/\sqrt{\delta(Z)}\).
Also \((1-|(Z^{-1})_{nn}|^2)^{-1/2}
=|\det Z|/\sqrt{\delta(Z)}\). All the parameters needed for
Lemma~\ref{lem:gramchart683} are at most \(H^2\), so
\(\|g_Z\|=\|g_Z^{-1}\|\le H^{64}\).
Work in its fixed-coordinate \(2n\)-space \(V_*\), spanned by
the isotropic and dual arms, \(z\) and one positive middle
coordinate. On this space form the square matrix
\(D=(\mathcal F,g_Z^{-1}\mathcal F)\). Its joint Gram has
off-diagonal block \(Z^*\), so the determinant calculation of
Theorem~\ref{thm:rationalflags680} at \(N=1\) gives
\(|\det D|=\sqrt{\delta(Z)}\). We have \(\|D\|\le H^{65}\), whence
\(\|D^{-1}\|\le H^{131n}\) by the singular value product formula.

Choose \(h\in S\) with \(\|h g_Zy\|\le2\rho\), and write
\(h^{-1}\mathcal F=\mathcal F B\) with \(B\) upper unitriangular.
The two matrices \(y^{-1}\mathcal F\) and
\(y^{-1}g_Z^{-1}\mathcal F B\) have norms at most \(\sqrt n\rho\)
and \(2\rho\), respectively. Their cross Gram is \(Z^*B\),
so \(\|B\|\le 2\sqrt n\rho^2\|Z^{-1}\|\le H^4\). Since
\(\det B=1\), \(\|B^{-1}\|\le H^{4n}\). It follows that the
restriction of \(y^{-1}D\) has norm at most \(H^{5n}\) and
\(\|y^{-1}|_{V_*}\|\le H^{136n}\). The negative eigenspace of
\(J\) is contained in \(V_*\). Writing \(P_-\) for its orthogonal
projection, any \(A\in G\) satisfies
\(AA^*=J+2AP_-A^*\) because \(AJA^*=J\). This bounds
\(\|y^{-1}\|^2\) by \(1+2H^{272n}\), and proves the claim.
\end{proof}

\begin{lemma}[A set of shifts in the real frame stabilizer]
\leavevmode\label{lem:shifts685}
Let \(T\ge2\), \(p\ge64r^3(2T)^8\), and \(C_T=\{g:\|g\|\le T\}\).
There is a compact inversion-invariant set \(B_p\subset P\) of
positive finite measure, with \(\|s\|\le p^7\) on \(B_p\), such
that \(\mathcal A=B_p C_T\) satisfies, for every \(W\) transforming
on the left by a unitary representation of \(S\),
\[
 \frac1{\operatorname{vol}_P(B_p)}\int_{\mathcal A}\|W(y)\|^2\,dy
          \le 2(2T^2)^{n(n-1)}\|W\|_{S\backslash G}^2 .
\]
\end{lemma}
\begin{proof}
In the coordinates of Theorem~\ref{thm:realgram682}, take all
\(h\in\mathrm U(k)\) and balls
\(\|B\|_{\rm HS}\le p^3,\ \|K\|_{\rm HS}\le p^6\).
The matrix formula gives the norm bound; inversion takes
\(B\) to \(-h^{-1}B\) and \(K\) to \(-K\). To bound the integral
by quotient disintegration, it suffices to bound the measure of
\(\{s\in S:sy\in\mathcal A\}\). On any left \(S\)-orbit meeting
\(\mathcal A\), choose a representative \(y_0\in C_T\).
Any \(s\) with \(sy_0\in B_pC_T\) is in
\(B_p(S\cap C_{T^2})\), because
\(\|y_0^{-1}\|=\|y_0\|\le T\).

For \(s_0\in S\cap C_{T^2}\), both its frame matrix \(A\in N_n\)
and \(A^{-1}\) have norm at most \(T^2\). A lift \(l(A)\in S\)
of norm at most \((2T^2+2)^3\) is obtained by extending
the leading \(a\)-block of \(A\) diagonally as in
Section~\ref{sec:gramtrace683} and following it by the shear
that adds the top part of its last column to \(z\). The diagonal
factor has norm at most \(1+T^2\); the shear has norm at most
\((1+\|A\|)^2+1\) by its displayed action there. The element
\(s_0 l(A)^{-1}\) is in \(P\); its \(B\)-coordinate is a matrix
block and its \(K\)-coordinate is the skew part of a block.
Each has Hilbert--Schmidt norm at most
\(\sqrt r\,T^2(2T^2+2)^3\le p\).

Multiplying \(B_p\) on the right by any such element of \(P\)
increases the two radii by at most \(p\) and \(p+p^4\).
Indeed the central cross term in the product is the skew
part of a product of the two \(B\)-coordinates, up to unitary
factors. Its Hilbert--Schmidt norm is at most the product of
their norms. The ratio of the enlarged \(P\)-volume to the
original is thus at most
\((1+p^{-2})^{2ka}(1+2p^{-2})^{a^2}\le2\) for the stated
\(p\). Integrate this bound over \(A\) using the Haar quotient
\(P\backslash S=N_n\). Each of its \(n(n-1)/2\) complex
coordinates has modulus at most \(T^2\), so their volume is
at most \((2T^2)^{n(n-1)}\). This bounds the fiber measure
and proves the lemma.
\end{proof}

We now allow all \(f\) real places. Use one copy of these groups
and measures for each place \(v\), and product Haar on
\(G_\infty\). Let \(W_v\) be any scalar multiple of the
translated discrete-series Bessel average of
Theorem~\ref{thm:archbessel684} and \(W=\bigotimes_v W_v\).
Let \(W_v^\flat\) multiply \(W_v\) by a bounded left
\(S\)-invariant cutoff supported on \(\ell(y)\le\rho_v\), and
\(W^\flat=\bigotimes_v W_v^\flat\). For
\(\gamma=(\gamma_v)_v\) define
\[
 K_\gamma(y)=W(\gamma y)\otimes W(y)^*,\qquad
 K_\gamma^\flat(y)=W^\flat(\gamma y)\otimes W^\flat(y)^* .
\]
Write \(Z_v=Z(\gamma_v)\) and \(H_v=\rho_v^{9n}/\delta(Z_v)\).

\begin{theorem}[A cut pairing at a rational Gram value]
\leavevmode\label{thm:realpair685}
Suppose \(n\ge32\), \(Z_v\in\Omega\), \(\rho_v\ge8r\), and
\(|(Z_v)_{ij}|\le\rho_v^2\). At each place assume the weight
factor \(4\sqrt r(L_v+3r+2n^2+1)\le H_v^2\). Let \(\phi\)
be a measurable scalar function on \(G_\infty\), bounded by
one in absolute value and invariant under left translation by
\(\prod_v P\). Suppose that whenever \(\phi(y)K^\flat_\gamma(y)\)
is nonzero, each \(y_v\) is in the frame box of radius \(\rho_v\).
Then this cut integrand is supported in a compact set and
\[
 \begin{split}
 \left\|\int_{G_\infty}\phi(y)K^\flat_\gamma(y)\,dy\right\|_{\rm HS}
 &\le
 \left(\prod_v
  \frac{H_v^{10^5 n^4}\,\|W_v\|_{S\backslash G}^{4}}
       {d_{\pi_v}\,c_G\,\delta(Z_v)^{k-1}}\right)^{1/2}\\
 &\quad+\int_{G_\infty}|\phi(y)|\,
             \|K_\gamma(y)-K_\gamma^\flat(y)\|_{\rm HS}\,dy .
 \end{split}
\]
The inequality is valid also for a nonfactorizable weight \(\phi\).
\end{theorem}
\begin{proof}
First work locally. By Theorem~\ref{thm:realgram682} we may
write \(\gamma=s_L g_Z s_R\) with \(s_L,s_R\in P\). Use
\(T=H^{140n}\), \(p=H^{2000n}\); the bounds in
Lemma~\ref{lem:pairbox685} and \(H\ge\rho^{3n}\) verify the
hypotheses for the chart, formal degree and
Lemma~\ref{lem:shifts685}. Choose \(s\) in \(B_p\) with its
normalized Haar restriction and set \(t(s)=s_R^{-1}s^{-1}\).
Any \(z\) in the frame box with \(\ell(\gamma t(s)z)\le\rho\)
has \(s^{-1}z\) in the frame box, as \(s\) fixes the frame;
also \(\ell(g_Zs^{-1}z)=\ell(\gamma t(s)z)\). Thus
\(s^{-1}z\in C_T\) by Lemma~\ref{lem:pairbox685}, so
\(z\in\mathcal A=B_p C_T\). In particular, at all places
together the nonzero cut integrand at \(t(s)z\) has
\(z\in\prod_v\mathcal A_v\), since each \(t(s)_v\) fixes the
frame and the cutoff enforces the length condition.
Taking any fixed shifts proves the compact support assertion.

For the full average consider the local map
\[
   \psi\longmapsto
       \left(s\longmapsto\int_{\mathcal A}
                              \psi(z)W(g_Zs^{-1}z)\,dz\right)
\]
from \(L^2(\mathcal A)\) to \(L^2(B_p;V_\tau)\) with unnormalized
measure there. Its squared operator norm is at most
\(p^{1000n^2}\|W\|_{S\backslash G}^2/
(d_\pi c_G\delta(Z)^{k-1})\) by
Theorem~\ref{thm:archbessel684}. Indeed \(\mathcal A\) has
norm at most \(p^7 T\le p^{20}\), \(B_p\) lies in the allowed
ball of shifts, and inversion preserves its measure. Apply
the tensor product of these operators to the scalar functions
\[
       \phi(z)\ \overline{\left(\bigotimes_v W_v(z_v)\right)_j}
               \quad\text{on }\prod_v\mathcal A_v
\]
for each value-space coordinate \(j\). The operator bound
multiplies: this follows first on Hilbert tensor products
from the positive-operator bound for each factor, and then
on product \(L^2\) by density of finite sums of rectangle
functions. Thus it holds for the indicated general \(\phi\).
Sum the squared output norms in \(j\). Left equivariance
under \(s_L\) and \(t(s)\) only inserts unitaries on the two
value spaces. Since \(\phi(t(s)z)=\phi(z)\), the resulting
quantity is the integral over the shifts of
\[
 \left\|\int_{\prod_v\mathcal A_v}
          \phi(t(s)z)K_\gamma(t(s)z)\,dz\right\|_{\rm HS}^2 .
\]
On the input side the sum of squared norms is at most the
product of \(\int_{\mathcal A_v}\|W_v(z_v)\|^2\,dz_v\).
Divide by the product of the shift volumes and apply
Lemma~\ref{lem:shifts685}. At each place the scalar loss,
apart from \(\|W_v\|_2^4/(d_{\pi_v}c_G\delta(Z_v)^{k-1})\),
is bounded by
\(2(2T^2)^{n(n-1)}p^{1000n^2}\le H^{10^5n^4}\):
with the chosen \(T,p\) its exponent to base \(H\) is at most
\(2{,}001{,}000\,n^3\), and \(n\ge32\).

For every shift the cut integral over \(\prod_v\mathcal A_v\),
after the indicated translation, is the original full-domain
cut integral by its support and Haar invariance. Its difference
from the corresponding integral of the full kernel is bounded
by the displayed tail on \(G_\infty\). Integrate this triangle
inequality against the probability measure of the shifts and
use Cauchy--Schwarz and the preceding second-moment estimate.
\end{proof}

For the shorter construction take \(\rho_v=b^{37n}\). Its
cut vanishes past length \(b^{6n+2}\), and the prescribed
large-node scales have \(n\ge32\) and \(r+L_v+1\le b^4\);
these bounds imply the numerical hypotheses above. For a
near rational term the two frames at a common point are both
in the box, so Cauchy--Schwarz there gives
\(|(Z_v)_{ij}|\le\rho_v^2\), since \(J\) has operator norm one.
The adjoint kernel \(W(y)\otimes W(\gamma^{-1}y)^*\) is covered
by replacing \(\gamma\) with \(\gamma^{-1}\) and conjugating
\(\phi\). The two Gram matrices are adjoints and have the same
\(\delta\). Corollary~\ref{cor:realmarked682} evaluates the
product of the density denominators on these very terms. To
bound the global second moment one must also sum the rational
orbits, their stabilizer weights, and the cut errors.


\subsection{Summing over rational frame products and one normal vector}
\label{sec:firstnormal686}

We return to the rational Hermitian space and to the integral
class lattices of Section~\ref{sec:discriminant681}. In this
subsection \(P\) is the rational group fixing every column of
\(\mathcal F\), and \(U_0\) is the kernel of \(P\to\mathrm U(X_0)\)
on the positive middle space. Thus \(P=U_0\rtimes\mathrm U(X_0)\).
Write \(D={\rm N}_{F/\mathbb Q}(\Delta)\) when \(\Delta\ne0\).
It is a positive integer for any integral class position. We
describe which rational orbits are counted by keeping the cross
Gram matrix \(C\) and the first positive normal \(q=q_\gamma\)
of Proposition~\ref{prop:normalvector681}.

\begin{proposition}[Rational frame-fixer orbits on the open cell]
\leavevmode\label{prop:rationalorbits686}
Each realized cross Gram matrix \(C\) with \(\Delta\ne0\)
determines a single \(P(F)\)-double orbit. The orthogonal
complement \(Y_\gamma=(T+\gamma T)^\perp\) is totally positive
of dimension \(k-1\), and the simultaneous pointwise stabilizer
of the two frames is \(\mathrm U(Y_\gamma)\). Its projection to
\(\mathrm U(X_0)\) is the full stabilizer of \(q_\gamma\).
The orbits under \(U_0(F)\) on the left and \(P(F)\) on the
right inside this double orbit are therefore distinguished
exactly by \(q_\gamma\).
\end{proposition}
\begin{proof}
The joint Gram is nonsingular, with negative index \(n\)
at every real embedding, by Theorem~\ref{thm:rationalflags680}.
These are all the negative directions in the ambient space.
Thus its orthogonal complement is totally positive of dimension
\(r-2n=k-1\). Matching the two frames for terms of the same
\(C\) gives an isometry of their joint spans. We recall an
extension argument over \(E/F\). Choose an orthogonal basis of
the common abstract nondegenerate span with nonzero diagonal
entries, by Hermitian diagonalization. Suppose two corresponding
next vectors \(u,v\) have the same nonzero norm and are both
orthogonal to the vectors already matched. For \(y=u-v\) and
\(\beta=h(y,u)\ne0\), the map
\(t\mapsto t-yh(y,t)/\beta\) preserves the form since
\(\beta+\bar\beta=h(y,y)\). It sends \(u\) to \(v\) and fixes
the preceding vectors. If \(\beta=0\), then
\(h(u+v,u)=2h(u,u)\ne0\); apply the same formula with target
\(-v\), and then negate on the nondegenerate line \(Ev\).
Induction gives the desired rational extension. Since it
matches both frames, it fixes \(\mathcal F\) and matches the
two images of \(\mathcal F\). This proves the double-orbit
assertion.

An element fixing both frames is the identity on their span,
and is exactly an arbitrary unitary transformation on its
orthogonal complement. The projection of \(Y_\gamma\) to
\(X_0\) is an isometry: \(Y_\gamma\subset T^\perp=L\oplus X_0\)
and its form is positive definite. Its image is orthogonal to
\(q_\gamma\), by pairing with
\(t=l+q_\gamma-(D_n/N)z\in\gamma T\) in
Proposition~\ref{prop:normalvector681}. Since
\(h(q_\gamma,q_\gamma)=\Delta/N\ne0\), dimension forces the
image to be the whole hyperplane \(q_\gamma^\perp\).
Every unitary transformation of this hyperplane extends,
through \(Y_\gamma\), to the simultaneous stabilizer just
described. Its action on \(X_0\) fixes \(q_\gamma\): it fixes
\(t\) and \(z\), and the middle projection on \(L^\perp/L\)
intertwines this action. Its image is therefore exactly
\(\mathrm U(q_\gamma^\perp)\), fixing \(q_\gamma\) pointwise.

Right multiplication by \(P(F)\) fixes the frame image. On
left multiplication by \(p\in P(F)\), \(C\) stays fixed and
\(q_\gamma\) changes to its image by the projection of \(p\).
The kernel is \(U_0(F)\). Quotienting the left copy of \(P(F)\)
by this kernel and by the simultaneous stabilizer gives the
stated classification.
\end{proof}

Let \(\Lambda_c\) be the integral class lattice in \(V\), and
write \(Q_c\) for its projection to \(X_0\) in the rational
decomposition. It is a full \(\mathcal O_E\)-lattice there.
The two hyperbolic arms are primitive direct lattice summands,
also after the class change. Removing their coordinates does
not affect \(Q_c\).

\begin{lemma}[Counting integral positive normals without a loss at \(N\)]
\leavevmode\label{lem:normalcount686}
Fix \(C\) and \(c\). Let \(\mathfrak a\) be the ideal generated
by the coordinates of \(\operatorname{adj}(C)e_n^{\rm std}\).
For an integral term on the \(p_i\)-slice,
\(q_\gamma\in p_i\mathfrak a Q_c\). More generally, if the
normals being counted belong to \(\mathfrak b Q_c\) for a
nonzero integral ideal \(\mathfrak b\), their number, when
nonzero, is at most
\[
      9^{kf}\big(D/{\rm N}_{E/\mathbb Q}\mathfrak b\big)^k .
\]
On the \(p_i\)-slice one may take
\({\rm N}_{E/\mathbb Q}\mathfrak b\ge p_i^{2nf}\).
\end{lemma}
\begin{proof}
For each middle lattice vector \(x+\alpha z\) one has
\(N\alpha\in\mathcal O_E\), by pairing with \(z\).
For two such vectors the difference between \(N h(x,x')\)
and \(N h(x+\alpha z,x'+\alpha' z)\) is
\(\overline{N\alpha}\,N\alpha'\). Hence
\(N h(Q_c,Q_c)\subseteq\mathcal O_E\). In particular if
\(0\ne u\in\mathfrak bQ_c\), then the totally positive
\(Nh(u,u)\) is in the \(F\)-ideal
\({\rm N}_{E/F}\mathfrak b\). Indeed
\(\overline{\mathfrak b}\mathfrak b\) is the extension of this
ideal, as one checks at split, inert and ramified primes.
Thus \({\rm N}_{F/\mathbb Q}(Nh(u,u))\ge
{\rm N}_{E/\mathbb Q}\mathfrak b\).

Each projected column of \(\gamma\mathcal F\) is in \(Q_c\),
and is in \(p_iQ_c\) at the places above \(p_i\) on the indicated
slice, by the frame-span preservation there. Hence
\(q_\gamma\in p_i\mathfrak a Q_c\). The first \(a\) rows of \(C\)
vanish modulo every component at \(p_i\); their \(a\)-minors
generate \(\mathfrak a\). Thus
\(p_i\mathfrak a\subseteq p_i^n\mathcal O_E\).

Set \(R=D^{1/(2f)}\) and
\(m=({\rm N}_{E/\mathbb Q}\mathfrak b)^{1/(2f)}\).
At embedding \(v\), use the Euclidean norm from \(Nh\)
and rescale it by \(R/\sqrt{\Phi_v(\Delta)}\). The scale
factors have product one. By the first paragraph, distinct
points of \(\mathfrak bQ_c\) are at distance at least \(m\)
in the maximum of the rescaled local norms: their squared
norms have product at least \(m^{2f}\). Each of the proposed
normals has local radius exactly \(R\) at every place. If a
normal exists then \(R\ge m\), by comparison with zero.
Disjoint open balls of radius \(m/2\) around these normals
lie in the ball of radius \(R+m/2\). Volume in this real
\(2kf\)-dimensional normed space bounds the number by
\((1+2R/m)^{2kf}\le(3R/m)^{2kf}\), proving the claim.
\end{proof}

For fixed archimedean vectors write \(t=10^5n^4\),
\(\rho=b^{37n}\), \(A_0=\prod_{v,j}b_{j,v}\) with the scales of
Corollary~\ref{cor:realmarked682}, and
\[
  B_x=\rho^{9nft} A_0^{-2(t+k-1)}
             \prod_v\frac{\|W_v\|_2^4}{d_{\pi_v}c_G}.
\]
By that corollary the square of the first term on the right
of Theorem~\ref{thm:realpair685} is exactly
\(B_xD^{-(t+k-1)}\).

\begin{theorem}[A sum over the first normal on either marked slice]
\leavevmode\label{thm:firstnormalsum686}
Fix \(c,c'\) and a marked slice \(p_i\).
Sum \(B_xD^{-(t+k-1)}\) once for each distinct pair
\((C,q_\gamma)\) from a near integral open-cell term on this
slice. With \(\rho=b^{37n}\), the sum is at most
\[
       B_x(1+2N\rho^2)^{2fn^2}\,9^{kf}\,p_i^{-2nfk}.
\]
If the entire residue matrix is fixed at both levels, the first
parenthesis can be replaced by \(1+2N\rho^2/M\).
For sufficiently large nodes of the shorter construction the
first bound is at most \(B_x b^{-18nkf}\).
\end{theorem}
\begin{proof}
A near term has normalized Gram entries of modulus at most
\(\rho^2\). Before the rescaling at \(v\), the \((j,l)\)-entry
therefore has modulus at most
\[
      R_{jl,v}=N\rho^2\kappa_v^{-(2n-j-l)} .
\]
The geometric mean over \(v\) is at most \(N\rho^2\): the
different generator has absolute norm \(|{\rm disc}(E)|\),
so \((\prod_v\kappa_v)^{1/f}
=M|{\rm disc}(E)|^{1/(2f)}/(2\pi)\ge1\).
To count possible integral entries of \(C\) rescale each
complex plane by \((\prod_v R_{jl,v})^{1/f}/R_{jl,v}\).
The factors have product one. Distinct algebraic integers
are separated by at least one in the rescaled maximum norm,
since the absolute field norm of their difference is a
nonzero integer. The same disjoint-ball argument as above
bounds the count by \((1+2N\rho^2)^{2f}\).
When the residue is fixed modulo both levels, all possible
entries in the same position are congruent modulo \(M\):
in class coordinates they are the entries of
\(\mathcal F^*Jg\mathcal F\) in the same integral Gram model.
The separation is then at least \(M\). Multiply these bounds
over the \(n^2\) positions.

For any chosen \(C\), Lemma~\ref{lem:normalcount686} counts
its first normals by \(9^{kf}D^k p_i^{-2nfk}\).
Multiplying by the summand leaves \(D^{-(t-1)}\le1\).
This proves both bounds. For the last estimate the separate
prime is chosen with \(N\le b^{60n}\), the level satisfies
\(p_i\ge b^{10}\), and \(b\ge9\). Thus the Gram count costs
at most \(b^{270fn^3}\). The family has
\(n=m/2\le\sqrt x\) and \(k=r-2n+1\ge x^5-2\sqrt x\), so
\(k\ge270n^2\) eventually. The Gram count and \(9^{kf}\)
then cost at most \(b^{2nkf}\), giving the assertion.
\end{proof}

For clarity, the finer rational orbit data can be described
without counting them. On a fixed \(C\), retain the second
normal \(q_{\gamma^{-1}}\) and the isometry between the two
normal hyperplanes induced by \(Y_\gamma\): send the positive
projection of \(t\in Y_\gamma\) to that of \(\gamma^{-1}t\).
These two projections are onto their respective hyperplanes
by Proposition~\ref{prop:rationalorbits686}. If both actions
are restricted to \(U_0(F)\), the remaining data are exactly
the two normals and this isometry. Indeed the stabilizer in
the two copies of \(\mathrm U(X_0)\) of the two normals and
the isometry is the graph of \(\mathrm U(Y_\gamma)\); each of
its elements lifts to the unitary on the full space acting
as the identity on the joint frame span. This is exactly the
image of the stabilizer of \(\gamma\). Passing further to
principal unipotent congruence groups can split these orbits
again. Theorem~\ref{thm:firstnormalsum686} counts its specified
pairs once; the second moment weights those finer lifts.


\subsection{The discriminant of two integral terms with a common link}
\label{sec:linkdisc687}

We prove the arithmetic discriminant gain for two open-cell
terms in the same level position whose positive links agree.
Keep the original class lattices \(\Lambda=\Lambda_c\) and
\(\Lambda'=\Lambda_{c'}\) in the indefinite space. Write
\(R=\mathcal O_E\), \(\Lambda_X=\Lambda\cap X_0\) and
\(\Lambda_L=\Lambda\cap L=\bigoplus_{i=1}^a R e_i\).
Use primes on subscripts for the corresponding sublattices of
\(\Lambda'\). A term \(\gamma\) carries \(\Lambda'\) isometrically
to \(\Lambda\). Let \(H_\gamma=q_\gamma^\perp\subset X_0\)
and \(j_\gamma:H_\gamma\to X_0\) be the isometry induced by the
two positive projections of \(Y_\gamma=(T+\gamma T)^\perp\)
in Proposition~\ref{prop:rationalorbits686}.

For \(x\in H_\gamma\) its unique lift to \(Y_\gamma\) has the
form \(t_\gamma(x)=x+l_\gamma(x)\) with \(l_\gamma(x)\in L\).
Its inverse image under \(\gamma\) has positive component
\(j_\gamma x\) and isotropic component in \(L\).
The integral lattice of the term is
\[
 L_\gamma=\{x\in\Lambda_X\cap H_\gamma:
                         l_\gamma(x)\in M\Lambda_L\}.
\]
For a positive Hermitian lattice \(A\) of rank \(e\) define
\(d(A)\) as its real covolume at the complex embeddings, divided
by \((2^{-f}\sqrt{|{\rm disc}(E)|})^e\). Use Euclidean volumes
for its positive form at each embedding. This gives \(d(R)=1\)
for the usual form, and passage to an index \(i\) sublattice of
full rank multiplies \(d\) by \(i\). Equivalently, in any
rational basis with Gram \(G_0\) and determinant ideal
\(\mathfrak d\) of the lattice coordinates, one has
\(d(A)={\rm N}_{E/\mathbb Q}(\mathfrak d)
       {\rm N}_{F/\mathbb Q}(\det G_0)\).
Indeed this holds in the basis lattice by the real Jacobian
at the complex embeddings, and the lattice index gives the
general formula after clearing denominators. It agrees with
the square root of the index of the local Gram determinant.

\begin{theorem}[An exact discriminant formula for the level lift lattice]
\leavevmode\label{thm:termdisc687}
For an integral term with \(\Delta_\gamma\ne0\), define the
finite index
\[
 I_\gamma=[R^n:\{(h(\gamma f_j,y))_{j=1}^n:
                         y\in M\Lambda_L+\Lambda_X\}].
\]
Then the lattice \(L_\gamma\) has rank \(k-1\) and
\[
                  d(L_\gamma)
              =\frac{M^{2af}\,{\rm N}_{F/\mathbb Q}(\Delta_\gamma)}
                       {I_\gamma}.
\]
\end{theorem}
\begin{proof}
Let \(P\) here denote the \(n\)-by-\(a\) matrix
\(P_{ji}=h(\gamma f_j,e_i)\), and put
\(B x=(h(\gamma f_j,x))_j\). All these values are integral on
the lattices in question. The matrix \(P\) is the first \(a\)
columns of \(C^*\) and has rank \(a\), since \(C\) is
invertible. Set \(u=\operatorname{adj}(C)e_n^{\rm std}\) and
let \(\mathfrak a=(u_1,\ldots,u_n)\). The functional
\(u^*:R^n\to\overline{\mathfrak a}\) is onto. Its kernel is
\(Y=R^n\cap\operatorname{span}_E(P)\), since \(Cu=D_n e_n^{\rm std}\).
The \(a\)-minors of \(P\) generate \(\overline{\mathfrak a}\).
Smith normal form at each prime gives
\([Y:PR^a]={\rm N}_{E/\mathbb Q}\mathfrak a\).
Moreover \(u^*B x=h(q_\gamma,x)\) by the construction of the
normal. Write \(J=h(q_\gamma,\Lambda_X)\), a nonzero integral
ideal contained in \(\overline{\mathfrak a}\).

Set \(A=MPR^a+B\Lambda_X\). It has rank \(n\) because
\(u^*A=J\ne0\). For \(x\in\Lambda_X\cap H_\gamma\) one has
\(B x\in Y\), and \(l_\gamma(x)\) has coordinates solving
\(P l=-Bx\). Therefore \(x\mapsto Bx\bmod MPR^a\) identifies
\((\Lambda_X\cap H_\gamma)/L_\gamma\) with
\((A\cap Y)/MPR^a\). Surjectivity follows because the
\(B\)-part of any element in \(A\cap Y\) is annihilated by
\(u^*\). The exact sequences for \(u^*\) give
\[
 \begin{split}
 [R^n:A]&=[Y:Y\cap A]\,[\overline{\mathfrak a}:J],\\
 [\Lambda_X\cap H_\gamma:L_\gamma]
     &=\frac{M^{2af}\,{\rm N}_{E/\mathbb Q}J}{I_\gamma}.
 \end{split}
\]
This also proves the asserted rank.

It remains to compute the discriminant of the saturated
hyperplane lattice. Project \(\Lambda_X\) orthogonally to
\(Eq_\gamma\). At embedding \(v\), in a unit vector coordinate
on this line, its projected coordinates are the embeddings
of \(J\) divided by \(\sqrt{\Phi_v h(q_\gamma,q_\gamma)}\).
Covolume in an exact sequence of real Euclidean lattices is
the product of the kernel covolume and the projected
covolume; this follows by lifting a basis of the image and
wedging with a basis of the kernel. Hence
\[
 d(\Lambda_X\cap H_\gamma)
     =d(\Lambda_X)\,
       \frac{{\rm N}_{F/\mathbb Q} h(q_\gamma,q_\gamma)}
            {{\rm N}_{E/\mathbb Q}J}.
\]
Finally \(d(\Lambda_X)=N^f\) in this construction. Before
the class change all middle basis vectors except the distinguished
one are ordinary. For that coordinate the integral middle
lattice uses \(x_*/\bar\pi+z/N,\ z\), where the finite idele
\(\pi\bar\pi=N\). A combination has zero \(z\)-part precisely
when the coefficient of the first vector is divisible by \(N\),
so the pure positive coordinate lattice there is \(\pi\widehat R\).
Its ideal norm is \(N^f\) by \(\pi\bar\pi=N\).
A finite unitary class change in \(X_0\) has determinant of
absolute ideal norm one, since its conjugate times itself is
one. Thus it preserves this discriminant. Inserting
\(h(q_\gamma,q_\gamma)=\Delta_\gamma/N\) cancels the factors of
\(N\). Multiplying the last two index formulas finishes the proof.
\end{proof}

Let
\[
   t_\gamma=[\,\Lambda\cap(L+\gamma L):
                         \Lambda_L+\gamma\Lambda'_L\,].
\]
For a fixed residue class write \(z(s)\) for the product of the
norms of all marked components where the two reduced full frame
spans meet. If character agreement has only been checked above
one rational mark, one may form the product just over the
components there. In particular on the \(p_i\)-Bessel slice
the factor \(z_i=p_i^{2f}\) is always available.

\begin{lemma}[A surviving factor after both isotropic arms are removed]
\leavevmode\label{lem:augmentedindex687}
For an agreeing finite residue class, \(t_\gamma z(s)\) divides
\(I_\gamma\). On a \(p_i\)-Bessel slice \(t_\gamma z_i\) divides
\(I_\gamma\) even without the agreement at the other mark.
\end{lemma}
\begin{proof}
The lattice \(\Lambda\) is Hermitian unimodular. Its hyperbolic
blocks are unimodular, the changed two-dimensional middle
Gram has determinant \(-1\), and the remaining block is the
identity. The frame lattice \(\Lambda\cap T\) is generated by
the columns of \(\mathcal F\). Its orthogonal lattice is
\(K=\Lambda\cap T^\perp=\Lambda_L+\Lambda_X\).
By the perfect pairing and saturation of \(K\), the dual of
\(K\) identifies with \(\Lambda/(\Lambda\cap T)\) via the
Hermitian pairing. This is a local statement about free
modules at a prime and its conjugate: an integral linear
functional on a saturated direct summand extends over the
whole ambient lattice. Its kernel in the pairing is the
annihilator of \(K\), whose rational span is \(T\).

Enlarging \(M\Lambda_L+\Lambda_X\) to \(K\) gives an index
dividing \(I_\gamma\). In the indicated perfect dual bases
the evaluation map on \(K\) is the conjugate transpose of
the columns of \(\gamma\mathcal F\) in
\(\Lambda/(\Lambda\cap T)\). The index of its image in \(R^n\)
is therefore the product of residue norms to the index lengths
\(\ell_w\) of the full joint frame columns at all finite primes.
Indeed append the primitive first frame in adapted lattice
coordinates; saturation of the combined columns is saturation
of the projected columns in the quotient. Hermitian transpose
exchanges conjugate places, which preserves the product of
norms. Likewise \(t_\gamma\) is the product using the lengths
\(\ell_{0,w}\) of the two isotropic arms. At every place
\(\ell_w-\ell_{0,w}\ge0\), by wedging in the two terminal
columns after saturation as in Theorem~\ref{thm:terminalmeet681};
this part of that argument works over any local discrete
valuation ring. At each agreeing meeting component the
difference is at least one by that theorem. On a Bessel slice
the terminal congruence \(g z\equiv z\bmod L\) in both
components gives the required dependence after saturation
directly. Multiplication yields the divisibility assertions.
\end{proof}

\begin{theorem}[The common-link discriminant gain for integral pairs]
\leavevmode\label{thm:samelink687}
Suppose \(\gamma,\eta\) are open-cell terms in the same integral
class position and with the same residue at both marked levels.
Suppose their positive link graphs agree: \(H_\gamma=H_\eta\)
and \(j_\gamma=j_\eta\) there. Put \(L_{12}=L_\gamma\cap L_\eta\).
Then
\[
 \big(z\,d(L_{12})\big)^2
    \le M^{4af}\,
       {\rm N}_{F/\mathbb Q}(\Delta_\gamma)\,
       {\rm N}_{F/\mathbb Q}(\Delta_\eta)
\]
for \(z=z(s)\) under finite character agreement, or \(z=z_i\)
on a \(p_i\)-Bessel slice. In particular the bound applies to
open terms with the same full link: imposing zero \(z\)-coordinate
in both middle components of the full link gives precisely
the positive graph above.
\end{theorem}
\begin{proof}
We first show \([L_\gamma:L_{12}]\mid t_\eta\). For
\(x\in L_\gamma\), put \(x'=j_\gamma x=j_\eta x\), and write
\[
 t_\gamma=x+l_\gamma,\quad \gamma^{-1}t_\gamma=x'+l'_\gamma,
 \qquad
 t_\eta=x+l_\eta,\quad \eta^{-1}t_\eta=x'+l'_\eta .
\]
The subscripts refer to the unique rational lifts. The given
finite residue equality implies
\((1-\eta\gamma^{-1})\Lambda\subset M\Lambda\):
at all primes over \(M\) it is the congruence in the common
class basis, and at the other primes \(M\) is a unit. Hence
\[
  u(x)=\frac{t_\gamma-\eta\gamma^{-1}t_\gamma}{M}
       =\frac{l_\gamma-l_\eta}{M}
          -\eta\,\frac{l'_\gamma-l'_\eta}{M}
              \ \in\ \Lambda\cap(L+\eta L).
\]
The two summands in \(L+\eta L\) are uniquely determined, since
\(T\cap\eta T=0\). The class of \(u(x)\) modulo
\(\Lambda_L+\eta\Lambda'_L\) vanishes precisely when
\(l_\eta\in M\Lambda_L\). In one direction use uniqueness and
\(l_\gamma\in M\Lambda_L\); in the other, after removing the
integral first summand from \(u(x)\), the remaining integral
vector belongs to \(\Lambda\cap\eta L=\eta\Lambda'_L\).
This yields the required injection of the quotient and the
divisibility. Interchanging the terms gives
\([L_\eta:L_{12}]\mid t_\gamma\).

By Theorem~\ref{thm:termdisc687} the two resulting bounds are
\[
 \begin{split}
 d(L_{12})&\le t_\eta\,M^{2af}
                   {\rm N}_{F/\mathbb Q}(\Delta_\gamma)/I_\gamma,\\
 d(L_{12})&\le t_\gamma\,M^{2af}
                   {\rm N}_{F/\mathbb Q}(\Delta_\eta)/I_\eta .
 \end{split}
\]
Multiply and use Lemma~\ref{lem:augmentedindex687} on each term.
The two isotropic torsion indices cancel, leaving \(z^{-2}\).
\end{proof}

This comparison uses the same lattices as the open-cell
integral and the same two levels. When the two positive graphs
differ, their common lattice has smaller rank. Bounding and
counting those pairs, as well as the terms of zero frame
discriminant, is still required for the full second moment.


\subsection{The discriminant gain when two open-cell links differ}
\label{sec:differentlinks688}

We now allow distinct positive link graphs. We work first with all
middle coordinates and impose the two terminal equations only after
comparing the links. Write \(R=\mathcal O_E\), use the discriminant
\(d\) of Section~\ref{sec:linkdisc687}, and put \(A=\mathcal F_{[1,a]}\),
the first \(a\) frame columns. For an open-cell term \(\gamma\) set
\[
 P_\gamma=(A,\gamma A),\quad Z_\gamma=(z,\gamma z),\quad
 V_\gamma=(L+\gamma L)^\perp .
\]
The columns of \(P_\gamma\) are independent and
\(\dim_E V_\gamma=k+1\). Every \(t\in V_\gamma\) has coordinates
\(t=l+x+u z\) and \(\gamma^{-1}t=l'+x'+u'z\), with
\(l,l'\in L\), \(x,x'\in X_0\). On \(D=X_0\oplus E\) use the
positive form \(h_D((x,c),(y,d))=h(x,y)+N\bar c d\).
Then \((x,u')\) and \((x',u)\) have equal norm, by unitarity of
\(\gamma\). The first coordinate pair determines \(t\):
if it is zero then the other pair vanishes by positivity and
\(t\in L\cap\gamma L=0\). Dimension now shows that the pairs
are related by a unitary map
\(J_\gamma:D\longrightarrow D\), sending \((x,u')\) to \((x',u)\).
Its graph is the full link; its restriction with both last
coordinates zero is the positive graph of
Proposition~\ref{prop:rationalorbits686}.

Take \(\Lambda=\Lambda_c\), \(\Lambda'=\Lambda_{c'}\) as before.
Let \(K_\gamma^0=\Lambda\cap V_\gamma\) and let \(K_\gamma^+\)
be the sublattice imposing \(l\in M\Lambda_L\). We identify both
lattices with their first pair of coordinates in \(D\). Let
\(\ell_\gamma(t)=(h(z,t),h(\gamma z,t))=(-Nu,-Nu')\), and set
\[
 Q_\gamma=|{\rm N}_{E/\mathbb Q} D_n(\gamma)|,\quad
 D_\gamma={\rm N}_{F/\mathbb Q}(\Delta_\gamma),\qquad
 G_\gamma=M^{2af}\max\{Q_\gamma/N^f,D_\gamma\}.
\]

\begin{lemma}[The full link and its two terminal linear forms]
\leavevmode\label{lem:fullborder688}
For each open integral term,
\[
 d(K_\gamma^0)=\frac{Q_\gamma}{N^{2f}t_\gamma},\qquad
 d(K_\gamma^+)\le \frac{M^{2af} Q_\gamma}{N^{2f} t_\gamma}.
\]
At an embedding of \(F\) in \(\mathbb R\), the product of the
squares of the two nonzero singular values of \(\ell_\gamma\),
using \(h_D\) and the usual metric on \(\mathbb C^2\), is
\(N^2\Delta_\gamma/|D_n(\gamma)|^2\) at that embedding.
Each one of the two component functionals has norm \(\sqrt N\).
\end{lemma}
\begin{proof}
Write \(H\) for the ambient Hermitian matrix in rational
coordinates, \(P=P_\gamma\), \(Z=Z_\gamma\), \(E_0=P^*H\) and
\(E_1=Z^*H\). On \(\ker E_0=V_\gamma\) the positive graph metric
is exactly the restriction of
\[
                  H'=H+N^{-1}E_1^*E_1 .
\]
Indeed the indefinite norm is \(h(x,x)-N|u|^2\) and the added
term is \(N(|u|^2+|u'|^2)\). At an archimedean embedding form
\(B_0=\left(\begin{smallmatrix}H'&E_0^*\\E_0&0\end{smallmatrix}\right)\).
Its absolute determinant is \(|\det H|\,|D_n|^2/N^2\). To see
this, start with the Hermitian block matrix with first row of
blocks \(H,E_0^*,E_1^*\), the two lower diagonal blocks \(0\)
and \(-NI_2\), and zero between the two lower blocks. Eliminating
\(-NI_2\) gives \(N^2\det B_0\). Eliminating \(H\) gives, up to
sign, \(\det H\) times the determinant of the joint frame Gram
with \(N\) added at the two terminal diagonal entries. Its two
diagonal frame blocks are now zero and its off-diagonal blocks
are \(C,C^*\). Its absolute determinant is \(|D_n|^2\).

Here is the lattice consequence of this bordered calculation.
The map \(E_0\) sends \(\Lambda\) to a full lattice in \(R^{2a}\)
of index \(t_\gamma\), since \(\Lambda\) is Hermitian unimodular
and evaluation on the columns of \(P\) has the elementary
divisors of those columns at the conjugate primes. Choose a
rational basis of \(\ker E_0\) and then lifts of the standard
basis of the target. Call the resulting basis matrix \(T_0\).
In this basis \(E_0T_0=(0,I_{2a})\); expansion along the
border gives
\(\det(H'|_{\ker E_0})=|\det T_0|^2|\det B_0|\), where the
left side denotes the determinant in the chosen kernel basis.
In these coordinates the determinant ideal of \(\Lambda\) is
the product of those of \(K_\gamma^0\) and \(E_0\Lambda\),
by the short exact sequence, or by choosing local bases.
The basis change contributes the inverse ideal of \(\det T_0\).
The formula for \(d\) in the preceding subsection therefore
gives \(d(K_\gamma^0)=Q_\gamma/(N^{2f}t_\gamma)\):
the factor consisting of the norm of the original determinant
ideal of \(\Lambda\) times \(|{\rm N}_F\det H|\) is one by
unimodularity. Imposing the congruence on \(l\) has index at
most \(|\Lambda_L/M\Lambda_L|=M^{2af}\).

For the singular values border \(B_0\) a second time by \(E_1\)
on the original ambient coordinates and zero elsewhere; give
this new border zero diagonal block. Its absolute determinant
is \(|\det H|\Delta_\gamma\), because the term
\(N^{-1}E_1^*E_1\) can be removed by row operations, and the
remaining determinant after eliminating \(H\) is the joint
Gram determinant of the two frames. Divide by \(|\det B_0|\).
In coordinates \(E_0T_0=(0,I)\) the upper ambient part of
\(B_0^{-1}\) is the inverse kernel metric on the kernel and
zero on the remaining coordinates, as is seen by solving its
bordered linear equations. Thus this ratio is the determinant
of the Gram of the two functionals on \(V_\gamma\) in the dual
positive metric, proving the formula. Finally each functional
is \(N\) times the last coordinate of a unitary identification
with \(D\), whose last coordinate line has squared length \(N\).
\end{proof}

\begin{proposition}[Congruence on the quotient by the common graph]
\leavevmode\label{prop:crosslink688}
For open terms \(\gamma,\eta\) in one class position with the
same residue at both levels, put \(I=\ker(J_\gamma-J_\eta)\)
in \(D\), \(c=\operatorname{codim}_D I\) and
\(T_i=K_i^+\cap I\) for \(i=\gamma,\eta\). Then
\[
       d(T_\gamma)d(T_\eta)
           \le (2/M)^{2cf} d(K_\gamma^+)d(K_\eta^+).
\]
\end{proposition}
\begin{proof}
On the two source spaces consider the cross pairing
\[
 B(u,v)=h_D(u,v)-h_D(J_\gamma u,J_\eta v).
\]
Its left and right kernels are both \(I\), since the two maps
are unitary and defined on all of \(D\). Its operator norm is
at most two. On \(K_\gamma^+\times K_\eta^+\) it is integral
and divisible by \(M\). Indeed if \(t_1,t_2\in\Lambda\) are
the two lifts, expansion of the swapped terminal coordinates
gives
\(B(u,v)=h(t_1,t_2)-h(\gamma^{-1}t_1,\eta^{-1}t_2)\).
Now \(\gamma^{-1}t_1,\gamma^{-1}t_2\in\Lambda'\), and
\((\gamma^{-1}-\eta^{-1})t_2\in M\Lambda'\) by the shared
residue. This proves the congruence.

Let \(B_i\) denote the orthogonal projection of \(K_i^+\) to
\(I^\perp\) for \(h_D\). Rationality of the projection makes
these full lattices in the \(c\)-dimensional quotient. There
\(B\) is nondegenerate and \(B(B_\gamma,B_\eta)\subset MR\).
Taking the determinant of the pairing gives
\[
                      M^{2cf}\le 2^{2cf}d(B_\gamma)d(B_\eta).
\]
Explicitly, in rational bases the determinant, times the
conjugate determinant ideal of the first lattice and the
determinant ideal of the second, generates an ideal contained
in \(M^cR\), by taking exterior powers at every finite prime.
At an infinite place the absolute determinant of the pairing is
at most \(2^c\) times the square root of the product of the
two positive Gram determinants. Take the absolute field norm
and use the formula for \(d\). Finally
\(d(K_i^+)=d(T_i)d(B_i)\) by orthogonal projection. The empty
determinant is one if \(c=0\).
\end{proof}

\begin{theorem}[The discriminant of any pair of open terms in one block]
\leavevmode\label{thm:openpair688}
In the setting of Proposition~\ref{prop:crosslink688}, let
\(I_X=\{u\in I:u\in X_0,\ J_\gamma u\in X_0\}\).
The common integral lattice of the two positive links is
\(L_{12}=I_X\cap L_\gamma\cap L_\eta\). Then
\[
                 d(L_{12})^2\le(2/M)^{2cf}G_\gamma G_\eta .
\]
For \(c=0\) the finer bound of Theorem~\ref{thm:samelink687}
also applies with its meeting-place factor. All common
dimensions in this assertion are allowed, including zero.
\end{theorem}
\begin{proof}
The two terminal maps \(\ell_i\) coincide on \(I\), since the
input and output coordinates there agree. Their common kernel
in \(I\) is \(I_X\). Write \(T_{i,X}=T_i\cap I_X\) and \(s\)
for the rank of this terminal map on \(I\). It has integral
values on \(T_i\) by its definition on the integral lifts.
We give bounds \(d(T_{i,X})\le A_i d(T_i)\).

If \(s=0\), take \(A_i=1\). If \(s=1\), choose a nonzero one
of the two component functionals on \(I\), so its kernel there
is \(I_X\). Its image lattice is a nonzero ideal inside \(R\).
Its real Jacobian on the orthogonal line to its kernel at each
place is at most \(N\) by Lemma~\ref{lem:fullborder688}. Thus
one can take \(A_i=N^f\). If \(s=2\), use the two components
together. Their image lattice in \(R^2\) has discriminant at
least one for the standard positive metric. The squared
exterior norm on a subspace is bounded by that on the entire
space, and Lemma~\ref{lem:fullborder688} gives
\(A_i=N^{2f}D_i/Q_i\).
In both positive ranks the covolume step is the same:
the discriminant of the projected lattice orthogonal to the
kernel is the image discriminant divided by the product of
the real Jacobians on that complement. Divide \(d(T_i)\) by
the projected discriminant to get the kernel discriminant.
Together with Lemma~\ref{lem:fullborder688} all three cases give
\(A_i d(K_i^+)\le G_i/t_i\).

The lattice \(T_{i,X}\) is \(L_i\cap I_X\) in the positive
coordinates: both terminal components vanish and its lift is
integral with the indicated level condition on the isotropic
part. We claim \([T_{\gamma,X}:L_{12}]\mid t_\eta\), and the
reverse bound with the terms exchanged. For \(u\in T_{\gamma,X}\)
the positive images \(j_\gamma u,j_\eta u\) coincide. Therefore
the class of
\((1-\eta\gamma^{-1})t_\gamma(u)/M\) in
\((\Lambda\cap(L+\eta L))/(\Lambda_L+\eta\Lambda'_L)\)
is defined, with kernel exactly \(L_{12}\) by the decomposition
in the proof of Theorem~\ref{thm:samelink687}. This argument uses
agreement of images at \(u\), which holds for every \(u\in I_X\).
Multiplying the resulting index inequalities and applying
Proposition~\ref{prop:crosslink688} gives
\[
 d(L_{12})^2\le t_\gamma t_\eta A_\gamma A_\eta
             (2/M)^{2cf}d(K_\gamma^+)d(K_\eta^+)
          \le (2/M)^{2cf}G_\gamma G_\eta .
\]
At \(c=0\) the two full unitary maps are equal and their
positive graphs agree.
\end{proof}

The factors in this estimate come from congruence on the actual
two class lattices. It covers each pair with nonzero frame
discriminants. To apply it in a second moment one still needs
the multiplicities of the integral pairs and estimates for the
other cell ranks.


\subsection{The pair discriminant at every rational cell rank}
\label{sec:allpair689}

Keep the Hermitian space, the two class lattices and the level
\(M=p_0p_1\) of the preceding subsection. We treat the terms
for which the two frames have a common vector. Their rational
lifts can have a common isotropic summand. We first remove
that summand over \(R=\mathcal O_E\).

Suppose that the characters agree on the rational stabilizer.
For \(\gamma\) with \(\Delta_\gamma=0\), let \(j\) be its rank in
Theorem~\ref{thm:rationalflags680}, put \(U=L\cap\gamma L\), and
write \(U_R=U\cap\Lambda\). Let \(A_j=(e_1,\ldots,e_j)\). A bar on a
vector in \(U^\perp\) denotes its image in \(U^\perp/U\).
Define
\[
 \begin{split}
 t_\gamma&=[\,\Lambda\cap(L+\gamma L):
                         \Lambda_L+\gamma\Lambda'_L\,],\\
 i_\gamma&=[\,\Lambda_L/U_R:\bar A_j R^j\,],\qquad
 i'_\gamma=[\,\gamma\Lambda'_L/U_R:\overline{\gamma A_j}R^j\,].
 \end{split}
\]
The columns in each denominator span its numerator over \(E\).
In fact the common flag from Theorem~\ref{thm:rationalflags680}
gives \(L=\langle A_j\rangle\oplus U\) and
\(\gamma L=\langle\gamma A_j\rangle\oplus U\).
The index \(i_\gamma\) can equally be computed from the first
\(j\) isotropic rows of the frame product matrix inside the
lattice generated by all its isotropic rows. Indeed evaluation of
\(l\in\Lambda_L\) against the second frame annihilates \(U_R\),
because \(U\subset\gamma L\). Its restriction to \(A_j\) has
rank \(j\) by \(D_j\ne0\), so its kernel is exactly \(U_R\).
Evaluation against the first frame proves the corresponding
description of \(i'_\gamma\). All indices here are absolute
orders of finite quotients.

\begin{lemma}[The two terminal vectors in the isotropic quotient]
\leavevmode\label{lem:quotientframe689}
The lattice \(\bar\Lambda=(\Lambda\cap U^\perp)/U_R\) is Hermitian
unimodular. In \(\bar V=U^\perp/U\) the subspaces \(\bar L\) and
\(\overline{\gamma L}\) pair perfectly over \(E\). Their sum has
an orthogonal complement \(W\) of dimension \(k+1\).
If \(d_*=\det C_{\{1,\ldots,j,n\},\{1,\ldots,j,n\}}\), then
\[
                        d_*=-N\zeta D_j ,
\]
where \(\zeta\) is the scalar for this term in
Theorem~\ref{thm:rationalflags680}. We include \(j=0\) by
\(D_0=1\). The two terminal linear forms
\(t\mapsto h(z,t),h(\gamma z,t)\) on \(W\) have joint rank one.
\end{lemma}
\begin{proof}
The module \(U_R\) is saturated in \(\Lambda\). Locally it is
a direct summand. Since \(\Lambda\) is self-dual, evaluation
on \(U_R\) is onto its conjugate dual. Choose a local basis
of \(U_R\) and dual lifts. Their joint Gram matrix has blocks
\(\left(\begin{smallmatrix}0&I\\ I&B\end{smallmatrix}\right)\),
whose determinant is a unit. The summand they generate splits
off orthogonally by integral projection. Its contribution to
\(U_R^\perp/U_R\) is zero, and the remaining summand is
self-dual. At a split place take bases and dual lifts in the
paired components of the local ring \(R\otimes_{\mathcal O_F}
\mathcal O_{F_v}\). The same block matrix is invertible there.
This argument applies also above two. Local self-duality gives
the assertion about \(\bar\Lambda\).

In the displayed bases of the two arms the cross pairing has
determinant \(D_j\ne0\). The arms have zero intersection and
are each isotropic. Together they span a nondegenerate
\(2j\)-space in the \((2j+k+1)\)-space \(\bar V\).
Both terminal vectors are orthogonal to \(U\). Subtract their
orthogonal projections onto the sum of the two arms.
For \(\bar z\) the subtracted vector belongs to \(\bar L\),
since \(\bar z\perp\bar L\) and the two arms pair perfectly.
For \(\overline{\gamma z}\) it belongs to \(\overline{\gamma L}\).
The resulting vectors \(z_0,z_1\in W\) both have norm \(-N\).
The common vector at the terminal level of the two flags gives
\(\gamma z-\zeta z\in L+\gamma L\), whence \(z_1=\zeta z_0\).
Changing each terminal vector by its indicated arm is a row
or column operation of determinant one on the subframe matrix.
Its remaining terminal entry is \(h(z_0,z_1)=-N\zeta\).
This proves the formula and shows that the two linear forms
on \(W\) are proportional and nonzero.
\end{proof}

For any rank use the positive space \(D=X_0\oplus E\), with
squared norm \(h(x,x)+N\bar u u\). Representatives of a class
of \(W\) may be taken in \((L+\gamma L)^\perp\). As in the open
case write
\[
 t=l+x+u z,\qquad \gamma^{-1}t=l'+x'+u'z .
\]
Changing \(t\) by a vector of \(U\) leaves both middle components
fixed. The pairs \((x,u')\) and \((x',u)\) have equal positive
norm by unitarity. If one vanishes, both vanish and \(t\in U\).
Since \(\dim_E W=\dim_E D\), the first pair identifies \(W\)
with \(D\), and the second differs from it by a unitary map
\(J_\gamma\). Write \(K_\gamma^0\) for the image in \(D\) of
\(\Lambda\cap(L+\gamma L)^\perp\). Write \(K_\gamma^+\) for the
image of those representatives that satisfy \(l\in M\Lambda_L\).
The lattice \(L_\gamma\) of the positive link is obtained by
imposing \(u=u'=0\) on \(K_\gamma^+\).

\begin{proposition}[The determinant of the integral link in lower rank]
\leavevmode\label{prop:closedborder689}
For \(0\le j\le a\) in the preceding notation,
\[
 d(K_\gamma^0)=
   \frac{|{\rm N}_{E/\mathbb Q}d_*|}
        {N^{2f}t_\gamma i_\gamma i'_\gamma},\qquad
 [K_\gamma^0:K_\gamma^+]\le M^{2jf}.
\]
Each terminal functional has norm \(\sqrt N\) for the positive
metric at every real place.
\end{proposition}
\begin{proof}
In \(\bar V\) take \(P=(\bar A_j,\overline{\gamma A_j})\),
\(Z=(\bar z,\overline{\gamma z})\), and write \(H\) for the
Hermitian matrix. The lattice \(K_\gamma^0\) identifies with
\(\bar\Lambda\cap\ker(P^*H)\). Indeed any integral lift of
such a class is already orthogonal to both arms. The positive
metric on this kernel is the restriction of
\(H+N^{-1}(Z^*H)^*(Z^*H)\), as is seen by substituting the
two expressions for \(t\). The bordered determinant with
border \(P^*H\) therefore has absolute value
\(|\det H|\,|d_*|^2/N^2\) at each embedding. Explicitly,
adjoin the two rows \(Z^*H\) with diagonal block \(-NI_2\).
Eliminating that block contributes \(N^2\). Eliminating \(H\)
instead gives \(\det H\) times, up to sign, the Gram matrix
of the two subframes with both terminal diagonal entries
changed to zero. The latter determinant has absolute value
\(|d_*|^2\).

The image of \(P^*H\) on \(\bar\Lambda\) has index
\(t_\gamma i_\gamma i'_\gamma\) in \(R^{2j}\).
By self-duality, local elementary divisors of the evaluation
map are those of the columns of \(P\) at the conjugate prime.
Their product of absolute norms is the index of \(PR^{2j}\)
in its saturation. Saturating the sum of the two full arm
lattices costs precisely \(t_\gamma\), since \(U_R\) lies in
both arms and a lift of a class in their rational sum belongs
to \(L+\gamma L\). Inside the sum the projected arms are
disjoint. Their separate indices therefore multiply to
\(i_\gamma i'_\gamma\). Apply the basis-and-image calculation
of Lemma~\ref{lem:fullborder688} to this border. The ambient
determinant ideal times the absolute norm of the ambient
Gram determinant is one by unimodularity. This proves the
first identity, including the empty border when \(j=0\).

In the class lattice the top isotropic coordinates are an
integral direct summand, so for an integral representative
the component \(l\) belongs to \(\Lambda_L\). Assign its
class modulo \(M\Lambda_L+U_R\) to the element of \(K_\gamma^0\).
This is independent of the representative and has kernel
\(K_\gamma^+\). The target is the reduction modulo \(M\)
of the rank \(j\) projective module \(\Lambda_L/U_R\).
Its order is \(M^{2jf}\), by local freeness. Finally the
functionals have values \(-Nu,-Nu'\). Each is \(N\) times the
last coordinate of a unitary identification with \(D\), so
has squared norm \(N\).
\end{proof}

For an open term write \(j=n=a+1\), \(d_*=D_n\),
\(i_\gamma=i'_\gamma=1\). For a term of any cell rank put
\[
  G_\gamma=M^{2\min(j,a)f}
       \max\left\{
       \frac{|{\rm N}_{E/\mathbb Q}d_*|}
            {N^f i_\gamma i'_\gamma},
                     {\rm N}_{F/\mathbb Q}\Delta_\gamma\right\}.
\]
The second quantity in the maximum vanishes at every lower rank.

\begin{theorem}[The common lattice at the two split levels]
\leavevmode\label{thm:allpair689}
Take two integral terms \(\gamma,\eta\) between the same class
lattices, with the same residue at both levels, each satisfying
the rational frame conclusion of Theorem~\ref{thm:rationalflags680}.
Set \(I=\ker(J_\gamma-J_\eta)\), \(c^+=\operatorname{codim}_D I\)
and
\[
 I_X=\{y\in I:y\in X_0,\ J_\gamma y\in X_0\},\qquad
 L_{12}=I_X\cap L_\gamma\cap L_\eta .
\]
Then \(L_{12}\) has full rank in \(I_X\), and
\[
             d(L_{12})^2\le(2/M)^{2c^+ f}G_\gamma G_\eta .
\]
Under finite character agreement, if \(c^+=0\) and one term
is open, the left side can be multiplied by \(z(s)^2\), the
square of the product of the norms of the meeting places.
In all cases
\[
 k-\dim_E I_X\ \le\ c^++\min\big\{
       {\bf1}_{\Delta_\gamma\ne0},{\bf1}_{\Delta_\eta\ne0}\big\}.
\]
\end{theorem}
\begin{proof}
The comparison of full links in Proposition~\ref{prop:crosslink688}
applies with the quotient representatives just constructed.
Indeed its cross pairing is
\(h(t_1,t_2)-h(\gamma^{-1}t_1,\eta^{-1}t_2)\). Its expression
in middle coordinates is independent of the representatives,
and choosing integral ones puts its values in \(MR\) by the
two residue congruences. On the quotient by \(I\) this pairing
is nondegenerate and has operator norm at most two. Thus for
\(T_i=K_i^+\cap I\) one has the same determinant inequality
\[
 d(T_\gamma)d(T_\eta)\le(2/M)^{2c^+ f}
                                     d(K_\gamma^+)d(K_\eta^+).
\]

On \(I\) the two terms have the same terminal forms. If these
forms have joint rank \(s\), then
\(\dim I_X=k+1-c^+-s\). On all of \(D\) their rank is one
for a lower-rank term by Lemma~\ref{lem:quotientframe689}
and two for an open term by Lemma~\ref{lem:fullborder688}.
This proves the dimension bound. Put \(T_{i,X}=T_i\cap I_X\).
If either term has lower rank, then \(s\le1\). At rank zero
\(T_{i,X}=T_i\). At rank one choose a nonzero component on \(I\).
Its image on \(T_i\) is a nonzero integral ideal by evaluation
on integral lifts. Its real Jacobian on the line transverse to
the kernel is at most \(N\) at every embedding. The orthogonal
projection formula for covolume therefore yields
\(d(T_{i,X})\le N^f d(T_i)\). Multiplying the preceding
inequality by the costs of these two cuts gives
\[
       d(T_{\gamma,X})d(T_{\eta,X})
       \le (2/M)^{2c^+ f}\frac{G_\gamma G_\eta}{t_\gamma t_\eta}.
\]
Here Proposition~\ref{prop:closedborder689} supplies the
factor for each lower-rank term. An open term in a mixed pair
uses the first estimate of Lemma~\ref{lem:fullborder688},
so its minor term in the maximum suffices. If both terms
are open, the same inequality, allowing \(s=2\), was proved
in Theorem~\ref{thm:openpair688} before the final passage
to the common integral lattice.

It remains to justify that passage when the lifts are
nonunique. For \(x\in T_{\gamma,X}\), choose a level lift
\(t_\gamma=x+l_\gamma\), \(l_\gamma\in M\Lambda_L\). Choose
any rational \(\eta\)-lift \(t_\eta=x+l_\eta\) to the same
positive graph point over \(I_X\). The two inverse images
have the same positive coordinate \(x'\), with respective
top components \(l'_\gamma,l'_\eta\in L\). Then
\[
 v=\frac{(1-\eta\gamma^{-1})t_\gamma}{M}
   =\frac{l_\gamma-l_\eta}{M}
      -\eta\frac{l'_\gamma-l'_\eta}{M}
        \in\Lambda\cap(L+\eta L).
\]
The numerator is divisible by \(M\) on \(\Lambda\) because
the two reductions agree. Any two level lifts for \(\gamma\)
differ by a vector in \(M(\Lambda\cap L\cap\gamma L)\).
Consequently the class of \(v\) modulo
\(\Lambda_L+\eta\Lambda'_L\) is well defined and additive.

Suppose it is zero and choose
\(a_0\in\Lambda_L,b_0\in\Lambda'_L\) with \(v=a_0+\eta b_0\).
The difference \(d_0=(l_\gamma-l_\eta)/M-a_0\) lies in
\(L\cap\eta L\) by the displayed expression. Replacing
\(t_\eta\) by \(t_\eta+M d_0\) gives a lift with top
component \(l_\gamma-Ma_0\in M\Lambda_L\). It is integral,
since \(x\in\Lambda\cap X_0\) for the chosen \(\gamma\)-lift.
Conversely, if an \(\eta\)-level lift exists, its use in the
formula makes the first summand integral. The remaining
summand belongs to \(\Lambda\cap\eta L=\eta\Lambda'_L\).
The kernel of the class map is therefore precisely \(L_{12}\).
This proves \([T_{\gamma,X}:L_{12}]\mid t_\eta\), and
interchanging the terms gives \([T_{\eta,X}:L_{12}]\mid
t_\gamma\). These finite indices prove the rank assertion.
Multiplication cancels the two torsion indices and gives
the asserted inequality.

Finally \(c^+=0\) makes the two maps \(J\) equal everywhere.
Their joint terminal ranks coincide, so if either term is
open both are open. Their positive graphs agree, and
Theorem~\ref{thm:samelink687} gives the stated meeting factor.
\end{proof}

The theorem proves the arithmetic pair estimate used in
\path{New/Spanning/RemainingTerms/Gain.lean} on the class lattices
of the shorter construction. Summing the orbital integrals also
requires the number and the analytic weights of pairs with a
given common lattice.


\subsection{The weighted count of lattices in the marked classes}
\label{sec:latticemean690}

The common lattice in Theorem~\ref{thm:allpair689} lies in
\(\Lambda_X=\Lambda_c\cap X_0\). Its discriminant enters the
period estimate with weight \(d(L)^{-k}\). We sum this weight
over every lattice of a prescribed corank in the actual source
classes. The prime \(N\) in this calculation is the norm of the
negative terminal vector with sign reversed. It is distinct from
the two primes dividing \(M\).

Put \(R=\mathcal O_E\), \(f=[E:\mathbb Q]/2\), and
\(\rho_E=|\operatorname{disc}E|^{1/(2f)}\). We first work in an
arbitrary positive Hermitian unimodular lattice \(\mathcal U\)
of rank \(k\) over \(R\). For an \(E\)-subspace \(V\) put
\(d_V=d(\mathcal U\cap V)\), with the normalized covolume of
Section~\ref{sec:linkdisc687}.

\begin{lemma}[The sum over rational subspaces]\leavevmode
\label{lem:spacezeta690}
For \(0\le \ell\le k\) and \(k\ge1\),
\[
 \sum_{\operatorname{codim}_E V=\ell}d_V^{-k-1}
             \le (25 f k\rho_E)^{f(k+1)\ell}.
\]
\end{lemma}
\begin{proof}
For a lattice \(P\) in its rational span, write \(P^\#\) for
its Hermitian dual there. The determinant-ideal formula for
normalized covolume gives \(d(P^\#)=1/d(P)\).
Orthogonal projection of \(\mathcal U\) onto \(V^\perp\)
has dual \(\mathcal U^\#\cap V^\perp=\mathcal U\cap V^\perp\).
Indeed pairing a vector of \(V^\perp\) with a projected
vector gives its original pairing. The projection is rational
and its image is a lattice, so taking its double dual
recovers it. Multiplicativity of covolumes on the kernel
and the orthogonal projection now gives \(d_V=d_{V^\perp}\).
Also \(d_V\ge1\), because the integral lattice in \(V\) is
contained in its dual. For \(\ell=0\) the asserted sum is one.

Use the Euclidean norm
\(\|y\|^2=\sum_{v\mid\infty}h_v(y,y)\) on all the complex
embeddings in the CM type. Put \(P=\mathcal U\cap V^\perp\)
and \(m=2f\ell>0\). Its real covolume is
\((2^{-f}\rho_E^f)^\ell d_V\).
Let \(\lambda_1,\ldots,\lambda_m\) be its Euclidean successive
minima. We record the ball case of Minkowski's inequality:
\(\omega_m\prod_i\lambda_i\le2^m\operatorname{covol}(P)\),
where \(\omega_m\) is the unit ball volume. Choose successively
a shortest lattice vector outside the real span of the ones
already chosen, and orthonormalize their flag. The lengths
at these choices are the successive minima, since \(i\)
independent vectors of length at most \(\lambda_i\) cannot
lie in a space of dimension \(i-1\). In these orthonormal
coordinates dilate axis \(i\) of the unit ball by \(\lambda_i\).
A nonzero lattice vector whose last nonzero coordinate is \(i\)
has length at least \(\lambda_i\), while a vector in the open
ellipsoid with that last coordinate has length strictly less
than \(\lambda_i\). The half-sized open ellipsoid thus has
disjoint translates by \(P\). Comparing volume with a
fundamental domain proves the inequality.\footnote{For
the successive-minima inequality for general symmetric
convex bodies, see R.~P. Bambah, A.~Woods and H.~Zassenhaus,
\href{https://doi.org/10.1017/S1446788700028482}{Journal of the
Australian Mathematical Society \textbf{5} (1965), 453--462}.}

Next choose \(x_1,\ldots,x_\ell\) in \(P\) by successively
minimizing length outside their previous \(E\)-span. The
real span at all embeddings of an \(E\)-space of dimension
\(i-1\) has dimension \(2f(i-1)\). Hence
\(\|x_i\|\le\lambda_{2f(i-1)+1}\). The unit ball contains a
cube of side \(2/\sqrt m\). Multiplying the minima in groups
of \(2f\) therefore gives
\[
          \prod_{i=1}^{\ell}\|x_i\|^{2f}
             \le m^{m/2}(2^{-f}\rho_E^f)^\ell d_V
             =(f\ell\rho_E)^{f\ell}d_V .
\]
The tuple determines \(V\) by taking its orthogonal complement.

For nonzero \(y\in\mathcal U\), the absolute norm to
\(\mathbb Q\) from \(F\) of the totally positive algebraic
integer \(h(y,y)\) is at least one. The arithmetic-geometric
mean inequality gives \(\|y\|^2\ge f\). Distinct lattice
points are thus at least one apart. In dimension \(D=2fk\)
the number in a ball of radius \(t\ge1\) is at most
\((2t+1)^D\), by packing balls of radius \(1/2\).
For \(p=2f(k+1)\), integrate \(p t^{-p-1}\) against the
count of nonzero vectors of norm at most \(t\). This gives
\[
     \sum_{0\ne y\in\mathcal U}\|y\|^{-p}
          \le 3^{2fk}(k+1)\le25^{f(k+1)} .
\]
Choose one of the preceding tuples for each \(V\).
Raising the product inequality to \(k+1\) bounds \(d_V^{-k-1}\)
by \((f\ell\rho_E)^{f(k+1)\ell}\prod_i\|x_i\|^{-p}\).
Sum over all tuples of nonzero lattice vectors and use
\(\ell\le k\).
\end{proof}

At a node in the shorter construction \(N\equiv1\bmod q_1d\),
so \(N\) splits completely in \(E\subseteq\mathbb Q(\mu_{q_1d})\).
At its \(f\) pairs of split places the distinguished middle
coordinate imposes one row congruence at one member of each pair.
More explicitly, in the middle basis \(x_*/\bar\pi+z/N,z\)
of Section~\ref{sec:linkdisc687}, a pure \(X_0\)-vector with
distinguished coordinate \(v_*\) has those two coefficients
\(\bar\pi v_*\) and \(-v_*/\pi\), where \(\pi\bar\pi=N\).
Their integrality is equivalent to \(v_*\in\pi R\) locally.
The other pure coordinates are integral. For a class of a
definite isometry \(g\), this condition is read in \(g^{-1}\)
coordinates.

\begin{proposition}[The average over the terminal rows]
\leavevmode\label{prop:rowmean690}
Let \(\mathcal C\) be the finite set of the definite classes
at the source level. Forget the level of a class and retain
its underlying unimodular lattice \(\mathcal U\), up to
isometry. To average an isometry-invariant sum on the classes
over this \(\mathcal U\), one can average uniformly over
\((N^k-1)^f\) tuples of nonzero residue covectors when the
summand depends only on their common kernel in \(\mathcal U\).
This common kernel is the marked lattice \(\Lambda_X\).
For a fixed rank \(r=k-\ell\) sublattice \(L\) of \(\mathcal U\),
put \(P=\mathcal U\cap E L\), \(i=[P:L]\). The fraction of
row tuples for which \(L\) belongs to the kernel is at most
\[
             \frac{i N^{f\ell}}{(N^k-1)^f}.
\]
\end{proposition}
\begin{proof}
Fix a finite adelic representative of the coarse class and
let \(K\) be its standard integral level in lattice coordinates.
The finer classes over it are the orbits on
\(K/K_{\rm src}\) under the finite automorphism group of
\(\mathcal U\) in these coordinates. This follows by writing
two representatives as \(gk,gk'\) in the double quotient
and cancelling \(g\). Finiteness of the automorphism group
follows from discreteness on a real Euclidean lattice and
compactness of its isometry group. It acts freely on these
cosets. In fact stabilizing a finer coset gives a finite-order
integral matrix congruent to the identity at the odd split
prime \(p_0\), where the source level is principal. The
torsion calculation in Theorem~\ref{thm:definitelevels678} applies:
if \(1+p^a B\) has a nonzero first residue \(B\bmod p\),
prime-to-\(p\) powers preserve its depth and the \(p\)-th
power increases the depth by exactly one. Such a matrix
has infinite order. Thus every orbit has the same size and
a uniform average over finer classes equals the coset
average for an invariant sum.

The image of \(K\) in the split component at each chosen
place over \(N\) is \(\mathrm{GL}_k(\mathbb F_N)\): lift an
invertible residue matrix integrally and use its inverse
transpose in the conjugate component. These choices are
independent at different places. Under \(k\in K\) the
distinguished row becomes the distinguished row of \(k^{-1}\)
in the fixed coarse coordinates. The subgroup \(K_{\rm src}\)
fixes this row pointwise modulo the chosen prime. Indeed its
embedded elements preserve the middle lattice and fix \(z\).
For an integral \(X_0\)-matrix \(A\), the coefficient along
\(z\) in the image of a middle basis vector \(x_j\) with
\(j\ne *\) is \(-A_{*j}/\pi\). In the image of
\(x_*/\bar\pi+z/N\) the coefficient is
\((1-A_{**})/N\). Hence on the chosen component the required
row is the original unit row modulo \(\pi\). Preservation at
the conjugate place is compatible with the inverse transpose.
The coordinate description before the proposition identifies
the pure lattice with the row kernel. Transitivity on nonzero
rows and equal fiber sizes now give the uniform row average.

At a chosen place choose a local basis of \(\mathcal U\)
beginning with a basis of the saturation \(P\). By local
elementary divisors take \(L\) to be generated there by
\(\varpi^{a_j}e_j\), \(1\le j\le r\). A residue covector
annihilating \(L\) must have zero \(j\)-th entry if \(a_j=0\).
Allowing also the zero covector bounds its number by
\(N^{\ell+\#\{j:a_j>0\}}\). The product of
\(N^{\#\{j:a_j>0\}}\) over the chosen places is at most the
global index \(i\), by \(\#\{j:a_j>0\}\le\sum_j a_j\).
Division by the total number of rows proves the bound.
\end{proof}

\begin{theorem}[The corank sum on the period classes]
\leavevmode\label{thm:latticemean690}
Write \(n=a+1\), \(H=b^{n^3 f}\), \(\tau=b^{kf/4}\). At
every sufficiently large node of the shorter construction
on \(V_0\le x\le V_0^{3/2}\), and for every \(0\le\ell\le k\),
\[
  \sum_{c\in\mathcal C}
    \ \sum_{\substack{L\subseteq\Lambda_{c}\cap X_0\\
                  \operatorname{rank}_R L=k-\ell,\ d(L)\le H}}
               d(L)^{-k}
          \le \frac{|\mathcal C|}{N^{kf}}\tau^{\ell+1}.
\]
Here \(L\) ranges over all \(R\)-lattices of the indicated rank
in their rational spans, with the zero lattice of discriminant
one at rank zero.
\end{theorem}
\begin{proof}
First fix a coarse lattice and average as in
Proposition~\ref{prop:rowmean690}. Group the lattices \(L\)
by \(V=EL\), and write \(d=d_V\) and \(i=[\mathcal U\cap V:L]\).
Then \(d(L)=di\), \(d,i\ge1\). Put \(\mu=2a+4\),
\(r=k-\ell\). When \(di\le H\) one has
\[
    i(di)^{-k}\le H^{\mu+1} d^{-k-1}i^{-r-\mu}.
\]
Indeed after division the right side is at least
\(d^\mu i^\ell\) times the left side.

For a rank \(r\) projective \(R\)-lattice \(P\) the sum of
\([P:L]^{-s}\) over its full-rank sublattices equals
\(\prod_{j=0}^{r-1}\zeta_E(s-j)\) for \(s>r\). Locally at
a prime of norm \(q\), project onto the last coordinate and
fix the intersection with the first \(r-1\) coordinates.
If that intersection has index \(t\) and the last image
is \(\varpi^m R_w\), there are \(t\) choices for the first
coordinates of a lift of \(\varpi^m\), and the total index
is \(tq^m\). The local series is thus \((1-q^{-s})^{-1}\)
times the rank \(r-1\) series at \(s-1\).
Every finite index sublattice is determined by these choices
at its finitely many exceptional primes. Conversely their
intersection inside \(P\) realizes the choices, by the
Chinese remainder theorem. This proves the formula including
nonfree \(P\). At \(s=r+\mu\) it is at most \(2^r\) for our
fields. Indeed \(\zeta_E(t)\le\zeta(t)^{2f}\) by the Euler
factors, and
\(\log\zeta(t)\le 2^{-t}(1+2/(t-1))\) for \(t>1\), by
summing the tail by integration. For \(t=\mu+1=2a+5\) and
\((a+1)2f\le2^{2a}\) these inequalities give
\(\log\zeta_E(t)<\log2\).

Combining these bounds with Lemma~\ref{lem:spacezeta690}
and \((N^k-1)^{-f}\le 2^f N^{-kf}\) bounds the conditional
mean over each coarse class by
\[
    N^{-kf} 2^{f+k} H^{2a+5}
        \big\{N^f(25 f k\rho_E)^{f(k+1)}\big\}^{\ell}.
\]
All terms summed are nonnegative, so finite partial sums
justify the changes of order. The bound holds with the same
constant for every coarse class, hence for the unweighted
mean on \(\mathcal C\).

We verify the sizes on the indicated interval. Write \(m=2n\)
for the number of lifts and \(r_0\) for the target rank.
Here \(\sqrt x\le m\le2\sqrt x\), \(x^5\le r_0\le3x^5\),
\(k=r_0+1-m\), and \(2^s=m^{22}\). The cyclotomic field
\(\mathbb Q(\mu_{2^s})\subseteq E\) gives \(f\ge2^{s-2}\).
The inclusion \(E\subseteq\mathbb Q(\mu_{q_1d})\) and
\(q_1d\ll_{q_0,d}x^{11}\log(2x)\) also give
\(2f,\rho_E\le q_1d\). For the latter inequality the
cyclotomic discriminant formula bounds its root discriminant
by its conductor, and the relative discriminant formula
can only lower the root discriminant in a subfield.
Consequently \(k\le f\), \((a+1)2f\le2^{2a}\), and
\(k\ge1000(2n)^9\) for large \(x\). The construction imposes
\(x^{10^6}\le x_1\), \(x_1^{10^6}\le b\) and
\(\log N\le30m\log b\). Thus
\((200 f k\rho_E)^{1000}\le b\), \(N\le b^{60n}\) and
\(b\ge1000\) throughout the interval for large \(V_0\).
We obtain \(2^{f+k}\le b^f\) and
\[
 1+n^3(2a+5)\le k/4,\qquad
 60n+(k+1)/1000\le k/4 .
\]
The first inequality absorbs the factor \(2^{f+k}H^{2a+5}\)
in \(\tau\). The second absorbs the factor inside braces in
\(\tau\). Multiplication by \(|\mathcal C|\) proves the claim.
\end{proof}

The sum includes every common lattice of the preceding theorem.
To pass from lattices to orbital pairs one must also count
the distinct rational terms over a fixed lattice and residue
label. Multiple such terms can have the same lattice.


\subsection{Rational frames over a lattice and the weighted pair count}
\label{sec:framepairs691}

We count the multiplicity with which a lattice in
Theorem~\ref{thm:latticemean690} can occur. Fix the first class
\(c\) and write \(\Lambda=\Lambda_c\). For the image of the
frame under a term from \(c'\) write its blocks in the original
rational coordinates as
\[
             \gamma\mathcal F=
               \begin{pmatrix} A_\gamma\\ X_\gamma\\ v_\gamma\\
                                  B_\gamma\end{pmatrix}.
\]
The four blocks have rows along \(L,X_0,Ez,L^\vee\), respectively. The entries of
\(v_\gamma\) are the scalar coefficients along \(z\). We order
the last block by the duals of the basis of the first block.
The cross products \(C=\mathcal F^*h\gamma\mathcal F\) specify
both \(B_\gamma=C_{[1,a],*}\) and \(v_\gamma=-C_{n,*}/N\).
Each column of the frame is integral in \(\Lambda\), and its
top and bottom coordinates are integral. Two frames with the
same \(C\) have their positive blocks differing by columns
of \(\Lambda_X\), since the terminal and bottom parts cancel
and the top part can be subtracted integrally.

Let \(B_{\rm res}\) be the image of the integral finite
Bessel group in the product of the two level groups. Choose one
representative for each coset \(s B_{\rm res}\) and write a
residue \(s=s_0t\) with \(t\in B_{\rm res}\); call \(t\) its
label. A term is counted once per double coset under the two
rational pointwise unipotent frame fixers of principal level.
We count cosets admitting a near representative. At both
\(y\) and \(\gamma^{-1}y\), near means that the points lie in
the reduced frame domain and the Bessel length at every real
place is at most \(b^{6n+2}\). We take the rank \(j\) of a
surviving term as in Theorem~\ref{thm:rationalflags680}, with
\(j=n\) in the open case. Write \(\delta=D_j\ne0\) and form
the \(n\)-by-\(n\) matrix \(Y\) from the first \(j\) rows of
\(C\) followed by the remaining rows of the identity. Then
\(\det Y=\delta\).

\begin{lemma}[Bounded invariants of near frames]\leavevmode
\label{lem:frameheight691}
At sufficiently large nodes of the interval of
Theorem~\ref{thm:latticemean690} there are at most
\(b^{270n^3 f}\) possible matrices \(C\) for near integral
terms. For fixed \(j,C\), at most \(b^{502n^4 f}\) arrays occur
as the top coordinates of the columns
\(\gamma\mathcal F\operatorname{adj}(Y)e_i\), \(j<i\le n\).
These columns belong to \(T\), and are unchanged by every
pointwise frame fixer on the left. Moreover the gauge of
Theorem~\ref{thm:allpair689} satisfies \(G_\gamma\le H=b^{n^3 f}\).
\end{lemma}
\begin{proof}
Use the normalized frame at a real place, with scale factors
\(|s_h|=\kappa_v^{n-h}/\sqrt N\) from
Corollary~\ref{cor:realmarked682}. We have
\(1\le\kappa_v\le b^{23}\) for all these places. Indeed the
chosen different generator, up to a factor of absolute value
one, is \(\sqrt{-d}\,\Phi'_{q_1}(\zeta_{q_1})\), where here
\(\Phi_j(T)\) denotes a cyclotomic polynomial. If
\(u=\operatorname{rad}(q_1)\), then \(u\) is fixed and
\(\Phi_{q_1}(T)=\Phi_u(T^{q_1/u})\). This identity follows by
comparing the orders of the roots and the degrees. The
derivative at a primitive root therefore has magnitude
\((q_1/u)|\Phi'_u(\xi)|\) for some primitive \(u\)-th root
\(\xi\). These finitely many derivative values are nonzero.
The bounds on \(q_1,b,M\) in the preceding subsection now
give \(2\le|\Phi_v(\delta_E)|\le b\) and the stated bounds
on \(\kappa_v=M|\Phi_v(\delta_E)|/(2\pi)\).

For clarity, the reduced domain is defined by taking the
strict upper coefficients of the Gram--Schmidt triangular
matrix in rational frame coordinates in a set of residue
representatives modulo \(MR\) in \(\mathbb C^f\). One can
take all representatives to have maximum modulus at most
\(b^{26}\). Choose a rational basis from the spanning
integral elements \(\zeta_{q_1}^i,\sqrt{-d}\,\zeta_{q_1}^i\).
A half-open parallelepiped for \(M\) times their integer span
has maximum modulus at most \(2fM\max(1,\sqrt d)\le b^{26}\).
For each coset of \(MR\) choose the lexicographically first
of its finitely many points there. This gives a measurable
set by comparison among finitely many translates.

At a reduced point the normalized triangular coefficient at
\((i,l)\), \(i<l\), is the rational one times \(s_l/s_i\),
of modulus ratio at most one. The diagonal coefficients are
one. Let \(F_y\) denote the matrix of normalized columns
\(y^{-1}f_h\). For a point in the cut choose a Bessel element
\(u_0\) with \(\|u_0y\|\le b^{7n}\), using the infimum defining
the Bessel length and \(n\ge2\). Its action on the frame is
upper unitriangular. Column operations of that kind preserve
the successive orthogonalized columns of \(F_y\). Their
lengths are bounded by \(\|(u_0y)^{-1}\|=\|u_0y\|\), by
orthogonalization after the change of frame. Hence each
column of \(F_y\) has length at most
\(n b^{26}b^{7n}\le b^{37n}\). This applies to both boxed
points at the same place. Their cross products are bounded
by \(b^{74n}\) since the normalized Hermitian matrix has
operator norm one. Removing the two scale factors costs at
most \(N\le b^{60n}\), so \(|\Phi_v C_{il}|\le b^{134n}\).
Integral tuples of \(h\) entries of maximum modulus \(R_0\)
at all embeddings number at most \((1+2R_0)^{2fh}\).
In fact the product formula gives maximum distance at least
one between distinct tuples, and balls of radius \(1/2\)
pack inside the ball of radius \(R_0+1/2\) in the maximum
norm. With \(h=n^2\) this proves the first count.

For \(j<n\), the kernel of evaluation against \(A_j\) on
\(\gamma T\) equals \(T\cap\gamma T\). It contains that space
because \(A_j\subset L\), and it has dimension \(n-j\) by
\(D_j\ne0\). This puts each indicated column in \(T\). Its
top coordinates are integral since \(Y\) is integral. To
bound them, the unitriangular matrix \(U\) of \(u_0\) on the
normalized frame satisfies
\(\|U\|\le b^{7n}\|F_y\|\le b^{44n+1}\). Here the normalized
fixed frame consists of Euclidean orthonormal basis vectors
and \(u_0y F_y\) has coefficient matrix \(U\). Since
\(\det U=1\), for all \(w\)
\[
                    \|w\|\le b^{46n^2}\|F_y w\|.
\]
Indeed \(\|Uw\|\le b^{7n}\|F_yw\|\) and
\(\|U^{-1}\|\le\|U\|^{n-1}\).
All entries of \(\operatorname{adj}Y\) and \(\det Y\) have
modulus at most \(b^{135n^2}\), by determinant expansion and
\(n\le b\). Express the shared column both in the fixed and
in the second frame at \(y\). Its second-frame coefficients
before scaling are entries of \(\operatorname{adj}Y\).
The ratio of any two frame scales is at most \(b^{23n}\).
The preceding coefficient inequality and the column lengths
at \(\gamma^{-1}y\) bound every fixed rational coefficient by
\[
              b^{46n^2}\, n b^{37n+23n+135n^2}
                    \le b^{250n^2}.
\]
Packing applies to the at most \(an\) top entries, proving
the second count. For \(j=n\) the array is empty.

For the gauge, \(0\le\Phi_v(\Delta)\le|\Phi_v(D_n)|^2\).
In the lower ranks Lemma~\ref{lem:quotientframe689} gives
\(|\Phi_v(d_*)|=N|\Phi_v(D_j)|\). In the open case \(d_*=D_n\).
All finite row indices are at least one. Since \(M\le b^{22}\),
the bound for either term in the gauge is
\(G_\gamma\le b^{(44a+60n+270n^2)f}\le b^{n^3f}\) for
sufficiently large \(n\).
\end{proof}

\begin{proposition}[Multiplicity over a fixed positive lattice]
\leavevmode\label{prop:fibercount691}
Fix \(c,t\) and a corank \(\ell\) lattice
\(L_0\subseteq\Lambda_X\), \(d(L_0)\le H\).
The number of near terms \((c',\gamma)\) of these ranks and
this label with \(L_0\subseteq L_\gamma\) is at most
\[
                       b^{2000 n^4(\ell+1)f}.
\]
\end{proposition}
\begin{proof}
First, equality of frames up to a left level move identifies
terms of one label. Indeed, after the move, the relative
rational isometry fixes the first frame pointwise. Write it
as \(u h\) in \(P=U_0\rtimes\mathrm U(X_0)\) as in
Proposition~\ref{prop:rationalorbits686}. The relative finite
matrix after the two class changes fixes the same frame in
lattice coordinates. Its residue thus lies in \(B_{\rm res}\);
the two residues belong to the same coset \(s_0 B_{\rm res}\).
Equal labels imply equality of the residues. The relative
finite matrix is integral and congruent to one at \(M\).
On the middle quotient of the integral lattice
\((\Lambda\cap L^\perp)/\Lambda_L\) it induces the class
translate of \(h\oplus 1_{Ez}\), still integral and congruent
to one. Extending this middle action by the identity on the
two integral arms shows that the class translate of \(h\)
belongs to the definite source level, the preimage of the
principal level under this extension. Thus the two source
classes \(c'\) coincide. With their representatives equal,
\(h\) is an automorphism at that level and is the identity
by Proposition~\ref{prop:rowmean690}. The relative rational
element is consequently a pointwise unipotent level move
on the right. This proves the identification in the original
double cosets.

Fix \(j,C\) and the shared array of Lemma~\ref{lem:frameheight691}.
Write \(B=C_{[1,a],*}\). We give all the level moves used
below. In coordinates \((p,y_0,q)\) on the two isotropic
arms and the middle block, with middle Gram \(H_0\), one such
move is
\[
   (p,y_0,q)\longmapsto
   (p-D^*H_0 y_0-\epsilon D^*H_0Dq,\ y_0+Dq,\ q).
\]
Here \(D=M(\eta_1,\ldots,\eta_a)\) with columns in \(M\Lambda_X\)
and \(\epsilon=-\zeta_3\in R\), so that
\(\epsilon+\bar\epsilon=1\). Substitution in the Hermitian
form verifies unitarity. The move fixes \(L\) and \(z\)
pointwise. It preserves \(\Lambda\) and is one modulo \(M\),
since any \(\eta_i\) pairs integrally with the middle lattice.
Its inverse is integral by triangularity. We may also use the
central move \(p\mapsto p+Sq\) with \(S=-S^*\) in
\(M\,\operatorname{Mat}_{a}(R)\).

For each term choose an \(R\)-linear lift
\(x\mapsto x+M w(x)\), \(x\in L_0\), orthogonal to the
second frame, with \(w(x)\in\Lambda_L=R^a\).
Such lifts exist by the definition of \(L_\gamma\); they
can be chosen linearly because the surjection of the module
of level lifts onto \(L_0\) splits over the projective
module \(L_0\). Each row of \(w\) is a functional on \(L_0\).
The submodule of functionals \(x\mapsto h(\eta,x)\),
\(\eta\in L_0\), has index \(d(L_0)^2\), by Hermitian
duality in its rational span. Grouping the lifts by these
classes costs at most \(d(L_0)^{2a}\). Within one group the
first move makes \(w\) equal to a fixed representative,
because it changes its rows by \(-h(\eta_i,x)\).
For the resulting frames,
\[
                    X_\gamma^*h_X x=-M B^*w(x)
                           \qquad(x\in L_0),
\]
where \(h_X\) in matrix expressions denotes the positive
Gram matrix. Their positive columns therefore differ from
the fixed frame's positive columns by vectors of
\(P_0=\Lambda_X\cap(EL_0)^\perp\), a lattice of rank \(\ell\).
Record these differences modulo \(M\delta P_0\) in each
column. There are at most
\(|{\rm N}_{E/\mathbb Q}(M\delta)|^{\ell n}\) records.

Within one record put \(V=X_\gamma-X_{\gamma'}\).
All columns of \(VY^{-1}/M\) belong to \(\Lambda_X\), by
the adjugate formula and the congruence. If \(j<n\), its
columns \(i>j\) vanish because the corresponding combinations
of the two frames lie in \(T\). If \(j=n\), only the last
column needs attention. The vector
\(\gamma\mathcal F C^{-1}e_n\) has bottom coordinates zero
and \(z\)-coordinate \(-1/N\). It belongs to \(\gamma T\),
where the form is negative semidefinite at every embedding.
Thus its \(X_0\)-component has norm at most \(1/\sqrt N\)
at each embedding in the positive Hermitian space.
The same holds for the second term. The last column of
\(VC^{-1}\) has norm at most \(2/\sqrt N<1\) everywhere.
It is in \(\Lambda_X\). A nonzero such vector has a totally
positive integral norm with product at least one, so this
column also vanishes. These arguments apply after the
preceding moves, since the moves fix the first frame.
Multiplying \(VY^{-1}\) by \(Y\) now gives \(V=D B\), with
the columns of \(D\) in \(M\Lambda_X\): use the first
\(\min(a,j)\) columns and fill the others with zeros. The
first type of move, with the appropriate sign, thus makes
the entire positive block equal throughout this record.

Only the top coordinates can still differ. Record their
residues modulo \(M\delta\bar\delta\) entry by entry. This
costs at most \(|{\rm N}_{E/\mathbb Q}(M\delta\bar\delta)|^{an}\).
Let \(\Xi\) be the top difference of two frames with one
record. The equal Gram matrices of the frames give
\(B^*\Xi=-\Xi^*B\). Moreover \(\Xi\) vanishes on \(\ker B\).
For \(j<n\) this kernel is spanned by \(Y^{-1}e_i\) for
\(i>j\), since \(B\) has rank \(j\) by
Theorem~\ref{thm:rationalflags680} and the first \(j\)
rows are independent. The shared top entries are the same
and have been fixed through all moves. For \(j=n\), \(B\)
has rank \(a\), so \(B^*\) is injective and the Gram identity
gives \(\Xi\ker B=0\).

Let \(R_0\) be \(Y^{-1}\) times the \(n\)-by-\(a\) rectangular
diagonal matrix with ones only in the first \(\min(a,j)\)
positions. Then \(BR_0B=B\): the first \(\min(a,j)\)
rows are preserved and span the row space. Put \(P=BR_0\),
\(T_0=\Xi R_0\), so \(T_0B=\Xi\), since
\(1-R_0B\) has image in \(\ker B\). The matrix
\[
                        S=T_0-T_0^*-P^*T_0
\]
is skew-Hermitian since \(P^*T_0=R_0^*B^*\Xi R_0\) is.
Multiplying the Gram identity by \(R_0^*\) on the left
gives \(T_0^*B=-P^*\Xi\), hence \(SB=\Xi\).
Each entry of \(R_0\) is integral after multiplication by
\(\delta\). The congruence on \(\Xi\) therefore puts \(S\)
in \(M\operatorname{Mat}_a(R)\). The central move aligns
the two frames, proving they represent one term. This
argument also covers \(j=0\), when the shared columns are
all the columns.

There are at most \(n+1\) choices of \(j\). The two norms of
moduli used above are each at most \(b^{600n^2 f}\) by
\(|\Phi_v\delta|\le b^{135n^2}\) and \(M\le b^{22}\).
Multiplying the two counts in Lemma~\ref{lem:frameheight691},
the lift count, and the two residue counts bounds the base-\(b\)
logarithm, divided by \(f\), by
\[
    2+270n^3+502n^4+2an^3+600 n^2(\ell n+an)
                                  \le 2000n^4(\ell+1).
\]
The first term absorbs \(n+1\). All the choices were made
inside a fixed \(c,t,L_0\), as required.
\end{proof}

\begin{theorem}[The sum over near pairs in agreeing residue blocks]
\leavevmode\label{thm:paircount691}
On the interval and fields of Theorem~\ref{thm:latticemean690},
sum over ordered pairs of near surviving terms in the same
class position \((c,c')\) and with the same agreeing residue
at both levels. At corank \(\ell=k-\dim I_X\) the sum satisfies
\[
       \sum_{c,c'}\ \sum_{\text{such pairs of corank }\ell}
                    d(L_{12})^{-k}
       \le |B_{\rm res}|\,\frac{|\mathcal C|}{N^{kf}}\,
                                                   \tau^{\ell+2}.
\]
Every such pair satisfies \(d(L_{12})\le H\) and \(\ell\le n^3\).
\end{theorem}
\begin{proof}
The integral positive lattice \(L_{12}\) has \(d(L_{12})\ge1\).
By Theorem~\ref{thm:allpair689} and
Lemma~\ref{lem:frameheight691},
\[
         1\le d(L_{12})\le H(2/M)^{c^+ f}.
\]
Since \(M/2\ge b^{19}\), this gives \(c^+\le n^3/19\).
The dimension assertion of Theorem~\ref{thm:allpair689}
then gives \(\ell\le c^++1\le n^3\).
Assign to each ordered pair its first class, the common
residue label and the actual lattice \(L_0=L_{12}\).
Its two terms both contain this lattice. For each
\((c,t,L_0)\) their number is bounded by the square of the
bound in Proposition~\ref{prop:fibercount691}. This remains
an upper bound when the second class is required to be the
same for both. Since \(k\ge1000(2n)^9\) and \(\ell\le n^3\),
the square is at most \(b^{kf/4}=\tau\). Sum the weight
\(d(L_0)^{-k}\) using Theorem~\ref{thm:latticemean690}
and take the \(|B_{\rm res}|\) labels. This proves the claim.
\end{proof}

The discriminant estimate and the weighted count now apply
to the same pairs in the unfolded period. The absolute sizes
of their orbital integrals and the contribution of the
discarded cut tails still enter the second moment.


\subsection{Projected frame volumes and the integral over positive shifts}
\label{sec:schurshift692}

Let \(V\) denote the complex Hermitian space at one infinite
place of the shorter construction, in the normalized coordinates of
Corollary~\ref{cor:realmarked682}. Its form is \(h(v,w)=v^*Hw\),
where \(H=H^*=H^{-1}\) exchanges the two isotropic arms, is \(1\)
on the \(k\)-dimensional positive summand \(X_0\), and is \(-1\)
on the unit terminal vector \(z\). Put \(F=(e_1,\ldots,e_a,z)\),
\(T=\operatorname{im}F\), \(L=\langle e_1,\ldots,e_a\rangle\),
\(n=a+1\). For an isometry \(y\), let \(\Omega=y\ker(H+1)\).
It is a negative space of dimension \(n\). For another isometry
\(g\), put \(Z=F^*HgF\). The ratio of projected frame volumes is
\[
            \nu_g(y)=\frac{|\det(\Omega_b^*HgF)|}
                                {|\det(\Omega_b^*HF)|},
\]
where \(\Omega_b\) denotes any basis matrix of \(\Omega\).
A basis change multiplies both volumes by the same positive factor.
Both determinants are nonzero. Indeed \(T^\perp=L\oplus X_0\)
is positive semidefinite, as is \(gT^\perp\), whereas \(\Omega\)
is negative definite.

We use the flags of Theorem~\ref{thm:rationalflags680}. Let
\(j\) be the rank of a surviving term, \(i=\min(a,j)\),
\(p=i+1\), \(F_i=(e_1,\ldots,e_i,z)\), and \(B=F_i^*HgF_i\).
For \(j=n\) this is the full invertible Gram matrix. For
\(j\le a\) invertibility follows from
Lemma~\ref{lem:quotientframe689} and the nonzero leading minor.
Write \(X\) for projection on the positive middle summand.

\begin{lemma}[A determinant bound in a negative plane]\leavevmode
\label{lem:negdet692}
There are a matrix \(\omega\) with \(p\) columns in \(\Omega\)
and an invertible lower triangular matrix \(D=F_i^*H\omega\),
depending only on \(y,Z,i\), such that \(\omega^*H\omega=-I_p\).
Put \(\omega'=gF_iB^{-1}D,\ V_g=X(\omega'-\omega)\) and let
\(w\) be the last column of \(V_g\). Then
\[
       |\det B|\det(I_k+V_gV_g^*)\le 2^p|\det D|^2\nu_g(y).
\]
Suppose \(\|y^{-1}f\|\le\rho\) for each column \(f\) of \(F\)
and the Bessel length of \(y\) is at most \(\theta\).
For \(d=D_{pp}\) we then have
\[
        |\det D|\le\rho^p,\qquad
        \frac{|d|^2}{1+\|w\|^2}\le\theta^{2(j-i)} .
\]
\end{lemma}
\begin{proof}
The map \(F^*H:\Omega\to\mathbb C^n\) is an isomorphism by the
nondegeneracy just proved. Take the preimages \(\omega_0\) of
the columns of \(F^*HgF_iB^{-1}\). Their pairings with \(F_i\)
form the identity. Apply Gram--Schmidt in the positive metric
\(-h|_{\Omega}\), beginning with the last column. Each new
column is a combination of that column and later ones, with a
positive real diagonal coefficient. Thus \(\omega=\omega_0 D\)
has the stated triangular properties and depends continuously
on \(y\) for fixed \(Z\). Continuity follows also by using the
basis \(y\Omega_-\) for a fixed basis \(\Omega_-\) of
\(\ker(H+1)\) in the initial linear system.

Put \(U=L\cap gL\). In the open case \(U=0\) since \(Z\) is
invertible. In the lower ranks the common triangular flag gives
\(a-i\) independent vectors in \(U\), one at each index from
\(i+1\) to \(a\). The leading \(i\)-minor implies
\(\dim U\le a-i\) by pairing \(L\) with \(gL\).
The columns of \(\omega\) are orthogonal to \(U\).
Indeed \(F^*H\omega=F^*H\omega'\), and the vectors of
\(U\subset T\) pair to zero with \(gF_i\). In both frames
replace the omitted columns \(e_{i+1},\ldots,e_a\) and their
images by the same triangular common vectors. Each replacement
has determinant of modulus one, since the two leading
coefficients at every replaced index have modulus one. Extend
\(\omega\) to a basis of \(\Omega\). In the two pairing
matrices the common columns have a zero block in the first
\(p\) rows, and they have identical entries in the other rows.
Expansion along these columns cancels their common determinant,
which is nonzero by invertibility of the full matrices. Hence
\[
              \nu_g(y)=
           \frac{|\det(\omega^*HgF_i)|}{|\det D|}.
\]
This also covers the full frame, where the omitted block is empty.

The columns of \(\omega'-\omega\) belong to \(T^\perp\).
Their Gram matrix is \(V_g^*V_g\), because \(L\) is the
radical of \(T^\perp\). The form restricted to \(gT\) is
negative semidefinite. On writing \(K=-\omega^*H\omega'\),
expansion therefore gives
\[
        K+K^*=I+V_g^*V_g-\omega'^*H\omega'
                                          \succeq I+V_g^*V_g .
\]
If \(P=(K+K^*)/2\), then \(P^{-1/2}KP^{-1/2}=I+\mathrm i S\)
with \(S=S^*\). Each eigenvalue of \(I+\mathrm i S\) has
modulus at least one, so \(|\det K|\ge\det P\).
Monotonicity of determinant on positive definite matrices
follows by conjugating the larger matrix by the inverse
square root of the smaller one. Also
\(\det(I+V_g^*V_g)=\det(I+V_gV_g^*)\) by the singular
values. Substituting
\(|\det K|=|\det D|^2\nu_g(y)/|\det B|\)
proves the determinant bound.

The columns of \(y^{-1}\omega\) are Euclidean orthonormal in
\(\ker(H+1)\). Each row of \(D\) has length at most the
Euclidean length of the corresponding \(y^{-1}f\), by the
projection inequality. Hadamard's inequality gives the first
bound on \(D\). If \(j\le a\), take a common terminal vector
with coefficients \(\zeta,1\) along \(z,gz\), respectively,
where \(|\zeta|=1\). Subtract the common vectors of indices
\(a,\ldots,i+1\), in that order, to kill the coefficients along
\(ge_a,\ldots,ge_{i+1}\). This gives a vector \(w_0\in T\cap
g\operatorname{im}F_i\) with the same terminal coefficients.
Now \(F_i^*Hw_0=-\zeta e_p^{\rm std}\), so the last column
of \(\omega'\) is \(-(d/\zeta)w_0\) by invertibility of \(B\).
It has squared Hermitian norm \(-|d|^2\) and is orthogonal to
the corresponding column of \(\omega'-\omega\). Subtracting
the latter column gives the last column of \(\omega\) and hence
\(-1=-|d|^2+\|w\|^2\).

If \(j=n\), triangularity gives \(\omega_p\perp L\), where
\(\omega_p\) denotes the last column. The Bessel length is one
plus the infimum of \(\|u y\|\) over the Bessel unipotent
group, so some such \(u\) has \(\|u y\|\le\theta\).
As \(u^{-1}z-z\in L\), Cauchy--Schwarz after applying \(y^{-1}\)
gives \(|d|=|h(u^{-1}z,\omega_p)|\le\|(uy)^{-1}z\|
\le\theta\). Here \(\|(uy)^{-1}\|=\|uy\|\) since \(H\) is a
unitary involution. This proves the second bound in every rank.
\end{proof}

\begin{proposition}[The integral of the real majorant]\leavevmode
\label{prop:schurreal692}
Fix a representative \(g_0\), its \(Z,j\), and a point \(y\)
in the preceding bounds. For \(A\in\mathbb C^{k\times i}\)
consider all surviving \(g\) with this \(Z,j\) and
\[
          X(gF_i)=X(g_0F_i)+A B_{[1,i],*}.
\]
There is a single nonnegative smooth function \(\varphi_y(A)\)
majorizing \(\nu_g(y)^{-(k+i)}\) for every such \(g\), with
\[
 \int_{\mathbb C^{k\times i}}\varphi_y(A)\,dA
   \le \frac{(2\rho^2)^{pi}(2^p\theta^{2(j-i)})^k}
                 {|\det B|^{k+i}}\qquad(k\ge7).
\]
Here \(dA\) is the ordinary real Lebesgue measure in complex
orthonormal coordinates. The majorant is measurable in \(y\).
\end{proposition}
\begin{proof}
Write \(D=\left(\begin{smallmatrix}D_1&0\\ *&d\end{smallmatrix}\right)\).
The shift in \(X(\omega')\) is \(A B_{[1,i],*}B^{-1}D
=(AD_1,0)\). Thus \(V_g=(C_1+AD_1,w)\) with \(C_1,w\)
independent of \(A,g\). Put \(G=I_k+ww^*\), \(s=k+i\).
The lemma gives the stated domination with
\[
 \varphi_y(A)=
 \left(\frac{2^p|\det D|^2}{|\det B|}\right)^s
 \det\!\bigl(G+(C_1+AD_1)(C_1+AD_1)^*\bigr)^{-s}.
\]
All its ingredients depend continuously on \(y\) for fixed
frame products and fixed \(g_0\).
Changing variables to \(W=G^{-1/2}(C_1+AD_1)\) gives an
integral of the determinant factor equal to
\[
      |\det D_1|^{-2k}(\det G)^{-k}
           \pi^{ki}\prod_{t=0}^{i-1}\frac{t!}{(k+t)!}.
\]
Indeed the real Jacobian for left multiplication by \(G^{1/2}\)
is \((\det G)^i\), and right multiplication by \(D_1^{-1}\)
contributes \(|\det D_1|^{-2k}\). To integrate the remaining
factor, for any positive \(P\) use
\(\det(P+vv^*)=\det P(1+v^*P^{-1}v)\). At exponent \(s'>k\)
one column integrates to \(\pi^k\Gamma(s'-k)/\Gamma(s')\)
times \((\det P)^{1-s'}\). This follows by writing the
negative power as its gamma integral and integrating the
Gaussian on \(\mathbb C^k\); nonnegativity justifies the
exchange. Starting at \(s'=k+i\) and repeating for all
columns yields the displayed product. For \(k=7\),
\(\pi^k/k!<1\), and the ratio decreases thereafter because
\(\pi/(k+1)<1\). Each \(\pi^k t!/(k+t)!\le1\).
Now \(\det G=1+\|w\|^2\) and
\(|\det D|=|\det D_1|\,|d|\). The two bounds of the lemma
give the claimed integral after collecting the \(i\)-th and
\(k\)-th powers. For \(i=0\) the integral over the zero space
is evaluation and the empty product is one.
\end{proof}

\begin{theorem}[Mean over a translate of the marked column lattice]
\leavevmode\label{thm:schurmean692}
Take the fields, marked classes and interval of
Theorem~\ref{thm:latticemean690}. At every sufficiently large
node put \(\rho=b^{37n}\), \(\theta=b^{6n+2}\),
\(\tau=b^{kf/4}\). Fix a near surviving term of rank \(j\ge1\),
put \(e=j-i\), and set
\[
        P_j=\prod_{v:E^+\hookrightarrow\mathbb R}
                     |\Phi_v(D_j)|
             =|{\rm N}_{E/\mathbb Q}D_j|^{1/2}.
\]
Choose one embedding \(\Phi_v\) of \(E\) over each real place \(v\).
Use the positive columns \(Q\in X_{0,\infty}^i\) in the
original rational basis. A shift by \(M Q\) times the first
\(i\) isotropic rows of the rational cross-product matrix
induces a normalized shift \(A_v=Q_v S_v\) at place \(v\).
Let \(\Psi_y(Q)=\prod_v\varphi_{y_v}(Q_vS_v)\).
For \(\mathcal A=(\Lambda_c\cap X_0)^i\), any shift \(t\), and
any fundamental parallelepiped \(\mathcal P\) for \(\mathcal A\),
\[
 \frac{1}{\operatorname{vol}\mathcal P}
 \int_{\mathcal P}\sum_{\lambda\in\mathcal A}
                         \Psi_y(Q+t+\lambda)\,dQ
 \ \le\
 \frac{\tau}{c_E^{ki}N^{fi}}\!
 \left[\left(\frac{b^{13}}{M^2}\right)^{if}
                           \frac{N^{ef}}{P_j}\right]^k ,
 \qquad c_E=2^{-f}\sqrt{|{\rm disc}(E)|}.
\]
The bound holds uniformly for the reduced points \(y\) in the
Bessel length cut. The measure on each \(X_{0,v}\) is Euclidean
volume for its positive form.
\end{theorem}
\begin{proof}
The bounds on \(y_v\) required for the proposition hold at
every place by the proof of Lemma~\ref{lem:frameheight691}.
Write \(s_{h,v}\) for the frame scale of column \(h\), of
modulus \(\kappa_v^{n-h}/\sqrt N\). Put
\(r_v=\prod_{h=1}^i\kappa_v^{n-h}\). Row \(h\) of the
normalized cross products acquires a factor \(\bar s_{h,v}\),
and each column acquires its own scale. The latter scales also
multiply the shifted image columns. Thus
\[
 S_v=\operatorname{diag}_{h\le i}(M/\bar s_{h,v}),\qquad
 |\det S_v|=\frac{(M\sqrt N)^i}{r_v},\qquad
 |\det B_v|=\frac{r_v^2}{N^j}|\Phi_v(D_j)|.
\]
For the last equality, the unscaled subframe has determinant
of modulus \(N|\Phi_v(D_j)|\) in closed rank and
\(|\Phi_v(D_j)|\) in open rank. The product of the squared
subframe scales, including the terminal scale, is
\(r_v^2/N^{i+1}\).

The integral in the \(Q_v\) coordinates is multiplied by
\(|\det S_v|^{-2k}\). The factor \(r_v^{2k}\) cancels exactly
and \(r_v^{-2i}\le1\), since \(\kappa_v\ge1\). By the
proposition the resulting bound at \(v\) is at most
\[
 \left[\frac{(2\rho^2)^pN^j}{|\Phi_v(D_j)|}\right]^i
 \left[\frac{b^{13i}N^{e}}{M^{2i}|\Phi_v(D_j)|}\right]^k .
\]
Indeed \(2^p\theta^{2e}\le b^{13i}\). In closed rank use
\(i\ge1\) and \(b\ge4\). In open rank \(n=i+1\), and
\(2(6n+2)+(i+1)/2\le13i\) when \(i\ge33\), as holds at
these nodes. Since \(D_j\) is a nonzero algebraic integer,
\(P_j\ge1\). The product of the first brackets to the power
\(i\) is at most \(\tau\), since \(2\rho^2\le b^{75n}\),
\(N\le b^{60n}\), and
\((75np+60nj)i\le135n^3\le k/4\).

Tonelli's theorem and translation of each summand identify
the integral of the periodization over \(\mathcal P\) with
the integral of \(\Psi_y\) on the whole real vector space.
The translates tile it up to null sets for every \(t\).
By the covolume formula of Section~\ref{sec:linkdisc687}
and \(d(\Lambda_c\cap X_0)=N^f\), the volume of
\(\mathcal P\) is \(c_E^{ki}N^{fi}\). Dividing the product
of the local bounds by this volume proves the assertion.
\end{proof}

To pass from this translation mean to an orbital contribution
one must also count the integral frames with the same positive
shift and integrate their cutoffs against the Bessel functions.


\subsection{The orbital Schur bound and the contribution of the closed cells}
\label{sec:orbitalclosed693}

Return to the rational frame \(\mathcal F\) and the fields and
interval of Theorem~\ref{thm:schurmean692}. Write \(U\) for
the Bessel unipotent group and \(U^0\) for the subgroup that
fixes \(\mathcal F\) pointwise. For a class \(c\) let \(\Gamma_c\)
be the rational elements of \(U\) acting integrally on
\(\Lambda_c\) and as the identity modulo \(M\). Write
\(\Gamma_c^0=\Gamma_c\cap U^0\). The full unipotent covolume
\[
               v_U=\operatorname{vol}(\Gamma_c\backslash U_\infty)
\]
is independent of \(c\). Indeed the finite lattice in \(U\) of
class \(c\) is conjugate to the one at the identity by its
representative in \({\rm U}(X_0)(\mathbb A_{F,f})\). The determinant
of this conjugation on the independent positive coordinates of
the Lie algebra, over \(F=E^+\), is a power of
\({\rm N}_{E/F}\det_E h_c=1\); the remaining graded
coordinates are fixed. Exponential coordinates identify \(U\)
with its Lie algebra as a variety. Thus the finite rational
approximation and product measure argument in
Lemma~\ref{lem:covol127} applies with this \(U\) and this
principal level. It expresses the displayed volume as the
volume of the adelic unipotent quotient divided by the volume
of the finite open lattice. Both volumes are positive and finite.

We describe the cutoffs used in the orbital integral. The image
of \(\Gamma_c\) on \(\mathcal F\) is exactly the upper
unitriangular group with entries above the diagonal in \(MR\),
where \(R=\mathcal O_E\). The congruence gives the inclusion,
since the two isotropic arms are split summands. Every elementary
entry between two isotropic vectors lifts by the elementary
matrix on \(L\) and its inverse adjoint on the dual arm. To
lift an entry from the terminal vector \(z\) to \(e_h\), choose
a vector \(t\) of the middle lattice with \(h(t,z)\) any specified
element of \(R\). Indeed \(z\) is part of a basis locally in the
marked lattice, so evaluation on it maps the integral dual onto
\(R\); the middle lattice is Hermitian unimodular. Set
\(\xi=-\zeta_3\in R\), so \(\xi+\bar\xi=1\). The map
\[
 v\longmapsto v+M\{e_h h(t,v)-t h(e_h,v)\}
              -M^2\xi h(t,t)e_h h(e_h,v)
\]
is the required element of \(\Gamma_c\). In the pairing of
two transformed vectors the terms linear in \(M\) cancel.
The square of the first-order term contributes
\(M^2h(t,t)\overline{h(e_h,v)}h(e_h,w)\), canceled by the
two last terms since \(\xi+\bar\xi=1\). Higher terms vanish
by \(e_h\perp t,e_h\). The map and its isometric inverse
preserve the lattice and the level. Every required triangular
matrix is a product of the elementary ones by elimination.

Let \(R_{\rm red}(y)\) be the reduced-set indicator from
Lemma~\ref{lem:frameheight691}. In rational frame coordinates
the upper unitriangular Gram--Schmidt coefficient matrix of
\(y^{-1}\mathcal F\) changes by right multiplication by the
inverse of the frame matrix of \(u\) when \(y\) is replaced
by \(uy\), \(u\in\Gamma_c\). A unique matrix in the image
just described puts its strict upper entries into the chosen
representatives modulo \(MR\). Reduce in increasing order
of the superdiagonal, using elementary column additions.
Each addition affects only its chosen entry and larger
superdiagonals, and at the first superdiagonal of a difference
the two choices would differ by \(MR\). This proves uniqueness.
Choose a measurable nonnegative \(p_c\) whose translates
under \(\Gamma_c^0\) sum to one. For example, a discrete
subgroup acting on a Lie group by left multiplication has a
Borel fundamental set: cover its quotient by countably many
open sets admitting local sections and restrict the sections
to the successive disjoint differences. Set
\[
     \beta_c=p_cR_{\rm red},\qquad
     \beta_c^\Box=\beta_c\,{\bf1}_{\{\ell(y_v)\le\theta\ \forall v\}} .
\]
Summing first over \(\Gamma_c^0\) in each coset gives
\(\sum_{u\in\Gamma_c}\beta_c(uy)=1\). In particular the
argument below applies to the class cutoffs for any choice
of this partition on the pointwise fixer.

Let \(W^\flat=\bigotimes_v W_v^\flat\) be the cut Bessel
functions, with cut of absolute value at most one supported
at length at most \(\theta=b^{6n+2}\). Use the squared local
quotient norm bound \(\|W_v^\flat\|_{U_v\backslash G_v}^2\le d_v\),
where \(d_v\) is the dimension of the value space. This is
the normalization of Theorem~\ref{thm:archbessel684}: the
uncut average has squared norm \(d_v\) when scaled there,
and cutting decreases it. The quotient by \(U_v\) gives the
same norm as the quotient by the Bessel group since the
remaining compact group has Haar mass one and acts unitarily.
The zero function also satisfies the bound. Write
\(d_B=\prod_v d_v\).

For a rational term \(\gamma\) in a fixed block \((c,s,c')\),
let \(\mathcal O_\gamma=\Gamma_c^0\gamma\Gamma_{c'}^0\) as a
set of rational isometries. The sum of the absolute cut
amplitudes on this orbit is bounded by the nonnegative integral
\[
 A_\gamma=\sum_{\delta\in\mathcal O_\gamma}\int_{G_\infty}
      \beta_c(y)\beta_{c'}(\delta^{-1}y)
        \|W^\flat(y)\|\,\|W^\flat(\delta^{-1}y)\|\,dy.
\]
The product of the two norms is exactly the Hilbert--Schmidt
norm of the corresponding simple tensor. The finite residue
and finite character have already been fixed in this expression.

\begin{lemma}[Positive shifts and the multiplicity at each shift]
\leavevmode\label{lem:orbitfib693}
Let \(j\ge1\) be the rank of \(\gamma\), \(i=\min(a,j)\),
and let \(H_i\subset H_{\rm all}\) be the \(R\)-span of the
first \(i\) isotropic rows inside the span of all isotropic rows
of \(C=\mathcal F^*h\gamma\mathcal F\). Put
\(i_\gamma=[H_{\rm all}:H_i]\), as in
Lemma~\ref{lem:quotientframe689}, with \(i_\gamma=1\) in open
rank. The orbit breaks up into at most \(i_\gamma^k\) sets.
In each one choose a base term \(\delta_0\). The changes in
the positive frame block there are \(M\lambda C_{[1,i],*}\)
with unique parameters \(\lambda\in(\Lambda_c\cap X_0)^i\).

Fix a reduced point \(y\) in the length cut and
\(\widehat y=u y\), \(u\in U_\infty^0\). For a fixed parameter
in one of these sets, the sum of
\(\beta_{c'}^\Box(\delta^{-1}\widehat y)\) over its terms is at
most \(\tau=b^{kf/4}\).
\end{lemma}
\begin{proof}
Right multiplication by \(\Gamma_{c'}^0\) fixes the image
frame. A left pointwise unipotent adds to its positive block
a multiple of its bottom block \(C_{[1,a],*}\). In fact the
coefficient of its \(h\)-th row is the positive component of
\((u-1)e_h^\vee\), where \(e_h^\vee\) is the dual isotropic
basis vector. For \(u\in\Gamma_c^0\) this component belongs to
\(M\Lambda_X\), \(\Lambda_X=\Lambda_c\cap X_0\): the bottom
and terminal changes are zero since the frame is fixed, and
the integral top component can be subtracted. Thus all the
changes belong to \(M(\Lambda_X\otimes_R H_{\rm all})\)
in the positive block. The submodule formed with \(H_i\)
has index \(i_\gamma^k\). This can be checked prime by prime,
where \(\Lambda_X\) is free of rank \(k\); both tensor products
embed in \(X_0^n\) by extending to \(E\), using local freeness.
The first \(i\) rows are independent by the nonzero leading
minor (in open rank they are rows of an invertible matrix).
They are a basis of \(H_i\) and give the unique parameter
relative to any base representative of the coset. All other
blocks except the top are now fixed, since \(C\) fixes the
terminal and bottom blocks.

For a nonzero cutoff the normalized columns
\(\widehat y_v^{-1}\delta F_v\) have lengths at most
\(\rho=b^{37n}\). Compare any two contributing terms of this
fiber. Their rational frame difference is only in the top
block and has entries \(t_{hj}\in R\), since the top arm
splits off integrally. Moreover
\(\widehat y_v^{-1}F_v=y_v^{-1}F_v\), and the bound in the
proof of Lemma~\ref{lem:frameheight691} gives
\(\|y_v^{-1}F_v w\|\ge b^{-46n^2}\|w\|\). A difference of
normalized columns of index \(j'\) has coefficient
\(\Phi_v(t_{hj'})s_{j',v}/s_{h,v}\) on this fixed frame.
Hence \(|\Phi_v(t_{hj'})|\le2\rho\, b^{46n^2+23n}\le b^{47n^2}\)
at every embedding for sufficiently large \(n\). Relative to
one contributing frame, the packing bound for integral entries
of Lemma~\ref{lem:frameheight691} bounds the number of full
frames in the fiber by
\((1+2b^{47n^2})^{2fan}\le b^{96n^4 f}\le\tau\).
An empty fiber contributes zero.

Terms with identical full frame differ on the right by an
element of \(\Gamma_{c'}^0\). Indeed their quotient fixes
\(\mathcal F\) and is principal on \(\Lambda_{c'}\), since their
residues coincide. Its action on the positive middle quotient
preserves \(\Lambda_{c'}\cap X_0\) and is the identity modulo
\(M\) there: apply the quotient isometry to a positive vector,
then subtract its top component as above. The definite
automorphism group is finite, and the principal subgroup at
either odd level prime is torsion free by the argument in
Proposition~\ref{prop:rowmean690}. The positive middle action
is therefore trivial, leaving exactly the pointwise unipotent
fixer. The sum of the cutoffs for this full frame is bounded
by the sum of the translates of \(p_{c'}\) under \(\Gamma_{c'}^0\),
which is one. Multiplication by the frame count proves the claim.
\end{proof}

\begin{theorem}[Schur's inequality on a rational orbit]\leavevmode
\label{thm:orbitschur693}
For a near surviving orbit of rank \(j\ge1\) put \(e=j-i\),
\(\mathcal V_i=c_E^{ki}N^{fi}\), and let \(P_j\) be as in
Theorem~\ref{thm:schurmean692}. Then
\[
 \frac{A_\gamma}{v_U d_B}
 \le \frac{\tau^2}{\mathcal V_i}
       \left[
       \left(\frac{b^{13}}{M^2}\right)^{if}
                  \frac{N^{ef}\sqrt{i_\gamma i'_\gamma}}{P_j}
       \right]^k .
\]
Here \(i'_\gamma\) is the row index for \(\gamma^{-1}\).
\end{theorem}
\begin{proof}
Write \(s_0=k+i\) and define
\[
 F_\gamma(y)=\sum_{\delta\in\mathcal O_\gamma}
          \beta_{c'}^\Box(\delta^{-1}y)
                       \prod_v\nu_{\delta_v}(y_v)^{-s_0}.
\]
It is measurable, nonnegative and left \(\Gamma_c^0\)-invariant:
left multiplication reindexes the orbit and fixes the original
frame. At a reduced \(y\) in the cut, fix \(u\in U_\infty^0\).
For \(Q\in X_{0,\infty}^i\), choose a real point \(T(Q)\in
U_\infty^0\) whose inverse adds \(Q_v S_v\) times the first
\(i\) bottom rows to the positive block at \(v\). Such a
continuous choice follows from the explicit matrix before
Theorem~\ref{thm:realgram682}: take \(h=1,K=0\) and the
negative of \(Q_v S_v\), padded by zero columns, for \(B\).
At \(uT(Q)y\) the volume ratio for \(\delta\) equals the one
at \(y\) for \(T(Q)^{-1}u^{-1}\delta\), since \(u,T(Q)\)
fix the first frame. The Gram matrix is unchanged and all
the common vectors of the flags retain their coefficients.
In any one of the sets of Lemma~\ref{lem:orbitfib693}, take
\(u^{-1}\delta_0\) as the real base frame. The difference of
its positive block from that of \(T(Q)^{-1}u^{-1}\delta\),
with the latter minus the former, is
\((Q_v+\Phi_v\lambda)S_v\) times the first \(i\) Gram rows.
Indeed the shift coming from \(u^{-1}\) is common to the two
frames as their bottom rows are identical.

The majorants of Proposition~\ref{prop:schurreal692} can thus
be taken the same throughout the fiber over \(\lambda\).
Lemma~\ref{lem:orbitfib693} bounds the remaining cutoff sum
there by \(\tau\), also at \(uT(Q)y\). Summing over all
\(\lambda\), applying Theorem~\ref{thm:schurmean692}, and
then summing over the at most \(i_\gamma^k\) sets yields
\[
 \frac1{\mathcal V_i}\int_{\mathcal P} F_\gamma(uT(Q)y)\,dQ
 \le C_\gamma,\qquad
 C_\gamma=\frac{\tau^2}{\mathcal V_i}
       \left[\left(\frac{b^{13}}{M^2}\right)^{if}
                        \frac{N^{ef}i_\gamma}{P_j}\right]^k .
\]
The proof of the translation mean uses only the Gram value,
the flags, and the real base frame. It therefore applies to the
base \(u^{-1}\delta_0\) even when \(u\) is irrational.

Integrate first on \(U_\infty^0\)-orbits in \(G_\infty\) using
Haar disintegration. On such an orbit the reduced indicator,
the length cut and \(\|W^\flat\|^2\) are constant. The remaining
cutoff \(p_c\) integrates an invariant function by the quotient
measure on \(\Gamma_c^0\backslash U_\infty^0\). Right
multiplication by each \(T(Q)\) preserves this measure,
since a unipotent group is unimodular. Averaging this identity
over \(Q\in\mathcal P\) and using the displayed bound therefore gives
\[
 \int_{G_\infty}\beta_c^\Box(y)\|W^\flat(y)\|^2F_\gamma(y)\,dy
 \le C_\gamma\int_{G_\infty}\beta_c(y)\|W^\flat(y)\|^2\,dy
 \le C_\gamma v_U d_B .
\]
For the last inequality integrate first over \(U_\infty\).
The \(\Gamma_c\)-translates of \(\beta_c\) sum to one; its
integral on each such orbit is \(v_U\). The remaining quotient
norm factors over the places. All the changes of order apply
to nonnegative functions by Tonelli's theorem.

Apply Cauchy--Schwarz on the product of counting measure on
the orbit and Haar measure on \(G_\infty\). In one squared
factor insert \(\prod_v\nu_{\delta_v}(y_v)^{-s_0}\), and in the
other insert its reciprocal. Use the measure with density
\(\beta_c^\Box(y)\beta_{c'}^\Box(\delta^{-1}y)\); the cut on
both vectors permits these boxed cutoffs in \(A_\gamma\).
The first squared factor is the integral just estimated.
For the other change variables by \(y\mapsto\delta^{-1}y\)
and invert the orbit. The projected volume ratio is then
inverted by its determinant definition. The inverse has Gram
matrix \(C^*\), the same \(P_j,j\), and row index \(i'_\gamma\).
Thus \(A_\gamma^2\le(v_Ud_B)^2 C_\gamma C_{\gamma^{-1}}\).
Taking the square root proves the formula.
\end{proof}

\begin{corollary}[The squared contribution of all positive closed ranks]
\leavevmode\label{cor:closedmoment693}
In each agreeing residue block \((c,s,c')\), sum \(A_\gamma\)
once per surviving orbit with \(1\le j\le a\), and denote that
sum by \(A_{\rm cl}(c,s,c')\). At all sufficiently large nodes,
\[
       \sum_{c,c',\,s\ {\rm agreeing}}
                \left(\frac{A_{\rm cl}(c,s,c')}{v_U}\right)^2
         \le \tau^{-40}|B_{\rm res}|\,|\mathcal C|\,d_B .
\]
This controls also the squared norms of the signed sums on
these cells, by the triangle inequality.
\end{corollary}
\begin{proof}
Only near orbits can contribute, since both frames belong to
the reduced length cut wherever the integrand is positive.
The common lattice and gauge of Theorem~\ref{thm:allpair689}
do not change under independent level moves on either frame:
left moves preserve the positive component of every lift with
zero bottom component, changing its top component by an
integral multiple of \(M\); right moves have the identical
property on the second lift. Thus the pair estimate applies
to representatives of these orbits. Alternatively choose
the near representatives when first expanding the square.
For \(j\le a\) we have \(|{\rm N}_{E/\mathbb Q}d_*|=N^{2f}P_j^2\)
and \(\Delta_\gamma=0\). The bracket of the theorem becomes
\[
        \left(\frac{b^{13}}{M}\right)^{if}
                          \frac{N^{f/2}}{\sqrt{G_\gamma}} .
\]
Moreover \(c_E\ge1\): the relative discriminant formula for
\(E/\mathbb Q(\mathrm i)\) gives \(|{\rm disc}(E)|\ge4^f\).
In particular \(\mathcal V_i\ge1\). For two closed terms
of ranks \(i,i'\ge1\), the positive corank \(\ell\) satisfies
\(\ell\le c^+\) in Theorem~\ref{thm:allpair689}. Its estimate
therefore gives
\[
 \frac{A_\gamma A_\eta}{(v_Ud_B)^2}
 \le \tau^4\, b^{-14kf} N^{kf}
           (2/M)^{\ell kf}d(L_{12})^{-k},
\]
using \(M\ge b^{20}\). By Theorem~\ref{thm:paircount691}
the sum over all ordered pairs is at most
\[
 |B_{\rm res}|\,|\mathcal C|\,d_B^2\,\tau^6 b^{-14kf}
          \sum_{\ell\ge0}[\tau(2/M)^{kf}]^\ell
\]
after division by \(v_U^2\). The series is at most two for
large nodes. Finally \(d_B\le b^{2kf}\). To see this on the
actual fields, the chosen diagram for the source at \(v\)
has only columns of heights one and \(k\); the number of
single boxes is \(t_v=xb\) at the designated places and zero
elsewhere. Removing columns of full height removes only a
determinant twist. The remaining symmetric power has a basis
of degree-\(t_v\) monomials on \(k\) variables, so its
dimension is at most \((t_v+1)^k\le b^{2k}\), since \(b\)
exceeds every fixed power of \(x\). Multiply over \(v\).
As \(b^{-14kf}=\tau^{-56}\) and \(d_B\le\tau^8\), the
asserted inequality follows from \(2\tau^{-42}\le\tau^{-40}\).
\end{proof}

The rank-zero terms, the open terms for which the
discriminant exceeds the minor bound, and the removal of
the Bessel length cut still enter the full second moment.


\subsection{Compact rotation of the terminal vector in the open cell}
\label{sec:terminalrotation694}

In open rank the last column of the real Schur estimate lies
in a complex unit ball. We integrate its contribution in the
direction of the positive middle space. This integration takes
place after the arithmetic unipotent average, so it requires
no compact rotation of the integral lattice.

\begin{lemma}[The angular integral of the terminal column]\leavevmode
\label{lem:terminalrotation694}
At one real place in open rank use the notation of
Lemma~\ref{lem:negdet692}. Write \(d=D_{nn}\),
\(v_0=gF B^{-1}e_n^{\rm std}\), \(q=Xv_0\) and
\(x=X\omega_n/d\). Then
\[
 t^2:=\|q\|^2
     =\frac{\Phi_v(\Delta_\gamma)}{|\Phi_v(D_n)|^2}\le1,\qquad
 s^2:=\|x\|^2=1-|d|^{-2}<1 .
\]
As the point \(y\) varies to \(hy\), \(h\in{\rm U}(X_0)\), the
quantity \(D\) is fixed, \(x\) varies to \(hx\), and
the factor \(|d|^2/(1+\|w\|^2)\) in the integral of the
majorant becomes \((1+t^2-2\Re\langle q,hx\rangle)^{-1}\).
For probability Haar \(dh\) and
\[
 c=\tfrac12\bigl(1+t^2+\sqrt{(1+t^2)^2-4t^2s^2}\bigr),
 \qquad r_0=ts/c ,
\]
its \(k\)-th power has the exact mean
\[
 \int_{{\rm U}(X_0)}
       (1+t^2-2\Re\langle q,hx\rangle)^{-k}\,dh
       =\frac{c^{-k}}{1-r_0^2}\le |d|^2 .
\]
\end{lemma}
\begin{proof}
In open rank \(B=F^*HgF\). Thus \(F^*Hv_0=e_n^{\rm std}\).
The bottom components of \(v_0\) vanish and its normalized
terminal component is \(-1\). Its squared Hermitian norm
can also be computed in the \(gF\) frame. It is
\(-|(B^{-1})_{nn}|^2\), since only the last column of that
frame has nonzero norm. The cofactor formula and the scale
of the terminal vector give
\[
            |(B^{-1})_{nn}|=
                 \frac{N|\Phi_v(D_a)|}{|\Phi_v(D_n)|}.
\]
Equating the two norms gives the formula for \(t^2\).
The subframe is the whole frame, and \(F^*H\omega_n=d e_n^{\rm std}\).
The same computation using \(h(\omega_n,\omega_n)=-1\)
gives the formula for \(s^2\). The last column of \(\omega'\)
is \(dv_0\), hence \(w=d(q-x)\). Using \(1-s^2=|d|^{-2}\)
now gives the indicated denominator. The negative basis
construction commutes with \(h\), as \(h\) preserves both the
metric and the ordered frame. This proves the transformation
law for \(D,x\).

For \(st>0\), the unit vector in the direction of \(hx\) is
uniform on the sphere in \(\mathbb R^{2k}\). Indeed the
push-forward of probability Haar is the unique rotation-invariant
probability under the transitive unitary action, and normalized
Euclidean surface measure has this invariance. Choosing the
axis in the direction of \(q\) makes the denominator equal to
\(c(1+r_0^2-2r_0 u_1)\), where \(u_1\) is the real first
coordinate on this sphere. The identities
\(c(1+r_0^2)=1+t^2\), \(cr_0=ts\), and \(c\ge1\) follow from
the quadratic formula and \(t,s\le1\). In particular
\(r_0\le s<1\).

For completeness, at dimension \(m=2k>2\) the mean of
\(\|z-u\|^{2-m}\) over surface probability on the unit sphere
is one for \(\|z\|<1\). It is a radial harmonic function there,
by differentiation on compact subballs, and its value at zero
is one. The radial harmonic differential equation and regularity
at the origin force it to be constant. Differentiate at
\(z=r_0 e_1\) to obtain
\(\int (r_0-u_1)\|r_0 e_1-u\|^{-m}\,du=0\).
On multiplying \(\|r_0 e_1-u\|^{-m}\) by
\(1+r_0^2-2r_0u_1\) and integrating, the result is one.
Consequently its mean without this multiplier is
\((1-r_0^2)^{-1}\). This proves the formula. For \(st=0\)
the denominator is constant and gives the same value directly.
Finally \(c\ge1\), \(r_0\le s\) give the bound.
\end{proof}

\begin{theorem}[A compact average in the orbital integral]\leavevmode
\label{thm:rotatedorbit694}
For every near surviving orbit of rank \(j\ge1\) on the
interval of Theorem~\ref{thm:orbitschur693} one has
\[
 \frac{A_\gamma}{v_U d_B}
 \le \frac{\tau^2}{\mathcal V_i}
       \left[\frac{2^{(i+1)f}}{M^{2if}}\,
                  \frac{N^{(j-i)f}\sqrt{i_\gamma i'_\gamma}}{P_j}
       \right]^k .
\]
\end{theorem}
\begin{proof}
In the exact integral in Proposition~\ref{prop:schurreal692}
retain the factor
\(R_v(y)^k=(|d|^2/(1+\|w\|^2))^k\). The factors
\(\pi^{ki}\prod_{t=0}^{i-1}t!/(k+t)!\) are at most one,
and \(|\det D|^{2i}\le\rho^{2pi}\). After the change of
coordinates in Theorem~\ref{thm:schurmean692} the local bound
before estimating \(R_v\) is
\[
 \left[\frac{(2\rho^2)^pN^j}{|\Phi_v(D_j)|}\right]^i
 \left[\frac{2^p N^{j-i}}{M^{2i}|\Phi_v(D_j)|}\right]^k R_v(y)^k .
\]
For \(j\le a\) one has \(R_v=1\) by the common terminal vector.
In open rank \(q=X(gF B^{-1}e_n^{\rm std})\) is independent
of the orbit representative and of the shifts in the positive
columns used in Lemma~\ref{lem:orbitfib693}. In fact left
multiplication by a pointwise unipotent preserves the positive
component of any vector with zero bottom component, and right
multiplication fixes the input frame. The matrix \(B\) is
unchanged. Also \(D\) and \(x\) are constant on
\(U_\infty^0\)-orbits of \(y\), by the same equivariance of
the triangular basis and the zero bottom component of
\(\omega_n\). Thus at a fixed \(y\) the factor
\(\prod_v R_v(y_v)^k\) is common to every fiber in the forward
Schur integral.

To justify averaging this factor, first perform the integration
over \(\Gamma_c^0\backslash U_\infty^0\) exactly as in
Theorem~\ref{thm:orbitschur693}. The right translations by
\(T(Q)\) give the same bound with the displayed local factors.
After this step the remaining measure on
\(U_\infty^0\backslash G_\infty\) is quotient Haar times
\[
  \operatorname{vol}(\Gamma_c^0\backslash U_\infty^0)\,
         R_{\rm red}(y)\,{\bf1}_{\rm cut}(y)\,\|W^\flat(y)\|^2 .
\]
Indeed integration of \(p_c(u y)\) over \(u\) gives that
covolume for each \(y\), by its partition identity. This
covolume is finite by the same rational unipotent argument
used for \(v_U\). The displayed density is invariant under
left \(\prod_v{\rm U}(X_{0,v})\): the compact group fixes the
frame, preserves the length, and acts unitarily on the
values of the Bessel function. It preserves the quotient
Haar measure since it normalizes \(U_\infty^0\) and preserves
its Haar measure. For the latter assertion its modulus is a
continuous homomorphism from a compact group to the positive
reals, hence is one. We may therefore replace the factor
\(\prod_v R_v(y_v)^k\) in the integral by its compact mean.
Tonelli gives the product of the local means.

In open rank the lemma bounds each mean by \(|d|^2\le\theta^2\)
on the cut. The product of the first displayed brackets raised
to \(i\), times \(\theta^{2f}\), is at most
\(b^{(135n^3+12n+4)f}\le\tau\) at these nodes, since
\(P_j\ge1\). The same estimate applies in closed rank with
the terminal factor one. The cutoff multiplicity costs another
\(\tau\), the coset count gives \(i_\gamma^k\), and the
translation mean divides by \(\mathcal V_i\). The integral
of the remaining squared Bessel norm against \(\beta_c\)
is at most \(v_Ud_B\). This proves that the forward squared
factor divided by \(v_Ud_B\) is at most
\[
       \frac{\tau^2}{\mathcal V_i}
       \left[\frac{2^{pf}N^{(j-i)f}i_\gamma}
                         {M^{2if}P_j}\right]^k .
\]
Apply the same result to the inverse orbit and take the
geometric mean as in Theorem~\ref{thm:orbitschur693}.
\end{proof}

\begin{corollary}[Open terms with bounded product of normal lengths]
\leavevmode\label{cor:normalrange694}
In each agreeing block sum \(A_\gamma\) over open surviving
orbits for which the normalized real normal vectors satisfy
\[
       \prod_v\|q_v\|
       =\frac{\sqrt{{\rm N}_{F/\mathbb Q}\Delta_\gamma}}{P_n}
       \le \frac{(M/b)^{af}}{N^{f/2}} .
\]
Call this sum \(A_*(c,s,c')\). Then
\[
     \sum_{c,c',\,s\ {\rm agreeing}}
       \left(\frac{A_*(c,s,c')}{v_U}\right)^2
       \le \tau^{-3a}|B_{\rm res}|\,|\mathcal C|\,d_B .
\]
\end{corollary}
\begin{proof}
For an open orbit \(i=a,\ i_\gamma=i'_\gamma=1\). Write
\(G_0=M^{2af}P_n^2/N^f\). Then the gauge in
Theorem~\ref{thm:allpair689} satisfies
\[
       G_\gamma/G_0
          =\max\{1,N^f\,{\rm N}_{F/\mathbb Q}\Delta_\gamma/P_n^2\}.
\]
Our condition gives \(\sqrt{G_\gamma/G_0}\le (M/b)^{af}\).
By the theorem, \(\mathcal V_a\ge1\), and
\(2^{a+1}\le b^{a/4}\) for \(b\ge 2^8,\ a\ge1\), it follows that
\[
        \frac{A_\gamma}{v_Ud_B}
         \le \tau^2
            \left[b^{-3af/4}\frac{N^{f/2}}{\sqrt{G_\gamma}}\right]^k.
\]
Two open terms in the same block can have positive corank
\(\ell=c^++1\); the estimate of Theorem~\ref{thm:allpair689}
always gives \(c^+\ge\max(0,\ell-1)\). Applying that theorem
to the product of these bounds leaves the factor
\((2/M)^{\max(0,\ell-1)kf}d(L_{12})^{-k}\). The factor
\(N^{kf}\) cancels upon summing by Theorem~\ref{thm:paircount691}.
Thus the normalized second moment is at most
\[
 |B_{\rm res}|\,|\mathcal C|\,d_B^2\,
 \tau^{6-6a}
   \sum_{\ell\ge0}\tau^\ell(2/M)^{\max(0,\ell-1)kf}.
\]
The sum is \(1+\tau/(1-\tau(2/M)^{kf})\le3\tau\).
Since \(d_B\le\tau^8\) by Corollary~\ref{cor:closedmoment693},
the result follows from \(3\tau^{15-6a}\le\tau^{-3a}\).
This inequality holds at all sufficiently large nodes.
\end{proof}

For larger products of normal lengths the compact mean still
gives Theorem~\ref{thm:rotatedorbit694}. The common-lattice
estimate then involves the larger discriminant term in
\(G_\gamma\); bounding those blocks on the scale of the full
period requires further information.

\section{Stationary real coefficients and truncation}\label{nav:stationary-cutoff}

\subsection{A stationary position at every real place}
\label{sec:effectivestationary695}

Fix a node \(V_0\le x\le V_0^{3/2}\) and a real place of the
field \(F=E^+\) used in Theorem~\ref{thm:rotatedorbit694}.
We construct a translate for the local period \(m\) of
Proposition~\ref{prop:averagedschur684}. Use normalized real
coordinates as in Corollary~\ref{cor:realmarked682}, so the
last vector \(z\) has Hermitian norm \(-1\) and the character
on the frame is \(\exp(\sigma-\bar\sigma)\), where \(\sigma\)
is the sum of its entries immediately above the diagonal. Write
\(a=n-1,\ r=k+2n-1\). The basis of the normalized space consists
of \(a\) isotropic vectors \(f_j\), the positive middle space
\(X_0\) of dimension \(k\), the line of \(z\), and the dual
vectors \(f_j^\vee\). It is Euclidean orthonormal. The
Hermitian form has matrix \(H\) with \(H^2=1\) and
\(\langle f_i,H f_j^\vee\rangle=\delta_{ij}\).

\begin{proposition}[The compact weights at the given node]\leavevmode
\label{prop:stationaryweights695}
Order the coordinates by decreasing Harish-Chandra parameter.
The weight operator \(L\) of the highest vector in the lowest
compact type is zero on the first \(k-1\) positive coordinates.
On the remaining \(2n\) coordinates the signs are
\(+,-,\ldots,+,-\), and the weights are distinct integers
\(\mu_0>\cdots>\mu_{2n-1}\) in \([-3B,-B]\). Here
\(B=\kappa+r^8\), where \(\kappa\) is the weight at this place.
The compact value representation has weight operator
\(L'=\operatorname{diag}(0,\ldots,0,t)\) after a permutation,
where \(t=xb\) at the designated \(a\)-places and \(t=0\)
elsewhere. For all sufficiently large nodes,
\(D_0:=3B+t\le b^2\).
\end{proposition}
\begin{proof}
The undoubled row parameter has entries \(k/2-1-i\) for
\(0\le i<k-1\). At \(i=k-1+j\) its entry is
\(-\kappa-2(j+1)r^8\), \(0\le j<2n\), by the shifted weight
formula \(\kappa_j=\kappa+(2j+1)r^8\). The first block
has positive sign and the last block alternates starting with
a positive sign. The prescribed scales have \(r=O(x^5)\),
\(2n\le 2\sqrt{x}\), \(\log b\ge(\log x)^6\), and
\(\log x_1\le 10^{-6}\log b\). The weight \(\kappa\) is
either \(b\) or \(xb/x_1\), with \(x_1\mid b\). Thus
\(2(2n)r^8+r=o(\kappa)\), and the row is in decreasing
order. Also \(\kappa\le xb\), \(t\le xb\), whence \(D_0\le b^2\).

For signs \(\epsilon_i\) and row entries \(\lambda_i\), the
formula \(L=\lambda+\rho_{\rm n}-\rho_{\rm c}\) recalled in
Lemma~\ref{lem:casimirs684} gives twice the difference between
\(\lambda_i\) and the compact weight at position \(i\) as
\[
        \sum_{j>i}\epsilon_i\epsilon_j
                 -\sum_{j<i}\epsilon_i\epsilon_j .
\]
Indeed a compact root contributes with the opposite sign to a
noncompact root. The alternating block has sign sum zero, so
the difference on the first block is \((k-2-2i)/2=\lambda_i\).
On the last block the correction has absolute value at most
\((r-1)/2\). Its entries are integers since \(r\) is odd.
The gap \(2r^8\) gives strict decrease, the first entry is at
most \(-\kappa-r^8\), and the last is at least
\(-\kappa-4nr^8-(r-1)/2\ge-3B\). Finally the diagram defining
the compact Bessel value space counts columns strictly below
height \(k\) with last weight zero. Its only columns at those
heights have height one, with the indicated multiplicity \(t\).
Thus it is the symmetric power with the stated highest weight.
\end{proof}

\begin{theorem}[A stationary flag with a uniform length bound]\leavevmode
\label{thm:effectivestationary695}
Write the defining space in sign coordinates and let \(p\) be
a Euclidean unitary identification of its sign form with \(H\)
in the normalized space. There are independent vectors
\(v_0,\ldots,v_a\), a positive isometric embedding
\(I:\mathbb C^k\to\mathbb C^r\), and \(g\in\mathrm U(H)\)
such that
\[
 \begin{gathered}
  gpv_j=f_{j+1}\ (0\le j<a),\qquad gpv_a=z,\qquad
  gpI=\text{the prescribed inclusion of }X_0,\\
  Lv_j=v_{j+1}\ (j<a),\qquad
  LI-IL'\ \text{has image in }\mathbb Cv_0,\qquad
  \inf_{s\in S}\bigl(1+\|sg\|_{\rm op}\bigr)\le b^{5n}.
 \end{gathered}
\]
Here \(S\) is the subgroup acting by upper unitriangular
matrices on the ordered frame \((f_1,\ldots,f_a,z)\).
\end{theorem}
\begin{proof}
Set \(P(T)=\prod_{j=0}^{2n-1}(T-\mu_j)\), \(Q(T)=T-t\).
For a polynomial \(A\) define \([A]\) to be zero on the first
\(k-1\) coordinates and to have tail coordinates
\[
             [A]_j=\frac{A(\mu_j)}
                 {\sqrt{|P'(\mu_j)(\mu_j-t)|}} .
\]
The sign of \(P'(\mu_j)\) is \((-1)^j\), and \(\mu_j<t\).
Thus the Hermitian product of two such vectors is
\[
 -\sum_j\frac{\overline{A(\mu_j)}\,C(\mu_j)}
                   {P'(\mu_j)(\mu_j-t)} .
\]
Put \(v_h=[T^hQ]\) for \(0\le h\le 2a\) and
\(u=\sqrt{P(t)}[1]\). Lagrange interpolation shows that
\(\sum_j R(\mu_j)/P'(\mu_j)\) is the coefficient of
\(T^{2n-1}\) in \(R\) whenever \(\deg R\le 2n-1\).
Apply this to \(R=T^{h+l}Q\). For \(h+l<2a\) the Hermitian
product of \(v_h,v_l\) vanishes, and for \(h+l=2a\) it is
\(-1\). Applying the formula to \(T^h\) gives
\(\langle u,v_h\rangle_{\epsilon}=0\) for \(h\le2a\).
Partial fractions give
\(\sum_j [P'(\mu_j)(t-\mu_j)]^{-1}=1/P(t)\), so \(u\)
has Hermitian norm one. The first \(a+1\) vectors \(v_h\)
are independent since a polynomial of degree at most \(a\)
vanishing at all \(2n\) nodes is zero.

Let \(I\) take the first \(k-1\) standard vectors to themselves
and the last one to \(u\). Multiplication by \(L\) acts as
multiplication by \(T\) on \([A]\). In particular
\(Lu=tu+\sqrt{P(t)}v_0\), which proves the weight assertions.
The Gram matrices now agree on the indicated frames and middle
spaces. To extend their identification to an isometry, choose
vectors dual to the independent isotropic arms in the
orthogonal complement of \(I(\mathbb C^k)\oplus\mathbb Cv_a\).
This complement is nondegenerate. Subtract from the dual
vectors the arms times half the Hermitian Gram matrix of the
duals. Their pairings with the arms stay equal to the identity
and their mutual pairings become zero. The resulting bases
give \(g\) with the required exact images.

We bound the length using a representative of \(Sg\). Since
the nodes and \(t\) are separated integers, \(|P'(\mu_j)|\ge1\)
and \(1\le t-\mu_j\le D_0\). Thus
\(|(v_h)_j|\le D_0^{h+1/2}\) and \(|u_j|\le D_0^n\).
Put \(K_0=\sqrt r\, D_0^{2n}\). Apply Euclidean Gram--Schmidt
without rescaling to \(p v_0,\ldots,p v_a\) and call the
resulting columns \(b_h\). Their lengths are at most \(K_0\).
The Hermitian product of \(b_h\) with \(p v_{2a-h}\) is still
\(-1\), since all earlier columns pair to zero with it.
Consequently \(\|b_h\|\ge K_0^{-1}\). For \(h<a\) their
isotropic duals \(H b_h/\|b_h\|^2\) have lengths at most
\(K_0\). The arms and these duals are pairwise Euclidean
orthogonal, by Gram--Schmidt, \(H^2=1\), and isotropy of the
span of the arms. Their joint span is \(H\)-stable. The
terminal vector \(b_a\) lies in its Euclidean orthogonal
complement, by Gram--Schmidt and Hermitian orthogonality
to the arms. Its Hermitian norm is \(-1\).

The complement has signature \((k,1)\). Write \(b_a=b_++b_-\)
in the positive and negative eigenspaces of \(H\). A unitary
identification of those eigenspaces followed by the real
two-dimensional boost with diagonal entries \(\|b_-\|\) and
off-diagonal entries \(\|b_+\|\) sends the standard negative
unit vector to \(b_a\). Choose the directions of the two
components as axes; if \(b_+=0\) only a unitary identification
is needed. Since \(\|b_-\|^2-\|b_+\|^2=1\), this is an isometry
of norm at most \(2K_0\), extended by the identity in the
remaining positive directions. Combining it with the map of
the standard arms and duals to the orthogonal columns just
constructed gives an isometry \(g_1^{-1}\) of norm at most
\(2K_0\). Its inverse has the same norm since \(H\) is a
Hermitian involution. The frames \(g^{-1}(f_1,\ldots,f_a,z)\)
and \(g_1^{-1}(f_1,\ldots,f_a,z)\) differ by an upper
unitriangular change of columns. Hence \(g_1g^{-1}\in S\).
Finally \(1+2K_0\le b^{5n}\) by \(D_0\le b^2\) and the stated
growth of the ranks. This proves the uniform assertion.
\end{proof}

\begin{corollary}[An exact trace identity for the period phase]\leavevmode
\label{cor:exactphase695}
For \(s\in S\), set \(Y=p^*g^{-1}sgp-1\),
and \(Y'=h-1\), where \(h\) is the induced action on \(X_0\).
Then
\[
  \operatorname{tr}Y=\operatorname{tr}Y',\qquad
  \operatorname{tr}(LY)-\operatorname{tr}(L'Y')
        =\sigma(s)-\overline{\sigma(s)} .
\]
In particular the integrand for \(m\) in
Proposition~\ref{prop:averagedschur684}, at the translate
\(g\) and the indicated unit highest weight vectors, has
vanishing first derivative on \(S\) at the identity.
\end{corollary}
\begin{proof}
Use the three-block Witt basis \(A,W,A^\vee\), with \(A\)
the arms \(v_0,\ldots,v_{a-1}\), \(W=I(\mathbb C^k)
\oplus\mathbb Cv_a\), and the isotropic duals constructed
in the theorem. The middle form is \(H_W=\operatorname{diag}(1_k,-1)\).
Both \(Y\) and its Hermitian adjoint \(Y^\dagger\), which equals
\(p^*g^{-1}s^{-1}gp-1\), are upper block triangular. Indeed
the isometries preserve \(A\) and \(A^\perp\). The action on
the middle quotient fixes the terminal vector, and therefore
restricts to \(h\) on its orthogonal complement. The \(AA\)
blocks of \(Y,Y^\dagger\) are strictly upper triangular and the
middle block of \(Y\) is \(\operatorname{diag}(Y',0)\). The
last diagonal block of \(Y\) is the conjugate transpose of
the first block of \(Y^\dagger\). This gives the ordinary trace.

For the weighted trace, \(L_{AA}\) is the shift \(v_j\mapsto
v_{j+1}\) up to the last arm, whose image is \(v_a\in W\).
The only entry of \(L_{WA}\) is this last map. Also
\(L_{A^\vee A}=0\) and the compression of \(L\) to the
\(I\)-part of \(W\) is \(L'\). Thus the contractions
\(\operatorname{tr}(L_{AA}Y_{AA})\) and
\(\operatorname{tr}(L_{WA}Y_{AW})\) add to \(\sigma(s)\),
including the terminal simple root. Since \(L=L^\dagger\),
the corresponding contractions on the dual blocks are the
complex conjugates of the same contractions with \(Y^\dagger\).
Explicitly,
\(L_{A^\vee A^\vee}=L_{AA}^*\),
\(L_{A^\vee W}=L_{WA}^*H_W\), and
\(Y_{WA^\vee}=H_W^{-1}((Y^\dagger)_{AW})^*\).
They give \(\overline{\sigma(s^{-1})}\).
All remaining contractions vanish except the middle trace
\(\operatorname{tr}(L'Y')\). The entries immediately above the
diagonal add under multiplication of upper unitriangular
matrices, so \(\sigma(s^{-1})=-\sigma(s)\).

For the derivative, the normalized matrix coefficient of a
torus weight vector of weight \(L\) has derivative
\(Z\mapsto\operatorname{tr}(LZ)\) on the unitary Lie algebra.
On the diagonal torus this follows by differentiating its
character. The complex-linear extension vanishes on
off-diagonal matrix units by torus conjugation, since the
coefficient is invariant under that conjugation. The same
argument applies to \(L'\). Differentiating the trace identity
now cancels both coefficient derivatives against the derivative
of the inverse simple-root character.
\end{proof}

The point and the trace identity hold at every place with the
same upper bound on the length. To normalize the global test
vector one still needs a lower bound for the integral \(m\);
this requires controlling the coefficient away from its
stationary region.


\subsection{Localization and the positive contribution near the stationary point}
\label{sec:positivecore696}

At the translate of Theorem~\ref{thm:effectivestationary695}
write \(G_s=p^*g^{-1}sgp,\ h_s=\operatorname{pr}_{X_0}s\) and
\(J_s=G_s\oplus h_s\). Put \(Y=G_s-1,\ Y'=h_s-1\), and let
\(C_s\) be the matrix retaining the entries of \(G_s\)
between coordinates of opposite signs and zero elsewhere.
Let \(\varphi(s)=\sigma(s)-\overline{\sigma(s)}\), and let
\[
 \Theta(s)=\langle w,\pi(g^{-1}sg)w\rangle
              \overline{\langle e,\tau_0(h_s)e\rangle}
                                  \exp(-\varphi(s)).
\]
Here \(\tau_0\) is the compact value representation and \(w,e\)
are the unit highest vectors used above. The unsigned norm
\(\|\cdot\|_2\) on matrices in this subsection is the
Hilbert--Schmidt norm. The integral of \(\Theta\) on \(S\)
is \(m\) in Proposition~\ref{prop:averagedschur684}.

\begin{lemma}[The phase in exponential coordinates]\leavevmode
\label{lem:logphase696}
If \(\|J_s-1\|_{\rm op}<1\), then
\[
          \operatorname{tr}(L\log G_s)-
          \operatorname{tr}(L'\log h_s)=\varphi(s)
\]
for the power-series logarithms at the identity.
\end{lemma}
\begin{proof}
Let \(D(A,A')=\operatorname{tr}(LA)-\operatorname{tr}(L'A')\).
By Corollary~\ref{cor:exactphase695} and additivity of
\(\varphi\), \(D(G_s^j-1,h_s^j-1)=j\varphi(s)\) for
each integer \(j\ge0\). For every \(l\ge2\), expand
\((G_s-1)^l,(h_s-1)^l\) by the binomial theorem. The
coefficients sum to zero, so one may replace \(G_s^j,h_s^j\)
by \(G_s^j-1,h_s^j-1\). The \(l\)-th finite difference of the
linear function \(j\varphi(s)\) vanishes. This gives
\(D(Y^l,(Y')^l)=0\). Termwise application of \(D\) to the
two convergent series proves the formula.
\end{proof}

\begin{theorem}[Localization modulo the common compact subgroup]\leavevmode
\label{thm:angularlocalization696}
The common first \(k-1\) positive coordinates define a subgroup
\(K_c=\mathrm U(k-1)\) of \(S\) acting by the same unitary matrix
on those coordinates in both factors of \(J\). The function
\(\Theta\) is invariant under left translation by \(K_c\). If
\(\|C_s\|_2=c\le1/6\), there is \(u\in K_c\) with
\[
       \|G_{us}-1\|_2^2\le 13c^2,\qquad
       \|h_{us}-1\|_2^2\le 11c^2 .
\]
In particular \(K_c\) consists precisely of the elements
of \(S\) whose first component is in the maximal compact
subgroup of the sign model.
\end{theorem}
\begin{proof}
An element of \(\mathrm U(k-1)\) on the common coordinates
extended by the identity on their complement fixes every
\(v_h\) and has the stated action on \(I\), so it defines a
member of \(S\). It fixes \(w\): the highest weight is zero
on the common block, and each adjacent compact root there has
highest weight difference zero. Its root \(\mathrm{SU}(2)\)
therefore fixes the highest vector by the weight string
formula. These root groups and the diagonal torus generate
the block. It fixes \(e\) in the symmetric power, and
\(\sigma\) vanishes there since the frame is fixed pointwise.
This gives the asserted invariance.

Let \(P_c\) project to the common coordinates and put
\(P_o=1-P_c\), \(\nu_i=-L_{ii}\). For each row of the
sign-unitary matrix \(G_s\), preservation of the sign form
gives the identity
\[
   2\Re Y_{ii}
      =2\sum_j |(C_s)_{ij}|^2-\sum_j |Y_{ij}|^2 .
\]
In fact \(\sum_j\epsilon_i\epsilon_j|(G_s)_{ij}|^2=1\);
expanding the diagonal entry \(1+Y_{ii}\) gives the identity.
The same formula holds for columns. By the exact phase
identity,
\(\Re\operatorname{tr}(\nu Y)
=-\Re\operatorname{tr}(L'Y')=t(1-\Re(h_s)_{kk})\ge0\).
Multiply the row identities by \(\nu_i\) and sum. Since
\(B\le\nu_i\le3B\) off the common block and \(\nu_i=0\) on
it, the row and column estimates give
\[
       \|P_oY\|_2^2\le6c^2,\qquad \|YP_o\|_2^2\le6c^2 .
\]
Also the unweighted row identities summed for both factors
and the ordinary trace identity give
\(\|Y'\|_2^2=\|Y\|_2^2-2c^2\).

Write \(A=(G_s)_{cc}\). From the \(cc\)-block of
\(G_s^*\epsilon G_s=\epsilon\) we have
\(\|A^*A-1\|_2\le\|P_oYP_c\|_2^2\le6c^2\), using
\(\|D^*\epsilon D\|_2\le\|D\|_2^2\).
Let \(A=U|A|\) be its polar decomposition. The bound on
\(c\) implies invertibility, and each singular value satisfies
\(|s_j-1|\le|s_j^2-1|\). Thus
\(\||A|-1\|_2\le6c^2\). Left translate by the member of
\(K_c\) acting by \(U^*\). The new \(cc\)-block is \(|A|\)
and the noncompact part has unchanged norm. The preceding
row and column estimates apply at the new point. Bounding
the entries outside \(cc\) by the sum of those two estimates
gives a squared norm at most \(12c^2+36c^4\le13c^2\).
The unweighted trace identity gives the second bound.
Taking \(c=0\) shows that the indicated element is the
inverse of \(s\), since \(s\mapsto G_s\) is faithful.
\end{proof}

\begin{theorem}[Positive mass in a neighborhood of the identity]\leavevmode
\label{thm:positivecore696}
Put \(d_*=r+k\), \(\beta=b^3\), \(\eta=(200\beta^5)^{-1}\).
With Haar measure on the actual Bessel subgroup let
\[
    \Omega_\eta=\{s:\|J_s-1\|_{\rm op}\le\eta\},\qquad
    \mathcal B_\beta=\{s:\|J_s\|_{\rm op}\le1+\beta\}.
\]
At every real place, for all sufficiently large nodes,
\[
   \Re\Theta(s)\ge\tfrac12\quad(s\in\Omega_\eta),\qquad
   \int_{\Omega_\eta}\Theta(s)\,ds
        \ge\tfrac12\beta^{-16d_*^2}\operatorname{vol}(\mathcal B_\beta)>0 .
\]
On the product of the \(f\) real groups the integral on
\(\prod_v\Omega_{\eta,v}\) is at least
\(2^{-f}\beta^{-16fd_*^2}\prod_v\operatorname{vol}(\mathcal B_{\beta,v})\).
\end{theorem}
\begin{proof}
For \(s\in\Omega_\eta\), \(Z=\log G_s\) has
\(\|Z\|_{\rm op}\le2\eta\) and \(\|Z\|_2\le2\sqrt r\,\eta\).
It is in the real Lie algebra: conjugation of the logarithm
by the sign involution after taking adjoints gives the
logarithm of \(G_s^{-1}\), which equals \(-Z\) on this
neighborhood. In an orthonormal real Lie basis the sum of
the absolute values of its coordinates is at most \(r\|Z\|_2\).
Lemma~\ref{lem:casimirs684} at length one and the unitary path
\(\exp(uZ)\), \(0\le u\le1\), give
\[
  |\langle w,\pi(g^{-1}sg)w\rangle-1|
       \le 8r^2(\max_i|\lambda_i|+3r+2n^2+1)\eta .
\]
Here the sign basis is carried to the normalized group by
\(p\), which preserves the matrix norms. On the symmetric
power the compact coefficient is \((h_s)_{kk}^t\), so its
difference from one has modulus at most \(t\eta\), also for
\(t=0\). Finally
\[
     |\exp(-\varphi(s))-1|\le|\varphi(s)|
                            \le (3rB+t)\eta
\]
by the trace identity and the fact that \(\varphi\) is imaginary.
Each coefficient has modulus at most one. Their product
differs from one by at most the sum of these three errors.
The sum is less than \(1/2\) for all sufficiently large nodes
since \(r=O(x^5)\), \(n\le\sqrt{x}\), \(3B+t\le b^2\),
\(\max_i|\lambda_i|\le 3xb\), and \(\log b\ge(\log x)^6\).

We compare the two volumes by packing in complex matrices
of size \(d_*\) with the operator norm. The graph
\(s\mapsto J_s\) is proper: its first projection is the
conjugate of the inclusion of the closed subgroup \(S\).
For \(R=1+\beta\), take a maximal \(\eta/R\)-separated
set \(J_{s_i}\) in the image of \(\mathcal B_\beta\).
Disjoint ambient balls of radius \(\eta/(2R)\) about them
lie in a ball of radius \(R+\eta/(2R)\). Comparing
Lebesgue volumes in real dimension \(2d_*^2\) bounds their
number by \((1+2R^2/\eta)^{2d_*^2}\le\beta^{16d_*^2}\)
for \(\beta\ge10^4\). Every \(J_s\) in that image is within
\(\eta/R\) of some \(J_{s_i}\). Since a sign-unitary matrix
and its inverse have equal operator norm,
\(\|J_s J_{s_i}^{-1}-1\|_{\rm op}\le\eta\).
The right translates \(\Omega_\eta s_i\) cover
\(\mathcal B_\beta\). Both sets have finite volume, and
each contains a neighborhood of the identity. This proves
the asserted comparison of volumes.

The group \(S\) is unimodular: its unipotent radical is
unimodular and its compact factor preserves Haar on the
radical. The set \(\Omega_\eta\) is inversion invariant
by the sign adjoint identity for \(J_s-1\).
The two unitary coefficients and the character take
conjugate values at inverses, so the integral on
\(\Omega_\eta\) is real. The lower bound follows by
integrating the first assertion, and the product bound
follows by Fubini on the product of these compact sets.
\end{proof}

In the full period the complementary region can contribute
with either sign. The localization theorem describes where
the integrand can have small noncompact displacement; estimating
its phase and decay there and beyond the compact neighborhood
is needed to bound the complete integral.


\subsection{Decay of the period away from the maximal compact subgroup}
\label{sec:noncompacttail697}

We estimate the part of the actual local integral on which
the displacement \(C_s\) of Theorem~\ref{thm:angularlocalization696}
is large. Put \(D=(r+k)^2,\ T=D\log\beta\) and
\(\delta_*=(Tr^{1/4})^{-1}\), with \(\beta=b^3\) as before.
Use coordinates in decreasing order of \(\lambda\). Write
\(\alpha=\lambda+\rho\), \(\mu=\lambda+\rho-2\rho_c\), and
\(K=\mathrm U(p)\times\mathrm U(q)\) for the maximal compact
subgroup of the sign unitary group. For an integer weight \(v\)
and an invertible matrix \(A\) put
\[
 m_v(A)=(\det A)^{v_{\min}}
    \prod_{v_{\min}<j\le v_{\max}}\Delta_{\{i:v_i\ge j\}}(A),
\]
where \(\Delta_B\) denotes the principal minor on \(B\) and
the empty minor is one. We first fix precisely the coefficient
formula used in the estimate.

\begin{proposition}[The compact integral in matrix coordinates]\leavevmode
\label{prop:polarcoefficient697}
Denote the highest-vector coefficient of \(\pi\) in sign
coordinates by \(c_\lambda\). Write \(y=k_0z\) by polar
decomposition, \(k_0\in K,\ z>0\), and normalize Haar on
\(K\) to have mass one. Then
\[
  c_\lambda(y)=\frac{1}{M_\mu}
       \int_K\frac{m_\mu(k_0u)\overline{m_\mu(u)}}
                     {m_{-\alpha}(u^*zu)}\,du,\qquad
  M_\mu=\int_K |m_\mu(u)|^2du
       =(\dim V_\mu)^{-1}.
\]
The germ at the identity of \(c_\lambda\) on the complex
matrix group is
\[
           \mathcal F(Z)=\frac{1}{M_\mu}
             \int_K\frac{|m_\mu(u)|^2}
                           {m_{-\alpha}(u^*Zu)}\,du .
\]
In this germ \(\mathcal F(b_- Z b_+)=\mathcal F(Z)\) for
block diagonal unitriangular \(b_-,b_+\) that are lower and
upper triangular, respectively, in the chosen order.
\end{proposition}
\begin{proof}
Use the compact-Cartan discrete-series realization by
cohomological induction from the Borel for the positive
roots of \(\lambda\). The torus character has weight
\(\lambda-\rho\), an integral weight here, and the good-range
inequalities are \(\lambda_i-\lambda_j>0\) for \(i<j\).
The bottom compact type has weight \(\mu\), as in
Lemma~\ref{lem:casimirs684}. On positive Hermitian matrices
the Flensted-Jensen integral for this compact type is the
average of the conjugates of its highest-line projection
times the inducing character multiplied by the determinant
character of the nilradical. In Cartan polar coordinates its
value at \(k_0z\) acquires the left factor \(\sigma_\mu(k_0)\).
These are the compact-Cartan case of the bottom-layer
coefficient theorem and its duality formula. We use
\href{https://doi.org/10.1112/S0010437X06002508}{B.~Sun,
\emph{Matrix coefficients of cohomologically induced representations},
Compos.~Math.~143 (2007), Theorems~2.1, 6.3 and formula~(26)}.

Here are the characters and projections in that formula.
The determinant character adds \(2\rho\) to the inducing
weight, giving \(\alpha\). For a positive Hermitian matrix
\(A=u^*Zu\), take its unique factorization \(A=n_0 d n_0^*\)
with \(n_0\) upper unitriangular and \(d\) positive diagonal.
It follows from successive Schur complements starting at the
last entry. Products of the last entries of \(d\) are exactly
the trailing principal minors of \(A\). Since \(\alpha\) is
strictly decreasing, \(m_{-\alpha}(A)=\prod_i d_i^{-\alpha_i}\).
Thus the torus factor is \(1/m_{-\alpha}(A)\). The compact
highest coefficient is \(m_\mu(u)\): in either sign block take
the tensor product of the exterior highest vectors for its
columns and the determinant power of its last weight.
The individual coefficients are their principal minors, and
the product of unit highest vectors generates the component
of highest weight \(\mu\). Taking the matrix entry of the
conjugated highest-line projection, and replacing the
integration variable by its inverse, gives the indicated
numerator. Schur orthogonality gives \(M_\mu=1/\dim V_\mu\).
In the polar decomposition the positive factor and the unitary
factor belong to the sign group: the involution
\(y\mapsto\epsilon y^{-*}\epsilon\) fixes \(y\) and preserves
both parts, so uniqueness gives the assertion.

For the local formula consider the operator-valued integral
before taking the highest entry. It extends holomorphically
to \(\|Z-1\|_{\rm op}<1\), since every compressed
\(u^*Zu\) is then invertible and the compact integral is
locally uniformly holomorphic. On positive matrices the
duality formula gives covariance under \(Z\mapsto aZa^*\)
for \(a\in K_{\mathbb C}\), with factors
\(\sigma_\mu(a),\sigma_\mu(a)^*\) on the two sides.
Consequently its germ is covariant under independent left
and right \(K_{\mathbb C}\)-multipliers. Indeed both sides
are holomorphic in the left multiplier, the right multiplier
and the argument; the conditions that the right multiplier
be the adjoint of the left and that the argument be Hermitian
give a totally real set of full real dimension. Uniqueness
of the Taylor expansion gives the identity of germs.
At a nearby real point \(y=k_0z\) this covariance gives
precisely the left factor in the polar formula, so the
asserted scalar germ is the one for \(c_\lambda\).
Upper compact unitriangular matrices fix the highest vector.
Lower ones fix its dual covector on the right of that
covector. The two multipliers therefore leave its matrix
entry unchanged. The identity extends along connected
parameter paths on which the three arguments of the
covariance identity stay in the ball of holomorphy.
\end{proof}

\begin{lemma}[Decay from the opposite-sign parameter gaps]\leavevmode
\label{lem:coefficientdecay697}
Take \(H_0=r^8\) and set \(A_y=(y+\epsilon y\epsilon)/2\),
the block diagonal part of \(y\). For the parameters of
Proposition~\ref{prop:stationaryweights695},
\[
       |c_\lambda(y)|^2|\det A_y|^{H_0}\le1,\qquad
       |c_\lambda(y)|^4\le(1+\|y-A_y\|_2^2)^{-H_0},\qquad
       |c_\lambda(y)|\le (2/\|y\|_{\rm op})^{H_0}.
\]
\end{lemma}
\begin{proof}
First suppose \(y=z>0\). Then \(\epsilon z\epsilon=z^{-1}\),
\(\det z=1\), and \(A_z=(z+z^{-1})/2\) is positive.
Jacobi's complementary-minor identity gives
\(\Delta_B(z)=\Delta_{B^c}(z)\), since \(z^{-1}\) is the
conjugate of \(z\) by a diagonal sign matrix. Set
\(S_j=\{i:\alpha_i\ge j\}\). All coordinates in
\(D_j=S_{j-H_0}\setminus S_j\) have the same sign whenever
that set is nonempty. Indeed opposite signs have
\(\lambda\)-gap at least \(H_0\): consecutive tail entries
are \(2r^8\) apart and the first block is separated from the
negative tail by \(\kappa-O(r)>r^8\). Adding the decreasing
coordinates of \(\rho\) increases the gaps.

For disjoint \(B,C\) the positive principal minors satisfy
\(\Delta_{B\cup C}(Z)\le\Delta_B(Z)\Delta_C(Z)\) by the
Schur complement formula and monotonicity of positive
determinants. On \(D_j\) the compressions of \(A_z\) and
\(z\) coincide. Applying the inequality first to \(A_z\)
and then to the complementary minor of \(z\) gives
\[
 \frac{\Delta_{S_{j-H_0}}(A_z)}{\Delta_{S_j}(A_z)}
 \le\Delta_{D_j}(z)
 =\Delta_{D_j^c}(z)
 \le\Delta_{S_j}(z)\,\Delta_{S_{j-H_0}^c}(z).
\]
Multiply over all integers \(j\). All but finitely many
ratios are one. The numerator on the left shifts the
denominator by \(H_0\) positions, giving \((\det A_z)^{H_0}\).
On the right, complements have the same minor, and
\(\det z=1\). The product is \(m_\alpha(z)^2
=m_{-\alpha}(z)^2\), again by complementary minors.

Conjugation by \(u\in K\) preserves \(\det A_z\).
For \(y=k_0z\) it also gives
\(|\det A_y|=\det A_z\), since \(A_y=k_0A_z\).
The preceding inequality thus bounds the denominator in
Proposition~\ref{prop:polarcoefficient697} below by
\(|\det A_y|^{H_0/2}\) at every integration point.
Cauchy--Schwarz and invariance of Haar bound the integral
of the absolute numerator by \(M_\mu\). This proves the
first inequality.

Put \(C=y-A_y\). The diagonal blocks of
\(y^*\epsilon y=\epsilon\) give \(A_y^*A_y=1+C^*C\).
Hence \(|\det A_y|^2\ge 1+\operatorname{tr}(C^*C)\) by
the nonnegative eigenvalues of \(C^*C\). Finally the
eigenvalues of \(z\) occur in reciprocal pairs, and the
largest is \(\|y\|_{\rm op}\). Each contributes
\((s+s^{-1})/2\ge1\) to \(\det A_z\), with the largest and
its reciprocal together contributing at least
\((\|y\|_{\rm op}/2)^2\) when \(\|y\|_{\rm op}>1\).
For norm one the bound is immediate. These facts give the
last two estimates.
\end{proof}

\begin{theorem}[Absolute integral outside the compact tube]\leavevmode
\label{thm:noncompacttail697}
At every real place of the same fields and at every sufficiently
large node in the prescribed interval,
\[
   \int_{\|C_s\|_2\ge\delta_*} |\Theta(s)|\,ds
       \le \beta^{-100D}\operatorname{vol}(\mathcal B_\beta).
\]
\end{theorem}
\begin{proof}
On this set the compact coefficient and the character have
modulus at most one. The lemma gives
\(|\Theta|\le\exp(-H_0\delta_*^2/8)\), since
\(\log(1+\delta_*^2)\ge\delta_*^2/2\).
Here \(H_0\delta_*^2/T=r^{15/2}/T^3\to\infty\) uniformly:
\(r\in[x^5,3x^5]\), \(r+k\le2r\), and
\(\log\beta=3\log b\le3x^{1/100}\). The uniform upper bound
is thus at most \(e^{-300T}\). By taking the geometric
mean with the global bound from the lemma and unitarity,
we get on that set
\[
  |\Theta(s)|\le e^{-150T}
           \min\{1,(2/\|J_s\|_{\rm op})^{H_0/2}\}.
\]
The operator norm of \(J_s\) equals that of \(G_s\), since
\(\|h_s\|_{\rm op}=1\le\|G_s\|_{\rm op}\).

The integral of this minimum is at most
\(2\operatorname{vol}(\mathcal B_\beta)\).
Indeed the portion inside \(R=1+\beta\) costs at most that
ball's volume. For \(U\ge1\), the packing argument of
Theorem~\ref{thm:positivecore696} with separation \(1/U\)
covers \(\{\|J_s\|_{\rm op}\le U\}\) by at most
\((1+2U^2)^{2D}\) translates of \(\{\|J_s-1\|_{\rm op}\le1\}\),
which is contained in \(\mathcal B_\beta\).
On the shell \(2^jR<\|J_s\|_{\rm op}\le2^{j+1}R\), \(j\ge0\),
the cost divided by \(\operatorname{vol}(\mathcal B_\beta)\)
is at most
\[
  12^{2D} R^{4D}(2/R)^{H_0/2}\,2^{j(4D-H_0/2)} .
\]
The first factor, including the power of \(2/R\), is at most
\(1/2\) eventually since \(H_0/D\to\infty\) and \(R\to\infty\).
The common ratio is also at most \(1/2\). Summing proves the
asserted bound, using \(2e^{-150T}\le\beta^{-100D}\).
\end{proof}

This tail is smaller than the positive integral on
\(\Omega_\eta\) in Theorem~\ref{thm:positivecore696}.
To compare the entire period with that positive integral,
it remains to bound cancellation in the tube
\(\|C_s\|_2<\delta_*\). Theorem~\ref{thm:angularlocalization696}
places its integrand within \(4\delta_*\) of identity
coordinates after compact translation.


\subsection{A lower bound for the full real Bessel period}
\label{sec:fullrealperiod698}

We prove that the signed integral in the tube of
Theorem~\ref{thm:noncompacttail697} retains a fixed fraction
of the mass estimated in Theorem~\ref{thm:positivecore696}.
The estimates below have absolute constants independent of
both the dimension and the number of minors in a weight.
For a projection \(P\), \(\det_P A\) denotes the determinant
of \(PAP\) acting on its range, and \([P,A]=PA-AP\).

\begin{lemma}[Expansion of a compressed product of exponentials]\leavevmode
\label{lem:compressedproduct698}
Let \(P\) be an orthogonal projection and let
\(\|A_j\|_2\le d\le1/40\) for \(1\le j\le3\).
Set \(u_0=\max_j\|[P,A_j]\|_2\),
\(F_P(A,B)=\operatorname{tr}(PA(1-P)BP)\). Then the
logarithm at the identity satisfies
\[
 \begin{split}
 \log\det_P(e^{A_1}e^{A_2}e^{A_3})
  ={}&\operatorname{tr}\bigl(P\textstyle\sum_j A_j\bigr)
      +\tfrac12\sum_j F_P(A_j,A_j)
      +\sum_{j<l}F_P(A_j,A_l)+E,\\
 & |E|\le 16000\,d\,u_0^2 .
 \end{split}
\]
\end{lemma}
\begin{proof}
Introduce a scalar \(t\) in all three exponents and expand
the trace of the logarithm on \(P\mathbb C^r\).
The compression differs from its identity by at most
\(e^{3|t|d}-1\) in operator norm. Its logarithmic series
converges for \(|t|\le1\). Expand the exponentials and the
logarithm in powers of \(t\). The sum of absolute scalar
coefficients of words of total degree \(l\) before inserting
projections is bounded by the coefficient of \(z^l\) in
\(-\log(2-e^{3z})\). This follows by replacing the
coefficients of the logarithm by their absolute values:
the majorizing series for its argument minus the identity
is \(e^{3z}-1\). The coefficients are nonnegative and the
series has value less than one at \(z=1/10\); hence this
bound is at most \(10^l\).

Insert \(1=P+(1-P)\) at each unseparated junction between
letters. There are at most \(2^l\) resulting paths per word.
Those staying in the range of \(P\) give exactly the
trace logarithm of the product of the three compressed
exponentials, namely of \(e^{P A_1 P}e^{P A_2 P}e^{P A_3 P}\)
on that range. Its determinant is the exponential of the
sum of their traces. Their contribution is thus precisely
the linear term. Every other path crosses the boundary of
\(P\) at least twice, since it begins and ends in its range.
The two crossing factors have Hilbert--Schmidt norm at most
\(u_0\); all other factors have operator norm at most \(d\).
The absolute trace is bounded by \(u_0^2d^{l-2}\) by the
Hilbert--Schmidt Cauchy--Schwarz inequality. In degree
two the crossing terms are exactly the two sums displayed
in the statement, from a square in one exponential or
from two exponentials in their indicated order. The sum in
degrees \(l\ge3\) is bounded by
\(\sum_{l\ge3}20^l d^{l-2}u_0^2\le16000d u_0^2\).
This also gives absolute convergence of the remainders.
\end{proof}

\begin{proposition}[A Gaussian estimate in the shrinking neighborhood]\leavevmode
\label{prop:localgaussian698}
Use the notation of Proposition~\ref{prop:polarcoefficient697}.
For an integer weight \(v\) put
\(Q_v(A)=\sum_{i,j}|v_i-v_j||A_{ij}|^2\). Suppose \(Z\)
is in the sign unitary Lie algebra and \(\|Z\|_2\le8\delta_*\).
For \(q=Q_\alpha(Z)/4\) and all sufficiently large nodes,
\[
 \left| c_\lambda(e^Z)e^{-\operatorname{tr}(\mu Z)}
                    -e^{-q}\right|
            \le \tfrac1{20}e^{-q}+e^{-200T}.
\]
Here \(\mu\) in a trace denotes the diagonal matrix of that
weight. If \(Z'\) is skew-Hermitian on the compact value
space of dimension \(k\) in its defining representation,
with \(\|Z'\|_2\le8\delta_*\), the corresponding estimate is
\[
 \left| ((e^{Z'})_{kk})^t e^{-tZ'_{kk}}-e^{-q'}\right|
          \le\tfrac1{20}e^{-q'}+e^{-200T},\qquad
 q'=\tfrac t2\sum_{j\ne k}|Z'_{kj}|^2 .
\]
\end{proposition}
\begin{proof}
Put \(d=8\delta_*\), which tends to zero with \(dT\to0\).
Let \(U,V\) be the strict lower and strict upper parts of
the compact (equal-sign) blocks of \(Z\). For
\((A_1,A_2,A_3)=(-U,Z,-V)\) all norms are at most \(d\).
Write \(Y=e^{-U}e^Ze^{-V}\). Proposition~\ref{prop:polarcoefficient697}
permits using \(\mathcal F(Y)\) for \(c_\lambda(e^Z)\).
In fact all partial products with the exponents scaled
between zero and one differ from the identity by at most
\(e^{3d}-1<1\) in operator norm, so the germ identity extends
to the matrices used here.

Fix a trailing coordinate projection \(P\) in decreasing
order of \(\lambda\) and \(u\in K\). Put \(R=uPu^*\),
\(s_0=\|[P,Z]\|_2\), \(h_0=\|[P,u]\|_2=\|R-P\|_2\).
Since taking either triangular mask can only delete entries
from \([P,Z]\), one has \(\|[P,A_j]\|_2\le s_0\) and
\(\|[R,A_j]\|_2\le s_0+2dh_0\). Apply
Lemma~\ref{lem:compressedproduct698} with projection \(R\).
The first term is \(\operatorname{tr}(R\,\operatorname{diag} Z)\):
the sum of the three arguments retains just the diagonal and
the opposite-sign entries, and \(R\) is block diagonal by
sign. Its difference from \(\operatorname{tr}(PZ)\) is at most
\(dh_0^2\), since
\(\sum_i|R_{ii}-P_{ii}|=2\operatorname{tr}(P-PR)=h_0^2\).

The quadratic term may be evaluated at \(P\) with error
at most \(C d h_0(s_0+d h_0)\) for an absolute \(C\).
Indeed the crossing blocks \(RA_j(1-R)\) and
\(PA_j(1-P)\) differ in Hilbert--Schmidt norm by at most
\(2dh_0\), by changing one projection at a time. The
reverse crossing blocks have the same bound. Their
individual norms are at most \(s_0+2dh_0\) at \(R\)
and \(s_0\) at \(P\). Expanding each of the six trace
products gives the estimate.

For \(i\in P,\ j\notin P\), necessarily \(i>j\). If the
signs differ, only \(F_P(Z,Z)/2\) contributes this pair,
giving \(Z_{ij}Z_{ji}/2=|Z_{ij}|^2/2\). If the signs agree,
the two square terms from \(U,V\) vanish on this pair.
The four remaining coefficients of \(Z_{ij}Z_{ji}\) are
\(1/2,-1,1,-1\), in order from the square in \(Z\) and the
three mixed terms. Here \(Z_{ji}=-\overline{Z_{ij}}\), so
the result is again \(|Z_{ij}|^2/2\). Summing gives
\(s_0^2/4\), since \(Z_{ij},Z_{ji}\) have the same modulus.
The lemma and the preceding estimates yield, with an
absolute constant,
\[
  \log\det_P(u^*Yu)=\operatorname{tr}(PZ)+s_0^2/4
                                   +O\bigl(d(s_0^2+h_0^2)\bigr).
\]

For the weight \(-\alpha\) all the projections in its
definition are trailing. For any weight \(v\), summing
\(\|[P,A]\|_2^2\) over its level projections gives \(Q_v(A)\):
the number of levels separating \(i\) from \(j\) is
\(|v_i-v_j|\). Also the sum of the projections is
\(\operatorname{diag}(v)-v_{\min}1\).
The determinant factor restores the constant since
\(\det Y=e^{\operatorname{tr} Z}\). It follows that
\[
 \frac{1}{m_{-\alpha}(u^*Yu)}
     =\exp\{\operatorname{tr}(\mu Z)-q+E(u)\},\qquad
 |E(u)|\le C_1 d\{q+Q_\mu(u)+r^2\}
\]
with absolute \(C_1\). To pass from \(\alpha\) to \(\mu\),
the difference \(2\rho_c\) has entries of modulus at most
\(r\). Thus the difference of the linear traces is at most
\(r^2d\) and \(Q_\alpha(u)\le Q_\mu(u)+2r^2\) by unitarity
of \(u\).

The weight \(m(u)=|m_\mu(u)|^2\) satisfies
\(0\le m(u)\le e^{-Q_\mu(u)/2}\). For one of its projections,
the squared compressed determinant is the determinant of
\(1-Pu^*(1-P)uP\) on the range of \(P\). Its eigenvalues
lie in \([0,1]\), so it is at most
\(\exp(-\|(1-P)uP\|_2^2)\). Unitarity gives
\(2\|(1-P)uP\|_2^2=\|[P,u]\|_2^2\). Multiply over levels,
using \(|\det u|=1\). By the compact Weyl dimension formula
\(M_\mu^{-1}=\dim V_\mu\le (3B+r)^{r(r-1)/2}\le e^T\):
every factor is the ratio of a highest-weight difference
plus a root height to that height, at least one in the
denominator, and all differences are at most \(3B\).

After removing \(e^{\operatorname{tr}(\mu Z)}\) the error
from \(e^{-q}\) is consequently bounded by the integral of
\(m(u)e^{-q}|e^{E(u)}-1|/M_\mu\). Where
\(q+Q_\mu(u)\le1000T\), the bound on \(E\) is \(o(1)\)
uniformly since \(r^2\le T\) and \(dT\to0\). This portion
costs at most \(e^{-q}/20\) eventually. On the complement
\(|E|\le(q+Q_\mu(u))/4\) eventually. The unnormalized
integrand is at most \(2\exp(-(q+Q_\mu(u))/4)\).
Integration against probability Haar and division by
\(M_\mu\) costs at most \(2e^{-249T}\le e^{-200T}\).
This proves the first assertion.

For the compact value coefficient take the last coordinate
projection and the single exponential \(e^{Z'}\) in
Lemma~\ref{lem:compressedproduct698}. Its linear term,
after multiplication by \(t\), is \(tZ'_{kk}\), and its
quadratic term is \(-q'\) by skew-adjointness. The error
has modulus at most \(C_2d q'\), because
\(\|[P,Z']\|_2^2=2\sum_{j\ne k}|Z'_{kj}|^2\). If
\(q'\le1000T\) the error in the exponential is \(o(1)\)
relative to \(e^{-q'}\). Otherwise its absolute contribution
is at most \(2e^{-q'/2}\le e^{-200T}\). This includes
\(t=0\), when the coefficient is constantly one.
\end{proof}

\begin{theorem}[Uniform lower bound for the complete real period]\leavevmode
\label{thm:fullrealperiod698}
At the translate \(g\) of Theorem~\ref{thm:effectivestationary695}
and at each real place of the same fields, for all sufficiently
large nodes \(V_0\le x\le V_0^{3/2}\), the complete integral
of the Bessel coefficient satisfies
\[
 m=\int_S\Theta(s)\,ds
   =d_\pi\|W_w\|_{L^2(S\backslash G)}^2
   \ \ge\ \tfrac14\beta^{-16D}\operatorname{vol}(\mathcal B_\beta)>0 .
\]
In particular the product of these full local periods is at
least \(4^{-f}\beta^{-16fD}\prod_v\operatorname{vol}(\mathcal B_{\beta,v})\).
\end{theorem}
\begin{proof}
If \(\|C_s\|_2<\delta_*\), Theorem~\ref{thm:angularlocalization696}
gives \(u\in K_c\) with the same integrand and both components
of \(J_{us}-1\) of Hilbert--Schmidt norm at most
\(4\delta_*\). Their logarithms \(Z,Z'\) have
Hilbert--Schmidt norms at most \(8\delta_*\) by the
logarithmic series and submultiplicativity. Their compact
weight traces are imaginary. By Lemma~\ref{lem:logphase696}
the two linear phases in Proposition~\ref{prop:localgaussian698}
cancel the character in \(\Theta\). Its value is therefore
the product of the first normalized coefficient in that
proposition with the conjugate of the second.

Put \(A=e^{-q}, A'=e^{-q'}\), \(e_0=e^{-200T}\).
If the two factors are within \(A/20+e_0\) and
\(A'/20+e_0\) of \(A,A'\), expanding their product bounds
its distance to \(AA'\) by
\((2/20+1/400)AA'+4e_0\), since \(A,A',e_0\le1\).
In particular \(\Re\Theta(s)\ge-4e_0\) throughout the tube.
The tube lies in \(\mathcal B_\beta\): the compact translation
is unitary in both factors, so it preserves \(\|J_s\|_{\rm op}\),
and \(\|J_{us}\|_{\rm op}\le1+4\delta_*\).

The integral is real and equals the squared-norm expression
by Proposition~\ref{prop:averagedschur684}. Integrate its
real part first on \(\Omega_\eta\) using
Theorem~\ref{thm:positivecore696}. On the rest of the tube
the negative contribution is at most
\(4e^{-200T}\operatorname{vol}(\mathcal B_\beta)\). The
remaining set contributes at worst the negative of the
absolute integral in Theorem~\ref{thm:noncompacttail697}.
Consequently
\[
  \frac{m}{\operatorname{vol}(\mathcal B_\beta)}
       \ge \tfrac12 e^{-16T}-4e^{-200T}-e^{-100T}
       \ge \tfrac14 e^{-16T}
\]
eventually. Every constant above is independent of the
place, and the uniform growth bounds used are those of the
given nodes. Multiplication gives the final assertion.
\end{proof}

The real test vector can now be multiplied by a positive
scalar to give it squared quotient norm equal to the
dimension of its value space. Passage to an integral
projected theta period still requires the global comparison.


\subsection{Normalized test vectors and their truncation}
\label{sec:normalizedcut699}

Continue with the real translates constructed in
Theorem~\ref{thm:effectivestationary695}. Haar on the
compact factor of \(S=U\rtimes\mathrm U(k)\) is probability
Haar, so an \(S\)-invariant scalar function has the same
integral on \(U\backslash G\) and \(S\backslash G\).
All vector norms below descend in this way.

\begin{lemma}[Integrals of an operator norm bound]\leavevmode
\label{lem:cosetpacking699}
At one real place let \(g\) be its stationary translate and
let \(\phi(t)=(2/\|t\|_{\rm op})^{9r^2}\) on \(G\). Write
\(V=\operatorname{vol}(\mathcal B_\beta)\) for the \(S\)-volume
in Theorem~\ref{thm:fullrealperiod698}. Uniformly for all
\(y\in G\), at all sufficiently large nodes,
\[
  \int_S\phi(g^{-1}sy)\,ds\le b^{r^2}V,\qquad
  d_\pi\int_G\phi(t)\,dt\le b^{38r^2}.
\]
\end{lemma}
\begin{proof}
Put \(P=9r^2\) and \(R\ge1\). In the part of the matrix
coset \(g^{-1}Sy\) of operator norm at most \(R\), take
a maximal \(1/R\)-separated set. As in
Theorem~\ref{thm:positivecore696} its cardinality is at
most \((1+2R^2)^{2r^2}\). If \(A=g^{-1}sy\) is within
\(1/R\) of \(A_i=g^{-1}s_i y\) in this set, then
\[
   \|g^{-1}s s_i^{-1}g-1\|
       =\|(A-A_i)A_i^{-1}\|\le1 .
\]
Here and below in this proof undecorated matrix norms are
operator norms. The corresponding \(s s_i^{-1}\) belongs
to \(\mathcal B_\beta\), since its first component has
operator norm at most two and its second is unitary.
Thus the \(S\)-volume with \(\|g^{-1}sy\|\le R\) is at
most \((1+2R^2)^{2r^2}V\) by right invariance. For
\(2^j\le\|g^{-1}sy\|<2^{j+1}\), \(j\ge0\), multiply the
volume bound by \(2^{P-jP}\). The sum is at most
\[
       C_r V,\qquad C_r=2^{P+1}12^{2r^2}\le b^{r^2}
\]
eventually, since the common ratio is \(2^{4r^2-P}\).
This is uniform in \(y\).

Taking the whole group in the packing argument gives
\(\int_G\phi\le C_r\operatorname{vol}_G\{\|t-1\|\le1\}\):
the relative points \(AA_i^{-1}\) now lie in that set.
Cover this last set, of norm at most two, by right
translates of the ball \(\{\|t-1\|\le\eta\}\) using
separation \(\eta/2\). The number required is at most
\((1+8/\eta)^{2r^2}\le\beta^{12r^2}\), since
\(\eta=(200\beta^5)^{-1}\) and \(\beta\to\infty\).
The coefficient derivative estimate in the proof of
Theorem~\ref{thm:positivecore696} applies to every \(t\)
on this small ball. In particular \(|c_\lambda(t)|\ge1/2\),
with the coefficient transported by the sign frame.
By Schur orthogonality \(d_\pi\int_G|c_\lambda(t)|^2\,dt=1\).
It follows that \(d_\pi\) times the small ball's volume
is at most four. The claim follows by \(C_r\le b^{r^2}\)
and \(\beta=b^3\).
\end{proof}

\begin{theorem}[A normalized local vector and its cut]\leavevmode
\label{thm:normalizedcut699}
At each real place put \(d_v=\dim V_{\tau_v}\),
\(\theta=b^{6n+2}\), \(j=3r^3 f\), and
\(l_b(y)=\max\{1,\ell(y)/b^{6n}\}\).
There is a positive scalar multiple \(F_v\) of the
stationary Bessel average and a continuous cut \(F_v^\flat\)
such that
\[
 \|F_v\|_{U\backslash G}^2=d_v,\qquad
 \|F_v^\flat\|_{U\backslash G}^2\le d_v,\qquad
 \|F_v-F_v^\flat\|_{U\backslash G}^2\le d_v b^{-2r^3 f}.
\]
They agree for \(\ell\le\theta/2\), and \(F_v^\flat\) vanishes
for \(\ell\ge\theta\). Both have the prescribed \(S\)-equivariance.
There is a bounded, measurable, \(S\)-invariant nonnegative
function \(A_v\) such that
\[
    l_b(y)^j\max\{\|F_v(y)\|,\|F_v^\flat(y)\|\}\le A_v(y),
    \qquad
    (\sup A_v)\int_{U\backslash G}A_v\le b^{400r^2}d_v .
\]
The full vector satisfies the reproducing identity
\[
       F_v(y)=d_\pi\int_G\overline{c_\lambda(h)}F_v(yh)\,dh .
\]
\end{theorem}
\begin{proof}
Write \(m\) for the full integral in
Theorem~\ref{thm:fullrealperiod698}. By that theorem
\(c=(d_v d_\pi/m)^{1/2}\) is finite and positive,
and \(F_v=c W_w\) has the asserted norm.
Let \(F_v^\flat=\zeta(\ell(y))F_v(y)\), where \(\zeta\)
equals one up to \(\theta/2\), is linear from one to zero
on \([\theta/2,\theta]\), and is zero thereafter.
The length is continuous: for \(t\in G\),
\(\ell(yt)\le\ell(y)\|t\|\) and
\(\ell(y)\le\ell(yt)\|t^{-1}\|\). These follow by taking
the infima and using \(\|t\|,\|t^{-1}\|\ge1\).
They give continuity as \(t\to1\), and the length is
left \(S\)-invariant by definition. Equivariance,
continuity, cut support and the cut norm bound follow.

For \(h=g^{-1}sy\), the length bound on \(g\) gives
\(\ell(y)=\ell(gh)\le b^{5n}\|h\|\). Hence
\(\ell(y)/b^{6n}\le\|h\|/2\) for large nodes.
The gap \(H_0=r^8\) of Lemma~\ref{lem:coefficientdecay697}
exceeds \(9r^2+j\), uniformly, since
\(f\ll_{q_0,d} x^{11}\log(2x)\) and \(r\ge x^5\).
It follows that
\[
       l_b(y)^j |c_\lambda(h)|\le (2/\|h\|)^{9r^2}.
\]
When \(l_b(y)>1\) one has \(\|h\|\ge2\) and uses coefficient
decay with exponent \(9r^2+j\). When \(l_b(y)=1\) use decay
if \(\|h\|\ge2\), and unitarity if \(\|h\|\le2\).

Take \(A_v(y)=c\int_S\phi(g^{-1}sy)\,ds\). Integration
of the last estimate, using the unitarity of the action
on the values, gives the pointwise bound for \(F_v\) and
its cut. Tonelli gives measurability and disintegration
over the quotient gives
\(\int_{U\backslash G} A_v=c\int_G\phi\); right invariance
on \(S\) gives its left invariance as a function of \(y\).
By Lemma~\ref{lem:cosetpacking699} and the lower bound for
\(m\),
\[
   (\sup A_v)\int A_v
   \le c^2 b^{r^2}V\int_G\phi
   \le 4d_v\,\beta^{16(r+k)^2} b^{39r^2}
   \le b^{400r^2}d_v .
\]
Outside \(\ell\le\theta/2\) we have \(l_b(y)\ge b\).
There \(\|F_v-F_v^\flat\|\le b^{-j}A_v(y)\);
inside, the difference vanishes. Integrating its square
gives \(d_v b^{-6r^3 f+400r^2}\le d_v b^{-2r^3 f}\)
for all sufficiently large nodes.

Finally \(d_\pi\int_G\overline{c_\lambda(h)}\pi(h)\,dh\)
is the orthogonal projection onto the line of \(w\):
Schur orthogonality gives the matrix entries of this
operator against arbitrary vectors, and the coefficient
is integrable by Proposition~\ref{prop:averagedschur684}.
The map \(u\mapsto W_u\) in that proposition intertwines
\(\pi\) with right translation and is bounded into the
quotient \(L^2\)-space. Applying it to the projection
identity on \(w\) proves reproduction in \(L^2\).
The convolution is continuous by boundedness of \(F_v\),
integrability of its kernel, and dominated convergence.
The quotient Haar measure has full support, so the
continuous representatives agree everywhere.
\end{proof}

\begin{corollary}[The cut at all real places of the given fields]\leavevmode
\label{cor:globalrealcut699}
With \(d_B=\prod_v d_v,\ F=\bigotimes_v F_v\) and
\(F^\flat=\bigotimes_v F_v^\flat\), one has
\[
 \|F\|_{U_\infty\backslash G_\infty}^2=d_B,\qquad
 \|F^\flat\|_{U_\infty\backslash G_\infty}^2\le d_B,\qquad
 \|F-F^\flat\|_{U_\infty\backslash G_\infty}
                          \le f b^{-r^3 f}\sqrt{d_B}.
\]
These specific nonzero full local tests and their cuts can
be used in Theorems~\ref{thm:realpair685} and
\ref{thm:rotatedorbit694}.
\end{corollary}
\begin{proof}
Product Haar and Tonelli multiply the squared quotient norms.
In the difference of the tensors replace one factor at a time.
In each of the \(f\) terms the changed factor contributes at
most \(b^{-r^3f}\sqrt{d_v}\) and all the others at most
\(\sqrt{d_{v'}}\). The triangle inequality proves the estimate.
The integrability, equivariance, support, choice of coefficient
and local norm required for the cited estimates are those just
verified. The bounds, including the exponent containing \(f\),
hold at the nodes of Theorem~\ref{thm:fullrealperiod698}.
\end{proof}


\subsection{The truncation error after summation over rational translates}
\label{sec:rationalcut700}

Use the full tensor \(F\), its cut \(F^\flat\), and the functions
\(A_v\) from Theorem~\ref{thm:normalizedcut699}. Fix a source
class \(c\) and a residue \(s\) at the two marked primes. Denote by
\(\mathcal T(c,s,c')\) the rational isometries \(\gamma\) for which
\(h_c^{-1}\gamma h_{c'}\) preserves the finite lattice and reduces
to \(s\). Throughout this subsection write
\[
 L(y)=\prod_{v\mid\infty}l_b(y_v),\qquad
 J=r^3f,\qquad d_0=2nrf,\qquad \rho=b^{37n}.
\]
The function \(\beta_c\) is the reduced partition from the proof
preceding Lemma~\ref{lem:orbitfib693}. Thus it is nonnegative,
is supported on reduced frames, and its \(\Gamma_c\)-translates
sum to one.

\begin{lemma}[Counting all translates at a prescribed length]\leavevmode
\label{lem:framepacking700}
For \(u\ge1\), uniformly in \(c,s,y\), at every sufficiently large
node one has
\[
 \sum_{c'}\ \sum_{\substack{\gamma\in\mathcal T(c,s,c')\\
                         L(\gamma^{-1}y)\le u}}
       \beta_{c'}(\gamma^{-1}y)
             \ \le\ b^{6n^2d_0}u^{2d_0}.
\]
Consequently the sum of \(\beta_{c'}(\gamma^{-1}y)
L(\gamma^{-1}y)^{-J}\) over all these classes and all their
translates is at most \(2b^{12n^3rf}\).
\end{lemma}
\begin{proof}
Write \(e_j\) for the rational frame columns, including the
terminal column \(e_n=z_0\) of norm \(-N\). Write
\(\widehat e_{j,v}=s_{j,v}\Phi_v(e_j)\) for their unit Euclidean
normalizations in the real model of Corollary~\ref{cor:realmarked682}.
The transformed frame we count has columns
\(y_v^{-1}\gamma_v\widehat e_{j,v}\). At \(x=\gamma^{-1}y\)
in the indicated support, \(\ell(x_v)\le b^{6n}u\). Choose
\(w_v\in S_v\) with \(\|w_vx_v\|\le2b^{6n}u\). In the normalized
model these unitary isometries of the indefinite form satisfy
\(\|g^{-1}\|=\|g\|\). The successive orthogonalized columns of
\(x_v^{-1}\widehat e_{j,v}\) are unchanged by the unitriangular
column operations induced by \(w_v^{-1}\). Their lengths are
therefore at most \(2b^{6n}u\). All strict upper coefficients
recovering the original columns are bounded by \(b^{26}\) on
the reduced set, as shown in Lemma~\ref{lem:frameheight691}.
Every transformed column has length at most
\(2nb^{26+6n}u\le\rho u\).

Two distinct transformed frames have distance at least
\((N\rho u)^{-1}\) in the maximum of the column norms over the
places. To prove this, take two translates \(\gamma_1,\gamma_2\),
possibly from different right classes, set
\(x=\gamma_1^{-1}y,\ \delta=\gamma_1^{-1}\gamma_2\), and assume
the distance is strictly smaller. For each pair \(i,j\) the
entry \(h(e_i,(\delta-1)e_j)\) belongs to \(\mathcal O_E\).
Indeed the two finite class lattices have the same frame, and
\(\delta\) takes the second lattice isometrically onto the first,
where the form is integral. At \(v\) the normalized pairing has
modulus at most
\[
 \|x_v^{-1}\widehat e_{i,v}\|\,
 \|x_v^{-1}(\delta_v-1)\widehat e_{j,v}\|<1/N ,
\]
since the defining Hermitian matrix has operator norm one.
Restoring the two scale factors costs at most \(N\).
Thus the algebraic integer has absolute value less than one at
every embedding, and its norm shows that it is zero.

It follows that \((\delta-1)e_j\in T^\perp=L_0\oplus X_0\),
where here \(L_0\) denotes the isotropic arm spanned by
\(e_1,\ldots,e_a\). Its Hermitian norm is zero: expand
\(h(\delta e_j,\delta e_j)=h(e_j,e_j)\) and use its
orthogonality to \(e_j\). The definiteness of \(X_0\) then
puts the difference in \(L_0\). Write it as
\(\sum_{h\le a}t_{hj}e_h\); each \(t_{hj}\) is integral because
the arm is a direct lattice summand. In any nonzero column
choose the largest \(h\) with \(t_{hj}\ne0\). Multiplication
by \(w_v\), upper unitriangular on this arm, leaves that top
coefficient unchanged. In normalized coordinates its modulus
is \(|s_{j,v}/s_{h,v}|\,|\Phi_v(t_{hj})|\). On the other hand
\[
   \|w_v(\delta_v-1)\widehat e_{j,v}\|
       \le \|w_vx_v\|\,
           \|x_v^{-1}(\delta_v-1)\widehat e_{j,v}\|
       < {2 b^{6n}\over N\rho}.
\]
Both scale ratios are at most \(b^{23n}\). Thus all conjugates
of \(t_{hj}\) have modulus less than one, a contradiction.
This proves the separation.

The ambient real dimension of the tuple of columns is \(d_0\).
Balls of radius \((2N\rho u)^{-1}\) about distinct frames have
disjoint interiors in the maximum norm. Comparing their volumes
with the concentric enlargement of the ball of radius \(\rho u\)
bounds their number by \((1+2N\rho^2u^2)^{d_0}\).
This is at most \((b^{6n^2}u^2)^{d_0}\), since \(N\le b^{60n}\)
and \(n\) tends to infinity. The first paragraph of the proof
of Proposition~\ref{prop:fibercount691} identifies the
multiplicity at a single frame: equality of the residue implies
that the relative map induces a principal-level isometry on the
definite middle quotient. Hence the right classes agree, and
their definite principal-level automorphism is trivial. The
relative maps then lie in \(\Gamma_{c'}^0\). The weights in the
fiber add to at most one, being distinct terms in the partition
by \(\Gamma_{c'}\). This proves the first assertion without a
restriction on the left point \(y\). For the second use
\(L^{-J}=\int_L^\infty Ju^{-J-1}\,du\) and Tonelli's theorem.
The extra factor is \(J/(J-2d_0)\le2\) at large nodes.
\end{proof}

\begin{theorem}[Absolute truncation error in a full rational row]\leavevmode
\label{thm:globalamplitudetail700}
Write \(K_\gamma(y)=F(y)F(\gamma^{-1}y)^*\) for the matrix
outer product, and \(K_\gamma^\flat\) for its cut counterpart.
Both kernels are absolutely integrable after summation over
every row \((c,s)\) with weights
\(\beta_c(y)\beta_{c'}(\gamma^{-1}y)\). With the same weights put
\[
 E_{csc'}=\sum_{\gamma\in\mathcal T(c,s,c')}\int_{G_\infty}
  \beta_c(y)\beta_{c'}(\gamma^{-1}y)
  \|K_\gamma(y)-K_\gamma^\flat(y)\|_{\rm HS}\,dy .
\]
On the fields and interval of Theorem~\ref{thm:fullrealperiod698},
for every sufficiently large node,
\[
                         \sum_{c'}E_{csc'}
                         \le \tfrac12 v_U b^{-k^3f}.
\]
\end{theorem}
\begin{proof}
Put \(A(y)=\prod_v A_v(y_v)\), \(a_*=\prod_v\sup A_v\).
At \(z=\gamma^{-1}y\) each kernel norm is at most
\(A(y)a_*L(y)^{-3J}L(z)^{-3J}\) by
Theorem~\ref{thm:normalizedcut699}. The cut matrix is the
full matrix multiplied by the product of all the scalar cutoffs
at \(y,z\), a real number in \([0,1]\). If the difference is
nonzero, some local length exceeds \(b^{6n+2}/2\).
Thus \(L(y)L(z)\ge b\) for large \(b\), and
\[
       \|K_\gamma(y)-K_\gamma^\flat(y)\|_{\rm HS}
           \le b^{-2J} A(y)a_*L(z)^{-J}.
\]
For the full and cut kernels separately the same upper bound
holds with the factor \(b^{-2J}\) omitted. The last factor is
summable against the right weights by Lemma~\ref{lem:framepacking700}.
For the left weight one has
\[
       \int_{G_\infty}\beta_c(y) A(y)\,dy
           =v_U\prod_v\int_{U_v\backslash G_v}A_v .
\]
Indeed at fixed \(y\) integration of \(\beta_c(uy)\) over
\(U_\infty\) unfolds to the integral of one over
\(\Gamma_c\backslash U_\infty\), of volume \(v_U\).
The \(A_v\) are left \(U_v\)-invariant. Tonelli's theorem and
Theorem~\ref{thm:normalizedcut699} consequently bound the
sum of the error norms by
\[
        2v_U d_B\, b^{(-2r^3+12n^3r+400r^2)f}.
\]
This also proves the asserted absolute convergence, by deleting
the gain \(b^{-2J}\). Finally \(d_B\le b^{2kf}\),
\(k\le r\), and \(12n^3 r+400r^2+2k=o(r^3)\) on this interval
because \(n\le\sqrt x\) and \(r\ge x^5\). The stated bound
follows. The estimates also hold for arbitrary subsets of
rational terms by nonnegativity.
\end{proof}

\begin{corollary}[The error in the residue second moment]\leavevmode
\label{cor:squaredamplitudetail700}
Uniformly on the same interval the array of differences between
full and cut block sums, normalized by \(v_U\), has squared
Hilbert--Schmidt norm at most
\[
  v_U^{-2}\sum_{c,s,c'}E_{csc'}^2
          \le |\mathcal C|b^{-k^3f}
          \le \tau^{-40}|B_{\rm res}|\,|\mathcal C|d_B .
\]
The bound persists when each block is restricted to a subset
of its rational translates or multiplied termwise by scalars
of modulus at most one.
\end{corollary}
\begin{proof}
For each \(c,s\), the sum of the squares of the nonnegative
errors is at most the square of their sum. At a split place of
\(F\) above \(p_i\) the number of possible residue matrices
is at most \(p_i^{r^2}\); there are \(f\) such places above
each mark. Therefore there are at most
\(M^{fr^2}\le b^{22fr^2}\) residues. Apply
Theorem~\ref{thm:globalamplitudetail700} and use
\(k^3\ge22r^2\) for all sufficiently large nodes, since
\(k=r-2n+1\). The last inequality follows from
\(\tau=b^{kf/4}\) and \(|B_{\rm res}|,d_B\ge1\).
The triangle inequality gives the assertion for signed matrix
sums. It bounds their difference as an array and hence also
the absolute difference of their Hilbert--Schmidt norms.
\end{proof}

\section{Poincar\'e series and their marked second moment}\label{nav:poincare-spanning}

\subsection{The rank-zero contribution at the two marks}
\label{sec:besselresidue701}

Write \(H=\mathrm U(X_0)\), \(S=U\rtimes H\), \(K\) for the
finite integral stabilizer, and \(K(M)\) for its principal
subgroup. The \(H\)-level is \(K_H(M)=H(\mathbb A_{F,f})\cap
K(M)\), with \(H\) extended by the identity on the arms and the
terminal line. The finite set
\(\mathcal C=H(F)\backslash H(\mathbb A_{F,f})/K_H(M)\)
has representatives \(h_c\). Let \(B_{\rm res}\) be the reduction
of \(S(\mathbb A_{F,f})\cap K\). At the infinite places let
\(V=\bigotimes_v V_{\tau_v}\), \(\dim V=d_B\), and let
\(\sigma_\infty\) denote the unitary action of \(S_\infty\)
in the equivariance of \(F\). Its restriction to \(U_\infty\)
is the prescribed scalar character. Its restriction to
\(\prod_v\mathrm U(k)\) is the exterior tensor product of the
irreducible local value representations. The cut has the same
equivariance.

\begin{lemma}[Rational lifts of a Bessel residue]\leavevmode
\label{lem:residuelifts701}
For \(t\in B_{\rm res}\) and each \(c\), there is exactly one
class \(c\cdot t\) with an element of \(S(F)\) over
\((c,t,c\cdot t)\). For any such lift \(\gamma_{c,t}\), all the
lifts there form the right coset \(\gamma_{c,t}\Gamma_{c\cdot t}\).
The rule is a right action on \(\mathcal C\), and
\[
  \gamma_{c,t}\Gamma_{c\cdot t}\gamma_{c,t}^{-1}=\Gamma_c .
\]
\end{lemma}
\begin{proof}
Lift \(t\) to \(k_t\in S(\mathbb A_{F,f})\cap K\), and factor
\(h_ck_t=u h\), with \(u\in U(\mathbb A_{F,f})\) and
\(h\in H(\mathbb A_{F,f})\). Write \(h=\eta h_d a_0\), with
\(\eta\in H(F)\), \(a_0\in K_H(M)\). Choose \(x\in U(F)\)
such that \(h^{-1}u^{-1}xh\in K(M)\). For this choice one only
needs density of \(U(F)\) in the finite adeles of \(U\). Indeed
in characteristic zero the finite logarithm and exponential
polynomials are inverse rational-coordinate maps between this
unipotent matrix group and its nilpotent Lie algebra. They give
homeomorphisms at the finite adeles, being integral outside
finitely many primes in a rational basis. The diagonal copy of
\(F\) is dense in \(\mathbb A_{F,f}\): multiply the given
coordinates and congruences by a common rational integer to
make them integral, and use the Chinese remainder theorem in
\(\mathcal O_F\), also imposing divisibility at any additional
primes of that integer. This gives the required approximation
in every Lie algebra coordinate. Now \(x\eta\) is a lift since
\[
 k_t^{-1}h_c^{-1}x\eta h_d
       =(h^{-1}u^{-1}xh)a_0^{-1}\in K(M).
\]

If \(\gamma_1,\gamma_2\) are lifts for the same \(c,t\) with
right classes \(d_1,d_2\), write
\(\gamma_1^{-1}\gamma_2=u'\eta'\) in \(S(F)\).
Its class translate is integral and principal, by equality
of residues. It acts on the middle quotient of the lattice
in the isotropic flag by \(h_{d_1}^{-1}\eta'h_{d_2}\) extended
by one on the terminal line. This induced map preserves the
middle lattice and is principal there. Its extension by the
identity on the two direct integral arms belongs to \(K(M)\).
Thus \(h_{d_1}^{-1}\eta'h_{d_2}\in K_H(M)\). The two classes
are equal and \(\eta'=1\) by freeness of the definite level
in Proposition~\ref{prop:rowmean690}. The relative element
then belongs to \(\Gamma_{d_1}\). Conversely every such right
multiple is a lift. Products and inverses of lifts have the
product and inverse residues by cancellation of the intermediate
class representatives. This proves the right action assertion.
Conjugation by a lift takes the rational unipotent to itself.
At the finite places it takes one indicated principal condition
to the other because \(K(M)\) is normal in \(K\). Applying
this also to its inverse proves the displayed equality.
\end{proof}

For a set \(T\subset G(F)\), write \(\mathcal A_T(s)_{cd}\) for
the sum, over \(\gamma\in T\cap\mathcal T(c,s,d)\), of the
integrals of \(K_\gamma\) against the two class weights, divided
by \(v_U\). Regard these matrices as blocks on
\(\mathcal H_X=\bigoplus_{c\in\mathcal C} V\). Every such sum
converges absolutely by Theorem~\ref{thm:globalamplitudetail700}.
Write \(R(t)=\mathcal A_{S(F)}(t)\).

\begin{theorem}[The exact matrix of the Bessel terms]\leavevmode
\label{thm:besselmatrix701}
For \(t\in B_{\rm res}\) the only nonzero block in row \(c\) of
\(R(t)\) is
\[
                  R(t)_{c,c\cdot t}=\sigma_\infty(\gamma_{c,t}).
\]
These operators form a unitary representation of \(B_{\rm res}\)
on \(\mathcal H_X\). For \(t\) outside \(B_{\rm res}\) one has
\(R(t)=0\). In particular, with \(D_X=|\mathcal C|d_B\),
\[
  R(1)=I_{\mathcal H_X},\qquad
  \sum_t\|R(t)\|_{\rm HS}^2=|B_{\rm res}|D_X .
\]
\end{theorem}
\begin{proof}
A Bessel term between any two classes has finite translate
in \(S(\mathbb A_{F,f})\), hence has residue in \(B_{\rm res}\).
Fix a lift \(\gamma=\gamma_{c,t}\), \(d=c\cdot t\).
The real character equals one on \(\Gamma_d\): its finite factor
there equals one by the conjugation and conductor computation
in Proposition~\ref{prop:stabilizercancel680}, and the adelic
character is trivial on rational points. Thus \(F\) is left
\(\Gamma_d\)-invariant. At each \(y\), all the matrices in the
right coset equal \(F(y)F(y)^*\sigma_\infty(\gamma)\).
Their right weights sum to one by the partition for \(\Gamma_d\).
It remains to evaluate
\[
          v_U^{-1}\int_{G_\infty}\beta_c(y)F(y)F(y)^*\,dy
             =\int_{U_\infty\backslash G_\infty}F(y)F(y)^*\,d\dot y .
\]
The equality follows by disintegration and the left
\(U_\infty\)-invariance of the matrix. Its trace is \(d_B\)
by Corollary~\ref{cor:globalrealcut699}. It commutes with the
product of compact middle groups: left translation by any
of their elements conjugates the integrand by its value
representation and preserves quotient measure. Indeed conjugation
on \(U_\infty\) preserves Haar measure, since a compact group
has trivial modulus. Successive applications of Schur's lemma
to the irreducible tensor factors show that this integral is
scalar. Its trace makes it the identity.

The displayed block is independent of the lift, as the factor
on \(\Gamma_d\) is the identity. The product of the lifts for
\((c,t)\) and \((c\cdot t,t')\) gives the block at \(tt'\).
Hence \(R(t)R(t')=R(tt')\), and \(R(1)=I\). Each operator
permutes classes and acts unitarily on each component. Its
squared Hilbert--Schmidt norm is \(D_X\); summing gives the
formula.
\end{proof}

\begin{corollary}[The full contribution of the isotropic parabolic]\leavevmode
\label{cor:parabolicmass701}
Let \(P_L(F)=\{\gamma\in G(F):\gamma L=L\}\). For
\(\xi\in\mu(E)\) write \(\bar\xi_{\rm res}\) for its scalar
reduction at the marks and \(\omega_\infty(\xi)\) for its
unitary central character on the full real vector. Then
\[
       \mathcal A_{P_L(F)}(s)=
          \sum_{\xi\in\mu(E)}
           \omega_\infty(\xi)R(\bar\xi_{\rm res}^{-1}s).
\]
The supports of these summands are disjoint. The squared norm
summed over all residues is
\(|\mu(E)|\,|B_{\rm res}|D_X\). On either marked slice, where
\(s\) agrees with a Bessel residue above the specified prime,
the parabolic matrix is exactly \(R(s)\); its squared norm
summed on that slice is \(|B_{\rm res}|D_X\).
\end{corollary}
\begin{proof}
Group the sum into double cosets under \(\Gamma_c,\Gamma_d\).
This is legitimate by absolute convergence. For one coset
first sum the right weights over all right multiples, giving
one. Then summing the left weights unfolds the integral to
\((\Gamma_c\cap\gamma\Gamma_d\gamma^{-1})\backslash G_\infty\).
In detail the remaining left translates are indexed by the
cosets of that intersection in \(\Gamma_c\); after translation
of the variable their weights add to a partition for the
intersection. The kernels are equivariant under this change
because the two lattice characters are trivial. For disagreeing
characters the orbital integral therefore vanishes by
Proposition~\ref{prop:stabilizercancel680}.

An agreeing \(\gamma\) preserving \(L\) lies in \(\xi S(F)\)
for a unique \(\xi\in\mu(E)\). Indeed the equal coefficient
ratio on the common flag of Theorem~\ref{thm:rationalflags680}
is the eigenvalue on its first line. The proof of
Corollary~\ref{cor:parabolicslice680} gives the conclusion
before imposing its congruence on that eigenvalue. Conversely
\(\xi S(F)\) normalizes the Bessel unipotent and preserves its
character, so every one of its elements agrees. The scalar
\(\xi\) commutes with \(y\), and its right action on the real
test is the central character of the tensor of the discrete
series, by the intertwining in Proposition~\ref{prop:averagedschur684}.
Thus it can be taken out of the adjoint factor in the outer
product with multiplier \(\omega_\infty(\xi)\). For these
cosets the right weights again sum to one as in the theorem,
even though the argument of each weight has also been
translated by the scalar. The residue is multiplied by
\(\bar\xi_{\rm res}\). This proves the stated matrix identity.

The first diagonal entry of an element of \(B_{\rm res}\) is
one. At any one chosen component above either mark, the roots
of unity of \(E\) inject on reduction, since the field there
is \(\mathbb Q_{p_i}\) and \(1+p_i\mathbb Z_{p_i}\) is
torsion-free for odd \(p_i\). Thus distinct scalars give
disjoint residue cosets. Only the identity scalar survives
on the indicated slice, which contains all of \(B_{\rm res}\).
The norm formulas now follow from the theorem.
\end{proof}

\begin{proposition}[Transport of the full and cut block sums]\leavevmode
\label{prop:residuetransport701}
If \(T\) is stable under left and right \(S(F)\), then
\(\mathcal A_T(st)=\mathcal A_T(s)R(t)\) for every residue
\(s\) and every \(t\in B_{\rm res}\). The same identity holds
with \(\mathcal A_T\) formed from cut vectors and with \(R(t)\)
as in Theorem~\ref{thm:besselmatrix701}.
\end{proposition}
\begin{proof}
For fixed \(d\) take \(\delta=\gamma_{d,t}\) and \(e=d\cdot t\).
Right multiplication by \(\delta\) identifies the terms over
\((c,s,d)\) with those over \((c,st,e)\), including their
membership in \(T\). It also identifies their right unipotent
cosets by Lemma~\ref{lem:residuelifts701}. In either sum the
right weights on a coset add to one at fixed \(y\); every
kernel in that coset is equal because the lattice character
is trivial. Passing to right cosets and integration is justified
by Theorem~\ref{thm:globalamplitudetail700}. Equivariance gives
\(K_{\gamma\delta}(y)=K_\gamma(y)\sigma_\infty(\delta)\),
also for the cut. The only block entering the product on the
right in column \(e\) is \(R(t)_{de}\), which proves the
identity. This evaluates the parabolic part of the marked
second moment. The open rational terms in the same residues
still enter the integral comparison.
\end{proof}


\subsection{Summing the signed open orbits on a marked slice}
\label{sec:signedopenslice702}

Put \(D_{\rm id}=|B_{\rm res}|D_X\),
\(n_p=|B_{\rm res}|\,|\mathcal C|N^{-kf}\) and
\(\Omega_i=\{s:s_i=t_i\text{ for some }t\in B_{\rm res}\}\).
Write \(X_\gamma={\rm N}_{F/\mathbb Q}(\Delta_\gamma)\) for an
open term. This is a positive integer. We shall retain the
factor in the archimedean bound which is independent of the
rational term. With the measures of Theorem~\ref{thm:realgram682}
define
\[
 \begin{gathered}
 A=10^5n^4,\quad e=A+k-1,\quad
 \Xi=\prod_{v,j}|s_{j,v}|^2,\quad d_\pi=\prod_v d_{\pi_v},
 \quad c_G=V_k S_k/V_r,\\
 C_\infty^2=
  \frac{d_B^2\rho^{9nAf}}
       {v_U^2d_\pi c_G^f\Xi^{2e}},\qquad
 D_{\rm op}=n_p M^{2akf}C_\infty^2 .
 \end{gathered}
\]
Here \(\rho=b^{37n}\) and \(\tau=b^{kf/4}\), as before.

\begin{lemma}[Elimination of nonagreeing residues]\leavevmode
\label{lem:finiteagreement702}
Let \(T\) be stable under left and right \(S(F)\). Both its full
and cut block matrices vanish unless the finite unipotent
characters agree at the residue. Explicitly, for integral
finite unipotents of residues \(u,u'\) with \(us=su'\), their
finite character values must coincide for the block at \(s\)
to be nonzero.
\end{lemma}
\begin{proof}
The adjoint of the array at \(s\) is the array for \(T^{-1}\)
at \(s^{-1}\): invert each rational term and change variables
by \(y=\gamma z\), interchanging the left and right weights.
Absolute convergence justifies this computation. Applying
Proposition~\ref{prop:residuetransport701} to \(T^{-1}\) and
taking adjoints gives the left transport formula
\(\mathcal A_T(ts)=R(t)\mathcal A_T(s)\). Both formulas hold
with cut or full kernels. For a finite integral unipotent,
Lemma~\ref{lem:residuelifts701} in each row permits a unipotent
rational lift to the same class: take \(h=h_c,\ d=c,\eta=1\)
in its proof. Its archimedean character is the inverse of
the finite character of the chosen residue. Indeed class
conjugation fixes the simple-root sum, and on principal
unipotents the finite character is one by
Proposition~\ref{prop:stabilizercancel680}. Hence \(R(u)\)
is that one scalar on the common source. The equalities for
\(us=su'\) give the result. One may take the integral
unipotents supported at just one place, so agreement holds
at each mark where a block survives.
\end{proof}

\begin{proposition}[One signed integral per open orbit]\leavevmode
\label{prop:unfoldopen702}
In a fixed block take a union of
\(\Gamma_c^0\)-left and \(\Gamma_d\)-right orbits of open terms.
Choose one representative of each such orbit, choosing a near
one whenever possible. Its normalized cut block is a sum of
integrals
\[
       v_U^{-1}\int_{G_\infty}R_{\rm red}(y)K_\gamma^\flat(y)\,dy .
\]
Only near representatives need be kept. The norm of the
integral for each retained representative is at most
\(C_\infty X_\gamma^{-e/2}+\epsilon_\gamma\), where the
nonnegative errors, summed in any block, are at most
\(E_{csd}/v_U\) of Theorem~\ref{thm:globalamplitudetail700}.
\end{proposition}
\begin{proof}
First sum over right cosets as in
Proposition~\ref{prop:residuetransport701}. The action of
\(\Gamma_c^0\) on these right cosets is free. In fact if
\(l=\gamma u\gamma^{-1}\) with \(l\) fixing the first frame
and \(u\in\Gamma_d\), write \(B\) for the upper unitriangular
action of \(u\) on the second frame in its own basis. The
cross Gram satisfies \(C=CB\), by invariance under \(l\).
Since \(C\) is invertible, \(B=1\). Now \(l\) fixes both frames
pointwise. Their joint span is nondegenerate and contains
all negative directions in the sense that its negative index
equals that of the ambient space
(Theorem~\ref{thm:rationalflags680}). Its pointwise stabilizer
is compact at every real place. A unipotent element of a
compact unitary group is the identity by diagonalizability.
This proves freeness.

For \(l\in\Gamma_c^0\), \(K_{l\gamma}(ly)=K_\gamma(y)\).
The sum of \(\beta_c(ly)\) over this free left orbit is
\(R_{\rm red}(y)\), because the reduced indicator is invariant
under the pointwise frame fixer and \(p_c\) is a partition
for \(\Gamma_c^0\). Applying these changes also to the
nonnegative integrals of the error norms gives precisely the
same unfolding, so the sum of their reduced-weight versions
is at most \(E_{csd}\). This proves absolute convergence for
the regrouping of both matrix integrals.

If a cut integrand with this reduced weight is nonzero,
both lengths are in the cut at every place. A right level
move reduces \(\gamma^{-1}y\) by the partition in
Section~\ref{sec:orbitalclosed693}; it preserves the length.
Consequently the orbit contains a near representative. The
reduced-weight integral is independent of the representative
by the preceding unfolding.

For a near choice \(R_{\rm red}\) is invariant under the
product of the real pointwise frame fixers. Every reduced
point of cut length is in the box of radius \(\rho\) by
Lemma~\ref{lem:frameheight691}. A witnessing pair of boxed
frames gives \(|Z_{ij,v}|\le\rho^2\). The hypotheses of
Theorem~\ref{thm:realpair685} thus apply with precisely
this weight, using its adjoint-kernel version. Each
\(\|F_v\|_2^4=d_v^2\), and
\(\delta(Z_v)=(\prod_j |s_{j,v}|^2)^2\Phi_v(\Delta_\gamma)\)
by Corollary~\ref{cor:realmarked682}. Substitution and division
by \(v_U\) give \(C_\infty X_\gamma^{-e/2}\); the other term of
that theorem is the integral of the error norm just unfolded.
\end{proof}

\begin{theorem}[The remaining open contribution on a marked slice]\leavevmode
\label{thm:signedopen702}
Let \(T_>\) consist of open rational terms with
\(\sqrt{X_\gamma}/P_n>(M/b)^{af}/N^{f/2}\). For \(i=0,1\),
the normalized full array on \(\Omega_i\) satisfies
\[
 \left(\sum_{s\in\Omega_i}\|\mathcal A_{T_>}(s)\|_{\rm HS}^2\right)^{1/2}
 \le (b^3/M)^{kf/2}\sqrt{D_{\rm op}}
                              +2\tau^{-20}\sqrt{D_{\rm id}} .
\]
In particular the full array of all rational terms satisfies
\[
 \begin{split}
 \left(\sum_{s\in\Omega_i}
       \|\mathcal A_{G(F)}(s)-R(s)\|_{\rm HS}^2\right)^{1/2}
 &\le (b^3/M)^{kf/2}\sqrt{D_{\rm op}}\\
 &\quad+(4\tau^{-20}+\tau^{-3a/2})\sqrt{D_{\rm id}} .
 \end{split}
\]
\end{theorem}
\begin{proof}
The two leading determinants of orders \(a,n\) are unchanged
by left and right multiplication by Bessel elements: on the
frame these perform upper unitriangular operations preserving
every initial segment. Thus \(T_>\) is stable. Remove rational
terms with disagreeing characters, whose whole orbital
integrals vanish by Proposition~\ref{prop:stabilizercancel680}
and the unfolding in Corollary~\ref{cor:parabolicmass701}.
This removal also preserves stability, as Bessel elements
normalize the unipotent and preserve its character. By
Lemma~\ref{lem:finiteagreement702} only agreeing residues
need be summed.

For the cut array at those residues apply
Proposition~\ref{prop:unfoldopen702}. The error array has norm
at most \(\tau^{-20}\sqrt{D_{\rm id}}\) by
Corollary~\ref{cor:squaredamplitudetail700}. For the main
majorant use \(X_\gamma\ge1,\ e\ge k\). The inequality defining
\(T_>\) gives \(X_\gamma>P_n^2/N^f\), so the gauge in
Theorem~\ref{thm:allpair689} is \(G_\gamma=M^{2af}X_\gamma\).
The individual main bound is thus at most
\(C_\infty M^{akf}G_\gamma^{-k/2}\).

Expand the square of the sum of these bounds in each block.
Each selected representative determines a distinct near orbit
for the smaller left and right pointwise subgroups: equality
there would imply equality in the larger orbit used in the
proposition. Thus the ordered pairs form a subset of the
pairs counted in Theorem~\ref{thm:paircount691}. For a pair
of positive corank \(\ell=k-\dim I_X\) (allowing \(\ell=0\)),
put \(c^+=\operatorname{codim}_D\ker(J_\gamma-J_\eta)\).
When \(c^+\ge1\), the common-lattice estimate gives
\[
 (G_\gamma G_\eta)^{-k/2}
 \le (2/M)^{\max(1,\ell-1)kf}\,d(L_{12})^{-k}.
\]
If \(c^+=0\), then \(\ell\le1\) and the full links agree.
Theorem~\ref{thm:samelink687} gives the stronger factor
\(p_i^{-2kf}\) multiplying \(d(L_{12})^{-k}\), since the
reduced full frames meet in each of the two components at
every place above the slice prime. Summing the two bounds by
Theorem~\ref{thm:paircount691} and dividing the result by
\(D_{\rm op}\) gives at most
\[
  (\tau^2+\tau^3)p_i^{-2kf}
  +\sum_{\ell\ge0}\tau^{\ell+2}
                       (2/M)^{\max(1,\ell-1)kf}.
\]
The series is at most \(4\tau^4(2/M)^{kf}\): the first
two terms have the same gain, and from \(\ell=2\) onward the
ratio is \(\tau(2/M)^{kf}<1/2\). The sizes of the two primes
give \(p_i^2\ge M/b\). The entire display is then at most
\(6\tau^4(b/M)^{kf}\le(b^3/M)^{kf}\). Taking square roots
and adding the error proves the cut bound. A further
\(\tau^{-20}\sqrt{D_{\rm id}}\) suffices for passage to full
kernels by Corollary~\ref{cor:squaredamplitudetail700}.

On the slice, Corollary~\ref{cor:parabolicmass701} gives the
entire parabolic block. For all positive closed ranks together
and the complementary open range, the two cut estimates are
Corollaries~\ref{cor:closedmoment693} and~\ref{cor:normalrange694},
respectively. Their union costs at most one further cut error,
by Corollary~\ref{cor:squaredamplitudetail700}. Adding their
norms gives the second inequality.
\end{proof}


\subsection{The marked second moment and the identity block}
\label{sec:markedmoment703}

We compare \(D_{\rm op}\) directly with the rank-zero mass at
the two prescribed primes. In this comparison the covolume of
the positive coordinates of the unipotent group and the
factorials in the real formal degree must both be retained.
Write \(w_v\) for the weight denoted by \(\kappa\) in
Proposition~\ref{prop:stationaryweights695}, so that \(w_v\)
is distinct in notation from the frame ratio \(\kappa_v\).
Put \(w=(\prod_v w_v)^{1/f}\), \(m=2n=2a+2\),
\(\rho_E=|{\rm disc}(E)|^{1/(2f)}\) and
\(c_E=2^{-f}\sqrt{|{\rm disc}(E)|}=(\rho_E/2)^f\).

\begin{lemma}[The terminal prime and the real formal degree]\leavevmode
\label{lem:calibratedlevel703}
For the terminal level of the shorter construction one can use
a prime \(N\equiv1\bmod Q^2\), \(Q=(q_0d)^{s+6}\), with
\[
 N_-\le N\le N_-\exp(\sqrt{\log b}),\qquad
 N_-=w^m(M\rho_E/(2\pi))^{m-1}b^{-15/2}.
\]
There is an algebraic integer \(\pi_N\in E\) with
\(\pi_N\bar\pi_N=N\). For every such choice at the same nodes,
\[
         d_\pi\ \ge\
          \frac{w^{fm(k-1)}}{\big(\prod_{j=k-1}^{r-1}j!\big)^f}.
\]
\end{lemma}
\begin{proof}
Here \(1\le w_v\le b\), since \(w_v=b\) or \(xb/x_1\)
with \(x_1\ge x\). We have \(\log \rho_E=O_{q_0,d}(\log(2x))\)
and \(b^{20}\le M\le b^{22}\). Thus for large \(x\),
\(m\log b\le\log N_-\le24m\log b\). The specified modulus
\(Q^2\) is bounded by a fixed power of \(x\), since
\(s\log 2=22\log m\). It is consequently at most a fixed
power of \(\log N_-\), uniformly on the interval, since
\(m\ge\sqrt{x}\). Apply the classical Siegel--Walfisz theorem
for rational primes in arithmetic progressions to \(Y=N_-\)
and \(2Y\), in the class \(1\bmod Q^2\). In the range
\(q\le(\log Y)^C\) it gives
\(\vartheta(Y;q,1)=Y/\varphi(q)+O_{C,B}(Y/(\log Y)^B)\)
for every fixed \(B\), where \(\vartheta\) sums \(\log p\)
over primes. Choosing \(B>C+2\) gives a strictly positive
difference. Since \(2\le\exp(\sqrt{\log b})\), this produces
a prime in the prescribed window. It splits in \(E\), since
\(q_1d\mid Q\) and \(E\subseteq\mathbb Q(\mu_{q_1d})\).
Also \(4\mid Q\) and \(\mathbb Q(i)\subset E\). The prime
\(N\) splits in the Euclidean ring \(\mathbb Z[i]\), so one
of the two conjugate prime ideals is generated by an element
of norm \(N\). This gives the required \(\pi_N\).
The construction retains \(p_0,p_1\), the fields, the weights
and the interval, and \(\log N\le30m\log b\) as used above.

For the real assertion use the formal degree formula with
compact-dual Haar measure in Theorem~\ref{thm:norm91}. Its
denominator is \(\prod_{j=1}^{r-1} j!\). The first \(k-1\)
Harish--Chandra entries in Proposition~\ref{prop:stationaryweights695}
are consecutive, contributing \(\prod_{j=1}^{k-2}j!\) to the
numerator. The other \(m\) entries have mutual gaps at least
one. Each of the \(m(k-1)\) cross gaps is at least \(w_v\):
the largest entry of the last block is \(-w_v-2r^8\) and
the smallest entry of the first is \(1-k/2\). Multiplication
of the resulting lower bounds proves the assertion.
\end{proof}

\begin{lemma}[A lower bound for the unipotent covolume]\leavevmode
\label{lem:unipotentvolume703}
In the normalized coordinates on the same group the Haar
measures used above satisfy
\[
       v_U\ \ge\ b^{-100n^3f}\,
        \Xi^{-k} M^{2akf}c_E^{ak}N^{-kf}.
\]
\end{lemma}
\begin{proof}
Take the successive quotients given by the frame action
\(U\to N_n\) (upper unitriangular matrices), then the positive
block \(B\) on its kernel \(U^0\), then the skew-Hermitian
center. On \(N_n\) Haar is Lebesgue measure on the strict
upper entries and on \(U^0\) it is \(dB\,dK\) in the
coordinates of Theorem~\ref{thm:realgram682}. These are
the quotient measures in Lemma~\ref{lem:shifts685}. At
each step the covolume multiplies by that of the image
lattice. Indeed, a measurable fundamental set in the image
can be lifted and multiplied by a fundamental set for the
intersection with the kernel. Uniqueness of the image class
and then the kernel class gives a fundamental set, and the
quotient integration formula gives the product of volumes.
The image lattices here are discrete by the integral conditions
in the coordinates below. They are cocompact since the original
rational unipotent lattice is cocompact, by the proof of
Proposition~\ref{prop:stabilizercancel680}.

In rational frame coordinates the action on \(N_n\) has image
in the strict upper entries of \(N_n(M\mathcal O_E)\), that
is, unitriangular matrices with those entries in
\(M\mathcal O_E\). A fundamental set for this latter lattice
is given by reducing entries modulo \(M\mathcal O_E\) in
order of increasing superdiagonal. Its volume is
\((M^{2f}c_E)^{n(n-1)/2}\ge1\). In normalized coordinates
each complex entry scales at \(v\) by \(s_{j,v}/s_{i,v}\),
of modulus at least \(b^{-23n}\). The contribution of this
quotient is therefore at least \(b^{-23n^3 f}\).

For the projection \(U^0\to B\), let a rational dual isotropic
basis vector be denoted by \(e_j^\vee\). The positive part of
\((u-1)e_j^\vee\) lies in \(M\Lambda_X\) for every
\(u\in\Gamma_c^0\), as in Lemma~\ref{lem:orbitfib693}.
In fact the terminal and bottom changes are zero by the fixed
frame; after subtracting the integral top summand, the
principal congruence puts the positive change in that lattice.
The normalized dual vector at \(v\) is
\(\Phi_v(e_j^\vee)/\bar s_{j,v}\), since its pairing with the
scaled arm vector has value one. The positive space has its
ordinary unitary metric. Each of the \(a\) columns of \(B\)
therefore lies in the Minkowski lattice \(M\Lambda_X\) scaled
by \(1/\bar s_{j,v}\) in its \(k\) complex coordinates at \(v\).
The product lattice containing this image has covolume
\[
   M^{2akf}(c_E^k N^f)^a
          \prod_{v,\,j\le a}|s_{j,v}|^{-2k}
     =M^{2akf} c_E^{ak}N^{af}\Xi^{-k}N^{-kf}.
\]
Here \(d(\Lambda_X)=N^f\) in the discriminant convention of
Section~\ref{sec:linkdisc687}, and the last frame vector has
scale \(N^{-1/2}\). The image contributes at least this
covolume, without needing surjectivity onto the containing
lattice.

In the central kernel the map adds \(K e_j\) in the first
arm to each dual vector, with rational entries in
\(\mathcal O_E\). Take off-diagonal upper complex coordinates
and imaginary diagonal coordinates as in the real Haar
formula. Before scaling, any nonzero difference in this
lattice has a coordinate of modulus at least one at some
embedding by the algebraic norm. In the real and imaginary
coordinates it thus has maximum modulus at least
\(1/\sqrt 2\). The translates of a real cube of side \(1/2\)
have disjoint interiors, giving covolume at least
\(2^{-a^2 f}\). Normalization scales a complex entry by
\(1/(s_{i,v}\bar s_{j,v})\) and a real diagonal coordinate by
\(|s_{i,v}|^{-2}\); in either case the modulus is at least
\(b^{-46n}\). Thus the center contributes at least
\(2^{-a^2f}b^{-46na^2 f}\). Multiplying the three lower
bounds and using \(N^{af}\ge1\) proves the lemma.
\end{proof}

\begin{theorem}[The second moment on both prescribed slices]\leavevmode
\label{thm:markedsecond703}
At all sufficiently large nodes of this family,
\[
    D_{\rm op}\le (M/b^6)^{kf}D_{\rm id},\qquad
    \sum_{s\in\Omega_i}\|\mathcal A_{G(F)}(s)-R(s)\|_{\rm HS}^2
                      \le b^{-12}D_{\rm id}\quad(i=0,1).
\]
Here \(\mathcal A_{G(F)}\) is the full signed array, with every
rational rank and the untruncated test vectors.
\end{theorem}
\begin{proof}
We first track the quantities whose logarithms are of order
\(akf\). Dropping such factors before taking ratios would lose
the comparison. Since \(|\log\Xi|\le 60n^2 f\log b\),
\(\rho=b^{37n}\), \(A=10^5 n^4\), and \(n^6=o(k)\), the
quantity \(b^{200n^3f}\rho^{9nAf}\Xi^{-2(A-1)}\) is at most
\(b^{o(k)f}\). All the little-\(o\) estimates here are uniform
as the lower endpoint tends to infinity; \(n\le\sqrt{x}\),
\(k\asymp x^5\) and \(\log b\ge(\log x)^6\) suffice. Using
Lemma~\ref{lem:unipotentvolume703} in the definition of
\(D_{\rm op}\) and \(n_p/D_{\rm id}=N^{-kf}/d_B\) gives
\[
 \frac{D_{\rm op}}{D_{\rm id}}
  \le b^{o(k)f}
       \frac{d_B}{d_\pi c_G^f}\,
       \frac{N^{kf}}{M^{2akf}c_E^{2ak}} .
\]
The identity \(c_G^{-1}=\prod_{l=k+1}^{r}S_l/S_k\), with
\(S_l=2\pi^l/(l-1)!\), is essential for the factorials.
Put \(h=r-k=2a+1=m-1\). Multiplying this identity by the
factorials in Lemma~\ref{lem:calibratedlevel703} gives
\[
 \frac1{d_\pi c_G^f}
 \le w^{-fm(k-1)}
   \big\{2^{h-1}\pi^{(h-1)k+h(h+1)/2}((k-1)!)^2\big\}^f .
\]
Indeed all \(j!\) with \(k\le j<r\) cancel, and the numerator
\((k-1)!\) from \(1/S_k\) is multiplied by the first surviving
factorial of the degree bound. Also
\[
     N^{kf}\le
       w^{mkf}(M\rho_E/(2\pi))^{hkf}
       b^{-15kf/2}\exp(kf\sqrt{\log b}) .
\]
Use \(c_E^{2ak}=(\rho_E/2)^{2akf}\) and \(h-1=2a\).
All powers in these three displays combine to give
\[
 \frac{D_{\rm op}}{D_{\rm id}}\le
 b^{-15kf/2+o(k)f}M^{kf}d_B
 \left\{
   w^m ((k-1)!)^2
   2^{h-1-k}\pi^{h(h+1)/2-k}\rho_E^k
 \right\}^{f}.
\]
The bracket is at most \(b^{o(k)}\). In detail,
\(m\log w\le m\log b\), \(\log((k-1)!)\le k\log k\),
\(\log\rho_E=O_{q_0,d}(\log x)\), and the remaining positive
exponents are \(O(n^2)\). The exponential of
\(kf\sqrt{\log b}\) has been absorbed in \(b^{o(k)f}\).
For the symmetric-power value space at each place the number
of degree-\(t_v\) monomials is at most \((t_v+1)^k\), with
\(t_v\le xb\). Thus \(d_B\le b^{(1+o(1))kf}\), using
\(\log x=o(\log b)\). The upper bound is
\((M b^{-13/2+o(1)})^{kf}\), which proves the first
inequality with room to spare.

Insert it into Theorem~\ref{thm:signedopen702}. After division
by \(\sqrt{D_{\rm id}}\) the resulting norm is at most
\(b^{-3kf/2}+4\tau^{-20}+\tau^{-3a/2}\). Since
\(\tau=b^{kf/4}\) and \(kf,a\) tend to infinity, this is at
most \(b^{-6}\) on the stated interval. Squaring gives the
second inequality.
\end{proof}

\begin{corollary}[The full block at the identity residue]\leavevmode
\label{cor:identityblock703}
The Hermitian operator \(A_1=\mathcal A_{G(F)}(1)\) on the
\(D_X\)-dimensional source satisfies
\[
        \|A_1-I\|_{\rm HS}^2\le b^{-12}D_X .
\]
At most \(4b^{-12}D_X\) of its eigenvalues, counted with
multiplicity, lie outside \([1/2,3/2]\).
\end{corollary}
\begin{proof}
For \(t\in B_{\rm res}\),
\(\mathcal A_{G(F)}(t)-R(t)=(A_1-I)R(t)\) by
Proposition~\ref{prop:residuetransport701}. Right multiplication
by this unitary preserves the Hilbert--Schmidt norm. Both
slices contain \(B_{\rm res}\). Sum there and divide the bound
of Theorem~\ref{thm:markedsecond703} by \(|B_{\rm res}|\).
Inversion of the rational term and the Haar change of variable
in Lemma~\ref{lem:finiteagreement702} give \(A_1^*=A_1\).
The spectral theorem expresses the squared norm as the sum
of \(|\lambda-1|^2\); each indicated eigenvalue costs at
least \(1/4\). These conclusions concern the complex
Poincar\'e pairing. The comparison of integral projected
theta columns in the main argument requires further estimates.
\end{proof}


\subsection{Poincar\'e series and projection at a marked prime}
\label{sec:poincareprojection704}

We realize the block matrices in a single automorphic Hilbert
space, with the Haar measures and terminal prime just used.
In this subsection write \(\mathbb A=\mathbb A_F\),
\([G]=G(F)\backslash G(\mathbb A)\), \(J=K(M)\) and
\(\mathcal G=K/J=\mathcal G_0\times\mathcal G_1\), where
\(\mathcal G_i=\mathrm{GL}_r(\mathbb F_{p_i})^f\).
The finite level is unrestricted within \(K\) at the other
places. Each residue matrix at a mark lifts to an integral
isometry: lift it invertibly on the chosen component and
use the inverse-transpose action on its paired component in
dual lattice bases. This gives the indicated equality for
\(\mathcal G\). On \(L^2([G])^J\) write \(T(s)\) for right
translation by a lift of \(s\). Put \(j_M=\operatorname{vol}(J)\).
Use product Haar on the adelic group and its quotient measure
for counting measure on \(G(F)\). Use inner products conjugate
linear in the first variable.

\begin{theorem}[The Petersson products of the full series]\leavevmode
\label{thm:poincaregram704}
There is a linear map \(\mathcal P:\mathcal H_X\to L^2([G])^J\)
such that for every \(s\in\mathcal G\)
\[
              \mathcal P^*T(s)\mathcal P=\mathcal A_{G(F)}(s).
\]
Its series are fixed by the archimedean convolution with
\(\prod_v d_{\pi_v}c_{\lambda_v}\) in the convention below.
The number of singular values of \(\mathcal P\) outside
\([1/\sqrt 2,\sqrt{3/2}]\), including the zeros on its source,
is at most \(4b^{-12}D_X\).
\end{theorem}
\begin{proof}
Use the product orthonormal basis of the value space, with
coordinates \(F_i\) of the full Bessel function. For the
source basis vector at class \(c\), coordinate \(i\), use the
measurable seed
\[
 \Phi_{c,i}(y,g_f)=(j_Mv_U)^{-1/2}\,
          {\bf1}_{h_cJ}(g_f)\,\beta_c(y)\overline{F_i(y)} .
\]
Define \(\mathcal P_{c,i}(g)=\sum_{\alpha\in G(F)}
\Phi_{c,i}(\alpha g)\). We first justify the construction,
including its inner products. Replace both seeds by their
absolute values and initially allow the sums to be infinite.
Unfolding the product of their periodizations, with the second
translated by \(s\), is justified by Tonelli's theorem. More
explicitly, sum over all ordered translates of a measurable
fundamental set for \(G(F)\) on the left. With the first
translate absorbed into the variable, the relative translate
in the second seed is \(\gamma^{-1}\). One obtains the integral
over the adelic group of the first seed times the sum of the
second seeds at \(\gamma^{-1}g k_s\), where \(k_s\) lifts \(s\).
The first support requires \(g_f\in h_cJ\). The second support
then either holds throughout this coset or is empty. The
condition is \(h_c^{-1}\gamma h_d\in k_sJ\), since \(k_s\)
normalizes \(J\). Thus the finite integration contributes
\(j_M\), and the resulting nonnegative real sum is bounded
by the absolute row estimate in
Theorem~\ref{thm:globalamplitudetail700} divided by \(v_U\).
In particular applying this to a seed paired with itself at
\(s=1\) shows that the periodization of its absolute value
lies in \(L^2([G])\). The series is absolutely convergent
almost everywhere; its sum is left rational invariant and
right \(J\)-invariant. There are finitely many coordinates.

Apply the same change of variables to the complex integrals,
now justified by the absolute bound. The two real coordinate
factors are \(F_i(y)\overline{F_j(\gamma^{-1}y)}\); the
remaining weights are \(\beta_c(y)\beta_d(\gamma^{-1}y)\).
The resulting matrix entry is exactly
\(\mathcal A_{G(F)}(s)_{cd,ij}\). This proves the stated
identity on the whole source.

For the real convolution group the series first by
\(\Gamma_c\backslash G(F)\). Along each left coset the finite
support is fixed, the \(F_i\) are invariant under \(\Gamma_c\),
and the class weights sum to one. The grouped seeds therefore
use \(\overline{F_i}\) with no class weight. They satisfy
\(\int C(h)W(gh)\,dh=W(g)\) with
\(C=\prod_v d_{\pi_v}c_{\lambda_v}\), by conjugating the
identity in Theorem~\ref{thm:normalizedcut699}. Here \(h\)
acts only at infinity. The kernel belongs to \(L^1(G_\infty)\).
The grouped absolute series also equals the periodization of
the absolute seed. Right translating that \(L^2\) function
and integrating it against \(|C(h)|\,dh\) gives an \(L^2\)
majorant by the integral triangle inequality and unitarity.
It is finite almost everywhere, permitting termwise
convolution there. This proves the real assertion; the
use of conjugate coordinates in the seed accounts for the
conjugation convention in the matrix identity. Finally the
squares of the singular values are the eigenvalues of
\(\mathcal P^*\mathcal P=A_1\); apply
Corollary~\ref{cor:identityblock703}.
\end{proof}

For an element \(x=\sum_s x(s)s\) of the finite group algebra
put \(R(x)=\sum_s x(s)R(s)\); the matrices \(R(s)\) vanish
off \(B_{\rm res}\). Write \(T(x)=\sum_s x(s)T(s)\).

\begin{theorem}[Compression of a projection at one mark]\leavevmode
\label{thm:singleprojection704}
For either \(i\), every orthogonal group-algebra projection
\(x=x^*=x^2\in\mathbb C[\mathcal G_i]\) satisfies
\[
   0\le R(x)\le I,\qquad
   \|\mathcal P^*T(x)\mathcal P-R(x)\|_{\rm HS}
                      \le b^{-6}\sqrt{D_X}.
\]
\end{theorem}
\begin{proof}
Realize the unitary representation induced from \(R|_{B_{\rm res}}\)
as functions \(f:\mathcal G\to\mathcal H_X\) satisfying
\(f(tg)=R(t)f(g)\), with squared norm the sum over left
\(B_{\rm res}\)-cosets of the squared value norms. The isometry
\(Iv\) supported on \(B_{\rm res}\) has value \(R(t)v\) at
\(t\in B_{\rm res}\). Compression of right translation by \(s\)
to this copy gives \(R(s)\): its value at the identity is
\((Iv)(s)\). Thus \(R(x)\) is the compression of an orthogonal
projection, and is a positive contraction.

Let \(j=1-i\). On \(\mathcal G\) form the matrix-valued
function \(D(s)={\bf1}_{\Omega_j}(s)
(\mathcal A_{G(F)}(s)-R(s))\).
It has squared norm at most \(b^{-12}D_{\rm id}\) for counting
measure, by Theorem~\ref{thm:markedsecond703}. Right
convolution \(D\mapsto[s\mapsto\sum_g x(g)D(sg)]\) is an
orthogonal projection on this Hilbert space, since the right
translations are unitary and form a representation. If
\(t\in B_{\rm res}\) and \(g\in\mathcal G_i\), then
\(tg\in\Omega_j\). Left transport from the proof of
Lemma~\ref{lem:finiteagreement702}, and the same formula for
\(R\), give
\[
    \sum_g x(g)D(tg)
       =R(t)\big(\mathcal P^*T(x)\mathcal P-R(x)\big).
\]
Take squared norms, sum over all \(t\in B_{\rm res}\), and
bound by the squared norm of the entire convolved function.
Since \(R(t)\) is unitary, division by \(|B_{\rm res}|\)
proves the estimate.
\end{proof}

\begin{corollary}[The cuspidal projection of the series]\leavevmode
\label{cor:cuspprojection704}
Take the two specified central projections \(e_i\) of
Section~\ref{sec:levelmass678}, put \(e=e_0e_1\), and set
\(\alpha_i=D_X^{-1}\operatorname{tr}R(1-e_i)\).
The image of \(T(e)\mathcal P\) lies in the cuspidal
\(L^2\)-space with the real convolution of
Theorem~\ref{thm:poincaregram704}. Its dimension satisfies
\[
    \operatorname{rank}(T(e)\mathcal P)\ \ge\
    D_X\left(\frac{\max(0,1-\alpha_0-\alpha_1-3b^{-6})}
                      {1+b^{-6}}\right)^2 .
\]
The \(\alpha_i\) are finite numbers in \([0,1]\), explicitly
the traces of the compressed finite-group projections.
\end{corollary}
\begin{proof}
Every finite type selected by \(e_i\) has a primitive factor
at some split place above \(p_i\). This factor is cuspidal
for \(\mathrm{GL}_r(\mathbb F_{p_i})\): it is the character
of a regular Frobenius orbit of length \(r\) with the one-box
partition in Green's parametrization
\cite[\S8]{XXV-XX-XX-green678}. In particular averaging it over the
unipotent radical of any proper finite parabolic gives zero.
For a proper \(F\)-parabolic with radical \(N_P\), consider
the constant term at an adelic \(g\). At the indicated split
place the conjugated parabolic by \(g_v^{-1}\) is conjugate
to a standard parabolic by an integral matrix: intersect
its flag with the standard lattice and extend bases along
the resulting flag of direct summands over the valuation
ring. The integral points in its unipotent radical reduce
onto the corresponding finite radical by lifting the free
entries. Averaging the form on the right over this compact
subgroup kills it. Its constant term at \(g\) is unchanged
by that average, since the subgroup becomes a subgroup of
\(N_P(\mathbb A)\) on the left of \(g\), and the measure on
\(N_P(F)\backslash N_P(\mathbb A)\) is right invariant.
Thus the constant term vanishes. For \(L^2\) representatives
these identities hold almost everywhere: the unipotent
quotient is compact by successive additive quotients, and
integration on its fibers gives locally integrable constant
terms by Fubini. Only finitely many reduced flag subgroups
occur, so the averaging identities can be taken simultaneously
outside a null set. Sum over the finitely many selected
types. The finite projection commutes with real convolution.

Let \(D_e=\mathcal P^*T(e)\mathcal P\). The commuting projections
give \(0\le I-T(e)\le T(1-e_0)+T(1-e_1)\). The trace of
\(A_1\) is at least \((1-b^{-6})D_X\), by
Corollary~\ref{cor:identityblock703} and Cauchy--Schwarz for
the matrix trace. The same trace inequality in
Theorem~\ref{thm:singleprojection704} bounds each of the two
losses by \((\alpha_i+b^{-6})D_X\). This proves the lower
bound for \(\operatorname{tr}D_e\) used in the numerator.
Also \(0\le D_e\le A_1\), so the minimax characterization of
eigenvalues gives \(\|D_e\|_{\rm HS}\le\|A_1\|_{\rm HS}
\le (1+b^{-6})\sqrt{D_X}\). If \(D_e\ne0\), apply
\((\operatorname{tr}D_e)^2\le
\operatorname{rank}(D_e)\|D_e\|_{\rm HS}^2\) to its nonzero
eigenvalues. For \(D_e=0\) the nonnegative trace lower bound
is zero. The kernel of \(D_e\) is the kernel of \(T(e)\mathcal P\),
which proves the asserted rank bound. Evaluating the finite
traces on this induced representation is necessary to obtain
a numerical lower bound for the prescribed cuspidal projection.
\end{proof}


\subsection{A finite Fourier transform and the projected cuspidal rank}
\label{sec:finitefourier705}

We evaluate the loss from the two finite projections by
counting cyclic vectors of a matrix over a finite field.
Let \(V_0\) have dimension \(r=k+2a+1\) over \(\mathbb F_q\),
\(a\ge1\), and put \(h=2a+1\). Order a basis in blocks as
\[
       e_1,\ldots,e_a,\quad W=\langle z\rangle\oplus X,\quad
       u_a,\ldots,u_1,\qquad \dim X=k .
\]
Let \(P\) be the upper block parabolic, \(U\) its unipotent
radical, and \(\mathfrak p,\mathfrak u\) the corresponding
matrix spaces, with \(\mathfrak u\) strictly upper in blocks.
Define \(C\) to advance one step in the string
\(e_1,\ldots,e_a,z,u_a,\ldots,u_1\), to kill \(u_1\),
and to vanish on \(X\). Fix a nontrivial additive character
\(\psi\) of \(\mathbb F_q\). On \(U\) the function
\(\chi(1+D)=\psi(\operatorname{tr}(CD))\) is a character:
the trace reads only entries between neighboring blocks,
and vanishes on \(\mathfrak u^2\). Extend it by zero on
\(\mathrm{GL}(V_0)\setminus U\), and denote its average under
conjugation by \(\chi^\natural\).

\begin{lemma}[Counting bases with prescribed successive images]\leavevmode
\label{lem:cyclicframes705}
For an endomorphism \(Y\) of \(V_0\) set
\[
 n_h(Y)=\#\{x\in V_0:x,Yx,\ldots,Y^{h-1}x
                            \text{ are linearly independent}\}.
\]
Then
\[
 \#\{g\in\mathrm{GL}(V_0): g^{-1}Yg\in C+\mathfrak p\}
           =|U|\,|\mathrm{GL}_k(\mathbb F_q)|\,q^k n_h(Y).
\]
\end{lemma}
\begin{proof}
Write \(\eta_1\) for the coordinate functional on the last
line \(u_1\). A basis counted on the left gives
\(x=g e_1,\ \eta=\eta_1 g^{-1}\) satisfying
\[
     \eta Y^j x=0\quad (0\le j<2a),\qquad
     \eta Y^{2a}x=1 .
\]
Indeed in that basis the only blocks below the diagonal
are the specified maps to the immediate next block. The
leading components of \(Y^j x\) for \(j<a\) are the successive
top vectors, and the leading middle component at \(j=a\) is
\(z\). Dually \(\eta Y^j\) reads the successive bottom lines,
and at \(j=a\) its middle component is \(z^*\), the functional
killing \(X\) with \(z^*(z)=1\). Pairing these components proves
the equations: for \(j<2a\) write \(j=l+t\) with \(l,t\le a\)
and at least one strictly less than \(a\); at \(j=2a\) only
\(z^*(z)\) survives.

Conversely, the equations determine the partial flag and
the simple-block identifications of every possible basis
above \(x,\eta\). Put \(x_j=Y^j x,\ \eta_j=\eta Y^j\)
and form
\[
    E_j=\langle x_0,\ldots,x_{j-1}\rangle,\qquad
    D_j=\bigcap_{t<j}\ker\eta_t \qquad (0\le j\le a).
\]
The first \(h\) powers on \(x\) are independent. In fact a
dependence at an earlier power would put every later power,
including \(Y^{2a}x\), in the span of the earlier ones,
contradicting its pairing with \(\eta\). The same argument
with the transpose gives independence for the first \(h\)
functionals. The flag is
\(E_1\subset\cdots\subset E_a\subset D_a\subset\cdots
\subset D_1\subset V_0\). On its top one-dimensional
quotients prescribe \(x_{j-1}\), and on its bottom quotients
\(D_{j-1}/D_j\) prescribe the covector \(\eta_{j-1}\).
The middle quotient \(D_a/E_a\) contains the class of \(x_a\),
and \(\eta_a\) induces a functional there: these assertions
use the zero moments up to \(2a-1\). Their pairing is one.
Identify the class of \(x_a\) with \(z\) and the induced
functional with \(z^*\). A basis respecting these data exists
by first choosing graded bases and then lifting them. It
gives the required lower blocks: \(Yx_j=x_{j+1}\) prescribes
the maps out of the top lines, including the map into the
middle. The identities \(\eta_j Y=\eta_{j+1}\) prescribe
the maps into the bottom lines, including the map out of
the middle. For example \(YD_j\subset D_{j-1}\) for
\(1\le j\le a\), and the coordinate on the next bottom
quotient is read by \(\eta_j\) on \(D_j\). These inclusions
and the corresponding top inclusions force every lower block
further than one step away to vanish.

For fixed \(x,\eta\) the changes of such a basis are arbitrary
upper unipotent block changes and arbitrary changes of basis
in the \(k\)-dimensional kernel in the middle graded space.
The singleton graded bases are fixed, and a middle
automorphism fixing both a vector and its covector of pairing
one is exactly the general linear group of that kernel.
There are \(|U|\,|\mathrm{GL}_k(\mathbb F_q)|\) choices.
Each independent \(h\)-tuple of powers permits precisely
\(q^{r-h}=q^k\) extensions of the specified functional
values to \(V_0\). This proves the count.
\end{proof}

\begin{theorem}[The conjugacy average at a split place]\leavevmode
\label{thm:conjugacyfourier705}
With \(v_j(q)=\prod_{l=1}^j(1-q^{-l})\), one has
\[
  1\le \sum_{t\in\mathrm{GL}(V_0)}|\chi^\natural(t)|^2
              \le v_k(q)/v_r(q).
\]
\end{theorem}
\begin{proof}
On the additive matrix algebra use the Fourier transform
\[
 \widehat f(Y)=q^{-r^2/2}\sum_D f(D)\psi(-\operatorname{tr}(YD)).
\]
It preserves the counting \(L^2\) norm by orthogonality
of the additive characters, since the matrix trace pairing
is nondegenerate. Take \(f(D)={\bf1}_{\mathfrak u}(D)
\psi(\operatorname{tr}(CD))\). Its transform is
\(q^{-r^2/2}|U|{\bf1}_{C+\mathfrak p}\), since the trace
annihilator of the strictly upper block matrices is the upper
block parabolic algebra. Fourier transform commutes with
averaging conjugates by the invariance of trace. Applying
Lemma~\ref{lem:cyclicframes705}, the transform of \(f^\natural\)
therefore equals
\[
          q^{-r^2/2}\frac{v_k(q)}{v_r(q)}\,q^{-r}n_h(Y).
\]
Indeed \(|U|=q^{a(2k+2a+1)}\),
\(|\mathrm{GL}_l(\mathbb F_q)|=q^{l^2}v_l(q)\), and
\(2a(2k+2a+1)+k^2+k-r^2=-r\).

For uniform independent \(Y,x\), the probability that the
first \(h\) powers starting at \(x\) are independent is
\(\prod_{j=0}^{h-1}(1-q^{j-r})=v_r(q)/v_k(q)\). First
exclude \(x=0\); after each independent list, condition
on the vectors already exposed and the images of all but
the last one. A uniform linear map with these prescribed
images sends the last independent vector uniformly to
\(V_0\). Excluding the current span gives the next factor.
Thus \(q^{-r}n_h(Y)\in[0,1]\) has mean \(v_r(q)/v_k(q)\)
over \(Y\), which also bounds its second moment. Plancherel's
identity now gives the upper bound. Finally \(f^\natural(D)\)
is supported where \(1+D\) is unipotent and agrees there
with \(\chi^\natural(1+D)\). This identifies the two counting
norms, and \(\chi^\natural(1)=1\) gives the lower bound.
\end{proof}

\begin{proposition}[The trace lost at an actual marked prime]\leavevmode
\label{prop:marktrace705}
For the prescribed central projection at \(p_i\) let
\(\delta_i\) be its discarded Plancherel mass as in
Proposition~\ref{prop:levelmass678}, and let \(\alpha_i\) be
the lost trace in Corollary~\ref{cor:cuspprojection704}.
Then on the same interval
\[
       |\alpha_i-\delta_i|
           \le 4\sqrt{f\delta_i}\,p_i^{-(k+1)/2}
                         \qquad(i=0,1)
\]
for all sufficiently large nodes.
\end{proposition}
\begin{proof}
For \(t\in\mathcal G_i\), so \(t=1\) at the other mark,
\(D_X^{-1}\operatorname{tr}R(t)\) is the scalar unipotent
character extended by zero off the finite unipotent image.
To see this, a nonzero diagonal class block gives a rational
Bessel lift \(u\eta\) from that class to itself. Its middle
component \(\eta\) preserves the integral middle quotient
and its positive lattice \(\Lambda_X\), since \(u\) acts
trivially on the middle quotient and \(\eta\) fixes the
terminal line. The automorphism group of that positive
lattice is finite. At the other odd split mark \(\eta\)
is principal on the middle quotient, so is the identity by
the torsion argument in Proposition~\ref{prop:rowmean690}.
Thus the lift is in \(U(F)\) and \(t\) is in the finite
unipotent image. On that image the scalar on every class
is the same, by Lemma~\ref{lem:finiteagreement702}. If a
residue is outside \(B_{\rm res}\) its entire matrix is zero.
This proves the trace identity.

At the mark the image of \(U\) is the full product of the
block upper unipotent radicals just used. On a chosen
split component lift any strictly upper block entries and
use the unitary inverse transpose on its paired component;
the integral form is perfect there and \(p_i\nmid N\).
Off the mark one can use the identity. In coordinates for
the chosen component the character reads the entries just
above the diagonal along the left arm, ending at \(z\).
The conjugate component supplies the opposite arm in reverse
order, starting at \(z\): its matrix is obtained from the
inverse transpose by the integral Gram identification. On
adjacent blocks the inverse of \(1+D\) has entries \(-D\),
since \(D^2\) skips a block. Thus the two middle coefficients
are a vector on the terminal line and its Hermitian dual,
with nonzero pairing modulo the mark since its norm is
\(-N\). Each simple coefficient is nonzero: the conductor
has precisely one factor \(p_i\) from \(M\) and the different
is a unit at this split place. Choose the middle basis starting
with this vector and continuing in the kernel of its dual.
Diagonal rescaling along the single string makes all its
coefficients one, successively, including the nonzero terminal
pairing. Hence up to conjugation the character on this factor
is the one of Theorem~\ref{thm:conjugacyfourier705}, with a
nontrivial additive character. This argument applies also
to the inverse character dictated by the real convention.

Let \(b_i^\natural\) be the conjugacy average of the normalized
trace on \(\mathcal G_i\). The character and its average
factor at its \(f\) places. As \(b_i^\natural(1)=1\), the
theorem gives, with \(q=p_i\) and counting norm,
\[
  \|b_i^\natural-{\bf1}_{\{1\}}\|_2^2
                   \le (v_k(q)/v_r(q))^f-1
                   \le 16 f q^{-(k+1)} .
\]
Indeed \(-\log(1-q^{-l})\le2q^{-l}\); the tail sum begins
at \(l=k+1\). Its logarithm after multiplication by \(f\)
is at most \(4fq^{-(k+1)}=o(1)\). Exponentiation gives
the displayed inequality.

Write \(x=1-e_i\). Its coefficients are constant on conjugacy
classes, and \(\sum_t|x(t)|^2=x(1)=\delta_i\), by taking the
identity coefficient of \(x^*x=x\) and the normalized trace
in the regular representation. By the trace identity,
\(\alpha_i=\sum_t x(t)b_i^\natural(t)\). Subtract the
identity coefficient and apply Cauchy--Schwarz to get the
claimed comparison.
\end{proof}

\begin{theorem}[Rank after both cuspidal projections]\leavevmode
\label{thm:projectedrank705}
On the fields and at all sufficiently large nodes of
Theorem~\ref{thm:markedsecond703}, let
\(\mathcal P_e=T(e_0e_1)\mathcal P\). Its Petersson matrix
and the dimension of its cuspidal image satisfy
\[
  \|\mathcal P_e^*\mathcal P_e-I\|_{\rm HS}
        \le 2 b^{-4}\sqrt{D_X},\qquad
  \operatorname{rank}\mathcal P_e\ge(1-4b^{-8})D_X .
\]
\end{theorem}
\begin{proof}
Proposition~\ref{prop:levelmass678} gives
\(\delta_0+\delta_1\le b^{-8}\), and the error at either
mark in Proposition~\ref{prop:marktrace705} is at most
\(b^{-20}\) for large nodes, since \(p_i\ge b^{10}\),
\(f\) is polynomial in \(x\), and \(k\to\infty\).
Each \(R(1-e_i)\) is a positive contraction with trace
\(\alpha_iD_X\). Its squared Hilbert--Schmidt norm is at
most that trace. By Theorem~\ref{thm:singleprojection704},
\[
  \|\mathcal P^*T(1-e_i)\mathcal P\|_{\rm HS}
             \le (b^{-6}+\sqrt{\alpha_i})\sqrt{D_X}.
\]
The positive difference
\(\mathcal P^*\mathcal P-\mathcal P_e^*\mathcal P_e\)
is bounded above by the sum of these two positive operators.
The Hilbert--Schmidt norm is monotone on positive operators
by eigenvalue monotonicity. Add
\(\|\mathcal P^*\mathcal P-I\|_{\rm HS}\) from
Corollary~\ref{cor:identityblock703}. The resulting factor
is at most
\(3b^{-6}+\sqrt{2(b^{-8}+2b^{-20})}\le2b^{-4}\)
for large \(b\). Each zero eigenvalue of the projected
Petersson matrix contributes one to its squared distance
from the identity. The bound for the rank follows. The image
is cuspidal by Corollary~\ref{cor:cuspprojection704}. This
is a dimension estimate relative to the source; the span
of its translates in the target and the integral comparison
for theta lifts still require control.
\end{proof}


\subsection{The second moment over all residues and the span of translates}
\label{sec:translates706}

Let \(\mathcal G=K/K(M)\), \(B=B_{\rm res}\), and let
\(\mathcal P_e=T(e_0e_1)\mathcal P\) be the projected
Poincar\'e map of Theorem~\ref{thm:projectedrank705}. We use
counting measure on \(\mathcal G\). Recall that \(z(s)\) is
the product of the norms of the marked components where the
two reduced full frames meet, with empty product one.

\begin{proposition}[A bound summed over every finite residue]\leavevmode
\label{prop:allresidue706}
On the fields and intervals of Theorem~\ref{thm:markedsecond703},
at all sufficiently large nodes one has
\[
  \sum_{s\in\mathcal G}\|\mathcal P^*T(s)\mathcal P\|_{\rm HS}^2
                \le (M/b^5)^{kf}\,|B|D_X .
\]
More precisely, on the residues with \(z(s)\ge Z\), \(Z\ge1\),
the contribution from the open terms \(T_>\) of
Theorem~\ref{thm:signedopen702} has counting Hilbert--Schmidt
norm at most
\[
 \big\{(\tau^2+\tau^3)Z^{-k}+4\tau^4(2/M)^{kf}\big\}^{1/2}
          \sqrt{D_{\rm op}}+2\tau^{-20}\sqrt{D_{\rm id}} .
\]
\end{proposition}
\begin{proof}
Theorem~\ref{thm:poincaregram704} identifies the left side
of the first display with the moment of the full signed
array. The set \(T_>\) is invariant under the rational Bessel
group on both sides by invariance of its two defining leading
determinants. Lemma~\ref{lem:finiteagreement702} therefore
eliminates every residue with disagreeing characters before
the square is expanded. In each remaining block, unfold it
by Proposition~\ref{prop:unfoldopen702}. After removing the
cut error, the norm is bounded by the sum of
\(C_\infty M^{akf}G_\gamma^{-k/2}\) over its retained near
representatives. This follows from \(e\ge k\), \(X_\gamma\ge1\),
and \(G_\gamma=M^{2af}X_\gamma\) in the indicated open
range, as in Theorem~\ref{thm:signedopen702}. The passage
to the cut and the replacement of the cut integral by its
majorant each cost at most \(\tau^{-20}\sqrt{D_{\rm id}}\)
in counting norm by Corollary~\ref{cor:squaredamplitudetail700}.
That corollary already sums over all residues.

For the square of the majorant expand in ordered pairs
within each block. Distinct chosen representatives in that
block belong to distinct orbits for the smaller pointwise
frame fixers. The pairs are therefore among those counted
by Theorem~\ref{thm:paircount691}. If their full links agree,
their positive corank \(\ell\) is zero or one, and
Theorem~\ref{thm:samelink687} bounds the pair weight
\((G_\gamma G_\eta)^{-k/2}\) by
\(Z^{-k}d(L_{12})^{-k}\).
If the links differ, Theorem~\ref{thm:allpair689} bounds
it by \((2/M)^{\max(1,\ell-1)kf}d(L_{12})^{-k}\).
Summing the first bound gives \((\tau^2+\tau^3)Z^{-k}\)
in units of \(D_{\rm op}\). In the second, the terms for
\(\ell=0,1,2\) share the gain \((2/M)^{kf}\), and from
\(\ell=2\) on the ratio is \(\tau(2/M)^{kf}<1/2\).
Their sum is at most \(4\tau^4(2/M)^{kf}\). This proves
the restricted estimate.

Take \(Z=1\). Its expression in braces is at most \(4\tau^3\)
for large nodes. The parabolic terms have norm
\(\sqrt{|\mu(E)|D_{\rm id}}\) by
Corollary~\ref{cor:parabolicmass701}. The positive closed
terms and the complementary open range have cut norms
\(\tau^{-20}\sqrt{D_{\rm id}}\) and
\(\tau^{-3a/2}\sqrt{D_{\rm id}}\) as upper bounds by
Corollaries~\ref{cor:closedmoment693} and~\ref{cor:normalrange694}.
One application of Corollary~\ref{cor:squaredamplitudetail700}
to their union bounds its passage back to the full kernels.
Thus, divided by \(\sqrt{D_{\rm id}}\), the norm of the entire
array is at most
\[
  2\tau^{3/2}(M/b^6)^{kf/2}
        +\sqrt{|\mu(E)|}+4\tau^{-20}+\tau^{-3a/2},
\]
using Theorem~\ref{thm:markedsecond703}. The square of the
first term is \(4(M/b^{21/4})^{kf}\). The roots of unity
inject into \(\mathbb F_{p_0}^\times\) as in
Corollary~\ref{cor:parabolicmass701}, so their number is
at most \(b^{11}\). Since \(M\ge b^{20}\) and \(kf\to\infty\),
the whole display is at most \((M/b^5)^{kf/2}\) for large
nodes. This gives the first bound.
\end{proof}

Define a real number determined by the two finite levels by
\[
             L=\frac{|\mathcal G|}{|B|(M/b^5)^{kf}} .
\]
It grows with the nodes. Indeed, at a split component of
residue size \(q\) the finite image of the Bessel group
has order \(q^{a(2k+2a+1)}|\mathrm{GL}_k(\mathbb F_q)|\).
The upper block unipotents lift as in
Proposition~\ref{prop:marktrace705}. The middle isometry
fixes the terminal vector and its dual and acts arbitrarily
on their common \(k\)-dimensional complement. Integral lifts
exist because the terminal norm is a unit there. On the
paired component the action is determined by the perfect
pairing. These choices are independent at distinct finite
places. Consequently
\[
 L=b^{5kf}M^{[(2a+1)k+(2a+1)(a+1)]f}
             \prod_{i=0}^1\left(\frac{v_r(p_i)}{v_k(p_i)}\right)^f .
\]
Here \(v_j\) is the product in Theorem~\ref{thm:conjugacyfourier705}.
The final product tends to one, since the negative of its
logarithm is at most \(4f\sum_i p_i^{-k-1}\). In particular
\(L\ge1\) at every sufficiently large node.

\begin{theorem}[Dimension obtained by finite translates]\leavevmode
\label{thm:translates706}
The cuspidal subspace generated by all
\(T(s)\mathcal P_e\mathcal H_X\), \(s\in\mathcal G\),
has dimension at least \((1-2b^{-4})^2 L D_X\). Moreover,
for every integer \(J_0\ge1\), at most \(J_0\) translates
generate a space of dimension at least
\[
      \frac{(1-2b^{-4})^2D_X}{L^{-1}+(1+b^{-6})^2/J_0}.
\]
\end{theorem}
\begin{proof}
Put \(Z=\mathcal P\mathcal P^*\) and
\(A=|\mathcal G|^{-1}\sum_s T(s)ZT(s)^*\). All traces in
the proof are of finite-rank operators. Expanding the square
and counting ordered pairs of residues with each prescribed
relative translate gives
\[
 \operatorname{tr}A^2
   =|\mathcal G|^{-1}\sum_s\operatorname{tr}
                     (ZT(s)ZT(s)^*)
   =|\mathcal G|^{-1}\sum_s
                     \|\mathcal P^*T(s)\mathcal P\|_{\rm HS}^2
   \le D_X/L .
\]
Write \(E=T(e_0e_1)\), \(Z_e=EZE\), \(A_e=EAE\).
The projection \(E\) commutes with each \(T(s)\) by
centrality. Consequently \(A_e\) is the same average with
\(Z_e\) in place of \(Z\). Compression cannot increase
Hilbert--Schmidt norm, and
\[
 \operatorname{tr}A_e=\operatorname{tr}Z_e
       =\operatorname{tr}(\mathcal P_e^*\mathcal P_e)
                  \ge (1-2b^{-4})D_X .
\]
For the last step pair the difference from the identity
in Theorem~\ref{thm:projectedrank705} with the identity
matrix by Cauchy--Schwarz. The range of \(A_e\) is precisely
the required span: its kernel is the intersection of the
orthogonal complements of the finitely many images, by
positivity of every summand. The inequality
\((\operatorname{tr}A_e)^2\le\operatorname{rank}(A_e)
\operatorname{tr}A_e^2\), applied to its eigenvalues,
proves the first assertion.

For the second choose \(J_0\) independent uniform residues
and let \(A'\) be the empirical average of their conjugates
of \(Z_e\). Its trace is still \(\operatorname{tr}Z_e\).
Independence in the off-diagonal terms gives
\[
 \mathbb E\big[\operatorname{tr}((A')^2)\big]
   =(1-J_0^{-1})\operatorname{tr}(A_e^2)
                         +J_0^{-1}\operatorname{tr}(Z_e^2).
\]
We have
\(\|Z_e\|_{\rm HS}\le\|Z\|_{\rm HS}=
\|\mathcal P^*\mathcal P\|_{\rm HS}\le(1+b^{-6})\sqrt{D_X}\)
by Corollary~\ref{cor:identityblock703}. Some tuple has
trace of its squared average at most the displayed mean.
Apply the same rank inequality to this tuple. Repetitions
do not affect its span.
\end{proof}

\begin{corollary}[A translate modulo a prescribed subspace]\leavevmode
\label{cor:avoidspace706}
For every subspace \(W\) of finite dimension \(d\) in the
target, some \(s\in\mathcal G\) has image modulo \(W\)
of dimension at least
\[
 D_X\left(\frac{\max\{0,1-2b^{-4}-\sqrt{d/(LD_X)}\}}
                       {1+b^{-6}}\right)^2 .
\]
\end{corollary}
\begin{proof}
If \(Q\) is the orthogonal projection onto \(W\), the
average trace lost to \(W\) is \(\operatorname{tr}(QA_e)\),
which is at most \(\sqrt{d}\|A_e\|_{\rm HS}\le\sqrt{dD_X/L}\).
Choose a residue with loss at most this average. Compress
its conjugate of \(Z_e\) to \(W^\perp\). The trace of the
compression is the original trace less this loss and is
nonnegative. Its Hilbert--Schmidt norm is at most
\((1+b^{-6})\sqrt{D_X}\). The trace--rank inequality gives
the assertion. Its scale is the explicit lower dimension
\(LD_X\). Passing to a bound in terms of a prescribed
fraction of the full target requires a sharper comparison
with the automorphic mass.
\end{proof}

\section{Spanning by translates and the terminal norm sphere}\label{nav:spanning-and-sphere}

\subsection{The good source dimension and the number of translates}
\label{sec:fewtranslates868}
Use the shorter interval, the marked group \(\mathcal G=K/J\)
on \(Y_0\), and the map \(P=T(e_0e_1)\mathcal P\) of
Theorem~\ref{thm:projectedrank705}. Write \(d_X=\dim\mathcal H_X\)
for its domain dimension and \(d=\dim H_{Y_0}=D_x\) for the
common good source of the determinant. These are different
spaces. The finite-level operators \(T(g)\) preserve \(H_{Y_0}\)
and act unitarily. Throughout this subsection \(b=b_x\).
The terminal level auxiliary prime \(N\) may be chosen as in
Lemma~\ref{lem:calibratedlevel703}. Its choice does not enter
the definition of \(d\).

\begin{proposition}[Multiplicity on the common good source]
\leavevmode\label{prop:goodmass868}
At every sufficiently distant node one has \(d>0\).
On the positive definite rank-\(s\) space \(I_s\), let
\(\mathcal C_s\) be the set of maximal-level lattice classes
and \(\Gamma_L\) the automorphism group of class \(L\).
Let \(V_\lambda\) be the compact coefficient representation
of the common Harish--Chandra parameter, with unit highest
vector and \(d_\lambda=\dim V_\lambda\). Put
\(\operatorname{mass}_s=\sum_{L\in\mathcal C_s}|\Gamma_L|^{-1}\).
Then
\begin{equation}
 d=|\mathcal G|\sum_{L\in\mathcal C_s}\frac1{|\Gamma_L|}
        \sum_{\gamma\in\Gamma_L}
            \operatorname{tr}_{V_\lambda}(\gamma)\,
                      e(\bar\gamma)
 \ \le\ |\mathcal G|\,d_\lambda\,\operatorname{mass}_s .
 \label{eq:goodmass868}
\end{equation}
Here \(e=e_0e_1\) is expressed in the group algebra with
counting convolution, and \(\bar\gamma\) denotes reduction
in the local frames of \(L\). Use the coefficient representation
in the left automorphy rule for \(V_\lambda\).
\end{proposition}
\begin{proof}
The rank bound of Theorem~\ref{thm:projectedrank705} is at
least \((1-4b^{-8})d_X\). Its domain dimension is
\(|\mathcal C|d_B\) by Theorem~\ref{thm:besselmatrix701},
with the identity class present and a nonzero compact type.
The image is in the fixed real discrete-series line and the
principal finite level by Theorem~\ref{thm:poincaregram704},
and has already been projected by \(e\). The parameters at
the node agree: the line in Proposition~\ref{prop:stationaryweights695}
has the \(k-1\) consecutive entries and the \(m\) shifted
entries of Theorem~\ref{thm:commonsource806} on \(Y_0\).
Thus \(P\) takes values in that theorem's \(H_{Y_0}\), and
the positive rank gives \(d>0\).

For the dimension calculation realize this good space at
the same finite places on \(I_s\). The packet argument in
Theorem~\ref{thm:commonsource806} applies with the signs all
positive. Indeed the unimodular lattices are locally isometric
by its finite frame calculation, the same two finite
supercuspidal components give a simple generic parameter,
and the all-positive assignment of the same regular entries
is the compact member of its real packet. Keeping the finite
parts and making that assignment gives multiplicity one
by the same simple-parameter multiplicity formula, in both
directions. Projection to the specified real line leaves
the finite multiplicity unchanged. On the definite group
this line projection identifies scalar forms with
\(V_\lambda\)-valued functions on the finite classes, by
compact Schur orthogonality.

At a maximal-level class \(L\) these functions at principal
level are the invariants in
\(V_\lambda\otimes\CC[\mathcal G]\) for \(\Gamma_L\) acting
by its coefficient representation and by left translation
on residue arguments. Project onto the invariants by averaging
this action. It commutes with the right action \(R(e)\).
For the convention \(R(t)f(g)=f(gt)\), the trace of left
translation by \(h\) followed by \(R(e)\) on the regular
space is \(\sum_{g\in\mathcal G} e(g^{-1}hg)=|\mathcal G|e(h)\).
The last equality uses centrality. Taking the trace of the
product of the two commuting projectors gives the equality
in \eqref{eq:goodmass868}. For the upper bound discard \(e\).
No nonidentity \(\gamma\in\Gamma_L\) reduces to the identity:
the principal subgroup at an odd original mark is torsion
free by the proof of Lemma~\ref{lem:terminalintegral867}.
Left translation on \(\mathcal G\) has trace zero for every
other group element. Thus the whole principal-level dimension
is exactly \(|\mathcal G|d_\lambda\operatorname{mass}_s\).
Projection can only decrease dimension.
\end{proof}

Recall the explicit number
\[
     L=\frac{|\mathcal G|}
          {|B_{\rm res}|(M/b^5)^{kf}}\ge1
\]
of Section~\ref{sec:translates706}. A stronger moment estimate
can be used in the following argument by replacing \(L\) with
any number \(L_*\ge1\) for which the positive operator
\[
         A=|\mathcal G|^{-1}\sum_{g\in\mathcal G}
                          T(g)PP^*T(g)^*
\]
satisfies \(\operatorname{tr}(A^2)\le d_X/L_*\).
For the actual columns \(L_*=L\) is available:
Proposition~\ref{prop:allresidue706} gives this bound before
the projection by \(e\), by expanding the trace as in
Theorem~\ref{thm:translates706}. Compressing that average by
the orthogonal projector \(T(e)\) gives \(A\), since \(e\) is
central, and cannot increase its Hilbert--Schmidt norm.

\begin{theorem}[Independent input columns from few translates]
\leavevmode\label{thm:fewtranslates868}
For every positive integer \(j\), the sum of at most \(j\)
spaces of the form \(T(g)P\mathcal H_X\) inside \(H_{Y_0}\)
has dimension at least
\begin{equation}
         L_*d_X\big(1-e^{-j/(2L_*)}-8b^{-4}\big).
 \label{eq:fewtranslates868}
\end{equation}
In particular with the explicit \(L\) there is a finite
set \(\mathcal T\subseteq\mathcal G\) with
\[
 |\mathcal T|d_X\le5d\log(1+2b^2),\qquad
 \dim\sum_{g\in\mathcal T} T(g)P\mathcal H_X
                             \ge (1-2b^{-2})L d_X .
\]
Independent vectors counted by this bound can be selected from
the translated images of any spanning family of \(\mathcal H_X\).
They may be used as the first \(n\) columns of
Theorem~\ref{thm:longframe867}.
\end{theorem}
\begin{proof}
Set \(\epsilon=2b^{-4}\) and \(N_*=L_*d_X\).
Theorem~\ref{thm:projectedrank705} gives
\(\|P^*P-I\|_{\rm HS}\le\epsilon\sqrt{d_X}\). Consequently
\(\operatorname{tr}A=\operatorname{tr}P^*P
\ge(1-\epsilon)d_X\). If the currently selected span has
dimension \(v\) and orthogonal projector \(Q\), then
\[
 \operatorname{tr}(QA)\le\sqrt{v\,\operatorname{tr}(A^2)}
                        \le d_X\sqrt{v/N_*}.
\]
By averaging choose \(g\) for which, with \(P_g=T(g)P\),
the positive operator \(B=P_g^*(1-Q)P_g\) has trace at
least \(\operatorname{tr}A-\operatorname{tr}(QA)\).
The rank of \(B\) is the exact dimension gained by adjoining
the columns of \(P_g\), because it is the Gram operator of
\((1-Q)P_g\). If \(Q_B\) is the orthogonal projector onto
its range in the domain, \(0\le B\le P^*P\) gives
\[
 \operatorname{tr}B\le\operatorname{tr}(Q_BP^*P)
       \le\operatorname{rank}B+
             \|Q_B\|_{\rm HS}\|P^*P-I\|_{\rm HS}
       \le\operatorname{rank}B+\epsilon d_X.
\]
This last estimate uses the near-isometry on the source.
It retains a linear gain in the number of missing columns.

For \(v\le N_*\) use \(1-\sqrt{v/N_*}\ge(1-v/N_*)/2\).
The new dimension is at least
\(v+(N_*-v)/(2L_*)-2\epsilon d_X\).
Writing \(s_j=(N_*-v_j)_+\) for the shortfall after \(j\)
choices, monotonicity of the span therefore gives
\[
 s_{j+1}\le(1-1/(2L_*))s_j+2\epsilon d_X.
\]
The inequality also holds when \(v_j>N_*\), since then the
next shortfall is zero. Starting with \(s_0=N_*\), summation
of the geometric progression bounds \(s_j\) by
\(N_*(e^{-j/(2L_*)}+4\epsilon)\), proving the first assertion.

For \(L_*=L\) choose \(j=\lceil2L\log(b^2)\rceil\); then
\(e^{-j/(2L)}\le b^{-2}\) and \(8b^{-4}\le b^{-2}\) eventually.
Also \(d\ge(1-\epsilon)^2 L d_X\) by
\((\operatorname{tr}A)^2\le d\,\operatorname{tr}(A^2)\).
It follows that \(j d_X\le5d\log(1+2b^2)\) for all
sufficiently large nodes, using \(L\ge1\). Repetitions of
translates may be discarded. The images of the indicated
spanning family span the same sum, so elementary basis
selection gives as many independent columns as its dimension.
\end{proof}

The unconditional rank above is measured in units of \(Ld_X\).
For example, a bound
\(\operatorname{tr}(A^2)\le(1+b^{-2})d_X^2/d\)
allows one to take \(L_*=d/((1+b^{-2})d_X)\) whenever this is
at least one; \eqref{eq:fewtranslates868} then has the scale of
the full good source. The estimate in
Proposition~\ref{prop:allresidue706} uses absolute majorants on
the open orbits. Obtaining this precision for the second moment
on all residues is a remaining issue in applying the terminal
covolume inequality.
\clearpage

\subsection{The contribution of equal links to the full second moment}
\label{sec:equallinks869}

Let \(b=b_x\) on the shorter interval and use the terminal prime of
Lemma~\ref{lem:calibratedlevel703}. In this subsection \(r=k+2a+1\)
is the rank called \(s\) in Proposition~\ref{prop:goodmass868}.
Set \(h=2a+1=m-1\), write \(\mathfrak m_j\) for the mass of the
standard positive genus in rank \(j\), and put
\[
 \mathcal M=\frac{d_X^2}{d_\lambda\mathfrak m_r},\qquad
 D_{\rm id}=|B_{\rm res}|d_X,\qquad \rho_E=|\operatorname{disc}E|^{1/(2f)}.
\]
These normalizations use counting measure on the marked residue group.
Indeed, averaging the translates of \(\mathcal P\mathcal P^*\) divides
the squared moment of \(\mathcal P^*T(s)\mathcal P\) by
\(|\mathcal G|\), while
\(d\le|\mathcal G|d_\lambda\mathfrak m_r\) by
Proposition~\ref{prop:goodmass868}.

\begin{lemma}[The normalization on the shorter interval]\leavevmode
\label{lem:massscale869}
At every sufficiently large node,
\begin{align}
 \frac{\mathcal M}{D_{\rm id}}
 &= \frac{d_B (N^k-1)^f\mathfrak m_k}
          {M^{fa(2k+h)}d_\lambda\mathfrak m_r},
 \label{eq:massscale869}\\
 \log\frac{\mathcal M}{D_{\rm id}}
 &= kf\log(M b^{-15/2})+\log d_B+f\log((k-1)!)
       +o(kf\log b).
 \label{eq:masslog869}
\end{align}
The error is uniform on the interval. In particular
\[
  (M/b^8)^{kf}D_{\rm id}\le\mathcal M
          \le (M/b^6)^{kf}D_{\rm id},\qquad
  D_{\rm op}\le b^{o(1)kf}\mathcal M ,
\]
where \(D_{\rm op}\) is the majorant in
Section~\ref{sec:signedopenslice702}.
\end{lemma}

\begin{proof}
At the two original marks the \(H=\mathrm U(X_0)\) level fixes
a full frame. Write \(\bar H\) for the product of the groups
\(\mathrm{GL}_k(\mathbb F_{p_i})\) at their split places. At each
of the \(f\) places of \(F\) over \(N\), the level in the chosen
component of \(E\) fixes one nonzero covector. To check the exact
index, use the middle basis in the proof of
Proposition~\ref{prop:rowmean690}. At a component where \(\pi\)
divides \(N=\pi\bar\pi\), \(\bar\pi\) is a unit. Preservation of
the lattice by a middle matrix and its inverse requires integrality
and the distinguished row congruence \(e_*^t A=e_*^t\bmod\pi\).
Conversely these conditions make all the coefficients integral:
the possible denominators in the \(z\)-coordinate are those of
\(A_{*j}/\pi\) for \(j\ne *\) and \((1-A_{**})/N\). The other
coordinates have unit denominators. The same row condition holds for
\(A^{-1}\), and the inverse transpose preserves the dual lattice in
the paired component. Reduction of \(\mathrm{GL}_k\) acts
transitively on nonzero rows. At the other primes the middle basis
differs from the standard one by an integral invertible change,
apart from the two prescribed principal levels. It follows that a
standard class \(L\) acquires
\(|\bar H|(N^k-1)^f\) structures before quotient by its automorphisms.
Their action is free by the odd principal level argument of
Proposition~\ref{prop:rowmean690}. Therefore
\[
 d_X=d_B|\bar H|(N^k-1)^f\mathfrak m_k .
\]
At a split mark the unipotent radical is the upper block radical
with block sizes \(1,\ldots,1,k+1,1,\ldots,1\), with \(a\)
singletons on either side. Its strict off-block entries number
\((r^2-(k+1)^2-2a)/2=a(2k+h)\). The residual middle subgroup
is \(\mathrm{GL}_k\) on the complement of the terminal line.
All these entries lift integrally, as in
Section~\ref{sec:translates706}. Thus
\(|B_{\rm res}|=|\bar H|M^{fa(2k+h)}\). Substitution gives
\eqref{eq:massscale869}.

We evaluate the part of the mass that is of order \(kf\log b\).
Here \(E/F\) is unramified at the finite places. Indeed \(E\)
contains \(\QQ(i)\) and the imaginary quadratic field of odd
discriminant \(-d_K\); inertia at any finite prime fixes at least
one of them and complex conjugation fixes neither. Hence
\(|\operatorname{disc}F|=\rho_E^f\).
Lemma~\ref{lem:mass209} gives
\(\mathfrak m_{i+1}/\mathfrak m_i=2^{-f}|L_F(-i,\eta^{i+1})|\)
with \(\eta\) the quadratic character of \(E/F\).
For \(j=i+1\) its abelian factors are the \(f\) primitive
Dirichlet characters of parity \(j\) in the character group of
\(E\), taken with their primitive Euler products. The conductors
for even parity have product \(|\operatorname{disc}F|\)
by the conductor-discriminant formula applied to \(F\).
For odd parity they have product
\(|\operatorname{disc}E|/|\operatorname{disc}F|
=|\operatorname{disc}F|\) by the same formula for \(E\).
For a primitive factor of conductor \(q\), the functional
equation at \(1-j\) gives the absolute value
\[
   2(j-1)!\,(2\pi)^{-j}q^{j-1/2}|L(j,\bar\chi)|.
\]
Both \(|L(j,\bar\chi)|\) and its reciprocal are at most
\(\zeta(j)\) by the Euler product. Since
\(\log\zeta(j)\le\sum_{n\ge2}n^{-j}\ll 2^{-j}\) for \(j\ge3\),
the factor \(2^f\) cancels and summation gives
\begin{equation}
 \log\frac{\mathfrak m_r}{\mathfrak m_k}
 = f\sum_{i=k}^{r-1}
 \left(\log i!+(i+\tfrac12)\log\rho_E-(i+1)\log(2\pi)\right)
       +O(f2^{-k}).
 \label{eq:massgamma869}
\end{equation}

Let \(w_v\) be the local weight of Lemma~\ref{lem:calibratedlevel703},
\(w_{\min}=\min_v w_v\), and \(w=(\prod_v w_v)^{1/f}\).
The first \(k-1\) entries of the compact Harish--Chandra parameter
are \(k/2-1-i\), \(0\le i\le k-2\); the last \(m\) entries are
\(-w_v-2jr^8\), \(1\le j\le m\), by
Proposition~\ref{prop:stationaryweights695}. In the Weyl degree
formula the first block cancels the factors \(1!,\ldots,(k-2)!\).
Each mixed gap divided by \(w_v\) lies between one and
\(1+(2mr^8+k)/w_{\min}\). The product of the gaps internal
to the last block is
\((2r^8)^{m(m-1)/2}\prod_{j=1}^{m-1}j!\).
Consequently
\[
 \log d_\lambda=fm(k-1)\log w-f\sum_{i=k-1}^{r-1}\log i!
             +\Delta,\qquad
 0\le\Delta\ll f\left(m^2\log(2r)
             +\frac{mk(2mr^8+k)}{w_{\min}}\right).
\]
Here \(w_{\min}\ge b^{1-10^{-6}}\), \(m=O(\sqrt{x})\),
\(k\asymp x^5\), and \(\log b\ge(\log x)^6\), so that
\(\Delta=o(kf\log b)\).

Write \(N=w^m(M\rho_E/(2\pi))^h b^{-15/2}e^\theta\),
\(0\le\theta\le\sqrt{\log b}\). Insert the last two formulas in
\eqref{eq:massscale869} and use
\(\log(N^k-1)=k\log N+O(N^{-k})\). The factorial sums leave
\(f\log((k-1)!)\). The coefficient of \(\log M\), before
splitting off its leading term, is \(f(hk-a(2k+h))=f(k-ah)\).
The remaining terms after the three displayed terms of
\eqref{eq:masslog869} are
\[
 f\left(m\log w-ah\log M-\tfrac12h^2\log\rho_E
       +\tfrac12 h(h+1)\log(2\pi)+k\theta\right)
       -\Delta+O\big(f(2^{-k}+N^{-k})\big).
\]
Each has size \(o(kf\log b)\), since \(m^2=o(k)\),
\(w\le b\), \(M\le b^{22}\), and \(\log\rho_E=O_{q_0,d_K}(\log x)\).
This proves \eqref{eq:masslog869}. The symmetric powers at infinity
give \(1\le d_B\le(xb+1)^{kf}\), and
\(\log((k-1)!)\le k\log k=o(k\log b)\), which give the stated
two-sided inequalities. Finally the proof of
Theorem~\ref{thm:markedsecond703}, before bounding \(d_B\), gives
\[
 D_{\rm op}/D_{\rm id}
 \le (M b^{-15/2})^{kf}d_B\exp(o(kf\log b)).
\]
Indeed the logarithms of the remaining factors there are
\(O(f(m\log w+k\log k+k\log\rho_E+h^2))+
o(kf\log b)\). Comparing with \eqref{eq:masslog869}, retaining
the same \(d_B\) in both expressions, proves the final bound.
\end{proof}

We now retain the signs of the orbital integrals. In each block
\((c,s,c')\) restrict to the large open range \(T_>\) of
Theorem~\ref{thm:signedopen702}, with agreeing characters.
Use the cut orbit representatives of
Proposition~\ref{prop:unfoldopen702} and let \(V_\gamma\) be
the actual signed matrix integral there, including \(v_U^{-1}\).
Choose one of the permitted near representatives for each orbit.
Its full link is the unitary map \(J_\gamma\) on \(X_0\oplus E\)
defined in Section~\ref{sec:differentlinks688}. For the product
\(z(s)\) of the norms of the meeting marked components, set
\begin{equation}
 \mathcal E=\sum_{\substack{c,s,c'\\z(s)<b^f}}
       \ \sum_J
       \left\|\sum_{\gamma:\,J_\gamma=J}V_\gamma\right\|_{\rm HS}^2 .
 \label{eq:linkenergy869}
\end{equation}
The inner sums are in a single block. The comparison below holds
for every such choice of near representatives.

\begin{theorem}[Removal of the pairs with different full links]
\leavevmode\label{thm:equallinks869}
Let \(\mathcal P\) be the unprojected Poincar\'e map on these same
fields, levels and intervals. At all sufficiently large nodes,
\[
 \left|
 \left(\sum_{s\in\mathcal G}\|\mathcal P^*T(s)\mathcal P\|_{\rm HS}^2
                         \right)^{1/2}-\sqrt{\mathcal E}\right|
       \le b^{-kf/10}\sqrt{\mathcal M}.
\]
In particular an upper estimate
\(\mathcal E\le(1+\tfrac12 b^{-2})\mathcal M\) suffices for
\(\operatorname{tr}(A^2)\le(1+b^{-2})d_X^2/d\) in
Theorem~\ref{thm:fewtranslates868}.
\end{theorem}

\begin{proof}
Theorem~\ref{thm:poincaregram704} identifies the full Gram array
with the sum of the orbital integrals. On all residues combined,
the parabolic contribution has norm
\(\sqrt{|\mu(E)|D_{\rm id}}\). The positive closed terms and the
remaining part of the open range have cut norms bounded by
\((\tau^{-20}+\tau^{-3a/2})\sqrt{D_{\rm id}}\), where
\(\tau=b^{kf/4}\), by Corollaries~\ref{cor:closedmoment693}
and~\ref{cor:normalrange694}. Their passage to full kernels costs
at most \(\tau^{-20}\sqrt{D_{\rm id}}\) by
Corollary~\ref{cor:squaredamplitudetail700}. Since
\(|\mu(E)|\le b^{11}\) and \(\mathcal M/D_{\rm id}\ge b^{12kf}\),
these terms together have norm \(O(b^{-5kf})\sqrt{\mathcal M}\).

For the full open array on \(z(s)\ge b^f\), use
Proposition~\ref{prop:allresidue706} with \(Z=b^f\).
The squared coefficient of \(D_{\rm op}\) in its majorant is
\[
 (\tau^2+\tau^3)b^{-kf}+4\tau^4(2/M)^{kf}
          \le b^{(-1/4+o(1))kf}.
\]
The other term in its norm bound is
\(2\tau^{-20}\sqrt{D_{\rm id}}\). By Lemma~\ref{lem:massscale869}
the whole norm is at most
\(b^{(-1/8+o(1))kf}\sqrt{\mathcal M}\).
On the complementary residues the error in using the cut orbit
integrals costs at most \(\tau^{-20}\sqrt{D_{\rm id}}\).
The triangle inequality now reduces the assertion to comparing
the norm of the signed sum of these \(V_\gamma\) with
\(\sqrt{\mathcal E}\).

Write \(B_\gamma=C_\infty M^{akf}G_\gamma^{-k/2}\) for the
individual majorant of Proposition~\ref{prop:allresidue706}.
There \(\|V_\gamma\|_{\rm HS}\le B_\gamma+\epsilon_\gamma\),
and the counting norm of the blockwise sums of
\(\epsilon_\gamma\ge0\) is at most
\(\tau^{-20}\sqrt{D_{\rm id}}\).
Project each \(V_\gamma\) radially onto the ball of radius
\(B_\gamma\) in its matrix space, obtaining \(W_\gamma\).
This changes its normed distance by at most \(\epsilon_\gamma\).
The error in the sum within each block is bounded by the sum of
these errors. The same bound on the total counting norm holds
after refining blocks by \(J\), since the sum of squares of
partial sums of nonnegative errors is at most the square of
their full sum.

Expand the squared norms for the \(W_\gamma\) before and after
refining blocks by \(J\). Their difference is a sum of inner
products over ordered pairs with \(J_\gamma\ne J_{\gamma'}\).
Its absolute value is bounded by the sum of
\(B_\gamma B_{\gamma'}\) over those pairs.
In the proof of Proposition~\ref{prop:allresidue706} this part
of the ordered-pair bound, even on all residues, is
\[
 D_{\rm op}\sum_{\ell\ge0}\tau^{\ell+2}
               (2/M)^{\max(1,\ell-1)kf}
 \le 4\tau^4(2/M)^{kf}D_{\rm op}.
\]
Indeed different links give positive codimension in
Theorem~\ref{thm:allpair689}; for \(\ell\ge2\) the successive
majorants have ratio \(\tau(2/M)^{kf}<1/2\).
Thus the squared-norm difference is at most
\(b^{(-19+o(1))kf}\mathcal M\).
For nonnegative \(u,v\) one has
\(|\sqrt u-\sqrt v|\le\sqrt{|u-v|}\).
Applying this to the two squared norms and then restoring the
\(V_\gamma\) on both sides proves the assertion, because
\(1/8>1/10\) and all error bounds are uniform.

Under the stated hypothesis on \(\mathcal E\), squaring gives
a full Gram moment at most \((1+b^{-2})\mathcal M\) for large
nodes. The average of \(\mathcal P\mathcal P^*\) has squared
Hilbert--Schmidt norm equal to that moment divided by
\(|\mathcal G|\), by expanding the two group averages and
using unitarity. Its compression by the central good projector
is \(A\), so its Hilbert--Schmidt norm dominates that of \(A\).
The inequality \(d\le|\mathcal G|d_\lambda\mathfrak m_r\) gives
the required conclusion.
\end{proof}

The estimate of \(\mathcal E\) with the displayed leading constant
is still required. Its summands retain interference between distinct
rational orbits with the same unitary map \(J\).
\clearpage

\subsection{Nonzero leading minors at a marked place}
\label{sec:openenergy870}

Keep \(\mathcal M\), the shorter interval and the unprojected
map \(\mathcal P\) of Section~\ref{sec:equallinks869}. In the
residue space at a place above \(p_0\), the two components give
the natural action on a vector space of dimension \(r=k+2a+1\)
and its contragredient. The pairings identify the frame in
the second component with \(n=a+1\) independent linear forms
on the first. Call a residue \(s\) open at \(w\) if for every
\(1\le l\le a\) the pairing matrix of the first \(l\) frame
forms against the \(s_w\)-images of the first \(l\) frame
vectors has nonzero determinant. Let \(\mathcal O\) be the
set of residues open at some \(w\) above \(p_0\). The first
\(a\) columns are the isotropic arm, and the extra column
is the terminal vector. The nonvanishing refers only to the arm.

\begin{lemma}[A weighted count of singular minors]\leavevmode
\label{lem:singularminor870}
Let a finite group \(H\) act on \(\mathbb F_q^r\) by a
representation \(\pi\). Suppose \(\beta\ge0\) on \(H\) and
\(\langle u_h,u_j\rangle=\beta(h^{-1}j)\) for vectors in a
complex Hilbert space. Write \(\mu=\sum_{h\in H}\beta(h)\).
Take independent \(l\)-tuples of vectors and forms, given by
embeddings \(B:\mathbb F_q^l\longrightarrow\mathbb F_q^r\)
and \(C:\mathbb F_q^l\longrightarrow(\mathbb F_q^r)^*\).
Suppose that each stabilizer of a nonzero \(Bv\) under \(\pi\),
or of a nonzero \(Cy\) under the contragredient, has
\(\beta\)-mass at most \(Y\). Then
\[
  q\sum_{\det(C^t\pi(h)B)=0}\beta(h)
       \le 2\mu+q^{l+r/2}Y .
\]
The columns of \(C\) here are written in dual coordinates.
\end{lemma}
\begin{proof}
Take a nontrivial additive character \(\psi\) of \(\mathbb F_q\).
In the counting space on \(V=\mathbb F_q^r\) use the point
masses \(e_z\) and the Fourier vectors
\(c_x(v)=q^{-r/2}\overline{\psi(x(v))}\).
Each family is orthonormal. For a nonzero \(x\) in the image of
\(C\) and \(z\) in the image of \(B\), form
\[
   U_x=\sum_{h\in H}c_{\pi(h)^{-t}x}\otimes u_h,\qquad
   V_z=\sum_{j\in H}e_{\pi(j)z}\otimes u_j.
\]
Orthogonality gives \(\|V_z\|^2=|H|\sum_{s:\pi(s)z=z}\beta(s)\)
by putting \(s=h^{-1}j\). The same computation for the dual
action gives \(\|U_x\|^2=|H|\sum_{s:\pi(s)^{-t}x=x}\beta(s)\).
With the first argument of the inner product conjugate-linear,
\[
 \langle U_x,V_z\rangle
   =|H|q^{-r/2}\sum_{s\in H}\psi(x(\pi(s)z))\beta(s).
\]
Cauchy--Schwarz bounds the absolute value of the character
sum by \(q^{r/2}Y\).

Put \(Z_s=C^t\pi(s)B\). The cardinality of its left kernel is
\(q^{-l}\sum_{y,v}\psi(y^t Z_s v)\), by summing over \(v\)
first. It is at least \(q\) when \(Z_s\) is singular.
In the weighted sum over \(s\), the pairs with \(y=0\) or
\(v=0\) contribute at most \(2\mu\) after division by \(q^l\).
Each other pair contributes at most \(q^{r/2}Y\) before that
division. There are at most \(q^{2l}\) pairs. This proves the
claim.
\end{proof}

A choice \(\Xi\) consists of a subset of the \(f\) places above
\(p_0\) and, at each chosen place, one nonzero vector in the
span of the full frame in either component. Let \(t=|\Xi|\)
denote the number of chosen places and \(\Gamma_\Xi\) the
subgroup fixing these vectors. Fixing a vector in the second
component means fixing its corresponding form by the
contragredient in the first. Define
\[
 \beta(s)=\|\mathcal P^*T(s)\mathcal P\|_{\rm HS}^2,\quad
 \mu(H)=\sum_{s\in H}\beta(s),\quad B_0=b^{6k}.
\]
For a choice \(\Xi\), write \(\mathcal E_\Xi\) for the sum in
\eqref{eq:linkenergy869} with residue range replaced by
\[
    s\in\Gamma_\Xi\cap\mathcal O,\qquad z(s)< b^{f+6t}.
\]
It uses the signed matrix integrals and exactly the same
partition by full link within each block as in that equation.

\begin{theorem}[Restriction to residues with nonzero leading minors]
\leavevmode\label{thm:openenergy870}
Uniformly for sufficiently large nodes on the shorter interval,
\begin{align}
 (1-2b^{-4})^2\mathcal M
     &\le \mu(\mathcal G),
 \label{eq:loweractual870}\\
  (1-3a/p_0)\,\mu(\mathcal G)
     &\le\left(
       \max_{\Xi:\,|\Xi|<f/4}B_0^{|\Xi|/2}\sqrt{\mathcal E_\Xi}
               +b^{-kf/10}\sqrt{\mathcal M}\right)^2 .
 \label{eq:upperactual870}
\end{align}
Every choice with \(t\ge f/4\) satisfies
\(B_0^t\mu(\Gamma_\Xi)\le b^{-kf/5}\mathcal M\).
In particular, on the right of \eqref{eq:upperactual870} it
suffices to prove
\(B_0^t\mathcal E_\Xi\le(1+\tfrac12b^{-2})\mathcal M\)
for \(t<f/4\) to obtain the full-source moment estimate in
Theorem~\ref{thm:fewtranslates868}.
\end{theorem}

\begin{proof}
For \(\mathcal P\mathcal P^*=S\), the positive operators
\(u_h=T(h)ST(h)^*\) are Hilbert--Schmidt. Cyclicity of trace
gives \(\langle u_h,u_j\rangle_{\rm HS}=\beta(h^{-1}j)\).
This proves the positivity hypothesis of
Lemma~\ref{lem:singularminor870} on every subgroup.
The average of the \(u_h\) has squared norm
\(\mu(\mathcal G)/|\mathcal G|\). Its compression by the
central good projector is \(A\) of Theorem~\ref{thm:fewtranslates868}.
The trace of \(A\) is at least \((1-2b^{-4})d_X\) by the
near-isometry proved there. Since \(A\) acts on a space of
dimension \(d\), its squared norm is at least
\((1-2b^{-4})^2d_X^2/d\). The mass upper bound for \(d\) in
Proposition~\ref{prop:goodmass868} now gives
\eqref{eq:loweractual870}.

We first track the error estimates under restriction to
\(\Gamma_\Xi\). Write \(Z=b^{f+6t}\). In units of \(\mathcal M\)
the term from \(D_{\rm op}\) in the squared majorant of
Proposition~\ref{prop:allresidue706} on \(z(s)\ge Z\), multiplied
by \(B_0^t\), is at most
\[
 b^{o(1)kf}\left(2b^{6kt+3kf/4}Z^{-k}
                        +4b^{6kt+kf}(2/M)^{kf}\right)
       \le b^{(-1/4+o(1))kf}.
\]
We used \(t\le f\), \(M\ge b^{20}\), and
Lemma~\ref{lem:massscale869}. The parabolic, closed, complementary
open and cutoff errors in the proof of Theorem~\ref{thm:equallinks869}
have combined norm \(O(b^{-5kf})\sqrt{\mathcal M}\) before the
weight. Their norm after multiplication by \(B_0^{t/2}\) is
\(O(b^{-2kf})\sqrt{\mathcal M}\). The bounds persist upon
restriction to any subset of residues. In comparing the cut sums
with their refinement by \(J\), the absolute cross terms between
different links after radial projection were bounded by
\(b^{(-19+o(1))kf}\mathcal M\). Restriction and multiplication
by \(B_0^t\) leave at most \(b^{(-13+o(1))kf}\mathcal M\).
The errors from the radial projections in both counting spaces
remain within the cutoff bound. Applying the same square-root
inequality as in Theorem~\ref{thm:equallinks869} therefore gives
\[
 \left|\sqrt{B_0^t\,\mu(\Gamma_\Xi\cap\mathcal O)}
            -\sqrt{B_0^t\,\mathcal E_\Xi}\right|
       \le b^{-kf/10}\sqrt{\mathcal M}.
\]
Only estimates of absolute errors were restricted here. No
positivity assertion about a signed orbital integral is needed.

If \(t\ge f/4\), every \(s\in\Gamma_\Xi\) satisfies
\(z(s)\ge p_0^t\): at each selected component the fixed vector
itself witnesses a nonzero intersection of the two frames.
Apply the same bound to all these residues using \(Z=p_0^t\).
Since \(p_0\ge b^{10}\), the exponent \(6kt+3kf/4-10kt\)
is at most \(-kf/4\). The terms independent of \(Z\) obey
the preceding estimates. After taking square roots and adding
the other orbital terms we obtain
\(B_0^{t/2}\sqrt{\mu(\Gamma_\Xi)}\le
 b^{-kf/10}\sqrt{\mathcal M}\), as claimed.

Choose \(\Xi\) to maximize \(B_0^t\mu(\Gamma_\Xi)\) among
all choices. They form a finite set, including the empty
choice of value \(\mu(\mathcal G)\). The lower bound
\eqref{eq:loweractual870} and the result for \(t\ge f/4\)
force \(t<f/4\) at every maximizer. Take an unselected place
\(w\) above \(p_0\). For every nonzero vector in either frame
span at \(w\), adding its stabilizer to the restrictions gives
mass at most \(\mu(\Gamma_\Xi)/B_0\), by maximality.

Apply Lemma~\ref{lem:singularminor870} on \(\Gamma_\Xi\) to the
first \(l\) vectors and forms at \(w\), for \(1\le l\le a\).
For all such \(l\),
\[
       p_0^{\,l+r/2}
       \le b^{11(a+r/2)}\le b^{6k}=B_0,
\]
since \(p_0\le b^{11}\), \(r=k+2a+1\), and \(a=o(k)\).
Each singular-minor locus therefore carries at most
\(3\mu(\Gamma_\Xi)/p_0\). A residue outside \(\mathcal O\)
has at least one of these minors zero at this particular \(w\).
Summing bounds its mass by \(3a\mu(\Gamma_\Xi)/p_0\).
Multiplying the bound for its complement by \(B_0^t\),
using maximality, and applying the weighted norm comparison
proves \eqref{eq:upperactual870}. Finally \(3a/p_0\le b^{-9}\)
eventually, and \(b^{-kf/10}=o(b^{-2})\). The asserted
sufficient estimate thus gives
\(\mu(\mathcal G)\le(1+b^{-2})\mathcal M\), which implies
the stated estimate for \(A\) by compression.
\end{proof}

For every surviving block there is a place where all these
minors are nonzero. The corresponding rational leading minors
are consequently nonzero in \(E\): reduction of their integral
matrix gives the displayed residue pairing since the two
class representatives belong to the middle group and fix the
frame. This supplies the domain on which successive elimination
of the frame coordinates is possible. The value of the signed
sum on that domain still enters \(\mathcal E_\Xi\).
\clearpage

\subsection{Factorization of a signed sum with a fixed link}
\label{sec:orbitfactors871}

On the rational space write \(Y=L\oplus W\oplus L^\vee\), where
\(L\) has ordered basis \(e_1,\ldots,e_a\), \(L^\vee\) has dual
isotropic basis \(e_i^\vee\) with \(h(e_i,e_j^\vee)=\delta_{ij}\),
and \(W=X_0\oplus Ez\) with \(h(z,z)=-N\). The Hermitian form
is linear in its second argument. The group \(U\) is the unipotent
radical of the stabilizer of the complete flag in \(L\). Extend
each \(\sigma\in\mathrm U(W)(F)\) by the identity on the arms.
For \(\mathbf p=(p_1,\ldots,p_a)\in(E^\times)^a\), define
\(w_{\mathbf p}e_i=p_i e_i^\vee\),
\(w_{\mathbf p}e_i^\vee=\bar p_i^{-1}e_i\) and
\(w_{\mathbf p}|_W=1\). Write \(D_j(g)\) for the leading minors
of the pairings of the frame \((e_1,\ldots,e_a,z)\) with its
image under \(g\). Put \(D_0=1\).

\begin{lemma}[The unipotent coordinates and the middle isometry]
\leavevmode\label{lem:pivots871}
If \(D_1(g),\ldots,D_a(g)\ne0\), there is a unique expression
\[
              g=u\,w_{\mathbf p}\sigma v^{-1},
       \qquad u,v\in U(F),\qquad \sigma\in\mathrm U(W)(F).
\]
One has \(p_i=D_i(g)/D_{i-1}(g)\), and the full link of \(g\)
determines \(\sigma\) and is determined by \(\sigma\).
Write \(\sigma z=\omega+\alpha z\), with \(\omega\in X_0\),
and \(t=h(z,\sigma z)\). Then
\begin{align}
 D_n(g)=D_a(g)t,\qquad
 \Delta_g&=D_a(g)\overline{D_a(g)}(t\bar t-N^2),\qquad
 N h(\omega,\omega)=t\bar t-N^2 .
 \label{eq:middlesphere871}
\end{align}
For \(\Delta_g\ne0\), one has \(|\Phi_v(t)|>N\) at every
real place. Membership in \(T_>\) then depends only on \(t\).
\end{lemma}
\begin{proof}
The ordinary upper unitriangular group on \(L\) lifts into \(U\)
by its inverse adjoint on the dual arm. Acting on the left
and right thus permits lower unitriangular row operations and
upper unitriangular column operations on the \(a\)-by-\(a\)
cross Gram matrix. Successive elimination makes it diagonal:
the first pivot is its upper left entry, and after clearing its
row and column the leading minors of the smaller block are
the next leading minors divided by that pivot. Induction gives
the claimed ratios. All leading minors are invariant under these
operations. More generally they are invariant under left and
right \(U\), including the minor with the last column \(z\),
because \(U\) acts on the entire ordered frame by upper
unitriangular matrices.

After diagonalization set \(q_i=g e_i/p_i\).
Then \(h(e_j,q_i)=\delta_{ji}\) and \(h(q_i,q_j)=0\).
There is an element of \(U\) taking \(e_i^\vee\) to \(q_i\).
Indeed fix the \(e_i\) and send \(y\in W\) to
\(y-\sum_i e_i h(q_i,y)\). These images have their original
pairings with one another and with \(L\), and are orthogonal to
the \(q_i\). The asserted map therefore preserves the whole
Gram matrix and induces the identity on each graded block of
the parabolic. Inverting it on the left now gives
\(g e_i=p_i e_i^\vee\). Consequently \(w_{\mathbf p}^{-1}g\)
fixes \(L\) pointwise and induces an isometry \(\sigma\) on
\(L^\perp/L=W\). After removing \(\sigma\) it lies in \(U\),
which proves existence.

To identify the middle isometry intrinsically, consider the
relation between the two images of \(y\) and \(g^{-1}y\) in
\(L^\perp/L\), for \(y\in(L+gL)^\perp\). Multiplication by \(U\)
on either side preserves this relation: change \(y\) by the left
factor and use that both unipotents act trivially on the middle
quotient. For \(g=w_{\mathbf p}\sigma\) the relation consists
of the pairs \((y,\sigma^{-1}y)\), \(y\in W\). Interchanging the
last scalar coordinates between the two copies of \(W\) is
exactly the passage from this relation to the graph of the full
link in Section~\ref{sec:differentlinks688}. It is an invertible
coordinate permutation, so equality of links is equivalent to
equality of \(\sigma\). For fixed \(\mathbf p,\sigma\), uniqueness
of the two unipotents follows from
\(U\cap(w_{\mathbf p}\sigma)U(w_{\mathbf p}\sigma)^{-1}=1\).
In fact \(\sigma\) normalizes \(U\), and interchanging both arms
takes the strictly upper block group to the strictly lower one
when the dual arm is ordered in reverse.

Calculation at \(w_{\mathbf p}\sigma\) gives \(D_n=D_at\)
and \(t=-N\alpha\). Isometry gives
\(h(\omega,\omega)-N\alpha\bar\alpha=-N\), which proves
\eqref{eq:middlesphere871}. Nonzero \(\omega\) has strictly
positive norm at every real place. Finally \(P_n=\prod_v
|\Phi_v(D_n)|\) in Corollary~\ref{cor:normalrange694}, so that
\[
 \frac{\sqrt{\mathrm N_{F/\QQ}\Delta_g}}{P_n}
   =\prod_v\left(1-\frac{N^2}{|\Phi_v(t)|^2}\right)^{1/2}.
\]
The right side, and hence its comparison with
\((M/b)^{af}/N^{f/2}\), is independent of the \(p_i\).
\end{proof}

Let \(J_w\) denote the local level at a finite place \(w\) of
\(F\), and take \(k_{s,w}\) to be the component of a lift of
the marked residue \(s\). For a class \(c\) put
\(U_{c,w}=U(F_w)\cap h_{c,w}J_wh_{c,w}^{-1}\).
Use Haar measure \(du_{c,w}\) on \(U(F_w)\) giving \(U_{c,w}\)
volume one. Write \(\chi_\infty\prod_{w<\infty}\chi_w\) for
the adelic Bessel character on \(U\), trivial on \(U(F)\).
With \(g=w_{\mathbf p}\sigma\), define
\begin{align*}
 K_w(c,s,c';g)
    &=\int_{U(F_w)^2}
         {\bf1}_{\,h_{c,w}k_{s,w}J_w h_{c',w}^{-1}}(u g v^{-1})
           \overline{\chi_w(u)}\,\chi_w(v)\,du_{c,w}\,dv_{c',w},\\
 I_v(g)
    &=\int_{G(F_v)} F_v^\flat(y)F_v^\flat(g^{-1}y)^*\,dy
              \qquad (v\mid\infty).
\end{align*}
The cut functions \(F_v^\flat\), their product and the real Haar
measures are those of Proposition~\ref{prop:unfoldopen702}.
The finite expression uses counting measure on left cosets of
the two local level groups. It can have complex values.

\begin{theorem}[Local factors for an actual open block]
\leavevmode\label{thm:orbitfactors871}
Fix \((c,s,c')\) occurring in \(\mathcal E_\Xi\) of
Theorem~\ref{thm:openenergy870}. A link \(J_*\) occurring in
that block has a unique middle isometry \(\sigma\). Its signed
sum in \eqref{eq:linkenergy869} equals
\begin{equation}
 \sum_{\gamma:\,J_\gamma=J_*} V_\gamma
   =\frac{1}{v_U}\sum_{\mathbf p\in\mathcal Q^a}
       \left(\prod_{w<\infty}K_w(c,s,c';w_{\mathbf p}\sigma)\right)
       \bigotimes_{v\mid\infty} I_v(w_{\mathbf p}\sigma).
 \label{eq:orbitfactors871}
\end{equation}
Here \(\mathcal Q\subset E^\times\) is a finite set common to the
blocks at the node. Each finite-place integral is a finite sum,
and almost all factors equal one. Formula~\eqref{eq:orbitfactors871}
applies when \(t\) has the range in Lemma~\ref{lem:pivots871}
defining \(T_>\); the sum restricted to \(T_>\) is zero for any
other link. Both class representatives and the original residue
enter the local factors.
\end{theorem}
\begin{proof}
Every term over a residue in \(\mathcal O\) has the required
nonzero leading minors by Theorem~\ref{thm:openenergy870}.
We first record the finiteness facts for the factorization.
The construction in Lemma~\ref{lem:pivots871} gives inverse
coordinates on the open set of matrices with all these minors
invertible. Gaussian elimination divides only by nonzero pivots.
The remaining changes of basis use the fixed nonsingular middle
form, isometric inverse, and the pivots and their conjugates.
In particular the inverse coordinates are continuous over each
completion on that open set, also in a split algebra by applying
the same operations in its two paired components. Where the
fixed form and bases are unimodular and the pivots are local
units, integral input with integral inverse gives integral
unipotent factors and integral middle isometry with integral
inverse. This follows step by step: the row and column additions,
the vectors \(q_i\), the correcting map on \(W\), and the
normalized Weyl map are all integral there, as are their inverses
by the unitriangular or isometry formulas.

For fixed \(g=w_{\mathbf p}\sigma\) its intersection
\(U(F_w)gU(F_w)\) with any compact set has relatively compact
unipotent coordinates. Indeed the leading minors of orders
\(1,\ldots,a\) are fixed nonzero constants on this orbit. The
closed locus fixing those minors in the compact set is contained
in the domain of the continuous inverse just described. Its
image is compact. Applied to the compact finite-place coset in
the definition of \(K_w\), this bounds both local unipotents.
The condition in that integral is unchanged by left
\(U_{c,w}\) on \(u\) and left \(U_{c',w}\) on \(v\).
The local characters on these subgroups are trivial by the
level and conductor computation in
Proposition~\ref{prop:stabilizercancel680}. Thus \(K_w\) is
a sum over a finite set of pairs of cosets, each of measure
one. At all but finitely many places the level is standard,
the class representatives are integral, \(g\) and its inverse
are integral and the pivots are units. The preceding integral
coordinate calculation gives exactly the two integral unipotent
subgroups as the preimage of the integral level. This gives
\(K_w=1\) there.

At a real place the cut functions are supported on \(U(F_v)C_v\)
for some compact \(C_v\). In fact their length bound is a bound
on the infimum of \(1+\|h y\|\) over the Bessel group, whose
quotient by \(U(F_v)\) is compact. Bounded matrices in the
unitary group lie in a compact set because their inverses are
obtained by adjunction with the fixed form. If both real factors
at \(y\) contribute, write \(y=u c_1\), \(g^{-1}y=v c_2\),
with \(c_i\in C_v\). Then \(u^{-1}g v=c_1c_2^{-1}\).
The compact-set argument of the preceding paragraph bounds
\(u,v\), so \(y\) lies in a compact set as well. The local
cut functions are bounded there. The real integral converges
absolutely and the product integral factors by Fubini.

At the node only finitely many pivot tuples can contribute to
\eqref{eq:orbitfactors871}. For a contributing term all \(D_j\)
are integral, since after the two class changes the isometry
preserves the integral lattice, the frame is integral, and the
class changes fix it. At each embedding the \(D_j\) are bounded
by their maxima on \(C_vC_v^{-1}\), by the same invariance
under \(U\). These are finitely many fixed bounds at the node.
The image of \(\mathcal O_E\) in its Minkowski space is a
lattice. Thus we may take a finite set \(\mathcal D\) comprising
\(1\) and all nonzero integers of \(E\) satisfying the bounds
for at least one of these minor orders at every embedding
(enlarging to the maximum bound for all orders is harmless).
Take \(\mathcal Q=\{d/d': d,d'\in\mathcal D\}\). Any pivot
of an orbit contributing to the finite and real integrals
then lies in \(\mathcal Q\).

It remains to unfold with the phases and measures. Write
\(U_c=\prod_{w<\infty}U_{c,w}\) and
\(\Gamma_c=U(F)\cap U_c\). The diagonal rational unipotent
group is dense in \(U(\mathbb A_{F,f})\): exponential and
logarithm are inverse polynomial maps with rational coefficients,
integral outside finitely many primes, and additive rational
points are dense in the finite adeles by approximation.
Consequently
\(\Gamma_c\backslash U(F)\longrightarrow
 U_c\backslash U(\mathbb A_{F,f})\) is a bijection.
Every left coset of the compact open subgroup is open and
contains a rational point. Two rational representatives differ
on the left by an element of \(U_c\) precisely when they differ
by an element of \(\Gamma_c\). The same assertions apply to \(c'\).

By uniqueness in Lemma~\ref{lem:pivots871}, the rational terms of
this double coset are parametrized without multiplicity by the
pairs \((u,v)\in U(F)^2\). Fix their two left level cosets.
For \(\lambda\in\Gamma_c,\lambda'\in\Gamma_{c'}\), the term
is \(\lambda u g v^{-1}\lambda'^{-1}\). In the integral with
the two class weights of Section~\ref{sec:besselresidue701}
substitute \(y=\lambda u y_0\). The second argument becomes
\(\lambda'v g^{-1}y_0\). Unipotent equivariance of the cut
function and triviality on the two rational level subgroups
leave the scalar \(\chi_\infty(u)\overline{\chi_\infty(v)}\).
Summing the two class weights independently gives one in each
variable, by their partition property. This unfolding is valid
in absolute value as well: the weights are nonnegative, and
after summing them the integrand is bounded by the product
of the two cut norms on the compact overlap just proved. Thus
the contribution of the fixed level cosets is the scalar times
\(v_U^{-1}\bigotimes_v I_v(g)\).

Finally \(\chi_\infty(u)=\prod_w\overline{\chi_w(u)}\) for
rational \(u\), and the conjugate relation holds for \(v\).
The two coset bijections identify the admissible pairs exactly
with the local condition \(u g v^{-1}\in
h_{c,w}k_{s,w}J_w h_{c',w}^{-1}\) at every \(w\).
Outside a finite set there is just the integral pair of cosets,
so their weighted count is the product of the \(K_w\).
Summing over the pivots gives the claimed formula. This
regrouping agrees with the cut integrals \(V_\gamma\) because
Proposition~\ref{prop:unfoldopen702} unfolds the same sum of
terms against the same two class weights. The full link and
the range \(T_>\) are invariant under every unipotent move,
as shown in Lemma~\ref{lem:pivots871}.
\end{proof}

In the energy \(\mathcal E_\Xi\) the squared norm is taken
after summation over \(\mathbf p\). Thus the finite character
sums with different pivots occur together with the corresponding
real matrix integrals. An estimate of this quadratic expression
with the leading constant in Theorem~\ref{thm:openenergy870}
remains necessary for this argument.
\clearpage

\subsection{The class sum on the terminal norm sphere}
\label{sec:spheremean872}

Let \(H=\mathrm U(X_0)\) and \(H_W=\mathrm U(W)\) in the
notation of Section~\ref{sec:orbitfactors871}. Thus \(H\) is the
stabilizer of \(z\) in \(H_W\). Put \(Q(x)=h(x,x)\) on \(X_0\).
Fix a choice \(\Xi\) of marked vectors and write
\(\Omega_\Xi=\{s\in\Gamma_\Xi\cap\mathcal O:z(s)<b^{f+6|\Xi|}\}\).
We compute the sum over the two definite classes and the middle
isometry in \(\mathcal E_\Xi\).
The level in \(H(\mathbb A_{F,f})\) is \(K_H=H(\mathbb A_{F,f})\cap J\).
In integrals over that group give \(K_H\) volume one.

Let \(\mathcal Q\) be the finite set of pivot values constructed in
Theorem~\ref{thm:orbitfactors871}. It also contains \(D_n/D_a=t\)
for any contributing term. Keep the \(t\in\mathcal Q\) for which
\[
       \beta_t=(t\bar t-N^2)/N
\]
is totally positive, is represented by \(Q\) over \(F\), and
satisfies the \(T_>\) condition of Lemma~\ref{lem:pivots871}
through its value of \(t\). Denote this finite set by \(\mathcal T\).
For each \(t\) choose one rational \(\sigma_t\in H_W(F)\) with
\(h(z,\sigma_t z)=t\). To obtain one, take
\(Q(\omega)=\beta_t\) and extend the isometry of nondegenerate
lines \(z\longmapsto\omega-(t/N)z\) by Witt's extension theorem.
For each \(\mathbf p\in\mathcal Q^a\) use the matrix
\[
        A_{\mathbf p}(t)=
              \bigotimes_{v\mid\infty}I_v(w_{\mathbf p}\sigma_t)
\]
in the Hilbert--Schmidt space of the coefficient representation.

For \(x\in H_W(\mathbb A_{F,f})\) write \(K(s;w_{\mathbf p}x)\)
for the product of the finite integrals in
Section~\ref{sec:orbitfactors871} with both class representatives
equal to one, using the standard \(U\)-level for both measures.
Define the nonnegative function
\[
       B_{\Xi,t}(x)=\sum_{s\in\Omega_\Xi}
          \left\|\sum_{\mathbf p\in\mathcal Q^a}
             K(s;w_{\mathbf p}x)A_{\mathbf p}(t)\right\|_{\rm HS}^2 .
\]
On the finite adelic sphere \(Q(\omega)=\beta_t\), set
\begin{equation}
    \Psi_{\Xi,t}(\omega)=\int_{H(\mathbb A_{F,f})}
                          B_{\Xi,t}(x h)\,dh,\qquad
               xz=\omega-(t/N)z .
 \label{eq:spheretest872}
\end{equation}
The integral uses one fiber of the map \(x\mapsto xz\).
Two choices of \(x\) differ on the right by \(H(\mathbb A_{F,f})\),
so its value is independent of that choice.

\begin{proposition}[Unfolding the right definite class]\leavevmode
\label{prop:sphereunfold872}
The test \(\Psi_{\Xi,t}\) is nonnegative, has compact support on
the finite adelic sphere, and is invariant under \(K_H\). The
identity for the actual energy is
\begin{equation}
 v_U^2\mathcal E_\Xi=
       \sum_{t\in\mathcal T}\ \sum_{c\in\mathcal C}\
       \sum_{\substack{\omega\in X_0(F)\\Q(\omega)=\beta_t}}
                       \Psi_{\Xi,t}(h_c^{-1}\omega).
 \label{eq:classsphere872}
\end{equation}
\end{proposition}
\begin{proof}
If \(\sigma,\sigma'\in H_W(F_v)\), \(v\mid\infty\), have the same
last entry \(t\), the \(X_0\)-parts of \(\sigma z\) and \(\sigma'z\)
have the same positive length by \eqref{eq:middlesphere871}.
An element \(h_1\in H(F_v)\) takes the first part to the second.
Then \((\sigma')^{-1}h_1\sigma\) fixes \(z\). Thus
\(\sigma'=h_1\sigma h_2^{-1}\) for some \(h_2\in H(F_v)\).
The arm interchange \(w_{\mathbf p}\) commutes with \(H\).
Changing the real integration variable by \(h_1\) shows that
all the matrices \(I_v(w_{\mathbf p}\sigma)\) are carried to
\(I_v(w_{\mathbf p}\sigma')\) by the same two unitary factors
on the coefficient spaces, by the left equivariance of \(F_v^\flat\).
Consequently any linear combination of the real tensor matrices
has the same squared norm when \(\sigma\) is replaced by
\(\sigma_t\), with its finite coefficients held fixed.

In the finite integrals conjugate the two unipotent variables
by \(h_c^{-1}\) and \(h_{c'}^{-1}\), respectively. These middle
elements preserve the Bessel character. The chosen Haar measures
go to the standard ones since the two normalizing open subgroups
go to \(U\cap J\). The resulting factor is
\(K(s;w_{\mathbf p}h_c^{-1}\sigma h_{c'})\). This expression as
a function of \(h_c^{-1}\sigma h_{c'}\) is invariant on both
sides by \(K_H\). Commuting a level element across the appropriate
unipotent variable preserves its character and its normalizing
subgroup, and \(j k_s J=k_sJ=k_sJj\) for \(j\in J\).
Residues with disagreeing characters give zero even for a fixed
pivot tuple. Indeed, if \(u_0,v_0\) are integral finite unipotents
with \(u_0 k_sJv_0^{-1}=k_sJ\), replacing the integration
variables by \(u_0u,v_0v\) preserves the coset condition and
multiplies the character factor by
\(\overline{\chi_f(u_0)}\chi_f(v_0)\).
Unequal character values therefore force the integral to vanish.

For fixed pivots, nonzero finite integral forces the middle
coordinate into a compact set at every place by the inverse
coordinate calculation of Theorem~\ref{thm:orbitfactors871}.
At almost all places this set is the standard integral middle
group. On a compact neighborhood of a middle element the
unipotent variables in the integral can be confined to a common
compact set, by the same inverse coordinate argument applied
to the integral coset and fixed pivots. The characteristic
function of that coset is locally constant. Its integral is
therefore a locally constant function there, as one also sees
by covering the compact integration set by finitely many cosets.
At all but finitely many places it is just the indicator of the
standard integral middle group. Hence \(B_{\Xi,t}\) has compact
support and is locally constant. Its fiber integral in
\eqref{eq:spheretest872} is measurable, supported in the image
of that compact set, and invariant by \(K_H\) on \(\omega\).
The fibers over the stated adelic sphere exist. Locally this is
Witt extension. At almost every finite place the prescribed
vector and \(z\) are integral of the same unit norm and an
integral isometry suffices. Over an unramified quadratic
valuation ring each such line splits off a unimodular
complement. To see that unimodular Hermitian lattices of the same
rank there are isometric, choose a vector of unit norm by
reduction, rescale by surjectivity of the norm on units, split
off its line and repeat on the complement. At a split place a primitive vector and its
paired covector with unit pairing likewise split off a direct
summand and its dual. These arguments give an integral
extension at every place required for the restricted product.

Now insert \eqref{eq:orbitfactors871} and group by
\(\sigma z=\omega-(t/N)z\).
For fixed rational \(\omega,t\) all choices of \(\sigma\) are
\(\sigma_0 h\), \(h\in H(F)\), for any one choice \(\sigma_0\);
every such choice gives one link. The pairs \((h,c')\) give
each right coset \(h h_{c'}K_H\) in \(H(\mathbb A_{F,f})\)
exactly once. Existence is the definition of the definite
classes, and uniqueness uses
\(H(F)\cap h_{c'}K_H h_{c'}^{-1}=1\) as in
Proposition~\ref{prop:rowmean690}. Each coset has volume one.
Their sum is exactly the integral in \eqref{eq:spheretest872}
with \(x=h_c^{-1}\sigma_0\).
Terms with no admissible finite representative give zero.
The resulting sums and integrals are of nonnegative squared
norms, so the regrouping is legitimate. This proves
\eqref{eq:classsphere872}.
\end{proof}

Normalize additive measure at a finite place \(w\) by giving
\(\mathcal O_{F_w}\) and \(\mathcal O_{E,w}^k\) volume one. On
the smooth fiber \(Q(\omega)=\beta\ne0\), let \(\nu_{\beta,w}\)
be the measure \(d\omega/dQ\) in these additive measures.
Equivalently, one can integrate a locally constant compactly
supported test over \(Q\in\beta+\mathfrak p_w^j\), divide by
the measure of \(\mathfrak p_w^j\) and pass to the limit.
In local coordinates with \(Q\) as one coordinate this formula
is the ordinary change of variables.

Choose a finite set \(S\) of places of \(F\) containing both
marks, the divisors of \(N\), and the places where any of the
pivots or entries in \(\mathcal Q\) is a nonunit. Include the
finitely many exceptions to the standard integral bases and
levels of Theorem~\ref{thm:orbitfactors871}. On the sphere,
write \(\Psi_{\Xi,t,S}\) for \eqref{eq:spheretest872} using
only the places of \(S\) in its finite integrals, product
measures and fiber. Write
\(\delta_w(\beta)=\nu_{\beta,w}(\mathcal O_{E,w}^k)\) and
\(c_E=2^{-f}\sqrt{|\operatorname{disc}E|}\).

\begin{theorem}[The mean on a prescribed terminal sphere]\leavevmode
\label{thm:spheremean872}
For every \(\Xi\) one has
\begin{align}
 v_U^2\mathcal E_\Xi
 &= |\mathcal C|\,
       \frac{\sqrt{|\operatorname{disc}F|}}{c_E^k}
       \left(\frac{\pi^k}{(k-1)!\,N^{k-1}}\right)^f
       \sum_{t\in\mathcal T}
          \prod_{v\mid\infty}\bigl(|\Phi_v(t)|^2-N^2\bigr)^{k-1}
          \ \mathcal D_{\Xi,t}, \label{eq:spheremean872}\\
 \mathcal D_{\Xi,t}
 &=\left(\prod_{w\notin S}\delta_w(\beta_t)\right)
       \int_{Q(\omega_S)=\beta_t}\Psi_{\Xi,t,S}(\omega_S)
                                \prod_{w\in S}d\nu_{\beta_t,w}.
 \notag
\end{align}
The product over \(w\notin S\) differs from one by
\(O(f2^{1-k})\), uniformly in \(t,S,\Xi\). In particular its
error is \(o(b^{-6})\) on the same shorter interval.
\end{theorem}
\begin{proof}
For \(w\notin S\) the local factor in each term of \(B_{\Xi,t}\)
is the indicator of the integral middle group, independent of
\(\mathbf p\). It is real and idempotent. Its fiber integral
equals one for integral \(\omega_w\) and zero otherwise.
Indeed \(\omega_w-(t/N)z\) is integral exactly then, and has
the unit norm \(-N\). The integral extension of the isometry
of lines proved above gives an integral point of the fiber.
The rest of the integral points form one coset of the integral
\(H\)-level, of volume one. Hence outside \(S\) the test on
the sphere is exactly the product of the standard lattice
indicators. This also shows \(\beta_t\) is integral there,
by its formula and the choice of \(S\).

We give the mean formula with its measure constant. Fix
\(t\in\mathcal T\), put \(\beta=\beta_t\), choose
\(x_0\in X_0(F)\) of norm \(\beta\), and put
\(H_0=\mathrm U(x_0^\perp)\). All rational vectors of this
norm are in the orbit of \(x_0\) by Witt extension. Integrate
\(\sum_{Q(x)=\beta}\Psi(h_f^{-1}x)\) over
\(H(F)\backslash H(\mathbb A_F)\) with probability Haar
measure, with \(\Psi=\Psi_{\Xi,t}\).
Its integral is the class sum divided by \(|\mathcal C|\).
The integrand is invariant under the full compact group at
infinity and \(K_H\) at finite places, and the finite-level
stabilizers are trivial, so that each class has equal measure.
Unfolding the rational orbit gives the integral of
\(\Psi(h_f^{-1}x_0)\) on \(H_0(\mathbb A_F)\backslash
H(\mathbb A_F)\) in quotient Tamagawa measure, multiplied by
\(\tau(H_0)/\tau(H)\). One sees the factor by first taking the
quotient by \(H_0(F)\) and then integrating the stabilizer
fiber of volume \(\tau(H_0)\). Both unitary groups have Tamagawa
number two, as both Hermitian dimensions \(k-1,k\) are positive
and both forms are definite at every real place.
See \href{https://web.math.princeton.edu/~jonhanke/Web-02/Mass-Formula-of-Shimura/shimura-mass.pdf}{Gan, Hanke and Yu, Section~4, equation~(4.3)}.

To identify the quotient measure, over \(F\) take the top
differential on the smooth sphere obtained by dividing a
constant volume form on \(X_0\) by \(dQ\). It is \(H\)-invariant
because the determinant of a unitary isometry has \(E/F\)-norm one.
The quotient of invariant rational top forms on \(H\) and \(H_0\)
is proportional to this form by a scalar in \(F^\times\),
since the tangent quotient at \(x_0\) is the tangent space of
the sphere and the action is transitive. The absolute values
of the scalar have product one. In the Tamagawa normalization
the character-module convergence factors of \(H\) and \(H_0\),
including the global leading \(L\)-value, cancel. Indeed over
an algebraic closure these groups are general linear groups,
whose characters are generated by determinant, the restriction
of determinant is determinant, and conjugation acts on that
generator by inversion in both cases. The discriminant
normalization for an \(F\)-group of dimension \(d'\) is
\(|\operatorname{disc}F|^{-d'/2}\) when rational-coordinate
finite measures give the base integer rings volume one.
The dimension difference here is \(2k-1\). Hence the quotient
is exactly the fiber measure obtained by disintegrating the
globally normalized additive measures on \(E^k\) and \(F\)
by \(Q\). This follows also place by place from the quotient
of top forms, since their underlying additive dimensions
differ by \(2k-1\). Local Witt extension and the integral
extension at almost all places identify the full adelic sphere
with the quotient being integrated.

For our finite conventions and ordinary Lebesgue measures on
\(\CC^f\) and \(\RR^f\), the covolumes of the rational additive
lattices in their adeles are respectively \(c_E\) and
\(\sqrt{|\operatorname{disc}F|}\). This follows by taking
the product of the finite integer rings and a fundamental
parallelotope of the Minkowski lattice. Thus disintegration
of the additive measures of covolume one multiplies the
product of local fiber measures by
\(\sqrt{|\operatorname{disc}F|}/c_E^k\).
At an infinite place the total fiber measure at
\(\beta_v>0\) is
\(\pi^k\beta_v^{k-1}/(k-1)!\), obtained by differentiating the real
\(2k\)-dimensional volume \(\pi^k u^k/k!\) of
\(\{\|x\|^2\le u\}\) at \(u=\beta_v\).
Use \(\beta_v=(|\Phi_v(t)|^2-N^2)/N\) and the factorization
outside \(S\). Inserting the resulting mean for every \(t\)
in \eqref{eq:classsphere872} yields \eqref{eq:spheremean872}.

It remains to estimate the standard local densities. At
\(w\notin S\) put \(q=|\mathcal O_{F_w}/\mathfrak p_w|\) and
\(\eta=1\) in the split case, \(\eta=-1\) in the unramified
field case. For a nontrivial additive character \(\psi\) of
\(\mathbb F_q\) and any \(d\ne0\), the sum of
\(\psi(d\,x\bar x)\) over the quadratic residue algebra is
\(\eta q\). In the split case sum first over the second
coordinate in \(x\bar x=u v\). In the field case every nonzero
residue norm has \(q+1\) preimages and zero has one.
Character orthogonality therefore gives
\[
 \#\{x: Q(x)=c\bmod\mathfrak p_w\}
   =q^{2k-1}+\eta^k q^{k-1}
       \begin{cases}-1,&c\ne0,\\q-1,&c=0.\end{cases}
\]
At each nonzero residue vector the differential of \(Q\) is
a nonzero linear form over \(\mathbb F_q\) by the perfect
residue trace pairing. Such a solution has precisely
\(q^{2k-1}\) lifts at each successive digit for any prescribed
target lift, including in residue characteristic two.
For unit \(\beta\) this proves \(\delta_w=1-\eta^k q^{-k}\).
For positive valuation the primitive solutions contribute
\(1+\eta^k(q-1)q^{-k}-q^{1-2k}\).
When the valuation is at least two, writing a nonprimitive
vector as \(\varpi y\) adds
\(q^{2-2k}\delta_w(\beta/\varpi^2)\) by additive scaling.
There is no such contribution at valuation one. This recursion,
ending at valuation zero or one, gives
\(|\delta_w(\beta)-1|\le 4q^{1-k}\) for \(k\ge3\) by summing
the geometric series with ratio \(q^{2-2k}\).
There are at most \(f\) places over any rational prime, hence
\(\sum_{w\notin S}q_w^{1-k}\ll f2^{1-k}\) as \(k\) tends to
infinity. The product converges absolutely and has the
asserted error, since \(f\ll_{q_0,d_K}x^{12}\),
\(k\asymp x^5\), and \(\log b\le x^{1/100}\).
\end{proof}

The integral in \(\mathcal D_{\Xi,t}\) still contains the
correlations between the finite character sums for different
pivots. The remaining upper estimate concerns this finite-place
integral together with the matrices \(A_{\mathbf p}(t)\) at the
infinite places.
\clearpage

\subsection{Simultaneous pivot congruences at the two original marks}
\label{sec:bothmarks873}

Let \(\mathcal W\) consist of the \(2f\) places of \(F\) above the
two split primes \(p_0,p_1\). At \(w\in\mathcal W\) the residue
algebra of \(E\) has two components, each \(\mathbb F_{q_w}\),
where \(q_w=p_i\) if \(w\mid p_i\). In the first component the
residue group acts on its natural space of dimension \(r\).
On the other it acts by the contragredient under the perfect
pairing. Say that \(s\) has \emph{unit pivots at \(w\)} if the
leading frame Gram minors of orders \(1,\ldots,a\) are nonzero
in both components. This condition can be tested on a lift to
the integral lattice group and depends only on its residue.

Permit a choice \(\Xi\) as in Section~\ref{sec:openenergy870}
to use places from either of the two primes, with one nonzero
vector in either frame component at each chosen place. Let
\(t=|\Xi|\). Define
\[
 \Omega_\Xi^\sharp
 =\{s\in\Gamma_\Xi:\ z(s)< b^{f+6t},\
              s\ \hbox{has unit pivots at every } w\notin\Xi
                   \hbox{ in }\mathcal W\}.
\]
For \(t<f/4\), let \(\mathcal E_\Xi^\sharp\) be the signed sum of
squares in \eqref{eq:linkenergy869} with this residue mask.
There is an unselected place over \(p_0\), so every residue in
this mask is in \(\mathcal O\). Both \eqref{eq:orbitfactors871}
and the sphere formula \eqref{eq:spheremean872} hold with this
mask. Indeed their proofs fix a set of residues in the open
cell and sum the squared norms there. The level \(K_H\), the
two class sums and the terminal sphere measure are the same.

\begin{theorem}[Unit pivots outside the selected places]\leavevmode
\label{thm:bothmarks873}
Put \(B_0=b^{6k}\) and \(q_*=\min(p_0,p_1)\). At all sufficiently
large nodes of the shorter interval the actual second moment
\(\mu(\mathcal G)=\sum_s\|\mathcal P^*T(s)\mathcal P\|_{\rm HS}^2\)
satisfies
\[
 (1-12af/q_*)\,\mu(\mathcal G)
 \ \le\
 \left(\max_{\Xi:\,|\Xi|<f/4} B_0^{|\Xi|/2}
           \sqrt{\mathcal E_\Xi^\sharp}
       +b^{-kf/10}\sqrt{\mathcal M}\right)^2 .
\]
In this bound \(12af/q_*\le b^{-9}\) for all sufficiently
large nodes. Every maximizing choice for
\(B_0^{|\Xi|}\mu(\Gamma_\Xi)\) among choices at both marks has
fewer than \(f/4\) selected places.
\end{theorem}
\begin{proof}
By \eqref{eq:masslog869}, \(M\ge b^{20}\), \(d_B\ge1\),
and \((k-1)!\ge1\), the identity scale satisfies the sharper
lower bound
\[
     \mathcal M/D_{\rm id}\ge b^{(25/2-o(1))kf}.
\]
This margin allows the stabilizer weights at up to \(2f\)
places. In the argument of Theorem~\ref{thm:equallinks869}, the
combined parabolic, closed, complementary open and cutoff norms
are bounded by
\[
  \big(\sqrt{|\mu(E)|}+O(\tau^{-20}+\tau^{-3a/2})\big)
                         \sqrt{D_{\rm id}} .
\]
After multiplying by \(B_0^{t/2}\) and dividing by
\(\sqrt{\mathcal M}\), the result is at most
\(b^{(-1/4+o(1))kf}\), since \(|\mu(E)|\le b^{11}\) and
\(3kt\le6kf\). The absolute contribution of cross terms
with different links after radial projection in that proof
was bounded by \(b^{(-19+o(1))kf}\mathcal M\) in the squared
norm. Its weighted bound for \(t\le2f\) is at most
\(b^{(-7+o(1))kf}\mathcal M\). Both statements hold on every
residue subset.

Every \(s\in\Gamma_\Xi\) has \(z(s)\ge q_*^t\) since a nonzero
fixed frame vector witnesses an intersection at that component.
Use Proposition~\ref{prop:allresidue706} with \(Z=q_*^t\).
In the squared majorant the term containing \(Z^{-k}\), in
units of \(\mathcal M\) and with weight \(B_0^t\), is bounded by
\(b^{o(1)kf}2b^{3kf/4+6kt}q_*^{-kt}\).
For \(t\ge f/4\) its exponent to base \(b\) is at most
\((-1/4+o(1))kf\), using \(q_*\ge b^{10}\).
The other term, \(4B_0^t\tau^4(2/M)^{kf}D_{\rm op}/\mathcal M\),
is at most \(b^{(-7+o(1))kf}\).
Taking square roots and adding the other orbital contributions
gives \(B_0^t\mu(\Gamma_\Xi)\le b^{-kf/5}\mathcal M\) for all
such choices. The empty choice has value at least
\((1-2b^{-4})^2\mathcal M\) by \eqref{eq:loweractual870}.
Consequently every maximizer has \(t<f/4\).

Fix a maximizer. For each unused place and each of its two
components, the stabilizer of every nonzero frame vector there
has mass inside \(\Gamma_\Xi\) at most \(\mu(\Gamma_\Xi)/B_0\).
Adding that one restriction to \(\Xi\) proves this directly.
To test nonvanishing of a minor at the first component apply
Lemma~\ref{lem:singularminor870} to the natural representation
and the primal and paired dual frame there. At the other
component apply it to the contragredient representation,
interchanging the two frames. Its dual is the natural
representation, so the same stabilizer bound applies on both
sides in each case. At both marks
\(q_w^{l+r/2}\le b^{6k}\) for \(1\le l\le a\), by the numerical
bounds in Theorem~\ref{thm:openenergy870}. A single singular
minor carries at most \(3\mu(\Gamma_\Xi)/q_*\). Summing over
at most \(2f\) places, two components, and \(a\) orders loses
at most \(12af\mu(\Gamma_\Xi)/q_*\). The stated upper bound
for this fraction follows from \(q_*\ge b^{10}\),
\(f\ll_{q_0,d_K}x^{12}\), and \(\log b\ge(\log x)^6\).

On the residues with unit pivots at all unused places remove
\(z(s)\ge b^{f+6t}\). In the preceding squared majorant the
term depending on \(Z=b^{f+6t}\) is at most
\(b^{(-1/4+o(1))kf}\mathcal M\), now for every \(t\).
The bound for different links and the cutoff bounds above
continue to hold under this restriction. The comparison of
square roots proved in Theorem~\ref{thm:equallinks869} thus
has error at most \(b^{-kf/10}\sqrt{\mathcal M}\) in the
weighted norm. Apply it to the subset of \(\Gamma_\Xi\)
with unit pivots at every unused place. Its weighted squared
mass is at least
\((1-12af/q_*)B_0^t\mu(\Gamma_\Xi)\), before the large-\(z\)
part is removed, and maximality makes this at least
\((1-12af/q_*)\mu(\mathcal G)\). This proves the result.
\end{proof}

\begin{proposition}[Evaluation of a marked factor with unit pivots]
\leavevmode\label{prop:localpivot873}
Use the finite factor of Theorem~\ref{thm:orbitfactors871} after
removing the two class representatives as in
Proposition~\ref{prop:sphereunfold872}. Fix an unselected \(w\)
and \(s\in\Omega_\Xi^\sharp\). Over the paired residue algebra
factor \(s\) uniquely as
\[
      s=\bar u_s w_{\bar{\mathbf p}_s}\bar x_s \bar v_s^{-1}.
\]
The bars in this display mean reduction. Here
\(\bar u_s,\bar v_s\) are unipotents, \(\bar x_s\) is a middle
isometry and the entries of \(\bar{\mathbf p}_s\) are units
in both components. The local factor vanishes unless
\(\mathbf p\) consists of local units, \(x\) is integral, and
\[
   \operatorname{red}_w(\mathbf p)=\bar{\mathbf p}_s,\qquad
   \operatorname{red}_w(x)=\bar x_s .
\]
When these conditions hold, the local Haar measures of
Section~\ref{sec:orbitfactors871} give
\[
     K_w(s;w_{\mathbf p}x)=
             \overline{\chi_w(\bar u_s)}\,\chi_w(\bar v_s).
\]
Here the character on a residue unipotent is evaluated on any
integral lift. In particular, let \(\mathfrak L_\Xi\) be the
product of both prime ideals of \(E\) over every unused
\(w\in\mathcal W\). If two rational pivot tuples contribute to
the same block at these places, then for every coordinate \(i\)
\[
  p_i-p_i'\in\mathfrak l\mathcal O_{E,\mathfrak l}
            \quad\text{for every }\mathfrak l\mid\mathfrak L_\Xi .
\]
For \(|\Xi|<f/4\), \(\mathrm N\mathfrak L_\Xi\ge b^{35f}\).
If \(p_i\ne p_i'\) and
\((p_i-p_i')=\mathfrak n/\mathfrak d\) in coprime integral
ideals, then \(\mathfrak L_\Xi\mid\mathfrak n\) and
\[
        \prod_{v\mid\infty}|\Phi_v(p_i-p_i')|^2
              \ \ge\ b^{35f}/\mathrm N\mathfrak d.
\]
\end{proposition}
\begin{proof}
At \(w\) the old level is principal of depth one, the middle
basis is integral, and \(N\) is a unit. The elimination in
Lemma~\ref{lem:pivots871} works over the product of the two
residue fields when each leading minor is a unit. Row operations,
the dual isotropic completion, and the upper and lower triangular
uniqueness use only inversions of these units and the nonsingular
middle Gram matrix. Over the local ring the same elimination
shows that an integral lift of \(s\) has integral factors with
the specified reductions. In particular, if
\(u w_{\mathbf p}x v^{-1}\) reduces to \(s\), the minors force
each rational pivot to be a local unit, and uniqueness over the
local algebra identifies \(u,x,v\) with the factors obtained
by integral elimination. Conversely integral factors of the
specified reductions give residue \(s\). Each residual
unipotent lifts because the radical at a split place is the
block upper unipotent group in one component, with the paired
component determined by inverse adjoint. The possible \(u\) form
one coset of \(U(F_w)\cap J_w\), as do the possible \(v\).
Each has volume one, and the character is constant on each.
This proves the local formula.

The ratios of minors determining \(\bar{\mathbf p}_s\) are fixed
by \(s\). Any two contributing rational tuples therefore have
the indicated congruences in both components, also in the
original class notation because the class changes fix the frame.
Each unused \(w\) contributes \(q_w^2\ge b^{20}\) to the
absolute ideal norm. There are \(2f-|\Xi|>7f/4\) such places.
A nonzero difference is integral at all their primes and its
numerator ideal contains each as a factor. Finally the absolute
field norm of the difference equals both the product of squared
complex absolute values and \(\mathrm N\mathfrak n/\mathrm
N\mathfrak d\). This gives the last inequality.
\end{proof}

\begin{corollary}[The integrated factor at an unused mark]\leavevmode
\label{cor:integratedmark873}
For \(|\Xi|<f/4\) enlarge the mask \(\Omega_\Xi^\sharp\) by
dropping the inequality on \(z(s)\), while retaining all the
fixed vectors and all the unit pivot conditions. The resulting
energy bounds \(\mathcal E_\Xi^\sharp\) from above. In the
sphere formula for this enlarged energy take an unused
\(w\in\mathcal W\), and write \(q=q_w\) and
\(A_w=\mathcal O_E\otimes_{\mathcal O_F}\mathcal O_{F_w}\).
For two pivot tuples \(\mathbf p,\mathbf p'\) and terminal entry
\(\theta\) the entire local factor in the expanded quadratic
expression is zero unless
\[
    \theta\in A_w,\qquad
    \mathbf p,\mathbf p'\in(A_w^\times)^a,\qquad
    \mathbf p\equiv\mathbf p'\pmod{\mathfrak p_w A_w}.
\]
When these conditions hold its value is
\[
       q^{\,2a(2k+2a+1)}
          |\mathrm{GL}_k(\mathbb F_q)|\,
            \delta_w\big((\theta\bar\theta-N^2)/N\big).
\]
The density is evaluated by the recursion in
Theorem~\ref{thm:spheremean872} with split sign \(+1\).
\end{corollary}
\begin{proof}
For a fixed link block the energy is a squared norm. Adding
residues therefore gives the stated upper bound and does not
alter any sum within a block. The enlarged mask is a product
of local conditions. In \eqref{eq:spheretest872} expand the
finite squared sum in the two pivot indices and sum the
residues at \(w\) before integrating the corresponding fiber.
This order is legitimate because all pivot and residue sums
are finite and each test at a finite place is bounded on
compact support.

At an integral middle matrix \(x\) the two local integrals
can both be nonzero for a given residue only when the pivot
tuples have the required unit reductions, by
Proposition~\ref{prop:localpivot873}. In that case the product
of one value with the conjugate of the other is one. Once
\(\bar{\mathbf p}\) and the reduction of \(x\) are fixed, the
residue \(s_w\) is specified uniquely by the two freely chosen
unipotents \(\bar u_s,\bar v_s\). The local residue sum is thus
\(|\bar U_w|^2=q^{2a(2k+2a+1)}\) times the indicator of an
integral middle matrix. For a nonintegral middle matrix every
term is zero.

For fixed \(\omega_w\) the fiber of interest consists of the
middle isometries carrying \(z\) to \(\omega_w-(\theta/N)z\).
Here \(N\) is a unit and the middle lattice is the standard
orthogonal sum. This fiber has an integral point exactly when
\(\theta,\omega_w\) are integral. Indeed the image vector has
the same unit norm as \(z\), so one can extend the isometry
of integral lines as in Proposition~\ref{prop:sphereunfold872}.
All integral points of the fiber then differ on the right by
the maximal integral subgroup of \(H(F_w)\). Its volume is
\(|\mathrm{GL}_k(\mathbb F_q)|\) because its principal kernel
has volume one in the measure of \eqref{eq:spheretest872}.
Finally the measure of integral \(\omega_w\) on its prescribed
norm sphere is precisely \(\delta_w((\theta\bar\theta-N^2)/N)\).
Multiplication gives the formula.
\end{proof}

The remaining terms now involve integral congruences between
the rational pivots at every unused mark. Their denominators
at the other places and their real orbital integrals still
enter the required leading-constant estimate.
\clearpage

\section{References}\label{nav:bibliography-XXV}
\begingroup\def\section*#1{}\begin{thebibliography}{99}\small
\bibitem{XXV-XX-XX-green678} J.~A. Green, The characters of the finite general linear groups, \emph{Trans. Amer. Math. Soc.} \textbf{80} (1955), 402--447. \href{https://doi.org/10.1090/S0002-9947-1955-0072878-2}{doi:10.1090/S0002-9947-1955-0072878-2}.
\bibitem{XXV-XX-XX-milicic684} D.~Mili\v ci\'c, Asymptotic behaviour of matrix coefficients of the discrete series, \emph{Duke Math. J.} \textbf{44} (1977), 59--88, pp.~59--60. \href{https://www.math.utah.edu/~ptrapa/math-library/milicic/milicic-asymptotics-duke-1977.pdf}{Article}.
\bibitem{XXV-XX-XX-blattner684} H.~Hecht and W.~Schmid, A proof of Blattner's conjecture, \emph{Invent. Math.} \textbf{31}, 129--154.
\end{thebibliography}\endgroup
\end{document}
